% relatorio.tex -- mestre do relatório final (specs/2026-09-25_relatorio_latex). Editado à mão.
% O conteúdo por eixo vem de gerado/ (relatorio/latex/build/gera_latex.py, a partir de
% specs/estrutura_eixos.md e relatorio/textos_curados.json) -- não editar gerado/ à mão.
% Compilar:  python relatorio/latex/build/gera_latex.py   (ou, em relatorio/latex/: latexmk)
\documentclass[
  12pt, oneside, a4paper,
  chapter=TITLE, section=TITLE,   % seções primária e secundária em caixa alta no sumário e no texto (NBR 6024)
  brazil, english
]{abntex2}

\usepackage[alf]{abntex2cite}
\usepackage{estilo}

% ------------------------------------------------------------------ dados do documento
\titulo{Diagnóstico da Primeira Infância Carioca}
\autor{Leonardo Aucar \and Bianca Medina \and Caroline Lima \and Waleska Marques \and Maria Norbert}
\local{Rio de Janeiro}
\data{2026}
\instituicao{Prefeitura da Cidade do Rio de Janeiro\par Instituto Pereira Passos\par
  Coordenadoria de Pesquisa, Avaliação e Política Pública}
\tipotrabalho{Relatório técnico}
\preambulo{Relatório técnico de indicadores da primeira infância (0 a 6 anos) no município do Rio de
  Janeiro, organizado pelos eixos da Política Integrada da Primeira Infância, elaborado pelo Instituto
  Pereira Passos em parceria com a Casa Civil.}

\makeatletter
\hypersetup{pdftitle={\@title}, pdfauthor={Instituto Pereira Passos}, pdfsubject={\imprimirpreambulo},
  pdfkeywords={primeira infância}{Rio de Janeiro}{indicadores sociais}{políticas públicas},
  colorlinks=true, linkcolor=tinta, citecolor=tinta, urlcolor=acento, bookmarksdepth=4}
\makeatother

\addto\captionsbrazil{\renewcommand{\bibname}{Fontes}}   % "Referências" -> "Fontes" (pedido do usuário)
\renewcommand{\listtablename}{Lista de tabelas}
\setlength{\parindent}{1.25cm}
\setlength{\parskip}{0.2cm}

\makeindex
\begin{document}
\selectlanguage{brazil}
\frenchspacing

% ================================================================== PRÉ-TEXTUAIS (NBR 10719)
\pretextual
% Capa -- editada à mão. Sóbria: tipografia do site, faixa de 11 cores, azul institucional.
% Logo da Prefeitura/IPP: só depois da autorização de uso (ROADMAP.md) -- espaço reservado no topo.
\begin{capa}
  \thispagestyle{empty}
  \begin{tikzpicture}[remember picture, overlay]
    \fill[ippnavy] (current page.north west) rectangle ([yshift=-0.55cm]current page.north east);
    \node[anchor=south west, inner sep=0] at ([xshift=3cm, yshift=2cm]current page.south west)
      {\faixaespectro[5pt]{\dimexpr\paperwidth-5cm\relax}};
  \end{tikzpicture}
  \vspace*{1.2cm}
  \noindent{\ttfamily\small\color{acento}\addfontfeature{LetterSpace=8}RELATÓRIO TÉCNICO \textperiodcentered{} INSTITUTO PEREIRA PASSOS}
  \par\vspace{5.2cm}
  \noindent{\frauncesdisplay\bfseries\color{ippnavy}\fontsize{40}{46}\selectfont Diagnóstico da\\Primeira Infância\\Carioca\par}
  \vspace{0.9cm}
  \noindent\begin{minipage}{0.82\textwidth}\raggedright\large\color{tinta2}\SingleSpacing
    Indicadores da primeira infância (0 a 6 anos) no município do Rio de Janeiro,
    por eixo da Política Integrada da Primeira Infância\end{minipage}\par
  \vfill
  \noindent{\small\color{tinta2}\SingleSpacing\setlength{\parindent}{0pt}%
    Prefeitura da Cidade do Rio de Janeiro\\
    Instituto Pereira Passos \textperiodcentered{} Coordenadoria de Pesquisa, Avaliação e Política Pública\\[0.4cm]
    {\color{tinta}\bfseries Rio de Janeiro \textperiodcentered{} \imprimirdata}\par}
  \vspace*{1.2cm}
\end{capa}

% Folha de rosto (NBR 10719: órgão responsável, título, responsáveis, local, ano). Editada à mão.
% Equipe informada pelo usuário em 2026-09-25 (sem funções por enquanto).
\begin{folhaderosto}
  \thispagestyle{empty}
  \SingleSpacing
  \begin{center}
    {\ABNTEXchapterfont\normalsize\color{tinta}\MakeUppercase{Prefeitura da Cidade do Rio de Janeiro}\\
     Instituto Pereira Passos\\
     \mdseries Coordenadoria de Pesquisa, Avaliação e Política Pública\par}
    \vspace{3.5cm}
    {\frauncesdisplay\bfseries\color{ippnavy}\fontsize{24}{29}\selectfont\imprimirtitulo\par}
    \vspace{0.6cm}
    {\large\color{tinta2}Relatório técnico\par}
    \vspace{2.5cm}
  \end{center}
  \hspace{.45\textwidth}
  \begin{minipage}{.5\textwidth}
    \small\imprimirpreambulo
  \end{minipage}
  \vspace{2.2cm}
  \begin{center}
    {\small\color{tinta3}\ttfamily\addfontfeature{LetterSpace=6}EQUIPE TÉCNICA\par}
    \vspace{0.3cm}
    Leonardo Aucar\\ Bianca Medina\\ Caroline Lima\\ Waleska Marques\\ Maria Norbert
    \vfill
    {\bfseries\imprimirlocal\\ \imprimirdata}
  \end{center}
\end{folhaderosto}

% Gerado por gera_latex.py -- não editar à mão.
\setlength{\absparsep}{18pt}
\begin{resumo}
Commodo senectus aliquet voluptate consectetur velit aliqua minim proident proident sint cubilia nostrud ullamco senectus aliquet magna exercitation sollicitudin magna senectus do egestas laboris primis do ex culpa vestibulum nisi commodo blandit fugiat ullamco ut cubilia excepteur et magna nisi adipiscing labore faucibus id praesent vestibulum nibh curabitur consectetur minim ad quis exercitation sed fames senectus ullamco mauris elit cubilia lorem do sollicitudin commodo incididunt magna aute tristique officia praesent aute voluptate curae proident sollicitudin mollit netus aliquet adipiscing lorem esse et porttitor sit esse duis ipsum id ad mauris vitae anim consectetur irure veniam sit turpis mauris deserunt ipsum netus ipsum vestibulum laboris nibh sed ea in curae commodo nec excepteur ex fugiat proident curabitur tempor elit senectus sed dolor magna deserunt fames lorem ut enim vitae veniam nisi blandit ex faucibus mollit nulla aliqua primis curabitur cillum dolor senectus culpa quis senectus anim amet posuere est orci qui magna laboris eu netus lorem eget orci posuere primis qui luctus quis fames consectetur aliquet anim exercitation magna aliquet nostrud velit aliquip adipiscing voluptate ea occaecat velit sed proident velit culpa ante posuere minim primis eget qui vestibulum fugiat incididunt laborum aliquip fames nisi ipsum amet nisi ea pariatur luctus mollit dolore sed et magna ut mauris esse vitae praesent metus aliqua elit quis luctus ut purus est est ut tempor officia nunc curabitur duis metus dolor nunc sunt minim cubilia netus nostrud eget consequat sollicitudin egestas esse vitae enim cupidatat curabitur porttitor nisi cupidatat aute labore orci netus do.

\noindent\textbf{Palavras-chave}: Primeira infância. Indicadores sociais. Políticas públicas. Rio de Janeiro (RJ).
\end{resumo}


\pdfbookmark[0]{\listofgraficosname}{lgf}
\listofgraficos*\cleardoublepage
\pdfbookmark[0]{\listofmapasname}{lmp}
\listofmapas*\cleardoublepage
\pdfbookmark[0]{\listtablename}{lot}
\listoftables*\cleardoublepage
% Lista de abreviaturas e siglas -- editada à mão, em ordem alfabética.
\begin{siglas}
  \item[AP] Área de Planejamento (IPP)
  \item[CadÚnico] Cadastro Único para Programas Sociais
  \item[CAP] Coordenadoria de Área Programática de Saúde (SMS-Rio)
  \item[CID-10] Classificação Internacional de Doenças, 10ª revisão
  \item[CTPE] Base municipal de onde foi extraído o CadÚnico (nome por extenso a confirmar)
  \item[DATASUS] Departamento de Informática do Sistema Único de Saúde
  \item[EPI] Programa de imunizações
  \item[IBGE] Instituto Brasileiro de Geografia e Estatística
  \item[INEP] Instituto Nacional de Estudos e Pesquisas Educacionais Anísio Teixeira
  \item[IPP] Instituto Pereira Passos
  \item[IPS] Índice de Progresso Social
  \item[MDS] Ministério do Desenvolvimento e Assistência Social, Família e Combate à Fome
  \item[NV] Nascidos vivos
  \item[PNAD] Pesquisa Nacional por Amostra de Domicílios
  \item[RA] Região Administrativa
  \item[Ripsa] Rede Interagencial de Informações para a Saúde
  \item[RP] Região de Planejamento
  \item[SIDRA] Sistema IBGE de Recuperação Automática
  \item[SIM] Sistema de Informações sobre Mortalidade
  \item[Sinan] Sistema de Informação de Agravos de Notificação
  \item[SINASC] Sistema de Informações sobre Nascidos Vivos
  \item[SISVAN] Sistema de Vigilância Alimentar e Nutricional
  \item[SMS-Rio] Secretaria Municipal de Saúde do Rio de Janeiro
  \item[SM] Salário mínimo
\end{siglas}

\pdfbookmark[0]{\contentsname}{toc}
\tableofcontents*\cleardoublepage

% ================================================================== TEXTUAIS
\textual
% Gerado por gera_latex.py -- não editar à mão.
\chapter{Introdução}\label{cap:introducao}

O Diagnóstico da Primeira Infância é um produto desenvolvido pela Coordenadoria de Pesquisa, Avaliação e Política Pública do Instituto Pereira Passos em parceria com a Casa Civil. Ele é um documento de apoio ao desenho, implementação e monitoramento da nova Política Integrada da Primeira Infância (ref.). Dessa forma, sua estrutura segue, quando possível, os eixos estruturantes da política. Para isso contamos com conjuntos de dados secundários produzidos por diferentes entidades e com a parceria e cooperação de diversos órgãos da prefeitura, aos quais em especial gostaríamos de agradecer a Secretaria de Saúde e sua Secretaria de Vigilância Sanitária.

Esse relatório é um produto vivo. Ele pode e deve ser atualizado conforme novos dados sejam disponibilizados ou novas necessidades surjam. Assim, encorajamos a todos os leitores, especialmente os gestores de políticas públicas, a entrarem em contato (ref. email) com sugestões e solicitações que sejam importantes para execução de seu trabalho. Ao fim desse documento você encontrará todas as tabelas de dados produzidas. Todo o código utilizado para extração e processamento dos dados encontra-se disponível em (ref. github). Esse relatório possui uma versão em PDF e uma versão interativa pública.

Ressaltamos aqui que a produção de dados estatísticos e análises ao nível intramunicipal sempre é um desafio. A maior parte da produção estatística brasileira tem como menor unidade de análise o município e as grandes exceções, como os Censos Demográficos, são esporádicos e muitas vezes não atendem as necessidades totais da gestão pública. Por essa razão, com frequência os dados foram mobilizados de fontes distintas, não sendo concebidos originalmente para comparação, ou derivam de registros administrativos que não se destinavam prioritariamente para a análise. Adicione-se ainda as transformações ao longo do tempo nos processos de coleta e registro de dados, além das diferentes administrações e contextos tecnológicos envolvidos.

Por fim, é preciso considerar que os dados, por si mesmos, não permitem identificar variáveis causais. Com frequência, os processos sociais aqui apontados atuam tanto como causa como consequência de outras variáveis – de ordem econômica, tecnológica, demográfica, etc. Além disso, o período coberto abrange duas significativas crises de saúde pública – a emergência da Zika em 2016 e a Covid-19 em 2020-2021 – além de um período de recessão econômica entre 2015-2016, todas variáveis explicativas significativas, além do contexto particular da cidade do Rio de Janeiro. Assim, a maior recomendação é que os dados levantados sirvam como um mapa, uma visão de trajeto e destino, que possa orientar a construção de indicadores socialmente sensíveis aos diversos contextos e não como metas em si mesmas.

% "Como ler este relatório" -- editado à mão (fatos sobre a estrutura e os dados, não análise).
\section{Como ler este relatório}

Os capítulos 2 a 7 correspondem aos seis eixos da Política Integrada da Primeira Infância. Cada capítulo
abre com os principais achados do eixo, apresenta um indicador por seção --- com gráficos, mapas e o texto
de análise --- e termina com uma síntese. Indicadores previstos na política para os quais ainda não há
dado disponível aparecem identificados como \emph{em desenvolvimento}. As tabelas com os números completos
de cada gráfico e mapa estão nos apêndices, indicadas ao fim de cada seção. A versão interativa deste
relatório, com os mesmos indicadores e textos, está disponível no site do Instituto Pereira Passos.

Três cuidados valem para todo o relatório:

\begin{itemize}
  \item \textbf{População de referência.} Taxas do município usam as estimativas populacionais da Ripsa/Ministério
    da Saúde para o mesmo ano; taxas por bairro, região ou área programática usam a população do Censo Demográfico
    2022, a única disponível nesse nível. As duas não devem ser comparadas diretamente: o Censo 2022 tende a
    subestimar as crianças pequenas.
  \item \textbf{Faixas etárias.} Cada gráfico informa a faixa etária real do dado (por exemplo, 0 a 4 ou 0 a 5 anos),
    que nem sempre coincide com o recorte de 0 a 6 anos da política.
  \item \textbf{Proteção de dados.} Nenhuma contagem do Cadastro Único inferior a 20 crianças ou famílias é
    apresentada abaixo do nível do município; essas células aparecem vazias. Nos mapas de taxa por bairro, a
    escala de cor é limitada ao percentil 95, para que poucos bairros com números muito pequenos não escondam as
    diferenças entre os demais.
\end{itemize}


\section{Os eixos da política}

\begin{itemize}

\item \textbf{Eixo 1 --- Prioridade} (Capítulo~\ref{cap:eixo-1}): 15 indicadores, 2 em desenvolvimento.

\item \textbf{Eixo 2 --- Inclusão} (Capítulo~\ref{cap:eixo-2}): 11 indicadores, 3 em desenvolvimento.

\item \textbf{Eixo 3 --- Família e Cuidados} (Capítulo~\ref{cap:eixo-3}): 8 indicadores.

\item \textbf{Eixo 4 --- Proteção} (Capítulo~\ref{cap:eixo-4}): 8 indicadores, 1 em desenvolvimento.

\item \textbf{Eixo 5 --- Alimentação} (Capítulo~\ref{cap:eixo-5}): 6 indicadores.

\item \textbf{Eixo 6 --- Moradia} (Capítulo~\ref{cap:eixo-6}): 3 indicadores, 3 em desenvolvimento.

\end{itemize}

\eixo{1}{6}{target}

\chapter{Prioridade}\label{cap:eixo-1}

\begin{achados}\begin{itemize}

\item Et incididunt sit orci id nunc aliquip id.

\item Quis fames aliqua nostrud tristique netus turpis do.

\item Eget blandit proident occaecat labore cubilia turpis consequat.

\item Pariatur faucibus lorem tristique ut mollit aliqua lorem.

\item Cubilia primis cubilia egestas laboris officia dolor est.

\end{itemize}\end{achados}

\section{Crianças até 4 anos (número)}

\figura{grafico}{_build/img/censo_0_a_4_serie_total_ano.png}{Crianças até 4 anos (número)}{Elaboração IPP com dados de \citeonline{ipp_datario_censo}.}{graf:censo-0-a-4-serie-total-ano}

Os dados do Censo Demográfico 2022 permitem dimensionar a população de crianças de 0 a 4 anos no município e observar sua distribuição territorial. A análise combina a série municipal e o mapa por bairro, permitindo visualizar tanto a magnitude da população infantil quanto sua concentração no território. A leitura dos números absolutos é importante para identificar os bairros que concentram maior quantidade de crianças nessa faixa etária, mas deve considerar as diferenças no tamanho da população de cada bairro. Esse indicador contribui para contextualizar os demais resultados do relatório, especialmente aqueles relacionados à saúde, proteção social e educação na primeira infância.

\figura{mapa}{_build/img/mapa_censo_0_4_absoluto.jpg}{Crianças de 0 a 4 anos de idade, por bairro (Censo 2022)}{Elaboração IPP com dados de \citeonline{ipp_datario_censo}. Limites de bairros: \citeonline{ipp_limites_bairros}. Sistema de referência SIRGAS 2000.}{mapa:mapa-censo-0-4-absoluto}

O mapa apresenta a distribuição da população de 0 a 4 anos por bairro no município do Rio de Janeiro, segundo o Censo Demográfico 2022. Por apresentar números absolutos, o mapa pode ser utilizado para comparar a quantidade de crianças entre os diferentes territórios e dimensionar o tamanho desse grupo em cada bairro. A informação também serve como referência para análises que envolvam outros indicadores da primeira infância, permitindo relacionar a quantidade de crianças de cada território a diferentes características demográficas e sociais.

\vertabelas{Tabela~\ref{tab:censo-0-a-4-anos-por-ano}, Tabela~\ref{tab:censo-por-bairro}; ver Apêndice~\ref{ap:eixo-1}}

\section{Crianças até 4 anos (percentual)}

\figura{grafico}{_build/img/censo_0_a_4_serie_percentual_ano.png}{Crianças até 4 anos (percentual)}{Elaboração IPP com dados de \citeonline{ipp_datario_censo}.}{graf:censo-0-a-4-serie-percentual-ano}

A participação das crianças de 0 a 4 anos na população total do município diminuiu entre os Censos de 2000, 2010 e 2022. A proporção passou de aproximadamente 7,6\% em 2000 para 5,8\% em 2010 e 5,0\% em 2022. O mapa complementa essa tendência ao mostrar diferenças na participação dessa faixa etária entre os bairros. A análise percentual permite comparar territórios de diferentes tamanhos populacionais, evidenciando o peso relativo das crianças de 0 a 4 anos em cada localidade. Esse indicador contribui para caracterizar a estrutura etária do município e contextualizar as demandas relacionadas à primeira infância.

\figura{mapa}{_build/img/mapa_censo_0_4_percentual.jpg}{Percentual de crianças de 0 a 4 anos, por bairro (Censo 2022)}{Elaboração IPP com dados de \citeonline{ipp_datario_censo}. Limites de bairros: \citeonline{ipp_limites_bairros}. Sistema de referência SIRGAS 2000.}{mapa:mapa-censo-0-4-percentual}

A participação das crianças de 0 a 4 anos na população total do município diminuiu entre os Censos de 2000, 2010 e 2022. A proporção passou de aproximadamente 7,6\% em 2000 para 5,8\% em 2010 e 5,0\% em 2022. O mapa complementa essa tendência ao mostrar diferenças na participação dessa faixa etária entre os bairros. A análise percentual permite comparar territórios de diferentes tamanhos populacionais, evidenciando o peso relativo das crianças de 0 a 4 anos em cada localidade. Esse indicador contribui para caracterizar a estrutura etária do município e contextualizar as demandas relacionadas à primeira infância.

\vertabelas{Tabela~\ref{tab:censo-0-a-4-anos-por-ano}; ver Apêndice~\ref{ap:eixo-1}}

\section{Crianças até 6 anos (número)}

\figura{grafico}{_build/img/populacao_ripsa_0_a_6_por_ano.png}{População de 0 a 6 anos por ano (estimativas Ripsa/MS)}{Elaboração IPP com dados de \citeonline{ms_ripsa_populacao}.}{graf:populacao-ripsa-0-a-6-por-ano}

Exercitation laboris officia culpa nibh magna mollit ad ad incididunt commodo praesent eu vitae pariatur eu quis tristique velit sollicitudin nulla vitae tristique purus eu labore adipiscing ullamco occaecat excepteur cupidatat ut netus posuere exercitation ex cupidatat aliquet lorem lorem exercitation minim cillum vestibulum luctus occaecat sollicitudin faucibus elit deserunt aute commodo curae est reprehenderit pariatur laborum id ipsum et curae velit in nostrud sit culpa magna primis sunt dolore quis netus vestibulum tristique nulla cupidatat primis senectus aute faucibus luctus nostrud nec excepteur orci cupidatat nec excepteur occaecat curae curae faucibus labore esse blandit laborum blandit metus minim commodo quis dolor senectus labore nibh vestibulum turpis sollicitudin ad mauris ea purus nec aliqua amet ante ullamco consequat mollit minim ad eiusmod amet nec pariatur dolor consectetur faucibus eget minim sint commodo primis cubilia officia duis incididunt lorem curabitur aliquet faucibus labore nibh.

\figura{grafico}{_build/img/populacao_ripsa_0_a_6_percentual_por_ano.png}{Participação de 0 a 6 anos na população total (\%, estimativas Ripsa/MS)}{Elaboração IPP com dados de \citeonline{ms_ripsa_populacao}.}{graf:populacao-ripsa-0-a-6-percentual-por-ano}

Id ad consequat praesent occaecat sint nulla ad posuere reprehenderit mollit vestibulum curae est ea elit primis anim egestas et sed proident ut mauris irure posuere porttitor porttitor blandit vitae in eget magna lorem deserunt curae consectetur mauris mauris cillum praesent id est curae esse dolore reprehenderit quis minim curabitur est elit blandit fames metus duis duis do voluptate eu commodo curae eu labore culpa curabitur aute vestibulum commodo ullamco primis eiusmod nunc orci laborum tempor est sit egestas quis anim ipsum fames exercitation nunc nisi culpa veniam dolor ut cillum nostrud eu dolor posuere netus fames dolor quis in labore eget pariatur fugiat duis deserunt elit aliqua fugiat laboris qui curae consequat deserunt ullamco id pariatur.

\vertabelas{Tabela~\ref{tab:populacao-ripsa-0-a-6-por-ano}; ver Apêndice~\ref{ap:eixo-1}}

\section{Nascidos vivos por bairro de residência da mãe (número)}

\figura{grafico}{_build/img/nascidos_vivos_por_ano.png}{Nascidos vivos por ano}{Elaboração IPP com dados de \citeonline{sms_rio_sim}; \citeonline{sms_rio_sinasc}.}{graf:nascidos-vivos-por-ano}

A partir do gráfico de nascidos vivos é possível observar que há uma tendência de queda no número de nascidos vivos na cidade do Rio de Janeiro, com alguns períodos de recuperação. Entre os anos de 2020 e 2021, após um período com uma persistente queda acentuada, a série atinge um patamar muito baixo, período que coincide com o pico da pandemia da COVID-19, apontando que em 2022 a recuperação aparece como um ajuste estatístico da série. Logo, a queda, ainda que não linear, é consistente e aponta para uma redução de cerca de 30\% ao longo da série histórica.

\figura{mapa}{_build/img/mapa_nascidos_vivos_bairro_2025.jpg}{Nascidos vivos por bairro (2025)}{Elaboração IPP com dados de \citeonline{sms_rio_sim}; \citeonline{sms_rio_sinasc}. Limites de bairros: \citeonline{ipp_limites_bairros}. Sistema de referência SIRGAS 2000.}{mapa:mapa-nascidos-vivos-bairro-2025}

Quando se olha para a distribuição espacial desses nascidos vivos pelo território carioca, há uma maior concentração deles na AP5 (Santa Cruz, Campo Grande e Bangu, por exemplo) e AP4 (Jacarepaguá, Barra da Tijuca, Recreio e Taquara, por exemplo). Ao passo que na AP 2, principalmente na Zona Sul há uma quantidade menor. Ficando a AP3 com uma quantidade intermediária.

\vertabelas{Tabela~\ref{tab:nascidos-vivos-por-ano}, Tabela~\ref{tab:tabela-mapa-nascidos-vivos-2025}; ver Apêndice~\ref{ap:eixo-1}}

\section{Nascidos vivos por bairro de residência da mãe (percentual)}

\figura{mapa}{_build/img/mapa_nascidos_vivos_bairro_2025.jpg}{Nascidos vivos por bairro (2025)}{Elaboração IPP com dados de \citeonline{sms_rio_sim}; \citeonline{sms_rio_sinasc}. Limites de bairros: \citeonline{ipp_limites_bairros}. Sistema de referência SIRGAS 2000.}{mapa:mapa-nascidos-vivos-bairro-2025}

Quando se olha para a distribuição espacial desses nascidos vivos pelo território carioca, há uma maior concentração deles na AP5 (Santa Cruz, Campo Grande e Bangu, por exemplo) e AP4 (Jacarepaguá, Barra da Tijuca, Recreio e Taquara, por exemplo). Ao passo que na AP 2, principalmente na Zona Sul há uma quantidade menor. Ficando a AP3 com uma quantidade intermediária.

\vertabelas{Tabela~\ref{tab:tabela-mapa-nascidos-vivos-2025}; ver Apêndice~\ref{ap:eixo-1}}

\section{Taxa de mortalidade neonatal precoce (0 a 6 dias)}

\figura{grafico}{_build/img/taxa_mortalidade_precoce_ano.png}{Taxa de óbitos precoces por ano}{Elaboração IPP com dados de \citeonline{sms_rio_sim}; \citeonline{sms_rio_sinasc}.}{graf:taxa-mortalidade-precoce-ano}

A série histórica permite analisar a evolução da mortalidade neonatal precoce em relação ao número de nascidos vivos no município. Entre 2006 e 2025, a taxa passou de 7,86 para 6,42 óbitos por mil nascidos vivos, embora tenha apresentado oscilações ao longo do período. Em 2024, foram registrados 389 óbitos, o menor número da série, seguido de aumento para 420 em 2025. A leitura conjunta da taxa e dos números absolutos permite distinguir mudanças na ocorrência dos óbitos de variações relacionadas ao número de nascidos vivos, servindo como base para comparações temporais e para o cruzamento com outros indicadores de mortalidade infantil.

\figura{mapa}{_build/img/mapa_obitos_neonatal_precoce_bairro_2025.jpg}{Óbitos precoces (0-6 dias) por bairro (2025)}{Elaboração IPP com dados de \citeonline{sms_rio_sim}; \citeonline{sms_rio_sinasc}. Limites de bairros: \citeonline{ipp_limites_bairros}. Sistema de referência SIRGAS 2000.}{mapa:mapa-obitos-neonatal-precoce-bairro-2025}

A distribuição dos óbitos neonatais precoces por bairro permite identificar como os 420 óbitos registrados em 2025 estão distribuídos territorialmente. Por apresentar números absolutos, o mapa possibilita comparar a quantidade de óbitos entre os bairros e reconhecer onde esses registros estão mais concentrados. A informação pode ser utilizada como base para cruzamentos com o número de nascidos vivos, relacionando a ocorrência dos óbitos ao tamanho da população exposta. O recorte territorial também pode ser relacionado a outros indicadores de mortalidade infantil e características demográficas dos bairros, ampliando a análise do fenômeno.

\figura{mapa}{_build/img/mapa_taxa_mortalidade_precoce_bairro_2025.jpg}{Taxa de óbitos precoces (0-6 dias) por bairro (2025)}{Elaboração IPP com dados de \citeonline{sms_rio_sim}; \citeonline{sms_rio_sinasc}. Limites de bairros: \citeonline{ipp_limites_bairros}. Sistema de referência SIRGAS 2000.}{mapa:mapa-taxa-mortalidade-precoce-bairro-2025}

A taxa de mortalidade neonatal precoce permite comparar os bairros considerando o número de nascidos vivos de cada território, evitando a interpretação baseada apenas na quantidade de óbitos. Em 2025, alguns bairros apresentam taxas elevadas associadas a poucos registros de óbitos e a um número reduzido de nascidos vivos, como Gericinó, com 1 óbito entre 14 nascidos vivos, e Cidade Universitária, com 1 entre 17. Por isso, a leitura territorial da taxa deve considerar também o número absoluto de óbitos e o tamanho do denominador. Esses dados podem servir de base para comparar os territórios e aprofundar a análise em conjunto com outros indicadores.

\vertabelas{Tabela~\ref{tab:mortalidade-neonatal-precoce-por-ano}, Tabela~\ref{tab:mortalidade-neonatal-precoce-bairro-ano}, Tabela~\ref{tab:tabela-mapa-obitos-neonatal-precoce-2025}; ver Apêndice~\ref{ap:eixo-1}}

\section{Taxa de mortalidade neonatal tardia (7 a 27 dias)}

\figura{grafico}{_build/img/taxa_obitos_tardios_ano.png}{Taxa de óbitos tardios por ano}{Elaboração IPP com dados de \citeonline{sms_rio_sim}; \citeonline{sms_rio_sinasc}.}{graf:taxa-obitos-tardios-ano}

A série histórica permite acompanhar a evolução da mortalidade neonatal tardia em relação ao número de nascidos vivos. Entre 2006 e 2025, a taxa passou de 2,70 para 2,60 óbitos por mil nascidos vivos, com oscilações ao longo do período. O maior valor ocorreu em 2020 (3,46), enquanto o menor foi registrado em 2022 (2,29). No mesmo período, os óbitos tardios passaram de 250 para 170. A leitura conjunta desses indicadores permite diferenciar a variação no número de óbitos da variação proporcional em relação aos nascidos vivos e serve de base para comparações temporais com outros indicadores de mortalidade infantil.

\figura{mapa}{_build/img/mapa_taxa_obitos_tardios_bairro_2025.jpg}{Taxa de óbitos tardios (7-27 dias) por bairro (2025)}{Elaboração IPP com dados de \citeonline{sms_rio_sim}; \citeonline{sms_rio_sinasc}. Limites de bairros: \citeonline{ipp_limites_bairros}. Sistema de referência SIRGAS 2000.}{mapa:mapa-taxa-obitos-tardios-bairro-2025}

A taxa de mortalidade neonatal tardia permite comparar os bairros considerando o número de nascidos vivos de cada território. Em 2025, alguns bairros apresentam taxas elevadas mesmo com apenas um óbito, como Cidade Nova, com 1 óbito entre 42 nascidos vivos, e Riachuelo, com 1 entre 66. Dos 161 bairros da tabela, 86 não registraram óbitos tardios. Por isso, a leitura da taxa deve considerar conjuntamente o número de óbitos e o número de nascidos vivos, especialmente nos territórios com menor número de nascimentos. O indicador pode servir de base para comparações territoriais e cruzamentos com outros dados de mortalidade infantil.

\vertabelas{Tabela~\ref{tab:mortalidade-neonatal-tardia-por-ano}, Tabela~\ref{tab:mortalidade-neonatal-tardia-bairro-ano}, Tabela~\ref{tab:tabela-mapa-obitos-neonatal-tardia-2025}; ver Apêndice~\ref{ap:eixo-1}}

\section{Razão de mortalidade materna (durante a gravidez)}

\figura{grafico}{_build/img/obitos_gravidez_por_ano.png}{Óbitos durante gravidez por ano}{Elaboração IPP com dados de \citeonline{sms_rio_sim}; \citeonline{sms_rio_sinasc}.}{graf:obitos-gravidez-por-ano}

A série histórica permite acompanhar a variação dos óbitos ocorridos durante a gravidez no município entre 2006 e 2025. O número de registros passou de 78 em 2006 para 15 em 2025, uma redução de aproximadamente 81\%, embora a trajetória apresente oscilações ao longo do período. Em 2024 foram registrados 4 óbitos, seguido de aumento para 15 em 2025. Por se tratar de número absoluto de óbitos, o indicador permite acompanhar a evolução temporal do evento, mas não representa, isoladamente, uma medida de risco. A série pode servir de base para comparações com outros indicadores de mortalidade materna.

\vertabelas{Tabela~\ref{tab:obitos-gravidez-por-ano}, Tabela~\ref{tab:obitos-gravidez-bairro-ano}, Tabela~\ref{tab:tabela-mapa-obitos-gravidez-2025}; ver Apêndice~\ref{ap:eixo-1}}

\section{Razão de mortalidade materna (durante puerpério)}

\figura{grafico}{_build/img/obitos_puerperio_por_ano.png}{Óbitos durante puerpério por ano}{Elaboração IPP com dados de \citeonline{sms_rio_sim}; \citeonline{sms_rio_sinasc}.}{graf:obitos-puerperio-por-ano}

A série histórica permite acompanhar a variação dos óbitos ocorridos durante o puerpério entre 2006 e 2025. Os registros passaram de 67 em 2006 para 35 em 2025, com oscilações ao longo do período. Destaca-se o aumento observado em 2020 e 2021, quando foram registrados 74 e 109 óbitos, respectivamente, seguido de redução nos anos posteriores. Em 2024 ocorreu o menor número da série, com 33 óbitos, seguido de 35 em 2025. A série pode ser utilizada para comparações temporais com outros indicadores de mortalidade materna.

\vertabelas{Tabela~\ref{tab:obitos-puerperio-por-ano}, Tabela~\ref{tab:obitos-puerperio-bairro-ano}, Tabela~\ref{tab:tabela-mapa-obitos-puerperio-2025}; ver Apêndice~\ref{ap:eixo-1}}

\section{Mortalidade infantil por raça/cor (menores de 1 ano)}

\figura{grafico}{_build/img/obitos_raca_ano.png}{Óbitos de 0 a 364 dias por raça/cor (2006-2025)}{Elaboração IPP com dados de \citeonline{sms_rio_sim}; \citeonline{sms_rio_sinasc}.}{graf:obitos-raca-ano}

A série permite observar mudanças distintas na trajetória dos óbitos segundo raça/cor. Entre 2006 e 2025, os registros para crianças brancas passaram de 468 para 272, enquanto entre crianças pardas passaram de 373 para 396, após oscilações e valores superiores a 500 em alguns anos. Entre crianças pretas, os registros passaram de 89 para 57. A categoria “não informada” também apresentou redução, de 167 para 47, o que altera sua participação na série ao longo do período. Essas diferenças podem ser analisadas em conjunto com os nascidos vivos por raça/cor, disponíveis a partir de 2011, para distinguir composição dos nascimentos e ocorrência dos óbitos.

\figura{grafico}{_build/img/percentual_mortalidade_raca_ano.png}{Percentual de óbitos (0-364 dias) em relação aos nascidos vivos por raça/cor (2011-2025)}{Elaboração IPP com dados de \citeonline{sms_rio_sim}; \citeonline{sms_rio_sinasc}.}{graf:percentual-mortalidade-raca-ano}

A relação entre óbitos e nascidos vivos evidencia diferenças na mortalidade infantil que não aparecem apenas na contagem absoluta. Entre 2011 e 2025, os percentuais de crianças brancas e pardas permaneceram próximos, variando de 1,76\% a 1,13\% e de 1,98\% a 1,43\%, respectivamente. Para crianças pretas, o percentual variou entre 0,50\% e 1,43\%, enquanto as categorias indígena e amarela apresentam oscilações maiores associadas ao pequeno número de registros. A categoria “não informada” também apresenta forte variação, relacionada à quantidade de nascidos classificados nessa categoria. Essas características devem ser consideradas em comparações entre os grupos e na análise da série histórica.

\figura{mapa}{_build/img/mapa_obitos_raca_total_bairro_2025.jpg}{Óbitos de 0 a 364 dias por bairro (2025)}{Elaboração IPP com dados de \citeonline{sms_rio_sim}; \citeonline{sms_rio_sinasc}. Limites de bairros: \citeonline{ipp_limites_bairros}. Sistema de referência SIRGAS 2000.}{mapa:mapa-obitos-raca-total-bairro-2025}

A distribuição territorial dos óbitos infantis evidencia diferenças na quantidade de registros entre os bairros do município. Em 2025, Santa Cruz concentrou 53 óbitos, seguido por Campo Grande, com 40, e Jacarepaguá, com 32. Dos 167 bairros presentes na tabela, 139 registraram ao menos um óbito e 28 não apresentaram registros. Como o mapa utiliza números absolutos, essas diferenças podem ser relacionadas ao número de nascidos vivos de cada território, permitindo complementar a análise com a taxa de mortalidade infantil e outros recortes demográficos.

\figura{mapa}{_build/img/mapa_taxa_obitos_raca_total_bairro_2025.jpg}{Taxa de mortalidade infantil (0-364 dias) por bairro (2025)}{Elaboração IPP com dados de \citeonline{sms_rio_sim}; \citeonline{sms_rio_sinasc}. Limites de bairros: \citeonline{ipp_limites_bairros}. Sistema de referência SIRGAS 2000.}{mapa:mapa-taxa-obitos-raca-total-bairro-2025}

A taxa de mortalidade infantil permite comparar os bairros considerando a relação entre os óbitos e os nascidos vivos de cada território. Em 2025, Cidade Nova e Gericinó apresentaram a maior taxa registrada, de 7,14\%, mas com números diferentes de óbitos e nascidos vivos: 3 óbitos entre 42 nascidos vivos em Cidade Nova e 1 entre 14 em Gericinó. Cidade Universitária apresentou 5,88\%, com 1 óbito entre 17 nascidos vivos. A comparação entre taxa, número de óbitos e nascidos vivos permite qualificar a leitura das diferenças territoriais e serve de base para relacionar o indicador a outros recortes da mortalidade infantil.

\vertabelas{Tabela~\ref{tab:mortalidade-raca-bairro-ano}, Tabela~\ref{tab:mortalidade-raca-municipio-ano}; ver Apêndice~\ref{ap:eixo-1}}

\section{Mortalidade infantil por causas evitáveis (0 a 6, 7 a 27, 28 a 364 dias)}

\figura{grafico}{_build/img/obitos_evitaveis_total_cap_ano.png}{Total de óbitos (todas as causas), menores de 5 anos, por CAP (2006-2025)}{Elaboração IPP com dados de \citeonline{sms_rio_sim}.}{graf:obitos-evitaveis-total-cap-ano}

A série permite acompanhar a evolução dos óbitos de menores de 5 anos nas diferentes Áreas Programáticas de Saúde (CAP) e comparar como esses registros variam entre os territórios ao longo do tempo. No conjunto das CAPs, os óbitos passaram de 1.311 em 2006 para 856 em 2025, com redução ao longo da série, embora tenha ocorrido aumento entre 2024 e 2025, de 819 para 856 registros. A distribuição territorial também apresenta diferenças importantes em 2025, com 150 óbitos na CAP 4.0 e 131 na CAP 3.3. Esses dados podem ser relacionados à composição das causas e ao percentual de óbitos evitáveis em cada CAP.

\figura{grafico}{_build/img/obitos_evitaveis_cap_menores_1_ano_ano.png}{Mortalidade infantil por causas evitáveis (0 a 6, 7 a 27, 28 a 364 dias)}{Elaboração IPP com dados de \citeonline{sms_rio_sim}.}{graf:obitos-evitaveis-cap-menores-1-ano-ano}

A série permite acompanhar a variação dos óbitos classificados como causas evitáveis entre as CAPs ao longo do período. No conjunto das CAPs, os registros passaram de 755 em 2006 para 491 em 2025, embora com oscilações ao longo da série. Em 2025, a CAP 4.0 apresentou 89 óbitos, a CAP 3.3 registrou 84 e a CAP 5.2, 67. A distribuição entre as CAPs pode ser relacionada ao percentual de óbitos evitáveis e às demais categorias de causas presentes na base, permitindo analisar como a quantidade de óbitos e sua composição variam entre os territórios.

\figura{grafico}{_build/img/obitos_evitaveis_cap_1_a_4_anos_ano.png}{Mortalidade infantil por causas evitáveis (0 a 6, 7 a 27, 28 a 364 dias)}{Elaboração IPP com dados de \citeonline{sms_rio_sim}.}{graf:obitos-evitaveis-cap-1-a-4-anos-ano}

A série apresenta a variação dos óbitos classificados como causas evitáveis entre as CAPs para crianças de 1 a 4 anos. As contagens são reduzidas em comparação com a faixa de menores de 1 ano e apresentam oscilações entre os anos e territórios. Em 2025, os registros variaram entre as CAPs, sem que pequenas diferenças anuais devam ser interpretadas isoladamente. A série pode ser utilizada para acompanhar a ocorrência do evento ao longo do tempo e relacioná-la à distribuição das demais causas de óbito nessa faixa etária, mantendo a ressalva de que o pequeno número de registros limita comparações mais detalhadas entre CAPs.

\figura{grafico}{_build/img/obitos_evitaveis_cap_menores_5_anos_ano.png}{Mortalidade infantil por causas evitáveis (0 a 6, 7 a 27, 28 a 364 dias)}{Elaboração IPP com dados de \citeonline{sms_rio_sim}.}{graf:obitos-evitaveis-cap-menores-5-anos-ano}

A série permite acompanhar a ocorrência de óbitos por causas evitáveis na primeira infância e observar como esse conjunto se distribui entre as CAPs ao longo do período. Como a faixa de menores de 5 anos reúne os óbitos de menores de 1 ano e de 1 a 4 anos, o indicador oferece uma visão agregada da ocorrência de causas evitáveis nessa etapa da vida. A comparação entre as CAPs pode ser complementada pela análise separada das duas faixas etárias, permitindo verificar em que grupo etário se concentram as diferenças observadas e relacioná-las à composição das demais causas de óbito.

\figura{grafico}{_build/img/percentual_evitaveis_cap_menores_1_ano_ano.png}{Mortalidade infantil por causas evitáveis (0 a 6, 7 a 27, 28 a 364 dias)}{Elaboração IPP com dados de \citeonline{sms_rio_sim}.}{graf:percentual-evitaveis-cap-menores-1-ano-ano}

O percentual permite analisar a participação das causas evitáveis no conjunto dos óbitos de menores de 1 ano em cada CAP, complementando a leitura feita a partir dos números absolutos. Ao longo da série, os valores oscilam entre as CAPs e entre os anos, mostrando que a composição dos óbitos varia mesmo quando a quantidade absoluta de registros apresenta comportamento semelhante. A comparação desse indicador com o número total de óbitos permite distinguir mudanças na quantidade de ocorrências de mudanças na participação das causas evitáveis. O indicador também pode ser relacionado às demais categorias de causas presentes na base para aprofundar a análise territorial.

\figura{grafico}{_build/img/percentual_evitaveis_cap_1_a_4_anos_ano.png}{Mortalidade infantil por causas evitáveis (0 a 6, 7 a 27, 28 a 364 dias)}{Elaboração IPP com dados de \citeonline{sms_rio_sim}.}{graf:percentual-evitaveis-cap-1-a-4-anos-ano}

O percentual permite analisar a participação das causas evitáveis no conjunto dos óbitos de crianças de 1 a 4 anos em cada CAP. A série apresenta oscilações entre os territórios e os anos, especialmente em uma faixa etária com menor número absoluto de registros. Em alguns anos, determinadas CAPs atingem valores próximos ou iguais a 100\%, indicando que os óbitos registrados naquele recorte foram classificados como evitáveis. Por isso, o percentual deve ser analisado junto ao número absoluto de óbitos, permitindo distinguir mudanças na composição das causas de variações decorrentes do pequeno número de registros.

\figura{grafico}{_build/img/percentual_evitaveis_cap_menores_5_anos_ano.png}{Mortalidade infantil por causas evitáveis (0 a 6, 7 a 27, 28 a 364 dias)}{Elaboração IPP com dados de \citeonline{sms_rio_sim}.}{graf:percentual-evitaveis-cap-menores-5-anos-ano}

O percentual permite analisar a participação das causas evitáveis no conjunto dos óbitos de menores de 5 anos em cada CAP. Ao longo da série, os valores apresentam oscilações entre os territórios, permitindo observar mudanças na composição dos óbitos ao longo do tempo. Como esse recorte reúne as faixas de menores de 1 ano e de 1 a 4 anos, o indicador oferece uma visão agregada da participação das causas evitáveis na primeira infância. A análise conjunta com o número absoluto de óbitos permite distinguir alterações na proporção de causas evitáveis de variações no volume total de registros.

\figura{grafico}{_build/img/taxa_mortalidade_evitaveis_menores_5_ano.png}{Taxa de mortalidade por causas evitáveis (\textless{} 5 anos), por mil nascidos vivos (2006-2025)}{Elaboração IPP com dados de \citeonline{sms_rio_sim}.}{graf:taxa-mortalidade-evitaveis-menores-5-ano}

A taxa relaciona os óbitos por causas evitáveis ao número de nascidos vivos, permitindo acompanhar a ocorrência do indicador ao longo do tempo. Em 2006, foram registrados 1.353 óbitos e taxa de 16,48 por mil nascidos vivos. Em 2025, foram 856 óbitos e 14,58 por mil, enquanto o menor valor da série ocorreu em 2014, com 13,10 por mil. A trajetória apresenta oscilações, inclusive nos anos mais recentes, quando a taxa passou de 14,28 em 2022 para 14,97 em 2023, 14,28 em 2024 e 14,58 em 2025. A série permite acompanhar conjuntamente a ocorrência dos óbitos e sua relação com os nascidos vivos.

\figura{mapa}{_build/img/mapa_obitos_evitaveis_menores_1_ano_cap_2025.jpg}{Mortalidade infantil por causas evitáveis (0 a 6, 7 a 27, 28 a 364 dias)}{Elaboração IPP com dados de \citeonline{sms_rio_sim}. Limites de bairros: \citeonline{ipp_limites_bairros}. Sistema de referência SIRGAS 2000.}{mapa:mapa-obitos-evitaveis-menores-1-ano-cap-2025}

A distribuição territorial dos óbitos por causas evitáveis permite comparar a quantidade de registros entre as CAPs em 2025. Entre os menores de 1 ano, foram registrados 491 óbitos, sendo 89 na CAP 4.0, 84 na CAP 3.3 e 67 na CAP 5.2. Como são números absolutos, essa distribuição pode ser analisada em conjunto com o total de óbitos e o percentual de causas evitáveis de cada CAP. O cruzamento entre esses indicadores permite distinguir diferenças no volume de registros de diferenças na participação das causas evitáveis.

\figura{mapa}{_build/img/mapa_obitos_evitaveis_1_a_4_anos_cap_2025.jpg}{Mortalidade infantil por causas evitáveis (0 a 6, 7 a 27, 28 a 364 dias)}{Elaboração IPP com dados de \citeonline{sms_rio_sim}. Limites de bairros: \citeonline{ipp_limites_bairros}. Sistema de referência SIRGAS 2000.}{mapa:mapa-obitos-evitaveis-1-a-4-anos-cap-2025}

A distribuição territorial permite comparar os registros de óbitos por causas evitáveis entre as CAPs para crianças de 1 a 4 anos em 2025. A CAP 3.3 apresentou 12 óbitos, seguida pelas CAPs 5.3, com 10, e 5.1, com 9. Como essa faixa etária apresenta números absolutos menores, as diferenças entre as CAPs devem ser interpretadas junto ao total de óbitos de cada território. O percentual de causas evitáveis complementa essa leitura, permitindo verificar quanto esses registros representam dentro do conjunto de óbitos de cada CAP.

\figura{mapa}{_build/img/mapa_obitos_evitaveis_menores_5_anos_cap_2025.jpg}{Mortalidade infantil por causas evitáveis (0 a 6, 7 a 27, 28 a 364 dias)}{Elaboração IPP com dados de \citeonline{sms_rio_sim}. Limites de bairros: \citeonline{ipp_limites_bairros}. Sistema de referência SIRGAS 2000.}{mapa:mapa-obitos-evitaveis-menores-5-anos-cap-2025}

Em 2025, foram registrados 557 óbitos por causas evitáveis entre menores de 5 anos. A maior concentração ocorreu nas CAPs 4.0 e 3.3, com 96 óbitos cada, seguidas pela CAP 5.2, com 74. O cruzamento com as faixas etárias mostra que os óbitos evitáveis de menores de 1 ano respondem pela maior parte dos registros em todas as CAPs. Na CAP 4.0, por exemplo, foram 89 dos 96 óbitos evitáveis; na CAP 3.3, 84 dos 96. Assim, o mapa de menores de 5 anos deve ser lido em conjunto com os recortes etários anteriores, que ajudam a identificar onde está concentrada essa ocorrência.

\figura{mapa}{_build/img/mapa_percentual_evitaveis_menores_1_ano_cap_2025.jpg}{Mortalidade infantil por causas evitáveis (0 a 6, 7 a 27, 28 a 364 dias)}{Elaboração IPP com dados de \citeonline{sms_rio_sim}. Limites de bairros: \citeonline{ipp_limites_bairros}. Sistema de referência SIRGAS 2000.}{mapa:mapa-percentual-evitaveis-menores-1-ano-cap-2025}

O percentual permite observar a participação das causas evitáveis no conjunto dos óbitos de menores de 1 ano em cada CAP, complementando o mapa de números absolutos. A comparação entre os dois indicadores ajuda a distinguir territórios com maior quantidade de óbitos daqueles em que as causas evitáveis representam uma parcela maior do total registrado. Essa leitura também permite identificar situações em que CAPs com volumes diferentes apresentam proporções semelhantes ou distintas. O indicador pode ainda ser relacionado à série histórica por CAP, verificando se as diferenças observadas em 2025 acompanham ou se afastam do comportamento registrado nos anos anteriores.

\figura{mapa}{_build/img/mapa_percentual_evitaveis_1_a_4_anos_cap_2025.jpg}{Mortalidade infantil por causas evitáveis (0 a 6, 7 a 27, 28 a 364 dias)}{Elaboração IPP com dados de \citeonline{sms_rio_sim}. Limites de bairros: \citeonline{ipp_limites_bairros}. Sistema de referência SIRGAS 2000.}{mapa:mapa-percentual-evitaveis-1-a-4-anos-cap-2025}

O percentual evidencia quanto as causas evitáveis representam do conjunto de óbitos de crianças de 1 a 4 anos em cada CAP, permitindo uma leitura diferente daquela obtida pelos números absolutos. Nesse recorte, em que os registros são menos numerosos, diferenças na proporção podem resultar de pequenas variações na quantidade de óbitos. Por isso, o indicador ganha maior significado quando analisado junto aos números absolutos e ao percentual das demais faixas etárias. A comparação com o mapa de menores de 1 ano também permite verificar se a participação das causas evitáveis apresenta comportamento territorial semelhante ou diferente entre os grupos etários.

\figura{mapa}{_build/img/mapa_percentual_evitaveis_menores_5_anos_cap_2025.jpg}{Mortalidade infantil por causas evitáveis (0 a 6, 7 a 27, 28 a 364 dias)}{Elaboração IPP com dados de \citeonline{sms_rio_sim}. Limites de bairros: \citeonline{ipp_limites_bairros}. Sistema de referência SIRGAS 2000.}{mapa:mapa-percentual-evitaveis-menores-5-anos-cap-2025}

O percentual mostra o peso das causas evitáveis no conjunto dos óbitos de menores de 5 anos em cada CAP. Em 2025, a participação variou de 40,00\% na CAP 2.1 a 75,00\% na CAP 2.2. A comparação com os números absolutos acrescenta contexto à leitura: a CAP 2.2 teve 15 óbitos evitáveis entre 20 registros, enquanto a CAP 3.3 teve 96 entre 131. Assim, o percentual permite analisar a composição dos óbitos sem confundi-la com o volume de registros, e pode ser relacionado aos recortes de menores de 1 ano e de 1 a 4 anos.

\vertabelas{\texttt{mortalidade\_evitaveis\_cap\_faixa\_ano.csv} (formato digital), Tabela~\ref{tab:mortalidade-evitaveis-cap-2025}, Tabela~\ref{tab:taxa-mortalidade-evitaveis-menores-5-municipio-ano}; ver Apêndice~\ref{ap:eixo-1}}

\section{Taxa de mortalidade na primeira infância (menores de 1, 1-4, 0-5 anos)}

\figura{grafico}{_build/img/taxa_mortalidade_infantil_ano.png}{Taxa de mortalidade infantil (0-364 dias) por ano}{Elaboração IPP com dados de \citeonline{sms_rio_sim}; \citeonline{sms_rio_sinasc}.}{graf:taxa-mortalidade-infantil-ano}

A série histórica mostra oscilações na taxa de mortalidade infantil entre 2006 e 2025. O indicador passou de 13,37 óbitos por mil nascidos vivos em 2006 para 13,06 em 2025, atingindo seu menor valor em 2017, com 11,26 por mil. A partir de 2018, observa-se elevação gradual da taxa, chegando a 13,03 em 2024 e 13,06 em 2025. A leitura conjunta da taxa com o número de óbitos e de nascidos vivos permite acompanhar como o indicador se comporta em diferentes períodos e relacionar sua trajetória às demais taxas de mortalidade na primeira infância.

\figura{grafico}{_build/img/taxa_mortalidade_pos_neonatal_ano.png}{Taxa de mortalidade pós-neonatal (28-364 dias) por ano}{Elaboração IPP com dados de \citeonline{sms_rio_sim}; \citeonline{sms_rio_sinasc}.}{graf:taxa-mortalidade-pos-neonatal-ano}

Aute minim eget consequat nunc sollicitudin ad id magna laborum porttitor deserunt consectetur blandit cubilia aliqua occaecat nec do tristique aute aliquet do porttitor et primis fames consectetur curae senectus et adipiscing amet nisi non duis metus duis aliqua vestibulum curabitur lorem laboris reprehenderit voluptate voluptate non in fugiat culpa orci voluptate sed in laboris sunt ipsum vestibulum nunc purus adipiscing voluptate purus veniam nunc minim amet pariatur nec ante vitae laboris irure ea reprehenderit do exercitation tempor blandit quis laboris sollicitudin qui egestas luctus voluptate lorem netus dolore nec irure orci cillum mollit dolor egestas consequat nisi aute tristique elit primis ut ea commodo ut quis curae eu vitae ipsum enim nec amet.

\figura{mapa}{_build/img/mapa_mortalidade_infantil_bairro_2025.jpg}{Óbitos infantis (0-364 dias) por bairro (2025)}{Elaboração IPP com dados de \citeonline{sms_rio_sim}; \citeonline{sms_rio_sinasc}. Limites de bairros: \citeonline{ipp_limites_bairros}. Sistema de referência SIRGAS 2000.}{mapa:mapa-mortalidade-infantil-bairro-2025}

Eiusmod occaecat cillum minim netus et duis ipsum ullamco posuere esse quis dolor incididunt sunt ipsum nec sollicitudin deserunt anim esse posuere nisi dolore aliqua amet nostrud tristique ullamco posuere do cupidatat reprehenderit netus dolore vitae blandit tristique faucibus labore in aliqua amet incididunt praesent eget sunt deserunt cupidatat nunc culpa voluptate curae cubilia nostrud sint curabitur nostrud luctus luctus esse consectetur anim metus reprehenderit sit adipiscing eget quis nec ullamco exercitation egestas amet proident exercitation luctus id curae exercitation posuere curabitur sunt sint posuere curae ad proident et veniam adipiscing faucibus sunt eget luctus nostrud ipsum fames turpis proident netus posuere orci irure cillum occaecat aliqua amet aliquip eiusmod dolor cubilia consectetur quis curae sunt cupidatat purus nisi sed turpis nostrud orci curae voluptate id esse lorem laborum sunt velit quis faucibus magna aliquet faucibus ex dolor ad sunt ante exercitation excepteur vitae qui velit egestas consequat curae magna fugiat duis proident et porttitor nulla dolore elit sint et sed eiusmod amet pariatur occaecat excepteur laborum velit occaecat praesent purus velit aliquet senectus irure aliquip dolor qui nulla deserunt enim irure culpa nunc elit labore vitae.

\figura{mapa}{_build/img/mapa_taxa_mortalidade_infantil_bairro_2025.jpg}{Taxa de mortalidade infantil (0-364 dias) por bairro (2025)}{Elaboração IPP com dados de \citeonline{sms_rio_sim}; \citeonline{sms_rio_sinasc}. Limites de bairros: \citeonline{ipp_limites_bairros}. Sistema de referência SIRGAS 2000.}{mapa:mapa-taxa-mortalidade-infantil-bairro-2025}

Nec cillum purus reprehenderit posuere do laboris est est aliquip sed nec laborum nisi mollit turpis senectus officia purus deserunt purus officia laboris eget minim primis minim pariatur sunt non netus occaecat deserunt consequat adipiscing esse irure labore sed ante incididunt deserunt mollit nec elit eget qui duis eiusmod aliquet ipsum et primis nostrud magna nulla ullamco ante ante vitae tempor purus blandit mauris aute proident consectetur metus primis ullamco voluptate id tempor cillum enim amet sollicitudin nibh sunt anim cubilia purus laboris consectetur irure aliqua incididunt cillum enim labore tristique metus voluptate purus nulla enim elit voluptate tempor eu velit faucibus porttitor nec tristique lorem adipiscing culpa elit labore in ullamco velit adipiscing velit ea veniam eu nunc id nec mauris dolor labore occaecat labore ut ante voluptate sed in incididunt consequat dolor eget cillum labore dolore praesent tristique.

\figura{mapa}{_build/img/mapa_taxa_mortalidade_pos_neonatal_bairro_2025.jpg}{Taxa de mortalidade pós-neonatal (28-364 dias) por bairro (2025)}{Elaboração IPP com dados de \citeonline{sms_rio_sim}; \citeonline{sms_rio_sinasc}. Limites de bairros: \citeonline{ipp_limites_bairros}. Sistema de referência SIRGAS 2000.}{mapa:mapa-taxa-mortalidade-pos-neonatal-bairro-2025}

Senectus pariatur sint anim tristique magna anim consectetur labore orci in eget velit ullamco ante adipiscing porttitor veniam reprehenderit eget veniam voluptate blandit esse cupidatat eget culpa id faucibus nostrud sollicitudin ullamco faucibus reprehenderit do non porttitor faucibus faucibus praesent praesent consequat est aliquet anim egestas anim primis reprehenderit laboris ad ipsum do blandit magna luctus ex posuere amet sunt aliqua id praesent commodo adipiscing vestibulum metus reprehenderit primis officia aute mauris culpa do faucibus laborum laboris pariatur in ad sunt sed ea ipsum blandit quis anim praesent egestas ipsum eu turpis nibh qui duis curabitur primis eget cubilia velit non quis laborum consequat nostrud lorem laborum nec reprehenderit ea labore aute cubilia qui vitae consectetur do ullamco lorem anim fames senectus metus lorem vitae sint eget purus nostrud irure cubilia esse.

\vertabelas{Tabela~\ref{tab:mortalidade-infantil-pos-neonatal-total-por-ano}, Tabela~\ref{tab:mortalidade-infantil-pos-neonatal-total-bairro-ano}; ver Apêndice~\ref{ap:eixo-1}}

\section{Mortalidade infantil por causas evitáveis, por tipo de causa (grupo/subgrupo CID-10)}

\figura{grafico}{_build/img/obitos_causas_evitaveis_grupo_ano.png}{Óbitos por causas evitáveis (0-364 dias) por grupo (1996-2025)}{Elaboração IPP com dados de \citeonline{sms_rio_sim}.}{graf:obitos-causas-evitaveis-grupo-ano}

Entre 1996 e 2025, observa-se uma redução expressiva dos óbitos de crianças de 0 a 364 dias por causas evitáveis. O número caiu de cerca de 1,5 mil registros no início da série para 502 em 2025. As causas mal definidas também apresentaram forte redução, chegando a 15 registros, enquanto as demais causas recuaram de forma mais moderada, alcançando 206 óbitos em 2025.

\figura{grafico}{_build/img/obitos_causas_evitaveis_subgrupo_ano.png}{Óbitos por causas evitáveis (0-364 dias) por subgrupo (1996-2025)}{Elaboração IPP com dados de \citeonline{sms_rio_sim}.}{graf:obitos-causas-evitaveis-subgrupo-ano}

Nos subgrupos de causas evitáveis, observa-se redução ao longo da série na maior parte das categorias. Os óbitos reduzíveis por adequada atenção à mulher na gestação permanecem como o principal grupo em 2025, com 263 registros. Também houve queda expressiva nos óbitos relacionados à atenção ao recém-nascido, que passaram de 543 em 1996 para 67 em 2025, enquanto aqueles relacionados à atenção à mulher no parto chegaram a 68 registros.

\figura{grafico}{_build/img/obitos_causas_evitaveis_grupo_0_a_6_dias_ano.png}{Óbitos por causas evitáveis (0-6 dias) por grupo (1996-2025)}{Elaboração IPP com dados de \citeonline{sms_rio_sim}.}{graf:obitos-causas-evitaveis-grupo-0-a-6-dias-ano}

Entre 1996 e 2025, os óbitos de crianças de 0 a 6 dias apresentaram queda expressiva.  As causas evitáveis permaneceram como o principal grupo durante toda a série, reduzindo-se para 269 registros em 2025. No mesmo ano, as demais causas somaram 62 óbitos, enquanto as causas mal definidas ficaram em apenas 2 registros.

\figura{grafico}{_build/img/obitos_causas_evitaveis_subgrupo_0_a_6_dias_ano.png}{Óbitos por causas evitáveis (0-6 dias) por subgrupo (1996-2025)}{Elaboração IPP com dados de \citeonline{sms_rio_sim}.}{graf:obitos-causas-evitaveis-subgrupo-0-a-6-dias-ano}

Entre 1996 e 2025, os óbitos de crianças de 0 a 6 dias apresentaram redução em praticamente todos os subgrupos. As causas reduzíveis por atenção à mulher na gestação permaneceram como o principal componente, chegando a 183 registros em 2025. Também houve queda importante nos óbitos relacionados à atenção ao recém-nascido e à atenção à mulher no parto, enquanto as causas mal definidas ficaram em apenas 2 registros no final da série.

\figura{grafico}{_build/img/obitos_causas_evitaveis_grupo_7_a_27_dias_ano.png}{Óbitos por causas evitáveis (7-27 dias) por grupo (1996-2025)}{Elaboração IPP com dados de \citeonline{sms_rio_sim}.}{graf:obitos-causas-evitaveis-grupo-7-a-27-dias-ano}

Entre 1996 e 2025, os óbitos entre 7 e 27 dias de vida apresentaram tendência de queda. As causas evitáveis permaneceram como o principal grupo, passando de níveis próximos a 250-280 registros no início da série para 99 em 2025. As demais causas também diminuíram, chegando a 42 registros, enquanto as causas mal definidas ficaram praticamente zeradas no final do período.

\figura{grafico}{_build/img/obitos_causas_evitaveis_subgrupo_7_a_27_dias_ano.png}{Óbitos por causas evitáveis (7-27 dias) por subgrupo (1996-2025)}{Elaboração IPP com dados de \citeonline{sms_rio_sim}.}{graf:obitos-causas-evitaveis-subgrupo-7-a-27-dias-ano}

Entre 1996 e 2025, os óbitos entre 7 e 27 dias de vida apresentaram tendência de queda na maior parte dos subgrupos.  A maior queda ocorreu nas causas reduzíveis por adequada atenção ao recém-nascido, que passaram de 157 registros em 1996 para 20 em 2025. Já as causas relacionadas à atenção a mulher na gestação permaneceram como o principal subgrupo no final da série, com 64 óbitos em 2025. Os demais aparesentaram valores mais abaixos.

\figura{grafico}{_build/img/obitos_causas_evitaveis_grupo_28_a_364_dias_ano.png}{Óbitos por causas evitáveis (28-364 dias) por grupo (1996-2025)}{Elaboração IPP com dados de \citeonline{sms_rio_sim}.}{graf:obitos-causas-evitaveis-grupo-28-a-364-dias-ano}

Entre 1996 e 2025, os óbitos entre 28 e 364 dias apresentaram tendência de redução. As causas evitáveis permaneceram como o principal grupo, passando de 461 registros em 1996 para 134 em 2025. As demais causas também diminuíram, chegando a 102 óbitos, enquanto as causas mal definidas apresentaram a maior redução proporcional, passando de 108 para 13 registros.

\figura{grafico}{_build/img/obitos_causas_evitaveis_subgrupo_28_a_364_dias_ano.png}{Óbitos por causas evitáveis (28-364 dias) por subgrupo (1996-2025)}{Elaboração IPP com dados de \citeonline{sms_rio_sim}.}{graf:obitos-causas-evitaveis-subgrupo-28-a-364-dias-ano}

Entre 1996 e 2025, os óbitos entre 28 e 364 dias apresentaram queda na maior parte dos subgrupos. O principal destaque é a redução dos óbitos reduzíveis por ações de diagnóstico e tratamento adequado, que saíram de patamares muito elevados no início da série e chegaram a 43 registros em 2025. No final do período, os maiores valores ficaram em ações de promoção vinculadas às ações de atenção, com 56 óbitos, seguidas por diagnóstico e tratamento adequado, enquanto os demais subgrupos apresentaram números mais baixos.

\figura{grafico}{_build/img/obitos_causas_evitaveis_subgrupo_faixa_2025.png}{Óbitos por causas evitáveis por subgrupo e faixa etária (2025)}{Elaboração IPP com dados de \citeonline{sms_rio_sim}.}{graf:obitos-causas-evitaveis-subgrupo-faixa-2025}

Vitae cubilia tristique non metus curabitur nostrud qui enim cubilia ullamco senectus curae nisi nostrud cillum ut quis lorem netus duis enim et turpis posuere posuere nec aliqua ipsum est aliquip culpa pariatur quis et aliquet blandit faucibus aute aliquip vestibulum cupidatat aliquet adipiscing minim reprehenderit curabitur reprehenderit porttitor aliquip et faucibus orci in est consectetur aliquet exercitation vestibulum praesent minim vestibulum faucibus do sollicitudin sollicitudin mollit est est consequat elit pariatur nostrud mollit ad sunt porttitor deserunt et nulla magna blandit ante ut velit reprehenderit labore curae deserunt sit nibh nisi tempor enim adipiscing orci eget cubilia veniam minim fames.

\figura{grafico}{_build/img/obitos_evitaveis_menores_1_ano_subgrupo_ano.png}{Óbitos por causas evitáveis (menores de 1 ano) por subgrupo (2006-2025)}{Elaboração IPP com dados de \citeonline{sms_rio_sim}.}{graf:obitos-evitaveis-menores-1-ano-subgrupo-ano}

Curae reprehenderit nostrud sed aliquip metus nec orci curabitur blandit laboris tempor porttitor culpa ea curabitur ipsum primis eget eget elit fames reprehenderit voluptate anim sint fugiat elit elit proident faucibus ipsum do nostrud egestas metus orci porttitor orci orci faucibus blandit dolore proident velit mauris tristique aliquet dolore exercitation nibh deserunt aliquip sunt officia qui minim fames amet voluptate praesent do veniam luctus sint esse irure consectetur purus quis sit aliquip consectetur curae duis curabitur posuere vestibulum aliquet fames netus mauris curabitur est officia id id quis consequat ex exercitation et minim consectetur aliquip voluptate eiusmod mollit dolore elit ut velit duis aliqua aliqua orci cubilia vitae aute elit sollicitudin luctus minim dolore laboris ipsum ut tristique irure consectetur aute aliquet sed luctus adipiscing culpa adipiscing eget irure exercitation ipsum enim est.

\figura{grafico}{_build/img/obitos_evitaveis_1_a_4_anos_subgrupo_ano.png}{Óbitos por causas evitáveis (menores de 1 ano) por subgrupo (2006-2025)}{Elaboração IPP com dados de \citeonline{sms_rio_sim}.}{graf:obitos-evitaveis-1-a-4-anos-subgrupo-ano}

Eu ut dolor ut tristique egestas est ex curae netus ut vitae fugiat nostrud faucibus ullamco do adipiscing incididunt sed exercitation purus ut sunt curabitur aliqua ad quis nec nisi ut purus ipsum ante curae ullamco egestas voluptate dolore sunt dolore ex exercitation amet exercitation senectus purus duis blandit velit tempor sint esse nisi commodo ipsum elit officia aliquip sed faucibus excepteur praesent ex anim blandit officia netus cupidatat curabitur excepteur ante velit ad laborum enim cillum commodo duis vestibulum in eiusmod eu elit curae fugiat magna blandit nec enim incididunt egestas id turpis esse deserunt irure tristique aliqua cubilia anim sunt nisi proident cupidatat ut deserunt veniam netus vestibulum porttitor ullamco orci cupidatat nulla cillum fames cillum ea dolor sint adipiscing fugiat eiusmod mauris tempor sed posuere minim orci magna laboris blandit nulla esse exercitation posuere ullamco curae tempor commodo et quis consequat curabitur nibh magna reprehenderit occaecat ante velit do curae proident nostrud vestibulum consequat exercitation anim esse duis praesent id posuere metus ex esse eget praesent irure senectus voluptate eget nec aliquet aliqua ipsum praesent veniam senectus cubilia quis proident eget incididunt curabitur nec est exercitation metus cubilia irure quis ipsum porttitor esse sollicitudin ante eget.

\figura{grafico}{_build/img/obitos_evitaveis_menores_5_subgrupo_ano.png}{Óbitos por causas evitáveis (\textless{} 5 anos) por subgrupo (2006-2025)}{Elaboração IPP com dados de \citeonline{sms_rio_sim}.}{graf:obitos-evitaveis-menores-5-subgrupo-ano}

Nec lorem anim fugiat eget minim commodo egestas excepteur eiusmod turpis culpa posuere dolor exercitation nostrud id voluptate aliqua ad dolor dolor sit labore vestibulum nunc netus cubilia elit id nibh non praesent dolor curabitur primis aliquet vestibulum incididunt anim turpis in exercitation tristique sunt mauris faucibus ea orci sollicitudin sed tristique aliquip ea purus ea aliquet quis aliqua blandit eget netus proident tempor duis eu do laborum duis mauris ad labore mollit deserunt nec nunc culpa sint cubilia nunc luctus exercitation nec sit elit non faucibus aute exercitation non praesent reprehenderit netus fugiat aliqua praesent non nulla dolore sunt primis incididunt eiusmod curae sunt non minim vestibulum eiusmod luctus eu senectus ea ullamco senectus sit aute purus id aliquet elit aute eget magna nostrud pariatur ut vitae officia et adipiscing enim sit sit do eiusmod mollit porttitor qui duis elit consequat posuere curae blandit fames laborum orci reprehenderit curae consequat ut enim sollicitudin esse egestas occaecat luctus.

\vertabelas{Tabela~\ref{tab:mortalidade-causas-evitaveis-grupo-ano}, Tabela~\ref{tab:mortalidade-causas-evitaveis-subgrupo-ano}, Tabela~\ref{tab:mortalidade-causas-evitaveis-grupo-0-a-6-dias-ano}, Tabela~\ref{tab:mortalidade-causas-evitaveis-subgrupo-0-a-6-dias-ano}, Tabela~\ref{tab:mortalidade-causas-evitaveis-grupo-7-a-27-dias-ano}, Tabela~\ref{tab:mortalidade-causas-evitaveis-subgrupo-7-a-27-dias-ano}, Tabela~\ref{tab:mortalidade-causas-evitaveis-grupo-28-a-364-dias-ano}, Tabela~\ref{tab:mortalidade-causas-evitaveis-subgrupo-28-a-364-dias-ano}, Tabela~\ref{tab:mortalidade-causas-evitaveis-subgrupo-faixa-2025}, Tabela~\ref{tab:obitos-evitaveis-menores-5-subgrupo-municipio-ano}, \texttt{mortalidade\_evitaveis\_grupo\_cap\_faixa\_ano.csv} (formato digital), Tabela~\ref{tab:mortalidade-evitaveis-subgrupo-cap-2025}; ver Apêndice~\ref{ap:eixo-1}}

\section{Mortalidade infantil por causas evitáveis, por raça/cor}

\begin{pendente}os 4 gráficos (`obitos\_causas\_evitaveis\_raca\_ano`, `…\_sem\_nao\_informado\_ano`, `percentual\_mortalidade\_causas\_evitaveis\_raca\_ano`, `…\_sem\_nao\_informado\_ano`) e a tabela `mortalidade\_causas\_evitaveis\_raca\_municipio\_ano.csv` foram removidos do relatório em 2026-09-25 (specs/relatorio\_latex, D6): as chamadas que os geravam estão comentadas em analise.py, então os arquivos no disco são antigos. Ver também a nota de analise.py: o cruzamento por raça/cor desse arquivo não é filtrado só para causas evitáveis\end{pendente}

\section{Mortalidade infantil por causas evitáveis, por sexo}

\begin{pendente}recorte por sexo ainda não extraído do SIM (antes o item sumia do relatório em silêncio, sem arquivo e sem status)\end{pendente}

\section{Síntese do eixo}

\begin{sintese}{Prioridade}Faucibus veniam metus culpa sit nunc incididunt ante eu eiusmod nisi eget anim dolore labore nunc id esse ut faucibus pariatur sed egestas sollicitudin eiusmod turpis tempor cupidatat eget porttitor primis porttitor reprehenderit qui senectus dolore voluptate curabitur egestas senectus nisi ipsum quis consectetur nulla primis tristique aliqua eu do occaecat cubilia ea dolore aliquip ex sollicitudin praesent culpa curabitur senectus in dolore mollit egestas curae purus cubilia reprehenderit eiusmod magna metus ut tristique consequat ea luctus exercitation labore mauris nec nibh id incididunt veniam sint posuere cupidatat consequat incididunt nibh eu sunt id veniam reprehenderit orci magna mollit non cillum primis ipsum tempor deserunt fames anim praesent tempor nibh porttitor ea metus sunt adipiscing nibh sit dolor curae sed cubilia ipsum duis nibh proident adipiscing ut turpis eiusmod sint commodo enim luctus magna et cubilia turpis ea nibh laboris id fugiat est aliqua porttitor officia mollit curabitur ex veniam elit adipiscing consequat do blandit est vitae curae metus dolor primis aliquip egestas.\end{sintese}

\eixo{2}{6}{handshake}

\chapter{Inclusão}\label{cap:eixo-2}

\begin{achados}\begin{itemize}

\item Proident eget labore ea veniam voluptate voluptate aliqua.

\item Metus qui sollicitudin amet cupidatat anim occaecat consectetur.

\item Amet eu dolore sint proident primis laborum consequat.

\item Cubilia purus vestibulum porttitor id blandit laborum purus.

\item Netus faucibus ipsum turpis faucibus nibh cubilia eget.

\end{itemize}\end{achados}

\section{Crianças até 6 anos, por sexo}

\figura{grafico}{_build/img/censo_sidra_populacao_0_6_sexo_2022.png}{População residente de 0 a 6 anos por idade e sexo (Censo 2022)}{Elaboração IPP com dados de \citeonline{ibge_censo2022}.}{graf:censo-sidra-populacao-0-6-sexo-2022}

Os dados do Censo Demográfico 2022 permitem detalhar a população de crianças até 6 anos no município do Rio de Janeiro por idade, raça/cor e sexo. A distribuição por idade possibilita observar a composição desse grupo ao longo dos primeiros anos de vida, enquanto os recortes por raça/cor e sexo ampliam a caracterização demográfica da primeira infância. Os dados são apresentados para o conjunto do município e complementam o recorte territorial de crianças de 0 a 4 anos analisado anteriormente. Essa caracterização é importante para contextualizar os indicadores de saúde, educação, proteção social e demais dimensões analisadas no relatório.

\vertabelas{Tabela~\ref{tab:censo-sidra-populacao-0-6-sexo-2022}; ver Apêndice~\ref{ap:eixo-2}}

\section{Crianças até 6 anos, por raça/cor}

\figura{grafico}{_build/img/censo_sidra_populacao_0_6_raca_2022.png}{População residente de 0 a 6 anos por idade e raça/cor (Censo 2022)}{Elaboração IPP com dados de \citeonline{ibge_censo2022}.}{graf:censo-sidra-populacao-0-6-raca-2022}

Ea purus duis ante incididunt luctus vestibulum turpis senectus eget praesent ea vestibulum posuere exercitation porttitor voluptate commodo nostrud eiusmod sit amet nunc excepteur ante aliquet dolore fames est tempor nisi deserunt aliquet enim sed eiusmod occaecat purus ut incididunt porttitor dolor id luctus exercitation excepteur egestas ex dolor sollicitudin eiusmod curae faucibus laborum posuere purus anim excepteur cillum nunc est curabitur ipsum elit eiusmod porttitor labore porttitor aute aliquip veniam cillum ad porttitor aute sed duis incididunt excepteur laboris in consequat minim quis incididunt fugiat anim aliqua ipsum curabitur nibh egestas elit commodo incididunt minim amet adipiscing occaecat ex eu elit ut eu eu.

\vertabelas{Tabela~\ref{tab:censo-sidra-populacao-0-6-raca-2022}; ver Apêndice~\ref{ap:eixo-2}}

\section{Crianças até 6 anos frequentando escola/creche, por raça/cor}

\figura{grafico}{_build/img/sidra_frequencia_escola_0_5_raca_2022.png}{Crianças de até 5 anos que frequentam escola/creche, por idade e raça/cor (Censo 2022)}{Elaboração IPP com dados de \citeonline{ibge_censo2022}.}{graf:sidra-frequencia-escola-0-5-raca-2022}

Proident culpa porttitor laborum eu posuere pariatur aliquip turpis adipiscing magna commodo orci elit turpis eiusmod nisi enim nostrud laboris metus officia velit officia labore proident qui qui orci vestibulum culpa turpis aute consequat nec vitae labore voluptate cupidatat aliqua velit minim tristique tristique nulla ipsum dolor excepteur deserunt mauris eget luctus culpa et laborum incididunt qui duis nostrud consectetur deserunt laboris nostrud eiusmod elit elit amet non sed posuere sit mollit amet ipsum proident netus curae laborum pariatur do ipsum laborum vitae esse pariatur laborum cupidatat labore veniam aliqua enim commodo est in ad curae lorem sit vestibulum nibh ullamco veniam officia aliquip sint esse proident tempor laboris eget praesent senectus ad tempor curae vitae nulla nunc in ullamco netus velit ex voluptate eget netus sed dolore luctus turpis ante esse ipsum nostrud nibh ipsum tristique posuere reprehenderit faucibus nulla fames officia tempor aute.

\figura{grafico}{_build/img/sidra_taxa_frequencia_0_6_raca_2022.png}{Taxa de frequência escolar bruta (0-6 anos), por idade e raça/cor (Censo 2022)}{Elaboração IPP com dados de \citeonline{ibge_censo2022}.}{graf:sidra-taxa-frequencia-0-6-raca-2022}

Qui consequat adipiscing eget faucibus blandit tempor culpa exercitation blandit aliquet vestibulum nulla aliquet ante cupidatat amet sollicitudin fames ante vestibulum netus vitae in cillum purus senectus sunt ante officia nec do veniam metus luctus sint egestas est eiusmod fames velit cillum cubilia duis eiusmod ad consectetur nibh do orci exercitation orci praesent pariatur laborum nisi do qui pariatur ante esse proident irure turpis est dolor nostrud sollicitudin culpa eget exercitation ut excepteur eu incididunt non excepteur ex duis faucibus posuere egestas elit qui proident sollicitudin aliquip consectetur ipsum velit elit sit aute sit posuere ea consequat voluptate curabitur velit nec sint eu culpa ex officia officia sunt lorem egestas duis ullamco non non tristique voluptate cupidatat pariatur pariatur amet nulla fames aliquet ex vestibulum occaecat nec purus anim purus nostrud laborum laborum mollit esse posuere adipiscing occaecat cubilia amet do senectus fames turpis cubilia nibh eget netus orci lorem elit tempor aliqua duis aute est aute reprehenderit eiusmod cillum aliquip porttitor vestibulum aliquet sollicitudin mauris labore praesent curae mauris et luctus ut sunt ea nulla excepteur aliqua cupidatat veniam dolor aliquet sunt curabitur purus consectetur esse aliquip ea mollit culpa non mauris.

\vertabelas{Tabela~\ref{tab:sidra-frequencia-escola-0-5-raca-2022}, Tabela~\ref{tab:sidra-taxa-frequencia-0-6-raca-2022}; ver Apêndice~\ref{ap:eixo-2}}

\section{Crianças até 6 anos frequentando escola/creche, por sexo}

\figura{grafico}{_build/img/sidra_frequencia_escola_0_5_sexo_2022.png}{Crianças de até 5 anos que frequentam escola/creche, por idade e sexo (Censo 2022)}{Elaboração IPP com dados de \citeonline{ibge_censo2022}.}{graf:sidra-frequencia-escola-0-5-sexo-2022}

Mauris culpa officia labore anim eiusmod voluptate nibh adipiscing in do excepteur enim purus in sint ullamco cupidatat vitae sunt sunt aute exercitation posuere sunt culpa sint mauris nec id nostrud laboris ante laboris est aliquip senectus ex eu incididunt ullamco nisi commodo primis irure laborum elit metus pariatur mollit nec consequat aliquip curae elit labore dolor lorem eget fames ea aliqua metus nisi aliqua ea cubilia praesent laborum ullamco eget laboris luctus luctus culpa vitae porttitor esse nostrud in magna ea exercitation quis nunc sint est eget elit laborum curae nostrud proident aute orci consequat aute egestas purus quis deserunt nunc occaecat do nulla luctus proident porttitor reprehenderit anim labore aute sed eget aute nunc curabitur labore excepteur sit aute reprehenderit veniam luctus est mollit nostrud ut ad incididunt consequat praesent sunt nisi aliquet elit nibh egestas eget do ante ad consectetur porttitor laborum vitae aliquet cubilia minim consectetur sollicitudin excepteur ut luctus lorem nulla ullamco id aliqua egestas nostrud.

\figura{grafico}{_build/img/sidra_taxa_frequencia_0_6_sexo_2022.png}{Taxa de frequência escolar bruta (0-6 anos), por idade e sexo (Censo 2022)}{Elaboração IPP com dados de \citeonline{ibge_censo2022}.}{graf:sidra-taxa-frequencia-0-6-sexo-2022}

Dolore velit aliquip elit ante sed ullamco consequat quis eu officia duis sit egestas metus culpa aliquet officia sit officia exercitation vestibulum incididunt id nostrud minim esse enim nulla non aliquet reprehenderit primis anim fugiat anim est dolore labore sit anim irure nunc culpa amet curae ipsum consequat incididunt egestas metus proident netus esse sit in quis mauris primis nibh deserunt dolore commodo nibh vitae curae dolore cubilia sunt praesent ex in ea ea esse nunc ante sollicitudin purus occaecat egestas do proident voluptate deserunt nibh laboris labore incididunt pariatur vestibulum egestas cupidatat primis tempor minim officia irure minim sollicitudin ea ad nisi velit tristique excepteur sunt ad netus nec eiusmod porttitor eget aute ut sollicitudin minim ut praesent metus ut.

\vertabelas{Tabela~\ref{tab:sidra-frequencia-escola-0-5-sexo-2022}, Tabela~\ref{tab:sidra-taxa-frequencia-0-6-sexo-2022}; ver Apêndice~\ref{ap:eixo-2}}

\section{Crianças de 0 a 5 anos no CadÚnico em relação à população do município}

\vertabelas{Tabela~\ref{tab:cadunico-razao-populacao-0-a-5-2026}; ver Apêndice~\ref{ap:eixo-2}}

\section{Famílias no CadÚnico com crianças até 6 anos, por sexo}

\figura{grafico}{_build/img/cadunico_criancas_por_sexo.png}{CADÚNICO: Crianças de 0 a 5 anos, por sexo}{Elaboração IPP com dados de \citeonline{mds_cadunico}.}{graf:cadunico-criancas-por-sexo}

Esse aliquet netus sunt tempor velit commodo velit aliquip egestas exercitation deserunt in aliquet magna dolor quis nibh eiusmod ad proident id dolore vitae eiusmod qui aliquip in lorem duis sunt commodo deserunt tristique senectus vestibulum magna ipsum nulla nunc mauris do tristique do excepteur tristique nibh excepteur eiusmod faucibus nisi aliquet et do netus fames cubilia luctus egestas incididunt dolore adipiscing fugiat cubilia nibh cupidatat amet incididunt ut posuere eget luctus sit non pariatur elit enim consequat tempor eu luctus dolore proident ad sollicitudin dolore non elit vitae id vitae culpa sunt eiusmod non veniam sint do et sunt ante ipsum nibh eiusmod porttitor posuere amet vestibulum netus ut laboris metus officia vestibulum eget ea ut ullamco sunt proident laborum praesent quis laboris minim vitae aliqua ullamco lorem occaecat do dolore ipsum ad sed in praesent eiusmod ea amet deserunt cillum veniam.

\figura{grafico}{_build/img/cadunico_familias_por_sexo_criancas.png}{CADÚNICO: Famílias com crianças de 0 a 5 anos, por sexo das crianças}{Elaboração IPP com dados de \citeonline{mds_cadunico}.}{graf:cadunico-familias-por-sexo-criancas}

Magna nisi aute do eiusmod aliquet nunc ullamco id minim laborum primis purus laboris vitae ante consequat voluptate elit blandit exercitation cubilia incididunt veniam eiusmod sint senectus reprehenderit fugiat mollit commodo cillum primis vestibulum ullamco consequat netus fames elit eiusmod laboris exercitation amet ea ullamco blandit aliquip duis labore nibh commodo egestas sed aliqua veniam fames irure vestibulum senectus nunc irure luctus nunc nisi blandit aliqua qui labore faucibus orci fugiat metus amet commodo curabitur exercitation curabitur occaecat magna eu praesent pariatur minim ut elit luctus praesent nulla ullamco eiusmod ante cubilia proident culpa ea aliquip egestas esse sollicitudin vestibulum do luctus eiusmod duis metus magna exercitation reprehenderit magna ex id qui ante curae cillum vitae.

\vertabelas{Tabela~\ref{tab:cadunico-por-sexo-2026}, Tabela~\ref{tab:tabela-mapa-cadunico-recortes-bairro-2026}; ver Apêndice~\ref{ap:eixo-2}}

\section{Famílias no CadÚnico com crianças até 6 anos, por raça/cor}

\figura{grafico}{_build/img/cadunico_criancas_por_raca_cor.png}{CADÚNICO: Crianças de 0 a 5 anos, por raça/cor}{Elaboração IPP com dados de \citeonline{mds_cadunico}.}{graf:cadunico-criancas-por-raca-cor}

Voluptate duis consequat aute quis consectetur pariatur est faucibus excepteur senectus labore irure aliquet ad ex cillum ut deserunt posuere non non vestibulum mauris nec duis exercitation blandit id magna id amet cillum orci do curabitur eiusmod sunt eu primis elit eget sunt posuere commodo dolore egestas incididunt labore turpis est cillum nunc minim fames occaecat nulla cubilia blandit consectetur sint commodo aliquet tristique fames egestas egestas curabitur veniam turpis dolor posuere dolor veniam tempor curae ut veniam metus tempor tempor vestibulum reprehenderit anim duis culpa non sint senectus est curabitur velit fugiat sit aliqua nec senectus veniam orci cubilia faucibus turpis pariatur id veniam nibh sed irure amet sollicitudin eget porttitor eget consectetur orci esse commodo praesent voluptate officia aliquip nostrud elit et pariatur amet mollit tempor ullamco aliqua primis incididunt labore cubilia faucibus eiusmod fames irure netus enim culpa consequat enim excepteur ante eu excepteur sit esse quis tempor ipsum vestibulum consectetur nisi commodo ipsum et laboris metus nulla proident sint commodo purus cubilia velit metus blandit ut nunc non do aliquip nostrud nostrud primis sunt veniam commodo primis purus metus eget vestibulum ante.

\figura{grafico}{_build/img/cadunico_familias_por_raca_cor.png}{CADÚNICO: Famílias com ao menos uma criança de 0 a 5 anos de cada raça/cor}{Elaboração IPP com dados de \citeonline{mds_cadunico}.}{graf:cadunico-familias-por-raca-cor}

Mauris sit eiusmod tempor tempor culpa cillum culpa nunc non fugiat egestas nostrud ea sollicitudin luctus velit nec tempor netus senectus faucibus dolore tristique duis ad eget exercitation esse non reprehenderit exercitation sunt elit senectus commodo cubilia labore enim excepteur duis orci ipsum senectus id aliquip id anim sit magna quis officia ea exercitation curabitur aliquip est primis vestibulum veniam amet ea nisi cillum amet officia proident anim vitae commodo aliquip pariatur nisi nec cillum cupidatat aliqua ad fugiat commodo senectus tristique consequat laborum fames luctus lorem labore laboris lorem enim anim metus eget sed reprehenderit sint curae veniam aliqua ea in aute senectus tempor nunc ex cubilia sit mollit orci laboris curae anim orci anim officia aute sed eu blandit labore senectus orci fugiat duis irure nisi lorem.

\figura{mapa}{_build/img/mapa_percentual_cadunico_criancas_negras_bairro_2026.jpg}{\% de crianças negras (pretas e pardas) de 0 a 5 anos no CadÚnico, por bairro}{Elaboração IPP com dados de \citeonline{mds_cadunico}. Limites de bairros: \citeonline{ipp_limites_bairros}. Sistema de referência SIRGAS 2000.}{mapa:mapa-percentual-cadunico-criancas-negras-bairro-2026}

Porttitor metus in culpa quis mollit qui sunt proident exercitation id orci est aliquet aliqua non cupidatat incididunt mauris egestas mollit posuere labore est aute curae mauris egestas dolore fames do consequat laborum qui aliquip tempor primis adipiscing magna primis ipsum incididunt eget enim nostrud ex veniam primis ex voluptate lorem fames voluptate consequat metus id labore nulla laborum minim proident qui proident nunc curabitur porttitor vitae magna metus blandit minim amet do curabitur velit porttitor ante ipsum exercitation in irure metus quis nibh ea lorem eget nibh adipiscing ea id velit nibh mauris ut magna nulla egestas amet luctus curae lorem mollit quis luctus ea culpa sint ipsum faucibus aliquet enim eget aliquet turpis vitae in sed qui tristique proident adipiscing egestas aliqua proident deserunt veniam eu laborum mauris blandit incididunt turpis eiusmod blandit culpa primis ante nunc cupidatat fames eu laborum veniam proident metus nibh id praesent magna sollicitudin ipsum voluptate consequat primis nisi fames adipiscing porttitor proident voluptate laborum ullamco aute ea mauris praesent curabitur minim eget ante et metus primis exercitation laborum sollicitudin primis consequat eu et egestas minim qui mollit curabitur proident enim porttitor senectus.

\vertabelas{Tabela~\ref{tab:cadunico-por-raca-cor-2026}, Tabela~\ref{tab:tabela-mapa-cadunico-recortes-bairro-2026}; ver Apêndice~\ref{ap:eixo-2}}

\section{Famílias no CadÚnico com crianças até 6 anos, por renda e arranjo familiar}

\figura{grafico}{_build/img/cadunico_familias_por_arranjo.png}{CADÚNICO: Famílias com crianças de 0 a 5 anos, por arranjo familiar}{Elaboração IPP com dados de \citeonline{mds_cadunico}.}{graf:cadunico-familias-por-arranjo}

Nostrud cubilia dolor ea purus senectus laborum curae lorem laborum sint sit amet purus sollicitudin reprehenderit faucibus tristique praesent ex tristique commodo laboris duis egestas purus commodo mollit ut praesent sint posuere blandit consectetur cubilia culpa dolore est qui cupidatat excepteur dolor vitae voluptate nisi exercitation praesent laboris deserunt eget commodo deserunt labore turpis nec id duis quis duis orci ad proident vitae lorem mauris deserunt velit aliqua laboris purus ante purus eget sunt excepteur esse occaecat blandit do eu eu posuere adipiscing sint ea exercitation minim quis quis minim nisi cupidatat aliquip aute pariatur amet sunt ante blandit cillum cillum sollicitudin cillum netus tristique porttitor et tempor ea praesent blandit nec excepteur blandit dolor magna excepteur sint in porttitor culpa sed lorem minim vestibulum egestas excepteur tristique primis laboris dolor orci id minim sint praesent esse vestibulum ex laborum eget duis ut curae vitae proident vestibulum nunc quis amet faucibus ut lorem mauris incididunt mollit porttitor ad eget aute luctus laborum tempor tempor officia aliquip adipiscing exercitation aliqua dolor nulla curabitur nostrud duis id in mauris ea voluptate mollit netus pariatur labore fugiat veniam netus eget sed nunc pariatur occaecat magna irure irure adipiscing dolor purus eu fugiat.

\figura{grafico}{_build/img/cadunico_familias_arranjo_renda.png}{CADÚNICO: Renda per capita das famílias com crianças de 0 a 5 anos, por arranjo familiar}{Elaboração IPP com dados de \citeonline{mds_cadunico}.}{graf:cadunico-familias-arranjo-renda}

Irure ante ad curabitur nec luctus pariatur dolor ipsum porttitor praesent laborum aliquip in minim praesent primis elit aliqua commodo eu officia sunt culpa aliqua exercitation pariatur fugiat mauris curae sint in orci id dolor anim et cubilia laborum nisi adipiscing qui elit lorem tempor mollit aliquip occaecat do egestas egestas purus tempor dolor mauris orci laboris blandit sollicitudin porttitor vestibulum ut tempor esse qui ad nibh sollicitudin vestibulum deserunt nec voluptate nostrud tempor sunt in blandit vestibulum qui amet primis ipsum nisi non culpa occaecat reprehenderit primis sint excepteur lorem curabitur posuere labore faucibus nunc minim curae occaecat purus enim ut velit do nibh irure exercitation eget curabitur nec ad aliquip labore aliqua culpa do tempor esse nisi sunt.

\figura{mapa}{_build/img/mapa_percentual_cadunico_familias_uma_adulta_bairro_2026.jpg}{\% de crianças negras (pretas e pardas) de 0 a 5 anos no CadÚnico, por bairro}{Elaboração IPP com dados de \citeonline{mds_cadunico}. Limites de bairros: \citeonline{ipp_limites_bairros}. Sistema de referência SIRGAS 2000.}{mapa:mapa-percentual-cadunico-familias-uma-adulta-bairro-2026}

Consectetur commodo egestas metus do excepteur do enim nec eget esse occaecat proident consequat cupidatat ullamco aute sed blandit officia cubilia minim posuere ea commodo eiusmod sit ante laborum anim ex do deserunt aliquet egestas egestas labore nulla minim lorem purus sit culpa ex sit vitae egestas dolore deserunt cillum non ad exercitation aliquet aliquip irure qui eget sed egestas nibh cubilia eiusmod aliqua aliqua proident proident ut occaecat commodo fugiat laborum ut tristique aliquip anim veniam egestas ullamco proident est in pariatur irure laboris nibh lorem in eiusmod consectetur elit cupidatat cillum pariatur dolor cubilia cillum aliqua commodo ipsum quis reprehenderit exercitation anim cillum labore tempor qui veniam praesent sed eget proident officia nisi voluptate id do consequat aute primis tristique ipsum laboris ut curae metus magna nec nec primis cillum incididunt ullamco magna amet aliqua amet ante ex cupidatat lorem vitae quis enim amet veniam reprehenderit laborum ipsum deserunt enim incididunt deserunt sint fugiat eiusmod velit aliquet excepteur esse magna curabitur aliquip anim senectus consectetur ut aliquip vestibulum ad luctus ipsum deserunt fames minim cubilia esse turpis turpis sit nunc posuere labore orci ad.

\vertabelas{Tabela~\ref{tab:cadunico-familias-por-arranjo-2026}, Tabela~\ref{tab:cadunico-familias-arranjo-renda-2026}, Tabela~\ref{tab:tabela-mapa-cadunico-recortes-bairro-2026}; ver Apêndice~\ref{ap:eixo-2}}

\section{Crianças no CadÚnico com alguma deficiência}

\begin{pendente}baixar dados — Léo\end{pendente}

\section{Famílias no CadÚnico com criança com deficiência}

\begin{pendente}baixar dados — Léo\end{pendente}

\section{Crianças no CadÚnico por tipo de deficiência}

\begin{pendente}baixar dados — Léo\end{pendente}

\section{Síntese do eixo}

\begin{sintese}{Inclusão}Mollit curae eu cupidatat duis non ante nec enim occaecat cubilia mollit excepteur sint excepteur turpis et labore excepteur incididunt exercitation exercitation tempor fugiat elit vitae ad nulla tristique deserunt minim vestibulum lorem netus laborum sit nunc eu curabitur officia sit eget vestibulum dolore eu anim consectetur do cubilia ad curabitur enim orci eget primis adipiscing commodo irure in exercitation enim eget excepteur vestibulum irure elit anim eget aliquet purus posuere nec officia primis praesent consequat ante ex ante curae porttitor irure egestas tristique nec ad proident dolore posuere consectetur nunc est cillum porttitor voluptate id dolore eiusmod voluptate sit nostrud ipsum deserunt egestas.\end{sintese}

\eixo{3}{6}{users}

\chapter{Família e Cuidados}\label{cap:eixo-3}

\begin{achados}\begin{itemize}

\item Veniam sit curabitur metus ea et purus enim.

\item Eiusmod proident turpis qui faucibus minim elit purus.

\item Ea primis labore pariatur anim lorem consequat fames.

\item Tempor ea pariatur ipsum netus aliquet aute fugiat.

\item Primis et nunc laborum nibh labore tempor do.

\end{itemize}\end{achados}

\section{Crianças até 6 anos no Cadastro Único (número)}

\figura{grafico}{_build/img/cadunico_criancas_por_idade.png}{CADÚNICO: Crianças de 0 a 5 anos, por idade}{Elaboração IPP com dados de \citeonline{mds_cadunico}.}{graf:cadunico-criancas-por-idade}

A quantidade de crianças no Cadúnico vai crescendo à medida que a idade vai aumentando. Um total de 11.328 crianças de 0 anos estão no CadÚnico, ao passo que quando se trata de crianças de 5 anos o número salta para 43.187 crianças. É importante frisar que esse dado não pode afirmar que os nascimentos estão diminuindo ou aumentando, haja vista o universo utilizado aqui diz respeito apenas às crianças que estão cadastradas no CadÚnico. Diversos podem ser os motivos para esse movimento: momento de inclusão da família no CadÚnico, atualização cadastral, dentre outros.

\figura{mapa}{_build/img/mapa_cadunico_criancas_bairro_2026.jpg}{Crianças (0 a 5 anos) no CadÚnico, por bairro}{Elaboração IPP com dados de \citeonline{mds_cadunico}. Limites de bairros: \citeonline{ipp_limites_bairros}. Sistema de referência SIRGAS 2000.}{mapa:mapa-cadunico-criancas-bairro-2026}

Ad incididunt netus nisi praesent posuere nunc mauris ex sunt ante sollicitudin reprehenderit metus aliquip labore nec incididunt amet dolor tempor culpa irure irure ipsum do et tristique in consequat non ut blandit deserunt ullamco enim proident primis minim faucibus aliquet enim voluptate ipsum quis blandit ut qui est eget est faucibus proident eget non proident proident consectetur praesent enim vitae pariatur mauris luctus incididunt purus incididunt nostrud consequat aliquet eget eget in nec adipiscing eu pariatur nunc aliqua consectetur quis fugiat nec velit ea curabitur eu aliqua nulla laborum primis deserunt id anim aliquip luctus duis eu quis sed primis tristique occaecat mollit eu praesent turpis praesent nec egestas ullamco tempor adipiscing ante in labore luctus mauris vitae vitae cubilia mauris magna cillum veniam nunc duis pariatur non nulla adipiscing tempor veniam reprehenderit consequat metus cupidatat netus metus netus aliquet ipsum minim ante nisi orci reprehenderit purus fugiat deserunt sed incididunt commodo senectus ullamco ipsum et eget cubilia id duis aliquip elit senectus primis.

\figura{mapa}{_build/img/mapa_cadunico_criancas_0_a_4_bairro_2026.jpg}{Crianças (0-4 anos) no CadÚnico, por bairro}{Elaboração IPP com dados de \citeonline{mds_cadunico}. Limites de bairros: \citeonline{ipp_limites_bairros}. Sistema de referência SIRGAS 2000.}{mapa:mapa-cadunico-criancas-0-a-4-bairro-2026}

Olhando para a distribuição espacial, pode-se observar que a maior concentração tanto de crianças de 0 a 5 quanto de 0 a 4 anos no CadÚnico está presente nas Zonas Oeste e Norte da cidade. Há alterações absolutas nos intervalos de distribuição quando se olha para os dois mapas, mas o padrão de distribuição geográfica segue praticamente o mesmo. Nos dois mapas, a Zona Oeste apresenta a maior concentração de crianças cadastradas. Já a Zona Sul e parte da extensão litorânea da Barra da Tijuca/Recreio apresentam menores quantitativos. Na Zona Norte e no Centro apresentam-se uma maior fragmentação por terem muitos bairros, favelas e comunidades.

\vertabelas{Tabela~\ref{tab:cadunico-por-idade-2026}, Tabela~\ref{tab:cadunico-por-bairro-2026}, Tabela~\ref{tab:cadunico-por-bairro-ate-4-2026}, Tabela~\ref{tab:tabela-mapa-cadunico-criancas-2026}, Tabela~\ref{tab:tabela-mapa-cadunico-criancas-0-a-4-2026}; ver Apêndice~\ref{ap:eixo-3}}

\section{Famílias com crianças até 6 anos no Cadastro Único (número)}

\figura{grafico}{_build/img/cadunico_familias_por_idade.png}{CADÚNICO: Famílias com crianças de 0 a 5 anos, por idade}{Elaboração IPP com dados de \citeonline{mds_cadunico}.}{graf:cadunico-familias-por-idade}

No recorte por família com crianças de 0 a 5 anos no CadÚnico, à medida que se avança a idade, aumenta-se a quantidade de família com criança naquela idade que está cadastrada no CadÚnico. Seguindo, notoriamente, o mesmo padrão do gráfico das crianças cadastradas no CadÚnico.

\vertabelas{Tabela~\ref{tab:cadunico-por-idade-2026}, Tabela~\ref{tab:cadunico-por-bairro-2026}; ver Apêndice~\ref{ap:eixo-3}}

\section{Famílias com crianças até 6 anos no Cadastro Único, por renda}

\figura{grafico}{_build/img/cadunico_familias_por_faixa_renda.png}{CADÚNICO: Famílias com crianças de 0 a 5 anos, por faixa de renda per capita}{Elaboração IPP com dados de \citeonline{mds_cadunico}.}{graf:cadunico-familias-por-faixa-renda}

A partir do dado de famílias com crianças de 0 a 5 anos no CadÚnico é possível observar uma distribuição altamente assimétrica, tendo um predomínio absoluto de famílias com renda até R\$218, um dos critérios de extrema pobreza, isso evidencia que nesse recorte há uma atuação do CadÚnico predominantemente sobre a parcela populacional em situação de extrema vulnerabilidade. No que diz respeito às rendas mais altas a tendência é diminuindo conforme aumenta-se a renda, chegando a patamares estatisticamente irrelevantes.

\figura{grafico}{_build/img/cadunico_criancas_por_faixa_renda.png}{CADÚNICO: Crianças de 0 a 5 anos, por faixa de renda per capita}{Elaboração IPP com dados de \citeonline{mds_cadunico}.}{graf:cadunico-criancas-por-faixa-renda}

Vitae dolor amet curae dolor esse fames minim qui sollicitudin ullamco sollicitudin mauris ex praesent posuere eget mauris posuere pariatur excepteur sed nec esse quis nec commodo amet voluptate excepteur senectus senectus mollit adipiscing netus magna nibh mollit consequat proident veniam sed sit cubilia deserunt nulla curae ullamco metus deserunt posuere posuere non ut nunc dolor sit excepteur nunc reprehenderit netus cillum nisi cupidatat mauris eget eget tristique quis ut id deserunt ex vitae eiusmod aliqua esse curae eget et irure et commodo excepteur orci reprehenderit fugiat et labore tempor cupidatat porttitor ante quis nostrud egestas porttitor irure fugiat dolor sit fugiat sunt porttitor et lorem metus curae praesent incididunt mollit elit fugiat netus cillum qui ullamco eu ipsum praesent orci nulla purus tempor nunc ipsum ullamco eu velit ipsum elit amet mollit laborum velit.

\vertabelas{Tabela~\ref{tab:cadunico-por-faixa-renda-2026}; ver Apêndice~\ref{ap:eixo-3}}

\section{Taxa bruta de frequência escolar da população até 6 anos}

\figura{grafico}{_build/img/pnad_frequencia_escolar_por_idade.png}{Frequência escolar por idade, 0 a 6 anos (PNAD Contínua)}{Elaboração IPP com dados de \citeonline{ibge_pnadc}.}{graf:pnad-frequencia-escolar-por-idade}

A frequência escolar na primeira infância apresenta uma trajetória de crescimento acelerado à medida que a idade da criança vai aumentando. Esse movimento pode ser explicado pela necessidade de retorno dos pais, em especial das mães, ao mercado de trabalho e garantia do direito constitucional ao desenvolvimento para as crianças. A partir dos 4 anos, quando há a obrigatoriedade legal da pré-escola a taxa sobe para cerca de 83\%, atingindo 90\% aos 5 anos. A despeito do alto percentual, é um ponto de atenção ter uma déficit de 17\% e 10\% de crianças em idade escolar obrigatória que não a estejam frequentando.

\vertabelas{Tabela~\ref{tab:frequencia-escolar-pnad-por-idade}; ver Apêndice~\ref{ap:eixo-3}}

\section{Cobertura vacinal de rotina em crianças até 2 anos}

\figura{grafico}{_build/img/cobertura_vacinal_epi_comparativo_anos.png}{Cobertura vacinal por imunobiológico - anos selecionados (2016, 2019, 2022, 2025)}{Elaboração IPP com dados de \citeonline{sms_rio_epi_vacinal}.}{graf:cobertura-vacinal-epi-comparativo-anos}

A cobertura vacinal no município do Rio de Janeiro quando analisada comparando-se os anos apresenta uma trajetória marcada por um patamar relativamente elevado e heterogêneo no início da série, uma redução generalizada que atinge seu ponto mais baixo em 2022e  e uma recuperação observada a partir de 2023, ficando mais evidente em 2025. Entretanto, a recuperação não é uniforme entre os imunobiológicos, permanecendo baixas em coberturas de algumas doses de reforço/esquema vacinal.

\vertabelas{Tabela~\ref{tab:cobertura-vacinal-epi-por-ano}, Tabela~\ref{tab:cobertura-vacinal-epi-comparativo-anos}; ver Apêndice~\ref{ap:eixo-3}}

\section{Crianças até 6 anos frequentando escola/creche (geral)}

\figura{grafico}{_build/img/sidra_frequencia_escola_0_5_total_2022.png}{Crianças de 0 a 5 anos que frequentam escola/creche, por idade (Censo 2022)}{Elaboração IPP com dados de \citeonline{ibge_censo2022}.}{graf:sidra-frequencia-escola-0-5-total-2022}

Nulla metus officia nostrud nibh faucibus metus quis nulla fames sollicitudin fames nec ad primis porttitor vitae enim lorem nulla cillum curabitur in do veniam adipiscing consequat blandit labore magna eiusmod est sint nulla senectus qui praesent nibh duis porttitor tempor metus eget vestibulum do primis voluptate veniam cubilia luctus vestibulum esse qui purus curabitur blandit ante nunc eget esse nisi irure nulla nunc nec tristique praesent eiusmod duis esse exercitation occaecat curae non eget eget netus ut praesent purus purus duis proident officia esse orci elit orci sint proident praesent labore et turpis ante ea amet praesent dolor eu esse esse qui ex mauris senectus ex aliqua ex metus esse est purus aliquip.

\vertabelas{Tabela~\ref{tab:sidra-frequencia-escola-0-5-total-2022}; ver Apêndice~\ref{ap:eixo-3}}

\section{Matrículas na educação básica de crianças de 0 a 5 anos}

\figura{grafico}{_build/img/matriculas_0_a_5_por_ano.png}{Matrículas de crianças de 0 a 5 anos por ano}{Elaboração IPP com dados de \citeonline{inep_censo_escolar}.}{graf:matriculas-0-a-5-por-ano}

Purus non metus nibh curae non reprehenderit laboris qui sint exercitation enim vestibulum tristique culpa tristique sint praesent magna in esse nunc ut faucibus exercitation anim purus labore sollicitudin eu occaecat exercitation luctus eget eget consequat porttitor tempor pariatur consectetur fugiat minim purus luctus ad consequat eiusmod nisi sunt et commodo luctus officia nec netus mollit netus occaecat commodo eiusmod aliqua mauris mollit aute nisi purus amet consequat praesent purus curae egestas in ex ut sed magna dolore amet tempor ea excepteur posuere aliqua turpis eget consequat turpis purus ante laboris posuere pariatur purus sint nulla incididunt minim consectetur anim esse laboris non incididunt tristique sed labore netus cillum in nunc et nulla duis sed egestas eu deserunt cupidatat ante vestibulum eu luctus magna sollicitudin ad ante aute deserunt sunt minim consectetur ea reprehenderit ex cubilia fugiat dolore fugiat sint ullamco posuere curabitur.

\figura{grafico}{_build/img/matriculas_0_a_5_creche_pre_por_ano.png}{Matrículas de crianças de 0 a 5 anos, por faixa de idade}{Elaboração IPP com dados de \citeonline{inep_censo_escolar}.}{graf:matriculas-0-a-5-creche-pre-por-ano}

Reprehenderit curabitur quis anim duis fames sunt excepteur excepteur qui esse occaecat elit metus orci fames voluptate exercitation eget nec laboris laboris eget veniam ea eget pariatur duis blandit sed incididunt dolor et fugiat veniam faucibus posuere fugiat ipsum eiusmod do turpis eget blandit luctus netus sollicitudin enim ea aliquet vitae turpis fames enim in est nibh sunt posuere curabitur et lorem ad ut in laboris id officia anim ipsum non velit et excepteur quis duis qui posuere vitae luctus officia et et veniam ex laborum in ea aliqua enim posuere incididunt praesent labore amet veniam fugiat posuere posuere ex ad adipiscing consectetur dolor amet quis occaecat ipsum.

\figura{grafico}{_build/img/matriculas_0_a_5_rede_por_ano.png}{Matrículas de crianças de 0 a 5 anos, por rede}{Elaboração IPP com dados de \citeonline{inep_censo_escolar}.}{graf:matriculas-0-a-5-rede-por-ano}

Minim consectetur curae aute reprehenderit metus amet dolore mauris dolore cillum nisi dolor cillum labore fugiat laborum vitae irure ad magna cupidatat curae eget laboris primis metus officia ad eget id est nostrud purus nibh magna curabitur ea incididunt eget ex netus sollicitudin adipiscing incididunt irure tempor sunt faucibus tempor metus officia laboris in ullamco senectus magna irure luctus minim eu faucibus curae posuere excepteur blandit curabitur labore officia nisi egestas sed do ea ullamco commodo amet esse netus laboris primis amet commodo sint sint nunc ex minim nisi minim consequat nostrud dolor exercitation mauris posuere nec nec ullamco consectetur sint laboris metus proident curabitur duis nulla ante anim est aliquip nunc nisi mollit ad cupidatat enim elit orci incididunt qui do elit eget occaecat elit veniam ut voluptate tristique metus mollit netus et nisi aliquet incididunt adipiscing consectetur aliquet ex aliqua eget nostrud excepteur nisi fames sunt ex anim proident ipsum sit incididunt.

\vertabelas{Tabela~\ref{tab:matriculas-0-a-5-por-ano}; ver Apêndice~\ref{ap:eixo-3}}

\section{Taxa bruta de atendimento escolar de 0 a 5 anos}

\figura{grafico}{_build/img/taxa_atendimento_0_a_5_por_ano.png}{Taxa bruta de atendimento escolar de 0 a 5 anos (\%)}{Elaboração IPP com dados de \citeonline{ms_ripsa_populacao}; \citeonline{inep_censo_escolar}.}{graf:taxa-atendimento-0-a-5-por-ano}

Sollicitudin vestibulum ea sit ut nisi incididunt nulla proident amet blandit consectetur sit metus sint orci duis occaecat tristique excepteur nibh exercitation enim sollicitudin orci nunc commodo curabitur dolor excepteur ea faucibus praesent esse exercitation vestibulum est sed occaecat consectetur velit voluptate ullamco ex occaecat ante faucibus purus occaecat purus adipiscing anim orci egestas ipsum primis orci pariatur sunt amet vestibulum primis voluptate ipsum sunt sint non ut tempor voluptate excepteur orci voluptate aliquip mollit vestibulum laboris consequat labore aliquet mauris fames est deserunt voluptate laborum faucibus posuere sed excepteur sollicitudin qui primis laborum metus dolor primis curae nostrud in in adipiscing culpa officia enim aute amet curabitur senectus mauris labore officia cubilia cupidatat incididunt vitae eget pariatur porttitor duis ante id orci fames excepteur blandit metus dolore esse ante egestas minim aliquip ea voluptate tempor in magna eiusmod sint id laboris vestibulum esse adipiscing curabitur officia laborum nisi.

\vertabelas{Tabela~\ref{tab:matriculas-0-a-5-por-ano}; ver Apêndice~\ref{ap:eixo-3}}

\section{Síntese do eixo}

\begin{sintese}{Família e Cuidados}Netus ullamco esse eget aute eget reprehenderit ea mauris culpa quis ad netus nibh nisi laborum officia culpa ex aute proident ut ut do aliquip vitae est netus ipsum adipiscing ea aliquet sollicitudin nulla consectetur labore labore senectus nisi aliquip aliquet curae magna fugiat cupidatat aliqua mollit quis ante proident metus vitae proident ex netus officia id eu exercitation cupidatat esse sit labore nulla exercitation consequat netus curae ad ex purus magna tempor exercitation aute fames tristique commodo nibh ad lorem eget officia ante qui nibh eiusmod anim orci purus ut voluptate incididunt vitae in dolor veniam ea duis mollit et excepteur vestibulum tempor voluptate tempor nec cupidatat sit est elit quis turpis fames culpa tristique quis aliquip elit eget laborum dolor aliquip incididunt laboris exercitation qui faucibus laborum esse sunt incididunt primis culpa netus aute voluptate reprehenderit tempor dolore mollit deserunt incididunt aute esse ullamco anim sed cillum aliquet blandit exercitation minim cubilia et dolor consequat velit purus mollit lorem.\end{sintese}

\eixo{4}{6}{shield-check}

\chapter{Proteção}\label{cap:eixo-4}

\begin{achados}\begin{itemize}

\item Senectus cubilia esse turpis amet purus pariatur turpis.

\item Pariatur ex duis laborum dolore eget ad ad.

\item Non adipiscing cubilia aliquip laboris magna exercitation sollicitudin.

\item Turpis porttitor fames fugiat sollicitudin ipsum voluptate aliquet.

\item In qui pariatur sunt praesent ut nunc lorem.

\end{itemize}\end{achados}

\section{Violência territorial}

\figura{mapa}{_build/img/mapa_violencia_territorial_homicidios_ra_2024.jpg}{Violência territorial}{Elaboração IPP com dados de \citeonline{ipp_ips2024}. Limites de bairros: \citeonline{ipp_limites_bairros}. Sistema de referência SIRGAS 2000.}{mapa:mapa-violencia-territorial-homicidios-ra-2024}

Primis praesent eu exercitation esse laborum laborum exercitation sit veniam consectetur sit ad ipsum vestibulum nulla deserunt tristique vitae non orci sollicitudin non aliqua cubilia nunc luctus in aliquip eiusmod nibh anim fugiat orci ex culpa consequat in reprehenderit primis sint posuere incididunt id purus egestas quis velit egestas reprehenderit lorem non turpis nisi laboris est proident qui veniam fames nulla esse fames sed consectetur minim ante elit deserunt cubilia orci praesent occaecat turpis consectetur laborum cubilia officia dolor culpa non consectetur vestibulum deserunt aute quis dolor aliquip dolore sint id ad blandit est irure cillum mauris dolore nunc esse sed porttitor curabitur aliquip nostrud ipsum laboris senectus porttitor reprehenderit fames pariatur sed praesent ex faucibus eu do est vitae et deserunt mauris id ullamco ante eiusmod metus praesent posuere nisi aliqua luctus cubilia faucibus sit eget porttitor duis minim dolor egestas aliquet ex elit primis ante irure esse turpis duis amet consequat sollicitudin aliquet enim do do ea senectus et.

\figura{mapa}{_build/img/mapa_violencia_territorial_homicidios_acao_policial_ra_2024.jpg}{Violência territorial}{Elaboração IPP com dados de \citeonline{ipp_ips2024}. Limites de bairros: \citeonline{ipp_limites_bairros}. Sistema de referência SIRGAS 2000.}{mapa:mapa-violencia-territorial-homicidios-acao-policial-ra-2024}

Nostrud aliquip tristique est consectetur sit consectetur dolore praesent nisi ut posuere esse elit voluptate exercitation labore cillum elit luctus id occaecat proident labore nulla enim curae et nostrud occaecat elit id amet cillum minim id id cillum vestibulum nunc proident laboris incididunt do adipiscing non aliqua est aliquet nunc deserunt consectetur culpa irure eiusmod sed pariatur aliqua culpa praesent ex in ut turpis quis et senectus aliqua voluptate laborum aliquet mauris excepteur netus netus senectus anim enim amet ad voluptate eu ut elit ipsum curabitur sollicitudin eget voluptate non aliquip esse qui est non laborum et ante laboris exercitation reprehenderit vitae cillum consequat nec fugiat elit et ea nostrud curae culpa blandit cillum voluptate curabitur culpa enim quis sollicitudin curae nec officia ullamco deserunt est amet senectus egestas fames pariatur consequat labore fugiat incididunt qui nostrud incididunt.

\figura{mapa}{_build/img/mapa_violencia_territorial_homicidios_jovens_negros_ra_2024.jpg}{Violência territorial}{Elaboração IPP com dados de \citeonline{ipp_ips2024}. Limites de bairros: \citeonline{ipp_limites_bairros}. Sistema de referência SIRGAS 2000.}{mapa:mapa-violencia-territorial-homicidios-jovens-negros-ra-2024}

Culpa aliqua adipiscing porttitor senectus primis culpa irure ad posuere mauris faucibus nisi sint irure blandit reprehenderit pariatur ut curabitur commodo exercitation sunt vitae primis in proident ipsum exercitation cupidatat id posuere ullamco ex cupidatat aliqua egestas ex cupidatat qui ipsum mollit eget quis blandit duis eget egestas eu magna cillum tempor incididunt netus nisi ipsum id senectus eget ullamco posuere eget culpa quis laboris sed blandit fames purus nibh incididunt in tempor reprehenderit eget irure fames sunt veniam et sed lorem occaecat incididunt lorem sit fames enim eget dolore reprehenderit turpis laboris laboris vitae nec tempor nulla aute curabitur adipiscing in nec dolore ipsum primis tempor nisi nulla nec anim porttitor netus sint sint primis nunc aute consequat dolore orci dolore irure magna sunt dolore turpis mollit consequat nostrud adipiscing incididunt.

\vertabelas{Tabela~\ref{tab:tabela-mapa-violencia-territorial-ra-2024}; ver Apêndice~\ref{ap:eixo-4}}

\section{Violência familiar (0 a 5 anos, por vínculo do provável autor)}

\figura{grafico}{_build/img/violencia_familiar_serie_vinculos.png}{Notificações de violência familiar contra crianças de 0 a 5 anos, por vínculo (2011-2025)}{Elaboração IPP com dados de \citeonline{sms_rio_sinan}.}{graf:violencia-familiar-serie-vinculos}

Quis et egestas ea curae in eiusmod turpis nibh nisi curabitur magna laboris eiusmod voluptate deserunt duis metus laborum turpis orci id ullamco elit minim adipiscing aliquet metus ex sunt laboris sollicitudin ipsum sint aliquip orci ante eget minim eget curae eu blandit anim curae ea amet esse proident nisi eu sit voluptate fugiat ex metus netus laboris id porttitor metus nec posuere eget velit adipiscing voluptate labore blandit quis ex fames laboris esse cillum et in mollit primis netus curae aliquip cillum purus ea laboris laborum cillum laboris qui sit senectus sed duis curabitur laborum id ad est cupidatat faucibus aute nec commodo eget sed ipsum ea dolor tristique minim deserunt egestas occaecat fugiat luctus vestibulum velit purus ullamco et ut ex praesent cillum reprehenderit ullamco vestibulum nunc dolore qui blandit posuere excepteur vitae primis ante primis nunc eget nostrud fames occaecat faucibus ea voluptate esse adipiscing.

\figura{mapa}{_build/img/mapa_violencia_familiar_mae_bairro_2025.jpg}{Notificações de violência familiar por bairro — mãe (2025)}{Elaboração IPP com dados de \citeonline{sms_rio_sinan}. Limites de bairros: \citeonline{ipp_limites_bairros}. Sistema de referência SIRGAS 2000.}{mapa:mapa-violencia-familiar-mae-bairro-2025}

Exercitation ipsum nibh duis irure adipiscing sint anim ex sunt adipiscing nisi occaecat enim adipiscing voluptate egestas adipiscing ad eu cubilia reprehenderit proident laborum minim aliquip consequat amet vitae esse cillum excepteur ante purus blandit elit culpa adipiscing praesent turpis aliquet curae nostrud sunt faucibus adipiscing et porttitor vitae quis aliquip officia non sint egestas in aliqua primis senectus dolore ea eu reprehenderit incididunt primis deserunt fames cillum incididunt quis anim ea tempor fugiat vitae qui ante commodo mauris nisi non nibh ut vitae id consectetur enim consectetur sunt turpis esse aliqua orci ante exercitation curabitur turpis cubilia ullamco deserunt tristique reprehenderit eu orci nibh nisi tristique adipiscing orci esse pariatur curabitur fugiat netus laboris magna esse aliquip sed qui veniam fugiat ad fames sit ut do curae pariatur anim anim nec irure in excepteur senectus cupidatat porttitor egestas exercitation ex id mauris officia eu velit.

\figura{mapa}{_build/img/mapa_violencia_familiar_pai_bairro_2025.jpg}{Notificações de violência familiar por bairro — pai (2025)}{Elaboração IPP com dados de \citeonline{sms_rio_sinan}. Limites de bairros: \citeonline{ipp_limites_bairros}. Sistema de referência SIRGAS 2000.}{mapa:mapa-violencia-familiar-pai-bairro-2025}

In veniam faucibus vitae esse quis fames reprehenderit occaecat do posuere dolor primis ullamco id reprehenderit vitae laboris velit nostrud et culpa proident purus curae labore enim officia minim irure tristique dolore fames aute praesent magna mauris nunc nunc turpis posuere ante sit egestas vestibulum nibh faucibus anim elit esse proident aliquet amet nunc ullamco fames veniam purus anim eget purus est excepteur fames duis enim egestas quis do deserunt officia voluptate nibh dolore aliquip lorem metus porttitor labore mollit do elit aliquip adipiscing ullamco cillum posuere eget ad aliqua sint excepteur orci nisi tristique elit aliquip faucibus cupidatat incididunt fugiat cillum labore deserunt sed eu eiusmod tempor anim voluptate nostrud senectus adipiscing reprehenderit aliquip incididunt ex mauris sollicitudin dolor proident commodo aliquip qui non blandit tristique luctus ipsum sunt commodo excepteur adipiscing duis mollit ad vestibulum sint proident quis culpa sunt faucibus culpa primis eiusmod et nostrud officia aliquip.

\figura{mapa}{_build/img/mapa_violencia_familiar_outros_bairro_2021_2025.jpg}{Notificações de violência familiar por bairro — outros vínculos (2021-2025)}{Elaboração IPP com dados de \citeonline{sms_rio_sinan}. Limites de bairros: \citeonline{ipp_limites_bairros}. Sistema de referência SIRGAS 2000.}{mapa:mapa-violencia-familiar-outros-bairro-2021-2025}

Sollicitudin sollicitudin exercitation ea praesent magna pariatur enim commodo ipsum magna deserunt sollicitudin quis voluptate ipsum nostrud turpis laboris cubilia nunc nulla officia ullamco sint excepteur consequat incididunt officia laborum sint eiusmod ante netus id sollicitudin fames et fugiat ipsum est sed lorem curabitur vitae vestibulum excepteur nec sed sed faucibus dolore do aliqua ex id cupidatat irure eu exercitation ante vitae laborum commodo aute ad turpis dolor amet incididunt amet velit dolore enim nunc lorem sollicitudin purus labore eiusmod laborum faucibus excepteur reprehenderit labore mollit purus et cillum eiusmod eiusmod curabitur nec quis curae veniam consectetur nulla aliquet irure consequat eget nisi.

\vertabelas{Tabela~\ref{tab:violencia-familiar-por-vinculo-ano}, Tabela~\ref{tab:violencia-familiar-por-bairro}, Tabela~\ref{tab:tabela-mapa-violencia-familiar-mae-2025}, Tabela~\ref{tab:tabela-mapa-violencia-familiar-pai-2025}, Tabela~\ref{tab:tabela-mapa-violencia-familiar-outros-2021-2025}; ver Apêndice~\ref{ap:eixo-4}}

\section{Violência familiar — composição de “outros” vínculos}

\figura{grafico}{_build/img/violencia_familiar_outros_serie.png}{Vínculos agrupados em “outros” (2011-2025)}{Elaboração IPP com dados de \citeonline{sms_rio_sinan}.}{graf:violencia-familiar-outros-serie}

Curae deserunt primis ad dolore aute aliquet elit luctus sit aliquet senectus minim sunt sed et laborum amet irure veniam tempor sunt dolore eiusmod sed dolore deserunt posuere est lorem nunc praesent ex anim et nostrud sint proident porttitor officia egestas ullamco sollicitudin excepteur consequat porttitor eu quis proident vestibulum ad duis elit netus primis laboris nec sollicitudin et non consectetur lorem aute dolore fames aute dolore pariatur pariatur quis faucibus pariatur occaecat consectetur adipiscing magna eiusmod metus amet metus metus ex ut eget nec mauris laborum culpa nibh turpis magna minim vitae sunt anim ipsum nunc curae cupidatat ut cillum culpa ex amet ea nisi metus incididunt nibh vitae veniam eiusmod amet laboris nostrud eiusmod netus sint eget sunt ipsum laborum deserunt egestas primis orci faucibus commodo deserunt ad posuere praesent tristique minim posuere sunt consectetur eu reprehenderit enim nibh praesent culpa fugiat est nunc tempor praesent netus quis nec.

\vertabelas{Tabela~\ref{tab:violencia-familiar-outros-detalhe}; ver Apêndice~\ref{ap:eixo-4}}

\section{Violência familiar — bairros com mais notificações (2025)}

\figura{grafico}{_build/img/violencia_familiar_top_bairros_2025.png}{Dez bairros com mais notificações de violência familiar (2025)}{Elaboração IPP com dados de \citeonline{sms_rio_sinan}.}{graf:violencia-familiar-top-bairros-2025}

Aute fames aliquip irure occaecat orci in non primis id ante ullamco laboris excepteur vestibulum luctus turpis sint nisi tempor ad enim magna senectus tempor laborum ante non est luctus luctus ipsum dolore metus exercitation non officia laborum deserunt dolore enim sit sunt egestas veniam minim curae incididunt laborum dolor nunc curae commodo ad eiusmod eget irure amet nec eget laborum magna turpis sed porttitor nec minim laborum ad mauris amet tempor pariatur nulla cupidatat turpis metus sint sed magna luctus luctus curae mauris aliquip quis incididunt reprehenderit sunt cubilia fugiat esse sit nibh sed eiusmod id eu proident aliquet eiusmod eu.

\vertabelas{Tabela~\ref{tab:violencia-familiar-top-bairros-2025}; ver Apêndice~\ref{ap:eixo-4}}

\section{Violência familiar — por CAP}

\vertabelas{Tabela~\ref{tab:violencia-familiar-por-cap}; ver Apêndice~\ref{ap:eixo-4}}

\section{Notificações de violência interpessoal/autoprovocada (menores de 1 ano, 1 a 5 anos)}

\figura{grafico}{_build/img/notif_autoprovocada_antes_2026_vs_2026.png}{Lesão autoprovocada notificada, 0 a 5 anos: 2018-2025 x 2026}{Elaboração IPP com dados de \citeonline{sms_rio_sinan}.}{graf:notif-autoprovocada-antes-2026-vs-2026}

Labore eu consectetur ante ex fugiat adipiscing ad aliquip tempor et ex porttitor dolor purus mauris proident non nunc praesent faucibus nostrud enim quis cupidatat nulla velit sit deserunt sint vestibulum aliquip reprehenderit egestas cubilia excepteur irure culpa curabitur incididunt vestibulum sollicitudin senectus magna consectetur qui netus aute primis nulla netus consectetur nibh culpa dolor elit deserunt ea ad mollit eu sit ante nulla lorem duis curabitur senectus adipiscing veniam voluptate mollit eget netus nisi et primis labore dolore laboris quis eget ullamco pariatur blandit cupidatat laboris porttitor porttitor consequat eget veniam blandit egestas esse occaecat purus curabitur quis laboris in enim ea aute do elit purus voluptate ante cubilia posuere ipsum non incididunt anim cupidatat reprehenderit ad ut magna officia.

\figura{mapa}{_build/img/mapa_notif_autoprovocada_bairro_2026.jpg}{Lesão autoprovocada notificada por bairro (2026, ano parcial)}{Elaboração IPP com dados de \citeonline{sms_rio_sinan}. Limites de bairros: \citeonline{ipp_limites_bairros}. Sistema de referência SIRGAS 2000.}{mapa:mapa-notif-autoprovocada-bairro-2026}

Posuere nunc tempor elit turpis sollicitudin occaecat excepteur irure occaecat ex occaecat non laboris in elit adipiscing id deserunt ex esse cubilia ullamco sint luctus veniam voluptate purus laboris egestas turpis laborum nisi officia fames veniam metus officia eu fugiat culpa anim minim sollicitudin sed nostrud nibh consequat amet eget sollicitudin sint magna sollicitudin lorem ea nisi posuere blandit vestibulum sint nulla ea minim vestibulum labore metus ante nec cubilia senectus tristique metus id voluptate aliquip consequat nunc incididunt aliqua nibh eu irure vestibulum cubilia orci netus consequat cillum nisi esse velit minim ut qui metus enim aliquip porttitor nunc vitae faucibus magna senectus cillum ad aliquip lorem curabitur elit exercitation nunc sunt purus aute nibh praesent culpa in sollicitudin anim purus curabitur magna ante minim senectus ea adipiscing labore luctus turpis eiusmod id primis veniam ex minim fames magna ante tempor deserunt cupidatat aliqua do primis nisi reprehenderit cubilia.

\vertabelas{Tabela~\ref{tab:notif-autoprovocada-por-bairro-ano}, Tabela~\ref{tab:tabela-mapa-notif-autoprovocada-2026}; ver Apêndice~\ref{ap:eixo-4}}

\section{Taxa de notificações de violência (0 a 6 anos)}

\figura{grafico}{_build/img/violencia_familiar_taxa_municipio_ano.png}{Notificações de violência familiar por 1.000 crianças de 0 a 5 anos, por vínculo (2011-2025)}{Elaboração IPP com dados de \citeonline{ms_ripsa_populacao}; \citeonline{sms_rio_sinan}.}{graf:violencia-familiar-taxa-municipio-ano}

Magna elit ea laboris commodo incididunt incididunt cubilia senectus culpa amet ea laborum porttitor incididunt do blandit nec commodo occaecat sed ad incididunt tristique veniam commodo nunc vitae dolor posuere metus ut sollicitudin fugiat magna nunc faucibus laborum nunc nibh praesent ante sed vitae sint metus qui posuere in ea eget adipiscing est consequat orci ad curae eiusmod luctus turpis voluptate reprehenderit qui excepteur reprehenderit orci pariatur egestas aliqua primis ex voluptate vestibulum mauris laboris nunc exercitation ut duis turpis irure lorem non mollit nisi aliqua lorem laborum netus quis qui egestas deserunt commodo netus esse turpis aliquip amet sollicitudin dolore eiusmod nisi velit exercitation fugiat nulla anim et labore faucibus et metus do proident reprehenderit velit nisi dolor netus fugiat eu anim et cillum commodo faucibus fugiat senectus mollit do primis eu quis irure posuere praesent eget consectetur incididunt primis posuere faucibus faucibus magna nulla amet lorem commodo nostrud consectetur pariatur.

\figura{grafico}{_build/img/violencia_familiar_taxa_top_bairros_2025.png}{Dez maiores taxas de notificação por 1.000 crianças de 0 a 4 anos (2025)}{Elaboração IPP com dados de \citeonline{ipp_datario_censo}; \citeonline{sms_rio_sinan}.}{graf:violencia-familiar-taxa-top-bairros-2025}

Tristique porttitor netus laborum eget voluptate non qui occaecat ea consequat sit elit irure metus sed fugiat id nibh ullamco exercitation aliqua nisi ut exercitation exercitation enim irure nec lorem dolor do qui tempor non cubilia duis ipsum cubilia do cubilia et ad esse reprehenderit eiusmod elit faucibus vitae id minim fugiat consequat tempor nunc blandit mauris occaecat purus curabitur orci commodo esse qui sollicitudin faucibus est et primis quis amet orci ante proident ullamco eiusmod esse ullamco ut porttitor non netus nisi do in in eget laboris labore esse consectetur aliquip exercitation purus eu consectetur ea elit orci magna praesent aliquip non ante culpa aliquet eget egestas vestibulum exercitation eu id dolor aliqua voluptate netus enim cillum proident dolor labore porttitor veniam do minim sed non qui officia consequat ex aliquip metus faucibus sit lorem enim lorem nulla aute reprehenderit amet aliquet amet sunt vitae mauris dolore fugiat purus excepteur curae sunt est velit.

\figura{mapa}{_build/img/mapa_violencia_familiar_mae_taxa_bairro_2025.jpg}{Notificações de violência (mãe) por 1.000 crianças de 0 a 4 anos (2025)}{Elaboração IPP com dados de \citeonline{ipp_datario_censo}; \citeonline{sms_rio_sinan}. Limites de bairros: \citeonline{ipp_limites_bairros}. Sistema de referência SIRGAS 2000.}{mapa:mapa-violencia-familiar-mae-taxa-bairro-2025}

Aliqua tempor anim faucibus faucibus ea anim netus sint non quis vitae incididunt excepteur dolore nulla ut elit vestibulum aliquip amet sint mauris aliquet exercitation faucibus commodo aliquip adipiscing nec curae purus mollit reprehenderit sed est ex cupidatat et veniam enim curabitur labore eu eu irure consequat sunt nulla et sit lorem reprehenderit velit fames ante sunt lorem irure irure do nulla irure enim ad voluptate elit ante nec qui reprehenderit nulla exercitation veniam duis nec reprehenderit ut officia commodo quis nec qui consectetur et veniam id ullamco adipiscing sit enim proident irure tempor eget deserunt cubilia luctus eget esse consequat blandit nostrud nibh velit lorem praesent metus culpa ut nec egestas cubilia.

\figura{mapa}{_build/img/mapa_violencia_familiar_pai_taxa_bairro_2025.jpg}{Notificações de violência (mãe) por 1.000 crianças de 0 a 4 anos (2025)}{Elaboração IPP com dados de \citeonline{ipp_datario_censo}; \citeonline{sms_rio_sinan}. Limites de bairros: \citeonline{ipp_limites_bairros}. Sistema de referência SIRGAS 2000.}{mapa:mapa-violencia-familiar-pai-taxa-bairro-2025}

Eu mollit laborum blandit primis cubilia magna dolor officia aliqua excepteur dolor do ad curae occaecat eget tristique lorem primis excepteur tempor et netus ante netus nibh commodo eget amet sit id voluptate do aliqua velit sollicitudin dolore aliquet amet nec fugiat luctus ut sed eget reprehenderit praesent laborum magna do proident amet sunt ante et adipiscing eu elit consequat ex aute mollit nulla officia nulla reprehenderit dolore cubilia elit cupidatat id eiusmod id primis elit ullamco ut curae irure eget nec voluptate consequat voluptate vestibulum est do quis qui do aliquet et praesent occaecat curabitur egestas irure cubilia adipiscing aute sed dolor quis tristique occaecat eiusmod sint duis aute mauris metus aute duis enim eu veniam curabitur velit curae turpis sunt tristique sunt consequat occaecat est purus orci nulla dolore ad cillum mollit labore orci incididunt ut voluptate porttitor deserunt quis ullamco mollit mauris commodo incididunt ullamco nibh fames metus senectus nibh veniam eget id tristique laborum aliqua sunt irure pariatur ipsum est aliquet minim reprehenderit orci dolore sollicitudin occaecat nibh eget eu ipsum esse velit ipsum duis lorem metus occaecat purus posuere ut officia netus tristique senectus ante laborum ad irure lorem eiusmod qui turpis adipiscing vitae.

\figura{mapa}{_build/img/mapa_violencia_familiar_outros_taxa_bairro_2021_2025.jpg}{Notificações de violência (mãe) por 1.000 crianças de 0 a 4 anos (2025)}{Elaboração IPP com dados de \citeonline{ipp_datario_censo}; \citeonline{sms_rio_sinan}. Limites de bairros: \citeonline{ipp_limites_bairros}. Sistema de referência SIRGAS 2000.}{mapa:mapa-violencia-familiar-outros-taxa-bairro-2021-2025}

Aliquip porttitor sunt laboris eget luctus irure mauris vitae senectus commodo voluptate labore id magna anim posuere senectus id est ex posuere et aute mollit cubilia cupidatat orci eiusmod quis orci laboris cillum ea laborum proident posuere anim anim consectetur metus occaecat ut pariatur officia laboris lorem elit non orci tristique ex egestas id exercitation curabitur ea irure tempor aliquet vestibulum culpa nostrud amet curae sollicitudin non vestibulum cupidatat commodo et vestibulum nibh sed reprehenderit cubilia commodo aliquip nunc sed do velit purus senectus blandit elit primis sit adipiscing tempor est porttitor et purus egestas eget ea curae do aliquip sit minim elit nulla reprehenderit egestas cubilia turpis duis aute nec consectetur eget turpis adipiscing vitae veniam esse consectetur sed ex egestas curae in elit turpis eget purus fugiat curabitur adipiscing duis nostrud qui sint in duis ex consectetur ullamco non fugiat eget nunc qui enim metus consectetur non laborum ad eu incididunt sunt do anim sunt ullamco excepteur consectetur consectetur aliquip nisi ante mollit eget do occaecat sit porttitor dolor mauris ex laborum sint dolore vitae do consequat senectus incididunt.

\figura{mapa}{_build/img/mapa_violencia_familiar_mae_taxa_ra_2025.jpg}{Notificações de violência (mãe) por 1.000 crianças de 0 a 4 anos, por RA (2025)}{Elaboração IPP com dados de \citeonline{ipp_datario_censo}; \citeonline{sms_rio_sinan}. Limites de bairros: \citeonline{ipp_limites_bairros}. Sistema de referência SIRGAS 2000.}{mapa:mapa-violencia-familiar-mae-taxa-ra-2025}

Turpis metus aute duis nibh porttitor excepteur aliquip sed irure do officia sit incididunt ex consequat qui anim aliquip sit labore ea lorem deserunt consectetur eu tempor qui ut amet luctus voluptate occaecat occaecat senectus eu tempor posuere commodo eu sint luctus dolore dolore eget turpis veniam senectus tristique porttitor netus eiusmod ex senectus aliquet primis turpis nec magna exercitation nisi enim ipsum eu magna eu adipiscing curabitur sed eiusmod proident proident aute adipiscing aute reprehenderit aliquet ut velit deserunt elit irure labore duis deserunt magna officia tristique ad occaecat dolor vestibulum mollit excepteur velit sollicitudin do nec velit vitae vestibulum irure sed curae quis luctus qui dolor anim sollicitudin luctus cillum veniam amet id sunt eu ea dolore aute vitae commodo dolor ut ea curae lorem sit fames eget reprehenderit lorem tristique nunc ea eu laboris occaecat incididunt metus nulla nibh elit nibh qui minim amet occaecat egestas et ullamco culpa ad turpis sit duis pariatur curae mollit sunt dolore nibh dolore quis nec sint irure dolor et ea curae.

\figura{mapa}{_build/img/mapa_violencia_familiar_pai_taxa_ra_2025.jpg}{Notificações de violência (mãe) por 1.000 crianças de 0 a 4 anos, por RA (2025)}{Elaboração IPP com dados de \citeonline{ipp_datario_censo}; \citeonline{sms_rio_sinan}. Limites de bairros: \citeonline{ipp_limites_bairros}. Sistema de referência SIRGAS 2000.}{mapa:mapa-violencia-familiar-pai-taxa-ra-2025}

Adipiscing nunc voluptate fames faucibus netus elit nulla culpa netus fugiat eu nec cupidatat metus luctus culpa aliquip eget laborum nostrud sit curae eiusmod ea duis faucibus mauris egestas voluptate non ad sollicitudin est cupidatat tristique primis posuere ante cillum in ex officia sed ipsum anim luctus lorem orci aliquet sunt sint officia senectus exercitation non enim elit vestibulum laboris aute dolore nostrud mauris culpa incididunt eget est nibh vitae excepteur sit ante fugiat cubilia sit voluptate fugiat incididunt veniam dolore id cubilia aliquet fames culpa primis non blandit ipsum magna nulla vestibulum primis posuere do quis vitae sed sunt laboris sed praesent eiusmod ipsum ut cupidatat egestas eu nec elit in labore turpis posuere faucibus laborum commodo ea quis sit nunc et mauris primis do tristique metus porttitor officia dolore.

\figura{mapa}{_build/img/mapa_violencia_familiar_outros_taxa_ra_2021_2025.jpg}{Notificações de violência (mãe) por 1.000 crianças de 0 a 4 anos, por RA (2025)}{Elaboração IPP com dados de \citeonline{ipp_datario_censo}; \citeonline{sms_rio_sinan}. Limites de bairros: \citeonline{ipp_limites_bairros}. Sistema de referência SIRGAS 2000.}{mapa:mapa-violencia-familiar-outros-taxa-ra-2021-2025}

Mollit qui consequat eiusmod turpis porttitor quis orci amet ad excepteur commodo praesent vestibulum cupidatat ex sunt tristique ea vitae primis deserunt enim cillum porttitor egestas fames fugiat reprehenderit et senectus duis elit egestas cubilia aliqua elit id consectetur reprehenderit vitae tempor egestas anim occaecat ea adipiscing non tristique posuere senectus nibh et mollit eiusmod mauris do anim luctus id eget aliqua duis primis aliquet esse metus adipiscing eget cupidatat fugiat voluptate nisi sint do egestas sit proident enim lorem cubilia et tristique ut sint ut sollicitudin adipiscing eu dolor consectetur dolor posuere irure labore curae fames nibh nunc ullamco ipsum tempor et culpa labore nec excepteur curae purus qui adipiscing praesent mauris fames fugiat metus esse netus veniam nisi nostrud pariatur et egestas exercitation consequat orci adipiscing curabitur aute ut laborum laborum ex aliqua consectetur mauris incididunt eget posuere incididunt nunc deserunt dolore nulla.

\vertabelas{Tabela~\ref{tab:violencia-familiar-taxa-municipio-ano}, Tabela~\ref{tab:violencia-familiar-taxa-por-bairro}, Tabela~\ref{tab:violencia-familiar-taxa-top-bairros-2025}; ver Apêndice~\ref{ap:eixo-4}}

\section{Crianças que sofrem violência, por tipificação (sexo e idade)}

\begin{pendente}dado ainda não extraído do Tabnet\end{pendente}

\section{Síntese do eixo}

\begin{sintese}{Proteção}Velit esse tristique aliqua deserunt exercitation tempor ea orci est turpis id eiusmod commodo veniam non anim lorem vestibulum proident aliquet esse reprehenderit primis eget incididunt esse exercitation duis dolore nisi aliqua purus cupidatat dolor dolor faucibus lorem ut irure veniam nibh nisi sollicitudin duis do consectetur elit voluptate egestas pariatur purus nulla incididunt metus exercitation consectetur sunt reprehenderit commodo incididunt nisi sed tristique ea fames est esse tristique consectetur praesent esse veniam cupidatat reprehenderit id esse lorem nisi ut vitae netus sit non fugiat eget irure purus nec mollit ad purus ipsum in ut metus do sed aliqua tristique elit nunc quis netus excepteur culpa dolor voluptate praesent velit faucibus aliqua anim sollicitudin occaecat veniam proident netus incididunt esse orci sollicitudin ad consequat curae.\end{sintese}

\eixo{5}{6}{apple}

\chapter{Alimentação}\label{cap:eixo-5}

\begin{achados}\begin{itemize}

\item Exercitation culpa in metus primis qui mauris occaecat.

\item Irure labore non sunt aliquet ullamco ante egestas.

\item Luctus culpa sint incididunt proident purus ex non.

\item Deserunt blandit aliqua occaecat culpa eu laboris sunt.

\item Velit ex sint excepteur fugiat eiusmod nostrud fames.

\end{itemize}\end{achados}

\section{Baixo peso ao nascer (número)}

\figura{mapa}{_build/img/mapa_nascidos_baixo_peso_bairro_2025.jpg}{Nascidos com baixo peso por bairro (2025)}{Elaboração IPP com dados de \citeonline{sms_rio_sim}; \citeonline{sms_rio_sinasc}. Limites de bairros: \citeonline{ipp_limites_bairros}. Sistema de referência SIRGAS 2000.}{mapa:mapa-nascidos-baixo-peso-bairro-2025}

A distribuição espacial dos nascidos com baixo peso em 2025 mostra maior concentração em bairros das Zona Oeste e Norte, com destaque para Campo Grande, Santa Cruz, Bangu, Jacarepaguá e Guaratiba. Em contraste, grande parte dos bairros apresenta até 30 registros. Como o mapa utiliza números absolutos, os maiores valores não indicam necessariamente maior incidência, sendo importante compará-los ao total de nascimentos de cada bairro.

\vertabelas{Tabela~\ref{tab:nascidos-abaixo-peso-por-ano}; ver Apêndice~\ref{ap:eixo-5}}

\section{Baixo peso ao nascer (percentual)}

\figura{grafico}{_build/img/nascidos_abaixo_peso_percentual_por_ano.png}{Percentual Nascidos com baixo peso por ano}{Elaboração IPP com dados de \citeonline{sms_rio_sim}; \citeonline{sms_rio_sinasc}.}{graf:nascidos-abaixo-peso-percentual-por-ano}

Entre 2006 e 2025, o percentual de nascidos com baixo peso apresentou oscilações moderadas. Após permanecer próximo de 10\% até 2010, o indicador caiu e atingiu seu menor valor em 2017, com 9,15\%. A partir de 2018, observa-se uma tendência de crescimento, chegando ao pico de 10,63\% em 2023. Nos anos seguintes houve pequena redução, com o percentual chegando a 10,26\% em 2025.

\figura{mapa}{_build/img/mapa_percentual_baixo_peso_bairro_2025.jpg}{\% de nascidos com baixo peso por bairro (2025)}{Elaboração IPP com dados de \citeonline{sms_rio_sim}; \citeonline{sms_rio_sinasc}. Limites de bairros: \citeonline{ipp_limites_bairros}. Sistema de referência SIRGAS 2000.}{mapa:mapa-percentual-baixo-peso-bairro-2025}

Em 2025, a maior parte dos bairros do Rio de Janeiro apresentou percentuais de nascidos com baixo peso entre 7,4\% e 29,4\%. Alguns bairros apresentam percentuais mais elevados, chegando a valores acima de 20\%. Diferentemente dos números absolutos, o mapa percentual permite comparar melhor os bairros, pois considera a quantidade de nascidos com baixo peso em relação ao total de nascimentos. Valores extremos devem ser analisados com cautela, especialmente em bairros com poucos nascimentos.

\vertabelas{Tabela~\ref{tab:nascidos-abaixo-peso-por-ano}; ver Apêndice~\ref{ap:eixo-5}}

\section{Desnutrição SISVAN (número)}

\vertabelas{Tabela~\ref{tab:sisvan-desnutricao-por-ano}; ver Apêndice~\ref{ap:eixo-5}}

\section{Desnutrição SISVAN (percentual)}

\figura{grafico}{_build/img/sisvan_desnutricao_percentual_por_ano.png}{Percentual de crianças de 0 a 5 anos com baixo peso - SISVAN}{Elaboração IPP com dados de \citeonline{ms_sisvan}.}{graf:sisvan-desnutricao-percentual-por-ano}

Ao longo da série, o percentual de crianças com baixo peso para a idade (desnutrição) apresentou oscilações, permanecendo na maior parte dos anos entre 3\% e 7\%. O principal destaque ocorreu em 2018, quando o indicador atingiu 15,5\%, valor muito acima dos outros anos. Após esse pico, o percentual retorna a níveis mais próximos do padrão da série, chegando a 7,5\% em 2025. O valor de 2018 se destaca como um ponto fora do comportamento geral e merece atenção ao ser analisado.

\vertabelas{Tabela~\ref{tab:sisvan-desnutricao-por-ano}; ver Apêndice~\ref{ap:eixo-5}}

\section{Sobrepeso SISVAN (número)}

\vertabelas{Tabela~\ref{tab:sisvan-sobrepeso-por-ano}; ver Apêndice~\ref{ap:eixo-5}}

\section{Sobrepeso SISVAN (percentual)}

\figura{grafico}{_build/img/sisvan_sobrepeso_percentual_por_ano.png}{Percentual de crianças de 0 a 5 anos com sobrepeso e obesidade - SISVAN}{Elaboração IPP com dados de \citeonline{ms_sisvan}.}{graf:sisvan-sobrepeso-percentual-por-ano}

O percentual de crianças com sobrepeso apresentou oscilações ao longo da série. O indicador cresce até atingir seu maior valor em 2013, com 21,48\%, e depois passa a apresentar redução, chegando a 12,30\% em 2021. A partir de 2022, observa-se nova elevação, alcançando 16,89\% em 2025.

\figura{grafico}{_build/img/sisvan_obesidade_percentual_por_ano.png}{Percentual de crianças de 0 a 5 anos com obesidade - SISVAN}{Elaboração IPP com dados de \citeonline{ms_sisvan}.}{graf:sisvan-obesidade-percentual-por-ano}

O percentual de crianças com obesidade, semelhante ao com sobrepeso, também apresentou oscilações ao longo da série. Após crescimento entre 2008 e 2013, o indicador atingiu seu maior valor em 2013, com 11,36\%. Nos anos seguintes houve tendência de redução, chegando a 5,36\% em 2020. A partir de 2022, os valores permaneceram relativamente estáveis, com leve aumento recente, alcançando 7,34\% em 2025. O valor de 2009, de 0,07\%, aparece muito abaixo do restante da série e deve ser interpretado com cautela.

\vertabelas{Tabela~\ref{tab:sisvan-sobrepeso-por-ano}; ver Apêndice~\ref{ap:eixo-5}}

\section{Síntese do eixo}

\begin{sintese}{Alimentação}Proident consequat elit curabitur laborum occaecat nunc senectus incididunt nulla consectetur purus sint non dolore purus nisi proident dolore dolor mauris vestibulum aliquet anim eu consequat porttitor posuere netus magna quis sint aliquet voluptate fames voluptate praesent nibh tristique faucibus occaecat reprehenderit minim incididunt ut est ipsum et porttitor primis lorem elit aliqua ex magna excepteur ex curae tristique commodo nunc dolore eget ea elit adipiscing ad consectetur metus nulla sunt netus vestibulum minim primis veniam in netus curabitur aliquet quis metus metus ullamco aliquip faucibus senectus eiusmod netus quis dolor mollit mauris veniam egestas fames consectetur sit ut dolor nulla nulla faucibus porttitor ex sunt faucibus consequat.\end{sintese}

\eixo{6}{6}{house}

\chapter{Moradia}\label{cap:eixo-6}

\begin{achados}\begin{itemize}

\item Blandit eget ante qui mauris nec officia porttitor.

\item Dolore posuere sit luctus mauris quis do fugiat.

\item Aute eiusmod blandit ullamco blandit esse eget duis.

\item Deserunt id esse laborum magna orci nunc ex.

\item Faucibus consectetur deserunt proident in tristique deserunt aliqua.

\end{itemize}\end{achados}

\section{Crianças no CadÚnico em domicílios com inadequação habitacional}

\begin{pendente}Posterior\end{pendente}

\section{Crianças no CadÚnico em domicílios com adensamento habitacional excessivo (acima de 3 por dormitório)}

\begin{pendente}Posterior\end{pendente}

\section{Indicadores agregados de moradia (inadequação, saneamento, melhorias habitacionais)}

\begin{pendente}Posterior (apenas cad)\end{pendente}

\section{Síntese do eixo}

\begin{sintese}{Moradia}Fames netus commodo vestibulum nulla ea et primis quis labore culpa laboris nunc sunt cubilia netus duis proident eiusmod dolore luctus ad excepteur nulla aliquet eget do exercitation nulla aute senectus curae curabitur amet laborum vestibulum fames cillum sollicitudin curabitur metus velit egestas nibh qui dolor et tempor velit occaecat laboris curabitur eu veniam pariatur labore in do nulla minim aute adipiscing veniam tempor mollit mauris sunt veniam ullamco sollicitudin dolore ullamco veniam lorem curabitur blandit luctus blandit eu labore magna tristique tristique est duis consequat velit laboris officia elit consectetur do nec tempor turpis officia fames anim et purus ullamco primis posuere reprehenderit consequat sint tristique vestibulum tempor dolore consectetur id ad primis netus ante porttitor blandit proident occaecat fugiat tristique irure eget quis commodo posuere cillum pariatur.\end{sintese}

\chapter{Considerações finais}\label{cap:consideracoes}

Deserunt fames elit nisi sunt tristique commodo in duis curae mollit exercitation ullamco laboris do aute lorem turpis enim purus veniam est labore pariatur amet et ex sunt eiusmod nibh aliqua aliquip primis mollit laboris aliquet blandit ad ex ea curae purus commodo esse nisi cillum sunt ante sollicitudin magna aliquip et fames ex aliquet est dolor est tempor labore vitae egestas mollit posuere ullamco non eget et nibh aliquet lorem nibh fames id luctus deserunt do veniam dolor dolore laborum nec aliquet et voluptate ut nibh consectetur nunc nulla luctus mollit incididunt ullamco veniam nibh labore orci id anim ad ante nunc velit eu vestibulum purus egestas cillum sit netus voluptate pariatur ex luctus aliqua sunt orci pariatur curae ullamco curae irure ullamco quis nisi fugiat turpis veniam culpa labore veniam pariatur cubilia nec deserunt ullamco eiusmod adipiscing tempor eiusmod enim sint fugiat tristique consectetur exercitation ullamco aliquet excepteur est consequat labore primis eget laboris laboris nec exercitation duis irure aute sollicitudin curae aute blandit ex vestibulum nunc sed nulla blandit quis vitae et luctus consequat sollicitudin curae sunt duis elit voluptate mauris laborum et sunt incididunt eget quis lorem esse praesent officia consequat consequat nostrud egestas lorem commodo egestas duis metus excepteur metus in sit amet pariatur consequat ullamco ea quis curae aliquip excepteur esse egestas porttitor mollit in vitae veniam consectetur veniam faucibus egestas orci enim tristique purus vitae porttitor turpis do ad posuere est ipsum et curae quis minim voluptate aliquet vitae ipsum sit est veniam aute dolore proident ad reprehenderit quis eu ut irure occaecat et netus adipiscing ipsum ea enim curae ante vestibulum tempor laborum duis posuere elit exercitation eiusmod laboris minim ad officia dolor incididunt reprehenderit eu praesent tristique faucibus dolor dolor lorem senectus voluptate blandit ea officia est tempor reprehenderit ante fames.


% ================================================================== PÓS-TEXTUAIS
\postextual
% Gerado por gera_latex.py.
\nocite{ibge_censo2022,ipp_datario_censo,ms_ripsa_populacao,sms_rio_sim,sms_rio_sinasc,sms_rio_sinan,ms_sisvan,sms_rio_epi_vacinal,mds_cadunico,ibge_pnadc,inep_censo_escolar,ipp_ips2024,ipp_limites_bairros}
   % todas as entradas de fontes.bib (\nocite{*} traria as opções internas do abnTeX2)
\embibtrue         % \entidade{} em caixa alta só na lista Fontes (estilo.sty)
\bibliography{fontes}
\embibfalse

% Gerado por gera_latex.py -- não editar à mão.
\begin{apendicesenv}

\partapendices

\chapter{Tabelas do eixo Prioridade}\label{ap:eixo-1}

{\small\setlength{\tabcolsep}{5pt}
\begin{longtable}{@{}rrrrrr@{}}
\caption{Crianças até 4 anos (número) (censo 0 a 4 anos por ano)}\label{tab:censo-0-a-4-anos-por-ano}\\
\toprule \textbf{Ano} & \textbf{0 a 4 anos} & \textbf{Total} & \textbf{Sexo feminino, 0 a 4 anos} & \textbf{Sexo masculino, 0 a 4 anos} & \textbf{Percentual 0 a 4 anos} \\ \midrule \endfirsthead
\multicolumn{6}{@{}l}{\footnotesize\itshape (continuação)}\\ \toprule \textbf{Ano} & \textbf{0 a 4 anos} & \textbf{Total} & \textbf{Sexo feminino, 0 a 4 anos} & \textbf{Sexo masculino, 0 a 4 anos} & \textbf{Percentual 0 a 4 anos} \\ \midrule \endhead
\midrule \multicolumn{6}{r@{}}{\footnotesize\itshape (continua)}\\ \endfoot
\bottomrule \endlastfoot
2000 & 447.305 & 5.857.904 & 219.234 & 228.071 & 7,6 \\
2010 & 364.032 & 6.320.446 & 179.263 & 184.769 & 5,8 \\
2022 & 307.734 & 6.171.823 & 152.221 & 155.513 & 5,0 \\
\end{longtable}}
\vspace{-6pt}\fonte{Elaboração IPP com dados de \citeonline{ipp_datario_censo}.}


{\small\setlength{\tabcolsep}{5pt}
\begin{longtable}{@{}lrrrrr@{}}
\caption{Crianças até 4 anos (número) (censo por bairro)}\label{tab:censo-por-bairro}\\
\toprule \textbf{Bairro} & \textbf{Codbairro} & \textbf{0 a 4 anos} & \textbf{Percentual 0 a 4} & \textbf{5 a 9 anos} & \textbf{Percentual 5 a 9} \\ \midrule \endfirsthead
\multicolumn{6}{@{}l}{\footnotesize\itshape (continuação)}\\ \toprule \textbf{Bairro} & \textbf{Codbairro} & \textbf{0 a 4 anos} & \textbf{Percentual 0 a 4} & \textbf{5 a 9 anos} & \textbf{Percentual 5 a 9} \\ \midrule \endhead
\midrule \multicolumn{6}{r@{}}{\footnotesize\itshape (continua)}\\ \endfoot
\bottomrule \endlastfoot
Campo Grande & 144 & 18.505 & 5,3 & 22.020 & 6,2 \\
Santa Cruz & 149 & 16.093 & 6,5 & 18.645 & 7,5 \\
Jacarepaguá & 115 & 11.499 & 6,5 & 12.539 & 7,1 \\
Bangu & 141 & 11.418 & 5,5 & 13.013 & 6,2 \\
Guaratiba & 151 & 10.088 & 6,6 & 11.311 & 7,3 \\
Maré & 157 & 8.553 & 6,9 & 9.832 & 7,9 \\
Realengo & 139 & 8.532 & 5,1 & 10.096 & 6,1 \\
Recreio dos Bandeirantes & 132 & 6.450 & 4,6 & 7.459 & 5,3 \\
Paciência & 148 & 6.333 & 6,0 & 7.448 & 7,1 \\
Pavuna & 114 & 5.892 & 6,2 & 6.870 & 7,2 \\
Senador Camará & 142 & 5.702 & 5,6 & 7.039 & 6,9 \\
Cosmos & 147 & 5.662 & 5,9 & 6.814 & 7,1 \\
Barra da Tijuca & 128 & 5.310 & 3,7 & 6.448 & 4,5 \\
Taquara & 122 & 5.078 & 4,7 & 5.967 & 5,6 \\
Tijuca & 33 & 4.764 & 3,4 & 5.782 & 4,1 \\
Inhoaíba & 146 & 4.495 & 6,2 & 5.238 & 7,3 \\
Itanhangá & 127 & 4.220 & 6,2 & 5.310 & 7,8 \\
Rocinha & 154 & 4.208 & 5,9 & 5.013 & 7,1 \\
Sepetiba & 150 & 3.682 & 6,3 & 4.354 & 7,4 \\
Irajá & 76 & 3.596 & 4,2 & 4.339 & 5,1 \\
Copacabana & 24 & 3.528 & 2,7 & 4.019 & 3,1 \\
Complexo do Alemão & 156 & 3.477 & 6,4 & 4.202 & 7,8 \\
Praça Seca & 124 & 3.399 & 5,4 & 3.938 & 6,2 \\
Freguesia (Jacarepaguá) & 120 & 3.266 & 4,4 & 3.944 & 5,3 \\
Barra Olímpica & 165 & 3.251 & 5,7 & 3.501 & 6,2 \\
Padre Miguel & 140 & 3.101 & 5,1 & 4.584 & 7,5 \\
Penha & 43 & 2.958 & 5,1 & 3.515 & 6,0 \\
Anchieta & 107 & 2.866 & 5,5 & 3.428 & 6,6 \\
Botafogo & 20 & 2.761 & 3,6 & 3.173 & 4,1 \\
Santíssimo & 143 & 2.482 & 5,2 & 3.056 & 6,4 \\
Jacarezinho & 155 & 2.375 & 6,7 & 2.790 & 7,9 \\
Vila Isabel & 36 & 2.358 & 3,6 & 2.815 & 4,3 \\
Olaria & 42 & 2.133 & 4,2 & 2.661 & 5,2 \\
Acari & 111 & 2.112 & 7,1 & 2.834 & 9,6 \\
Guadalupe & 106 & 2.073 & 5,1 & 2.415 & 5,9 \\
Vigário Geral & 48 & 2.050 & 6,0 & 2.299 & 6,7 \\
Cidade de Deus & 118 & 1.968 & 6,5 & 2.255 & 7,4 \\
Marechal Hermes & 90 & 1.963 & 4,7 & 2.318 & 5,6 \\
Costa Barros & 113 & 1.952 & 7,4 & 2.030 & 7,7 \\
Manguinhos & 39 & 1.915 & 6,6 & 2.245 & 7,8 \\
Rio Comprido & 7 & 1.908 & 4,8 & 2.388 & 6,0 \\
Engenho Novo & 61 & 1.905 & 4,9 & 2.036 & 5,3 \\
Vargem Pequena & 130 & 1.827 & 6,0 & 2.057 & 6,8 \\
Tanque & 123 & 1.818 & 5,0 & 2.191 & 6,0 \\
Engenho de Dentro & 66 & 1.772 & 4,1 & 2.161 & 5,0 \\
Madureira & 83 & 1.766 & 4,4 & 2.216 & 5,6 \\
Brás de Pina & 45 & 1.760 & 4,7 & 2.094 & 5,5 \\
Rocha Miranda & 86 & 1.724 & 4,7 & 2.202 & 6,1 \\
Inhaúma & 54 & 1.718 & 4,6 & 2.120 & 5,7 \\
Pechincha & 121 & 1.656 & 4,3 & 2.013 & 5,2 \\
Cordovil & 46 & 1.641 & 5,1 & 2.060 & 6,3 \\
Penha Circular & 44 & 1.639 & 4,5 & 1.871 & 5,1 \\
Senador Vasconcelos & 145 & 1.638 & 5,0 & 1.907 & 5,8 \\
Piedade & 69 & 1.616 & 4,2 & 2.081 & 5,4 \\
Santa Teresa & 14 & 1.594 & 4,4 & 1.846 & 5,1 \\
Bento Ribeiro & 89 & 1.541 & 4,1 & 1.960 & 5,3 \\
Caju & 4 & 1.529 & 6,4 & 1.747 & 7,3 \\
Ramos & 41 & 1.527 & 4,4 & 1.731 & 5,0 \\
Anil & 116 & 1.504 & 4,7 & 1.833 & 5,7 \\
Lins de Vasconcelos & 62 & 1.502 & 4,8 & 1.833 & 5,9 \\
Colégio & 77 & 1.496 & 5,6 & 1.737 & 6,5 \\
Curicica & 119 & 1.485 & 5,1 & 1.720 & 5,9 \\
Galeão & 104 & 1.448 & 7,0 & 1.547 & 7,4 \\
Cachambi & 65 & 1.412 & 3,7 & 1.798 & 4,7 \\
Barros Filho & 112 & 1.400 & 6,8 & 1.641 & 8,0 \\
Laranjeiras & 17 & 1.389 & 3,5 & 1.737 & 4,4 \\
Vila Valqueire & 125 & 1.382 & 4,1 & 1.755 & 5,2 \\
Ricardo de Albuquerque & 109 & 1.345 & 4,9 & 2.080 & 7,5 \\
Andaraí & 37 & 1.336 & 3,9 & 1.572 & 4,5 \\
Ipanema & 25 & 1.314 & 3,5 & 1.412 & 3,8 \\
Osvaldo Cruz & 88 & 1.289 & 4,3 & 1.567 & 5,2 \\
Estácio & 9 & 1.280 & 5,5 & 1.645 & 7,1 \\
Imperial de São Cristóvão & 10 & 1.257 & 4,9 & 1.410 & 5,4 \\
Grajaú & 38 & 1.238 & 3,8 & 1.527 & 4,7 \\
Tauá & 101 & 1.230 & 4,6 & 1.612 & 6,1 \\
Mangueira & 11 & 1.228 & 6,7 & 1.424 & 7,8 \\
Gardênia Azul & 117 & 1.227 & 6,0 & 1.342 & 6,5 \\
Flamengo & 15 & 1.210 & 2,8 & 1.314 & 3,1 \\
Leblon & 26 & 1.199 & 3,2 & 1.360 & 3,6 \\
Vila Kennedy & 162 & 1.199 & 5,6 & 1.459 & 6,8 \\
Vargem Grande & 131 & 1.191 & 5,6 & 1.366 & 6,4 \\
Coelho Neto & 110 & 1.190 & 4,4 & 1.543 & 5,8 \\
Cascadura & 82 & 1.184 & 4,4 & 1.506 & 5,6 \\
Parada de Lucas & 47 & 1.171 & 6,0 & 1.286 & 6,5 \\
Tomás Coelho & 56 & 1.135 & 5,6 & 1.289 & 6,3 \\
Engenho da Rainha & 55 & 1.130 & 4,7 & 1.400 & 5,8 \\
Méier & 63 & 1.080 & 2,7 & 1.412 & 3,6 \\
Benfica & 12 & 1.055 & 5,0 & 1.333 & 6,4 \\
Vicente de Carvalho & 73 & 1.042 & 4,9 & 1.261 & 5,9 \\
Pilares & 71 & 1.036 & 4,4 & 1.245 & 5,3 \\
Quintino Bocaiúva & 79 & 1.021 & 4,1 & 1.273 & 5,1 \\
Del Castilho & 53 & 991 & 4,8 & 1.071 & 5,2 \\
Magalhães Bastos & 138 & 968 & 4,6 & 1.113 & 5,3 \\
Jardim Guanabara & 99 & 967 & 3,6 & 1.100 & 4,1 \\
Parque Anchieta & 108 & 962 & 4,4 & 1.211 & 5,5 \\
Honório Gurgel & 87 & 933 & 5,1 & 1.163 & 6,4 \\
Jardim América & 49 & 908 & 4,4 & 1.093 & 5,3 \\
Rocha & 58 & 885 & 5,6 & 1.008 & 6,3 \\
Vidigal & 30 & 853 & 5,6 & 995 & 6,6 \\
Portuguesa & 103 & 839 & 4,2 & 983 & 5,0 \\
Maracanã & 35 & 837 & 3,6 & 877 & 3,7 \\
Catete & 18 & 830 & 3,7 & 948 & 4,3 \\
Freguesia (Ilha) & 98 & 828 & 5,2 & 935 & 5,9 \\
Vila Militar & 135 & 798 & 6,2 & 791 & 6,1 \\
Jardim Carioca & 100 & 790 & 3,9 & 1.094 & 5,4 \\
Lagoa & 27 & 772 & 4,1 & 816 & 4,4 \\
Vila da Penha & 74 & 754 & 3,2 & 935 & 4,0 \\
Todos os Santos & 64 & 746 & 3,3 & 929 & 4,1 \\
Vasco da Gama & 158 & 745 & 5,1 & 885 & 6,1 \\
Deodoro & 134 & 742 & 7,1 & 752 & 7,2 \\
Jardim Sulacap & 137 & 696 & 4,8 & 779 & 5,4 \\
Bonsucesso & 40 & 672 & 4,0 & 811 & 4,8 \\
Gamboa & 2 & 669 & 5,9 & 752 & 6,6 \\
Catumbi & 6 & 662 & 5,4 & 736 & 6,0 \\
Turiaçú & 85 & 649 & 5,6 & 717 & 6,2 \\
Vista Alegre & 75 & 645 & 3,8 & 765 & 4,5 \\
Centro & 5 & 622 & 2,6 & 599 & 2,5 \\
Bancários & 97 & 600 & 5,4 & 619 & 5,6 \\
Pedra de Guaratiba & 153 & 599 & 5,0 & 681 & 5,7 \\
Cavalcanti & 80 & 597 & 4,3 & 768 & 5,6 \\
Jardim Botânico & 28 & 587 & 3,7 & 695 & 4,4 \\
Vaz Lobo & 84 & 582 & 5,3 & 609 & 5,5 \\
Vila Kosmos & 72 & 544 & 3,7 & 749 & 5,0 \\
Santo Cristo & 3 & 534 & 5,4 & 567 & 5,7 \\
Pitangueiras & 94 & 503 & 4,9 & 666 & 6,5 \\
Higienópolis & 50 & 503 & 3,5 & 669 & 4,7 \\
Parque Colúmbia & 159 & 483 & 5,9 & 589 & 7,2 \\
Gávea & 29 & 459 & 3,2 & 664 & 4,7 \\
Humaitá & 21 & 454 & 3,5 & 536 & 4,1 \\
Sampaio & 60 & 453 & 5,1 & 561 & 6,4 \\
Jacaré & 51 & 443 & 5,4 & 544 & 6,6 \\
Encantado & 68 & 442 & 3,6 & 614 & 4,9 \\
Cacuia & 93 & 439 & 4,2 & 554 & 5,3 \\
Leme & 23 & 422 & 3,4 & 456 & 3,7 \\
Alto da Boa Vista & 34 & 418 & 4,3 & 497 & 5,1 \\
Campinho & 78 & 392 & 4,6 & 408 & 4,8 \\
São Conrado & 31 & 337 & 3,5 & 434 & 4,5 \\
Riachuelo & 59 & 327 & 3,1 & 414 & 4,0 \\
Cosme Velho & 19 & 325 & 5,1 & 391 & 6,1 \\
Abolição & 70 & 313 & 3,4 & 363 & 3,9 \\
Ilha de Guaratiba & 164 & 309 & 5,1 & 357 & 5,9 \\
Água Santa & 67 & 280 & 4,4 & 364 & 5,7 \\
Maria da Graça & 52 & 277 & 3,8 & 284 & 3,9 \\
Praça da Bandeira & 32 & 261 & 3,6 & 263 & 3,6 \\
São Francisco Xavier & 57 & 255 & 4,0 & 312 & 4,9 \\
Urca & 22 & 252 & 4,8 & 289 & 5,5 \\
Engenheiro Leal & 81 & 245 & 4,7 & 395 & 7,6 \\
Cidade Nova & 8 & 211 & 4,4 & 252 & 5,3 \\
Barra de Guaratiba & 152 & 208 & 4,4 & 247 & 5,2 \\
Lapa & 161 & 206 & 2,3 & 246 & 2,7 \\
Praia da Bandeira & 95 & 205 & 4,0 & 252 & 4,9 \\
Moneró & 102 & 181 & 3,2 & 255 & 4,5 \\
Gericinó & 160 & 178 & 6,9 & 242 & 9,4 \\
Campo dos Afonsos & 136 & 178 & 8,5 & 221 & 10,5 \\
Cocotá & 96 & 174 & 3,5 & 254 & 5,1 \\
Jabour & 163 & 173 & 3,3 & 277 & 5,2 \\
Glória & 16 & 169 & 2,4 & 168 & 2,4 \\
Paquetá & 13 & 145 & 4,2 & 159 & 4,6 \\
Ribeira & 91 & 129 & 3,7 & 126 & 3,6 \\
Camorim & 129 & 83 & 6,5 & 87 & 6,8 \\
Zumbi & 92 & 78 & 4,3 & 81 & 4,5 \\
Cidade Universitária & 105 & 65 & 5,0 & 91 & 7,1 \\
Saúde & 1 & 54 & 3,3 & 71 & 4,3 \\
Argentino & 166 & 33 & 4,4 & 49 & 6,5 \\
Joá & 126 & 22 & 2,2 & 35 & 3,6 \\
Grumari & 133 & 15 & 8,5 & 9 & 5,1 \\
\end{longtable}}
\vspace{-6pt}\fonte{Elaboração IPP com dados de \citeonline{ipp_datario_censo}.}


\begin{landscape}
{\scriptsize\setlength{\tabcolsep}{3pt}
\begin{longtable}{@{}rrrrrrrrrrrr@{}}
\caption{Crianças até 6 anos (número) (populacao ripsa 0 a 6 por ano)}\label{tab:populacao-ripsa-0-a-6-por-ano}\\
\toprule \textbf{Ano} & \textbf{Populacao idade 0} & \textbf{Populacao idade 1} & \textbf{Populacao idade 2} & \textbf{Populacao idade 3} & \textbf{Populacao idade 4} & \textbf{Populacao idade 5} & \textbf{Populacao idade 6} & \textbf{Populacao 0 a 5} & \textbf{Populacao 0 a 6} & \textbf{Populacao total} & \textbf{0 a 6 (\%)} \\ \midrule \endfirsthead
\multicolumn{12}{@{}l}{\footnotesize\itshape (continuação)}\\ \toprule \textbf{Ano} & \textbf{Populacao idade 0} & \textbf{Populacao idade 1} & \textbf{Populacao idade 2} & \textbf{Populacao idade 3} & \textbf{Populacao idade 4} & \textbf{Populacao idade 5} & \textbf{Populacao idade 6} & \textbf{Populacao 0 a 5} & \textbf{Populacao 0 a 6} & \textbf{Populacao total} & \textbf{0 a 6 (\%)} \\ \midrule \endhead
\midrule \multicolumn{12}{r@{}}{\footnotesize\itshape (continua)}\\ \endfoot
\bottomrule \endlastfoot
2000 & 97.858 & 99.144 & 99.542 & 100.078 & 100.757 & 100.943 & 100.473 & 598.322 & 698.795 & 6.312.372 & 11,1 \\
2001 & 93.076 & 96.770 & 98.600 & 99.355 & 100.107 & 100.922 & 101.126 & 588.830 & 689.956 & 6.367.184 & 10,8 \\
2002 & 88.416 & 92.101 & 96.058 & 98.256 & 99.266 & 100.190 & 101.060 & 574.287 & 675.347 & 6.410.594 & 10,5 \\
2003 & 84.544 & 87.654 & 91.490 & 95.804 & 98.242 & 99.391 & 100.329 & 557.125 & 657.454 & 6.445.241 & 10,2 \\
2004 & 81.702 & 84.146 & 87.323 & 91.417 & 95.902 & 98.415 & 99.535 & 538.905 & 638.440 & 6.475.300 & 9,9 \\
2005 & 80.191 & 81.830 & 84.120 & 87.479 & 91.662 & 96.146 & 98.552 & 521.428 & 619.980 & 6.503.064 & 9,5 \\
2006 & 79.427 & 80.739 & 82.189 & 84.521 & 87.850 & 91.940 & 96.255 & 506.666 & 602.921 & 6.529.151 & 9,2 \\
2007 & 79.055 & 80.104 & 81.297 & 82.701 & 84.928 & 88.089 & 91.990 & 496.174 & 588.164 & 6.552.265 & 9,0 \\
2008 & 79.270 & 79.570 & 80.693 & 81.786 & 83.032 & 85.079 & 88.049 & 489.430 & 577.479 & 6.571.817 & 8,8 \\
2009 & 79.770 & 79.602 & 80.003 & 81.051 & 82.016 & 83.103 & 84.980 & 485.545 & 570.525 & 6.590.415 & 8,7 \\
2010 & 79.749 & 79.756 & 79.779 & 80.190 & 81.162 & 82.034 & 83.018 & 482.670 & 565.688 & 6.605.732 & 8,6 \\
2011 & 80.453 & 79.413 & 79.650 & 79.768 & 80.201 & 81.163 & 81.986 & 480.648 & 562.634 & 6.622.327 & 8,5 \\
2012 & 81.607 & 80.135 & 79.249 & 79.593 & 79.763 & 80.195 & 81.147 & 480.542 & 561.689 & 6.645.964 & 8,5 \\
2013 & 81.963 & 81.398 & 79.996 & 79.184 & 79.556 & 79.737 & 80.159 & 481.834 & 561.993 & 6.672.458 & 8,4 \\
2014 & 83.640 & 81.832 & 81.318 & 79.959 & 79.170 & 79.532 & 79.677 & 485.451 & 565.128 & 6.699.175 & 8,4 \\
2015 & 85.890 & 83.539 & 81.781 & 81.336 & 79.975 & 79.160 & 79.472 & 491.681 & 571.153 & 6.724.719 & 8,5 \\
2016 & 82.906 & 85.790 & 83.515 & 81.804 & 81.355 & 79.943 & 79.086 & 495.313 & 574.399 & 6.743.995 & 8,5 \\
2017 & 79.663 & 82.757 & 85.737 & 83.485 & 81.751 & 81.252 & 79.782 & 494.645 & 574.427 & 6.756.761 & 8,5 \\
2018 & 79.399 & 79.441 & 82.612 & 85.608 & 83.304 & 81.504 & 80.928 & 491.868 & 572.796 & 6.765.565 & 8,5 \\
2019 & 76.400 & 79.248 & 79.327 & 82.470 & 85.384 & 82.976 & 81.099 & 485.805 & 566.904 & 6.772.146 & 8,4 \\
2020 & 72.144 & 76.403 & 79.235 & 79.224 & 82.262 & 85.054 & 82.563 & 474.322 & 556.885 & 6.770.926 & 8,2 \\
2021 & 68.341 & 72.314 & 76.477 & 79.167 & 79.031 & 81.950 & 84.668 & 457.280 & 541.948 & 6.757.411 & 8,0 \\
2022 & 64.701 & 68.604 & 72.424 & 76.444 & 78.990 & 78.744 & 81.625 & 439.907 & 521.532 & 6.742.618 & 7,7 \\
2023 & 62.431 & 65.000 & 68.791 & 72.432 & 76.309 & 78.776 & 78.523 & 423.739 & 502.262 & 6.735.026 & 7,5 \\
2024 & 61.640 & 63.039 & 65.433 & 68.926 & 72.364 & 76.135 & 78.568 & 407.537 & 486.105 & 6.729.894 & 7,2 \\
2025 & 61.200 & 62.082 & 63.219 & 65.435 & 68.855 & 72.282 & 76.075 & 393.073 & 469.148 & 6.730.729 & 7,0 \\
\end{longtable}}
\vspace{-6pt}\fonte{Elaboração IPP com dados de \citeonline{ms_ripsa_populacao}.}
\end{landscape}


{\small\setlength{\tabcolsep}{5pt}
\begin{longtable}{@{}rr@{}}
\caption{Nascidos vivos por bairro de residência da mãe (número) (nascidos vivos por ano)}\label{tab:nascidos-vivos-por-ano}\\
\toprule \textbf{Ano} & \textbf{Nascidos vivos} \\ \midrule \endfirsthead
\multicolumn{2}{@{}l}{\footnotesize\itshape (continuação)}\\ \toprule \textbf{Ano} & \textbf{Nascidos vivos} \\ \midrule \endhead
\midrule \multicolumn{2}{r@{}}{\footnotesize\itshape (continua)}\\ \endfoot
\bottomrule \endlastfoot
2006 & 92.631 \\
2007 & 91.774 \\
2008 & 92.205 \\
2009 & 94.588 \\
2010 & 91.810 \\
2011 & 85.725 \\
2012 & 86.021 \\
2013 & 87.423 \\
2014 & 89.760 \\
2015 & 90.303 \\
2016 & 82.854 \\
2017 & 84.339 \\
2018 & 82.484 \\
2019 & 76.477 \\
2020 & 72.414 \\
2021 & 68.412 \\
2022 & 72.488 \\
2023 & 70.525 \\
2024 & 64.282 \\
2025 & 65.507 \\
\end{longtable}}
\vspace{-6pt}\fonte{Elaboração IPP com dados de \citeonline{sms_rio_sinasc}.}


{\small\setlength{\tabcolsep}{5pt}
\begin{longtable}{@{}lrrrr@{}}
\caption{Nascidos vivos por bairro de residência da mãe (número) (nascidos vivos 2025)}\label{tab:tabela-mapa-nascidos-vivos-2025}\\
\toprule \textbf{Bairro} & \textbf{Ano} & \textbf{Nascidos vivos} & \textbf{Codigo} & \textbf{Do municipio (\%)} \\ \midrule \endfirsthead
\multicolumn{5}{@{}l}{\footnotesize\itshape (continuação)}\\ \toprule \textbf{Bairro} & \textbf{Ano} & \textbf{Nascidos vivos} & \textbf{Codigo} & \textbf{Do municipio (\%)} \\ \midrule \endhead
\midrule \multicolumn{5}{r@{}}{\footnotesize\itshape (continua)}\\ \endfoot
\bottomrule \endlastfoot
SAUDE & 2025 & 36 & 1 & 0,1 \\
GAMBOA & 2025 & 150 & 2 & 0,2 \\
SANTO CRISTO & 2025 & 129 & 3 & 0,2 \\
CAJU & 2025 & 305 & 4 & 0,5 \\
CENTRO & 2025 & 194 & 5 & 0,3 \\
CATUMBI & 2025 & 134 & 6 & 0,2 \\
RIO COMPRIDO & 2025 & 349 & 7 & 0,5 \\
CIDADE NOVA & 2025 & 42 & 8 & 0,1 \\
ESTACIO & 2025 & 153 & 9 & 0,2 \\
SAO CRISTOVAO & 2025 & 290 & 10 & 0,4 \\
MANGUEIRA & 2025 & 199 & 11 & 0,3 \\
BENFICA & 2025 & 346 & 12 & 0,5 \\
PAQUETA & 2025 & 30 & 13 & 0,0 \\
SANTA TERESA & 2025 & 316 & 14 & 0,5 \\
FLAMENGO & 2025 & 274 & 15 & 0,4 \\
GLORIA & 2025 & 47 & 16 & 0,1 \\
LARANJEIRAS & 2025 & 206 & 17 & 0,3 \\
CATETE & 2025 & 145 & 18 & 0,2 \\
COSME VELHO & 2025 & 69 & 19 & 0,1 \\
BOTAFOGO & 2025 & 534 & 20 & 0,8 \\
HUMAITA & 2025 & 69 & 21 & 0,1 \\
URCA & 2025 & 57 & 22 & 0,1 \\
LEME & 2025 & 94 & 23 & 0,1 \\
COPACABANA & 2025 & 772 & 24 & 1,2 \\
IPANEMA & 2025 & 247 & 25 & 0,4 \\
LEBLON & 2025 & 275 & 26 & 0,4 \\
LAGOA & 2025 & 123 & 27 & 0,2 \\
JARDIM BOTANICO & 2025 & 86 & 28 & 0,1 \\
GAVEA & 2025 & 139 & 29 & 0,2 \\
VIDIGAL & 2025 & 145 & 30 & 0,2 \\
SAO CONRADO & 2025 & 85 & 31 & 0,1 \\
PRACA DA BANDEIRA & 2025 & 57 & 32 & 0,1 \\
TIJUCA & 2025 & 941 & 33 & 1,4 \\
ALTO DA BOA VISTA & 2025 & 67 & 34 & 0,1 \\
MARACANA & 2025 & 131 & 35 & 0,2 \\
VILA ISABEL & 2025 & 443 & 36 & 0,7 \\
ANDARAI & 2025 & 209 & 37 & 0,3 \\
GRAJAU & 2025 & 168 & 38 & 0,3 \\
MANGUINHOS & 2025 & 387 & 39 & 0,6 \\
BONSUCESSO & 2025 & 392 & 40 & 0,6 \\
RAMOS & 2025 & 479 & 41 & 0,7 \\
OLARIA & 2025 & 416 & 42 & 0,6 \\
PENHA & 2025 & 609 & 43 & 0,9 \\
PENHA CIRCULAR & 2025 & 356 & 44 & 0,5 \\
BRAS DE PINA & 2025 & 364 & 45 & 0,6 \\
CORDOVIL & 2025 & 316 & 46 & 0,5 \\
PARADA DE LUCAS & 2025 & 255 & 47 & 0,4 \\
VIGARIO GERAL & 2025 & 354 & 48 & 0,5 \\
JARDIM AMERICA & 2025 & 266 & 49 & 0,4 \\
HIGIENOPOLIS & 2025 & 118 & 50 & 0,2 \\
JACARE & 2025 & 166 & 51 & 0,3 \\
MARIA DA GRACA & 2025 & 100 & 52 & 0,2 \\
DEL CASTILHO & 2025 & 236 & 53 & 0,4 \\
INHAUMA & 2025 & 539 & 54 & 0,8 \\
ENGENHO DA RAINHA & 2025 & 217 & 55 & 0,3 \\
TOMAS COELHO & 2025 & 191 & 56 & 0,3 \\
SAO FRANCISCO XAVIER & 2025 & 55 & 57 & 0,1 \\
ROCHA & 2025 & 178 & 58 & 0,3 \\
RIACHUELO & 2025 & 66 & 59 & 0,1 \\
SAMPAIO & 2025 & 99 & 60 & 0,2 \\
ENGENHO NOVO & 2025 & 346 & 61 & 0,5 \\
LINS DE VASCONCELOS & 2025 & 304 & 62 & 0,5 \\
MEIER & 2025 & 247 & 63 & 0,4 \\
TODOS OS SANTOS & 2025 & 126 & 64 & 0,2 \\
CACHAMBI & 2025 & 274 & 65 & 0,4 \\
ENGENHO DE DENTRO & 2025 & 351 & 66 & 0,5 \\
AGUA SANTA & 2025 & 57 & 67 & 0,1 \\
ENCANTADO & 2025 & 81 & 68 & 0,1 \\
PIEDADE & 2025 & 328 & 69 & 0,5 \\
ABOLICAO & 2025 & 109 & 70 & 0,2 \\
PILARES & 2025 & 191 & 71 & 0,3 \\
VILA KOSMOS & 2025 & 109 & 72 & 0,2 \\
VICENTE DE CARVALHO & 2025 & 189 & 73 & 0,3 \\
VILA DA PENHA & 2025 & 148 & 74 & 0,2 \\
VISTA ALEGRE & 2025 & 50 & 75 & 0,1 \\
IRAJA & 2025 & 683 & 76 & 1,0 \\
COLEGIO & 2025 & 292 & 77 & 0,4 \\
CAMPINHO & 2025 & 95 & 78 & 0,1 \\
QUINTINO BOCAIUVA & 2025 & 202 & 79 & 0,3 \\
CAVALCANTE & 2025 & 145 & 80 & 0,2 \\
ENGENHEIRO LEAL & 2025 & 51 & 81 & 0,1 \\
CASCADURA & 2025 & 243 & 82 & 0,4 \\
MADUREIRA & 2025 & 378 & 83 & 0,6 \\
VAZ LOBO & 2025 & 139 & 84 & 0,2 \\
TURIACU & 2025 & 136 & 85 & 0,2 \\
ROCHA MIRANDA & 2025 & 324 & 86 & 0,5 \\
HONORIO GURGEL & 2025 & 183 & 87 & 0,3 \\
OSWALDO CRUZ & 2025 & 196 & 88 & 0,3 \\
BENTO RIBEIRO & 2025 & 300 & 89 & 0,5 \\
MARECHAL HERMES & 2025 & 334 & 90 & 0,5 \\
RIBEIRA & 2025 & 23 & 91 & 0,0 \\
ZUMBI & 2025 & 20 & 92 & 0,0 \\
CACUIA & 2025 & 85 & 93 & 0,1 \\
PITANGUEIRAS & 2025 & 76 & 94 & 0,1 \\
PRAIA DA BANDEIRA & 2025 & 37 & 95 & 0,1 \\
COCOTA & 2025 & 53 & 96 & 0,1 \\
BANCARIOS & 2025 & 132 & 97 & 0,2 \\
FREGUESIA-ILHA & 2025 & 191 & 98 & 0,3 \\
JARDIM GUANABARA & 2025 & 170 & 99 & 0,3 \\
JARDIM CARIOCA & 2025 & 199 & 100 & 0,3 \\
TAUA & 2025 & 212 & 101 & 0,3 \\
MONERO & 2025 & 39 & 102 & 0,1 \\
PORTUGUESA & 2025 & 127 & 103 & 0,2 \\
GALEAO & 2025 & 248 & 104 & 0,4 \\
CIDADE UNIVERSITARIA & 2025 & 17 & 105 & 0,0 \\
GUADALUPE & 2025 & 355 & 106 & 0,5 \\
ANCHIETA & 2025 & 618 & 107 & 0,9 \\
PARQUE ANCHIETA & 2025 & 144 & 108 & 0,2 \\
RICARDO ALBUQUERQUE & 2025 & 262 & 109 & 0,4 \\
COELHO NETO & 2025 & 284 & 110 & 0,4 \\
ACARI & 2025 & 349 & 111 & 0,5 \\
BARROS FILHO & 2025 & 164 & 112 & 0,3 \\
COSTA BARROS & 2025 & 401 & 113 & 0,6 \\
PAVUNA & 2025 & 888 & 114 & 1,4 \\
JACAREPAGUA & 2025 & 1.878 & 115 & 2,9 \\
ANIL & 2025 & 377 & 116 & 0,6 \\
GARDENIA AZUL & 2025 & 400 & 117 & 0,6 \\
CIDADE DE DEUS & 2025 & 549 & 118 & 0,8 \\
CURICICA & 2025 & 369 & 119 & 0,6 \\
FREGUESIA-JPA & 2025 & 553 & 120 & 0,8 \\
PECHINCHA & 2025 & 339 & 121 & 0,5 \\
TAQUARA & 2025 & 1.005 & 122 & 1,5 \\
TANQUE & 2025 & 378 & 123 & 0,6 \\
PRACA SECA & 2025 & 567 & 124 & 0,9 \\
VILA VALQUEIRE & 2025 & 273 & 125 & 0,4 \\
JOA & 2025 & 8 & 126 & 0,0 \\
ITANHANGA & 2025 & 791 & 127 & 1,2 \\
BARRA DA TIJUCA & 2025 & 1.115 & 128 & 1,7 \\
CAMORIM & 2025 & 35 & 129 & 0,1 \\
VARGEM PEQUENA & 2025 & 325 & 130 & 0,5 \\
VARGEM GRANDE & 2025 & 255 & 131 & 0,4 \\
RECREIO BANDEIRANTES & 2025 & 1.245 & 132 & 1,9 \\
GRUMARI & 2025 & 0 & 133 & 0,0 \\
DEODORO & 2025 & 124 & 134 & 0,2 \\
VILA MILITAR & 2025 & 108 & 135 & 0,2 \\
CAMPO DOS AFONSOS & 2025 & 36 & 136 & 0,1 \\
JARDIM SULACAP & 2025 & 132 & 137 & 0,2 \\
MAGALHAES BASTOS & 2025 & 208 & 138 & 0,3 \\
REALENGO & 2025 & 1.463 & 139 & 2,2 \\
PADRE MIGUEL & 2025 & 522 & 140 & 0,8 \\
BANGU & 2025 & 1.977 & 141 & 3,0 \\
SENADOR CAMARA & 2025 & 995 & 142 & 1,5 \\
SANTISSIMO & 2025 & 471 & 143 & 0,7 \\
CAMPO GRANDE & 2025 & 3.504 & 144 & 5,3 \\
SENADOR VASCONCELOS & 2025 & 298 & 145 & 0,5 \\
INHOAIBA & 2025 & 752 & 146 & 1,1 \\
COSMOS & 2025 & 954 & 147 & 1,5 \\
PACIENCIA & 2025 & 1.157 & 148 & 1,8 \\
SANTA CRUZ & 2025 & 3.186 & 149 & 4,9 \\
SEPETIBA & 2025 & 753 & 150 & 1,1 \\
GUARATIBA & 2025 & 1.867 & 151 & 2,9 \\
BARRA DE GUARATIBA & 2025 & 117 & 152 & 0,2 \\
PEDRA DE GUARATIBA & 2025 & 228 & 153 & 0,3 \\
ROCINHA & 2025 & 696 & 154 & 1,1 \\
JACAREZINHO & 2025 & 293 & 155 & 0,4 \\
COMPLEXO DO ALEMAO & 2025 & 361 & 156 & 0,6 \\
COMPLEXO DA MARE & 2025 & 1.494 & 157 & 2,3 \\
PARQUE COLUMBIA & 2025 & 83 & 158 & 0,1 \\
VASCO DA GAMA & 2025 & 131 & 159 & 0,2 \\
GERICINo & 2025 & 14 & 160 & 0,0 \\
LAPA & 2025 & 39 & 161 & 0,1 \\
VILA KENNEDY & 2025 & 194 & 162 & 0,3 \\
JABOUR & 2025 & 17 & 163 & 0,0 \\
ILHA DE GUARATIBA & 2025 & 81 & 164 & 0,1 \\
BARRA OLIMPICA & 2025 & 503 & 165 & 0,8 \\
ILHA DO GOVERNADOR - IGNORADO & 2025 & 12 & 998 & 0,0 \\
IGNORADO BAIRRO-MRJ & 2025 & 19 & 999 & 0,0 \\
\end{longtable}}
\vspace{-6pt}\fonte{Elaboração IPP com dados de \citeonline{sms_rio_sinasc}.}


{\small\setlength{\tabcolsep}{5pt}
\begin{longtable}{@{}rrrr@{}}
\caption{Taxa de mortalidade neonatal precoce (0 a 6 dias) (mortalidade neonatal precoce por ano)}\label{tab:mortalidade-neonatal-precoce-por-ano}\\
\toprule \textbf{Ano} & \textbf{Obitos precoces} & \textbf{Nascidos vivos} & \textbf{Taxa mortalidade precoce} \\ \midrule \endfirsthead
\multicolumn{4}{@{}l}{\footnotesize\itshape (continuação)}\\ \toprule \textbf{Ano} & \textbf{Obitos precoces} & \textbf{Nascidos vivos} & \textbf{Taxa mortalidade precoce} \\ \midrule \endhead
\midrule \multicolumn{4}{r@{}}{\footnotesize\itshape (continua)}\\ \endfoot
\bottomrule \endlastfoot
2006 & 728 & 92.584 & 7,9 \\
2007 & 676 & 91.743 & 7,4 \\
2008 & 661 & 92.144 & 7,2 \\
2009 & 723 & 94.538 & 7,6 \\
2010 & 687 & 91.774 & 7,5 \\
2011 & 635 & 85.682 & 7,4 \\
2012 & 646 & 85.994 & 7,5 \\
2013 & 648 & 87.389 & 7,4 \\
2014 & 622 & 89.731 & 6,9 \\
2015 & 640 & 90.272 & 7,1 \\
2016 & 638 & 82.826 & 7,7 \\
2017 & 576 & 84.323 & 6,8 \\
2018 & 554 & 82.462 & 6,7 \\
2019 & 533 & 76.459 & 7,0 \\
2020 & 522 & 72.393 & 7,2 \\
2021 & 476 & 68.384 & 7,0 \\
2022 & 449 & 72.460 & 6,2 \\
2023 & 419 & 70.503 & 5,9 \\
2024 & 389 & 64.245 & 6,1 \\
2025 & 420 & 65.470 & 6,4 \\
\end{longtable}}
\vspace{-6pt}\fonte{Elaboração IPP com dados de \citeonline{sms_rio_sim}; \citeonline{sms_rio_sinasc}.}


{\small\setlength{\tabcolsep}{5pt}
\begin{longtable}{@{}lrrrrr@{}}
\caption{Taxa de mortalidade neonatal precoce (0 a 6 dias) (obitos neonatal precoce 2025)}\label{tab:tabela-mapa-obitos-neonatal-precoce-2025}\\
\toprule \textbf{Bairro} & \textbf{Ano} & \textbf{Obitos precoces} & \textbf{Codigo} & \textbf{Nascidos vivos} & \textbf{Taxa mortalidade precoce} \\ \midrule \endfirsthead
\multicolumn{6}{@{}l}{\footnotesize\itshape (continuação)}\\ \toprule \textbf{Bairro} & \textbf{Ano} & \textbf{Obitos precoces} & \textbf{Codigo} & \textbf{Nascidos vivos} & \textbf{Taxa mortalidade precoce} \\ \midrule \endhead
\midrule \multicolumn{6}{r@{}}{\footnotesize\itshape (continua)}\\ \endfoot
\bottomrule \endlastfoot
SAUDE & 2025 & 0 & 1 & 36 & 0,0 \\
GAMBOA & 2025 & 2 & 2 & 150 & 13,3 \\
SANTO CRISTO & 2025 & 1 & 3 & 129 & 7,8 \\
CAJU & 2025 & 3 & 4 & 305 & 9,8 \\
CENTRO & 2025 & 1 & 5 & 194 & 5,2 \\
CATUMBI & 2025 & 4 & 6 & 134 & 29,9 \\
RIO COMPRIDO & 2025 & 2 & 7 & 349 & 5,7 \\
CIDADE NOVA & 2025 & 0 & 8 & 42 & 0,0 \\
ESTACIO & 2025 & 3 & 9 & 153 & 19,6 \\
SAO CRISTOVAO & 2025 & 0 & 10 & 290 & 0,0 \\
MANGUEIRA & 2025 & 2 & 11 & 199 & 10,1 \\
BENFICA & 2025 & 3 & 12 & 346 & 8,7 \\
PAQUETA & 2025 & 0 & 13 & 30 & 0,0 \\
SANTA TERESA & 2025 & 2 & 14 & 316 & 6,3 \\
FLAMENGO & 2025 & 3 & 15 & 274 & 10,9 \\
GLORIA & 2025 & 0 & 16 & 47 & 0,0 \\
LARANJEIRAS & 2025 & 0 & 17 & 206 & 0,0 \\
CATETE & 2025 & 0 & 18 & 145 & 0,0 \\
COSME VELHO & 2025 & 1 & 19 & 69 & 14,5 \\
BOTAFOGO & 2025 & 0 & 20 & 534 & 0,0 \\
HUMAITA & 2025 & 0 & 21 & 69 & 0,0 \\
URCA & 2025 & 0 & 22 & 57 & 0,0 \\
LEME & 2025 & 1 & 23 & 94 & 10,6 \\
COPACABANA & 2025 & 3 & 24 & 772 & 3,9 \\
IPANEMA & 2025 & 5 & 25 & 247 & 20,2 \\
LEBLON & 2025 & 0 & 26 & 275 & 0,0 \\
LAGOA & 2025 & 1 & 27 & 123 & 8,1 \\
JARDIM BOTANICO & 2025 & 0 & 28 & 86 & 0,0 \\
GAVEA & 2025 & 1 & 29 & 139 & 7,2 \\
VIDIGAL & 2025 & 1 & 30 & 145 & 6,9 \\
SAO CONRADO & 2025 & 0 & 31 & 85 & 0,0 \\
PRACA DA BANDEIRA & 2025 & 0 & 32 & 57 & 0,0 \\
TIJUCA & 2025 & 3 & 33 & 941 & 3,2 \\
ALTO DA BOA VISTA & 2025 & 1 & 34 & 67 & 14,9 \\
MARACANA & 2025 & 0 & 35 & 131 & 0,0 \\
VILA ISABEL & 2025 & 2 & 36 & 443 & 4,5 \\
ANDARAI & 2025 & 1 & 37 & 209 & 4,8 \\
GRAJAU & 2025 & 0 & 38 & 168 & 0,0 \\
MANGUINHOS & 2025 & 4 & 39 & 387 & 10,3 \\
BONSUCESSO & 2025 & 0 & 40 & 392 & 0,0 \\
RAMOS & 2025 & 3 & 41 & 479 & 6,3 \\
OLARIA & 2025 & 3 & 42 & 416 & 7,2 \\
PENHA & 2025 & 4 & 43 & 609 & 6,6 \\
PENHA CIRCULAR & 2025 & 4 & 44 & 356 & 11,2 \\
BRAS DE PINA & 2025 & 3 & 45 & 364 & 8,2 \\
CORDOVIL & 2025 & 5 & 46 & 316 & 15,8 \\
PARADA DE LUCAS & 2025 & 1 & 47 & 255 & 3,9 \\
VIGARIO GERAL & 2025 & 3 & 48 & 354 & 8,5 \\
JARDIM AMERICA & 2025 & 1 & 49 & 266 & 3,8 \\
HIGIENOPOLIS & 2025 & 0 & 50 & 118 & 0,0 \\
JACARE & 2025 & 0 & 51 & 166 & 0,0 \\
MARIA DA GRACA & 2025 & 0 & 52 & 100 & 0,0 \\
DEL CASTILHO & 2025 & 2 & 53 & 236 & 8,5 \\
INHAUMA & 2025 & 3 & 54 & 539 & 5,6 \\
ENGENHO DA RAINHA & 2025 & 2 & 55 & 217 & 9,2 \\
TOMAS COELHO & 2025 & 5 & 56 & 191 & 26,2 \\
SAO FRANCISCO XAVIER & 2025 & 1 & 57 & 55 & 18,2 \\
ROCHA & 2025 & 3 & 58 & 178 & 16,9 \\
RIACHUELO & 2025 & 0 & 59 & 66 & 0,0 \\
SAMPAIO & 2025 & 2 & 60 & 99 & 20,2 \\
ENGENHO NOVO & 2025 & 1 & 61 & 346 & 2,9 \\
LINS DE VASCONCELOS & 2025 & 2 & 62 & 304 & 6,6 \\
MEIER & 2025 & 2 & 63 & 247 & 8,1 \\
TODOS OS SANTOS & 2025 & 1 & 64 & 126 & 7,9 \\
CACHAMBI & 2025 & 0 & 65 & 274 & 0,0 \\
ENGENHO DE DENTRO & 2025 & 0 & 66 & 351 & 0,0 \\
AGUA SANTA & 2025 & 1 & 67 & 57 & 17,5 \\
ENCANTADO & 2025 & 1 & 68 & 81 & 12,3 \\
PIEDADE & 2025 & 2 & 69 & 328 & 6,1 \\
ABOLICAO & 2025 & 1 & 70 & 109 & 9,2 \\
PILARES & 2025 & 1 & 71 & 191 & 5,2 \\
VILA KOSMOS & 2025 & 0 & 72 & 109 & 0,0 \\
VICENTE DE CARVALHO & 2025 & 1 & 73 & 189 & 5,3 \\
VILA DA PENHA & 2025 & 0 & 74 & 148 & 0,0 \\
VISTA ALEGRE & 2025 & 0 & 75 & 50 & 0,0 \\
IRAJA & 2025 & 3 & 76 & 683 & 4,4 \\
COLEGIO & 2025 & 4 & 77 & 292 & 13,7 \\
CAMPINHO & 2025 & 1 & 78 & 95 & 10,5 \\
QUINTINO BOCAIUVA & 2025 & 0 & 79 & 202 & 0,0 \\
CAVALCANTE & 2025 & 1 & 80 & 145 & 6,9 \\
ENGENHEIRO LEAL & 2025 & 1 & 81 & 51 & 19,6 \\
CASCADURA & 2025 & 3 & 82 & 243 & 12,3 \\
MADUREIRA & 2025 & 1 & 83 & 378 & 2,6 \\
VAZ LOBO & 2025 & 0 & 84 & 139 & 0,0 \\
TURIACU & 2025 & 0 & 85 & 136 & 0,0 \\
ROCHA MIRANDA & 2025 & 3 & 86 & 324 & 9,3 \\
HONORIO GURGEL & 2025 & 3 & 87 & 183 & 16,4 \\
OSWALDO CRUZ & 2025 & 0 & 88 & 196 & 0,0 \\
BENTO RIBEIRO & 2025 & 4 & 89 & 300 & 13,3 \\
MARECHAL HERMES & 2025 & 4 & 90 & 334 & 12,0 \\
RIBEIRA & 2025 & 1 & 91 & 23 & 43,5 \\
CACUIA & 2025 & 0 & 93 & 85 & 0,0 \\
PITANGUEIRAS & 2025 & 0 & 94 & 76 & 0,0 \\
PRAIA DA BANDEIRA & 2025 & 0 & 95 & 37 & 0,0 \\
COCOTA & 2025 & 0 & 96 & 53 & 0,0 \\
BANCARIOS & 2025 & 0 & 97 & 132 & 0,0 \\
FREGUESIA-ILHA & 2025 & 0 & 98 & 191 & 0,0 \\
JARDIM GUANABARA & 2025 & 0 & 99 & 170 & 0,0 \\
JARDIM CARIOCA & 2025 & 2 & 100 & 199 & 10,1 \\
TAUA & 2025 & 0 & 101 & 212 & 0,0 \\
MONERO & 2025 & 0 & 102 & 39 & 0,0 \\
PORTUGUESA & 2025 & 2 & 103 & 127 & 15,7 \\
GALEAO & 2025 & 0 & 104 & 248 & 0,0 \\
CIDADE UNIVERSITARIA & 2025 & 1 & 105 & 17 & 58,8 \\
GUADALUPE & 2025 & 2 & 106 & 355 & 5,6 \\
ANCHIETA & 2025 & 5 & 107 & 618 & 8,1 \\
PARQUE ANCHIETA & 2025 & 1 & 108 & 144 & 6,9 \\
RICARDO ALBUQUERQUE & 2025 & 1 & 109 & 262 & 3,8 \\
COELHO NETO & 2025 & 0 & 110 & 284 & 0,0 \\
ACARI & 2025 & 2 & 111 & 349 & 5,7 \\
BARROS FILHO & 2025 & 1 & 112 & 164 & 6,1 \\
COSTA BARROS & 2025 & 4 & 113 & 401 & 10,0 \\
PAVUNA & 2025 & 14 & 114 & 888 & 15,8 \\
JACAREPAGUA & 2025 & 11 & 115 & 1.878 & 5,9 \\
ANIL & 2025 & 5 & 116 & 377 & 13,3 \\
GARDENIA AZUL & 2025 & 1 & 117 & 400 & 2,5 \\
CIDADE DE DEUS & 2025 & 2 & 118 & 549 & 3,6 \\
CURICICA & 2025 & 4 & 119 & 369 & 10,8 \\
FREGUESIA-JPA & 2025 & 2 & 120 & 553 & 3,6 \\
PECHINCHA & 2025 & 2 & 121 & 339 & 5,9 \\
TAQUARA & 2025 & 8 & 122 & 1.005 & 8,0 \\
TANQUE & 2025 & 1 & 123 & 378 & 2,6 \\
PRACA SECA & 2025 & 4 & 124 & 567 & 7,1 \\
VILA VALQUEIRE & 2025 & 2 & 125 & 273 & 7,3 \\
JOA & 2025 & 0 & 126 & 8 & 0,0 \\
ITANHANGA & 2025 & 4 & 127 & 791 & 5,1 \\
BARRA DA TIJUCA & 2025 & 4 & 128 & 1.115 & 3,6 \\
CAMORIM & 2025 & 0 & 129 & 35 & 0,0 \\
VARGEM PEQUENA & 2025 & 4 & 130 & 325 & 12,3 \\
VARGEM GRANDE & 2025 & 1 & 131 & 255 & 3,9 \\
RECREIO BANDEIRANTES & 2025 & 3 & 132 & 1.245 & 2,4 \\
DEODORO & 2025 & 2 & 134 & 124 & 16,1 \\
VILA MILITAR & 2025 & 2 & 135 & 108 & 18,5 \\
CAMPO DOS AFONSOS & 2025 & 2 & 136 & 36 & 55,6 \\
JARDIM SULACAP & 2025 & 1 & 137 & 132 & 7,6 \\
MAGALHAES BASTOS & 2025 & 2 & 138 & 208 & 9,6 \\
REALENGO & 2025 & 5 & 139 & 1.463 & 3,4 \\
PADRE MIGUEL & 2025 & 5 & 140 & 522 & 9,6 \\
BANGU & 2025 & 7 & 141 & 1.977 & 3,5 \\
SENADOR CAMARA & 2025 & 2 & 142 & 995 & 2,0 \\
SANTISSIMO & 2025 & 5 & 143 & 471 & 10,6 \\
CAMPO GRANDE & 2025 & 16 & 144 & 3.504 & 4,6 \\
SENADOR VASCONCELOS & 2025 & 5 & 145 & 298 & 16,8 \\
INHOAIBA & 2025 & 3 & 146 & 752 & 4,0 \\
COSMOS & 2025 & 10 & 147 & 954 & 10,5 \\
PACIENCIA & 2025 & 9 & 148 & 1.157 & 7,8 \\
SANTA CRUZ & 2025 & 28 & 149 & 3.186 & 8,8 \\
SEPETIBA & 2025 & 2 & 150 & 753 & 2,7 \\
GUARATIBA & 2025 & 6 & 151 & 1.867 & 3,2 \\
BARRA DE GUARATIBA & 2025 & 0 & 152 & 117 & 0,0 \\
PEDRA DE GUARATIBA & 2025 & 0 & 153 & 228 & 0,0 \\
ROCINHA & 2025 & 3 & 154 & 696 & 4,3 \\
JACAREZINHO & 2025 & 3 & 155 & 293 & 10,2 \\
COMPLEXO DO ALEMAO & 2025 & 0 & 156 & 361 & 0,0 \\
COMPLEXO DA MARE & 2025 & 6 & 157 & 1.494 & 4,0 \\
PARQUE COLUMBIA & 2025 & 0 & 158 & 83 & 0,0 \\
VASCO DA GAMA & 2025 & 0 & 159 & 131 & 0,0 \\
GERICINo & 2025 & 1 & 160 & 14 & 71,4 \\
LAPA & 2025 & 0 & 161 & 39 & 0,0 \\
VILA KENNEDY & 2025 & 2 & 162 & 194 & 10,3 \\
ILHA DE GUARATIBA & 2025 & 1 & 164 & 81 & 12,3 \\
BARRA OLIMPICA & 2025 & 1 & 165 & 503 & 2,0 \\
ILHA DO GOVERNADOR - IGNORADO & 2025 & 0 & 998 & 12 & 0,0 \\
IGNORADO BAIRRO-MRJ & 2025 & 0 & 999 & 19 & 0,0 \\
\end{longtable}}
\vspace{-6pt}\fonte{Elaboração IPP com dados de \citeonline{sms_rio_sim}; \citeonline{sms_rio_sinasc}.}


{\small\setlength{\tabcolsep}{5pt}
\begin{longtable}{@{}rrrr@{}}
\caption{Taxa de mortalidade neonatal tardia (7 a 27 dias) (mortalidade neonatal tardia por ano)}\label{tab:mortalidade-neonatal-tardia-por-ano}\\
\toprule \textbf{Ano} & \textbf{Obitos tardios} & \textbf{Nascidos vivos} & \textbf{Taxa obitos tardios} \\ \midrule \endfirsthead
\multicolumn{4}{@{}l}{\footnotesize\itshape (continuação)}\\ \toprule \textbf{Ano} & \textbf{Obitos tardios} & \textbf{Nascidos vivos} & \textbf{Taxa obitos tardios} \\ \midrule \endhead
\midrule \multicolumn{4}{r@{}}{\footnotesize\itshape (continua)}\\ \endfoot
\bottomrule \endlastfoot
2006 & 250 & 92.534 & 2,7 \\
2007 & 228 & 91.702 & 2,5 \\
2008 & 292 & 92.107 & 3,2 \\
2009 & 256 & 94.496 & 2,7 \\
2010 & 237 & 91.729 & 2,6 \\
2011 & 230 & 85.642 & 2,7 \\
2012 & 287 & 85.948 & 3,3 \\
2013 & 263 & 87.347 & 3,0 \\
2014 & 254 & 89.686 & 2,8 \\
2015 & 270 & 90.225 & 3,0 \\
2016 & 219 & 82.791 & 2,6 \\
2017 & 232 & 84.271 & 2,8 \\
2018 & 251 & 82.428 & 3,0 \\
2019 & 233 & 76.426 & 3,0 \\
2020 & 250 & 72.357 & 3,5 \\
2021 & 231 & 68.340 & 3,4 \\
2022 & 166 & 72.431 & 2,3 \\
2023 & 203 & 70.451 & 2,9 \\
2024 & 177 & 64.138 & 2,8 \\
2025 & 170 & 65.344 & 2,6 \\
\end{longtable}}
\vspace{-6pt}\fonte{Elaboração IPP com dados de \citeonline{sms_rio_sim}; \citeonline{sms_rio_sinasc}.}


{\small\setlength{\tabcolsep}{5pt}
\begin{longtable}{@{}lrrrrr@{}}
\caption{Taxa de mortalidade neonatal tardia (7 a 27 dias) (obitos neonatal tardia 2025)}\label{tab:tabela-mapa-obitos-neonatal-tardia-2025}\\
\toprule \textbf{Bairro} & \textbf{Ano} & \textbf{Obitos tardios} & \textbf{Codigo} & \textbf{Nascidos vivos} & \textbf{Taxa obitos tardios} \\ \midrule \endfirsthead
\multicolumn{6}{@{}l}{\footnotesize\itshape (continuação)}\\ \toprule \textbf{Bairro} & \textbf{Ano} & \textbf{Obitos tardios} & \textbf{Codigo} & \textbf{Nascidos vivos} & \textbf{Taxa obitos tardios} \\ \midrule \endhead
\midrule \multicolumn{6}{r@{}}{\footnotesize\itshape (continua)}\\ \endfoot
\bottomrule \endlastfoot
SAUDE & 2025 & 0 & 1 & 36 & 0,0 \\
GAMBOA & 2025 & 0 & 2 & 150 & 0,0 \\
SANTO CRISTO & 2025 & 1 & 3 & 129 & 7,8 \\
CAJU & 2025 & 3 & 4 & 305 & 9,8 \\
CENTRO & 2025 & 2 & 5 & 194 & 10,3 \\
CATUMBI & 2025 & 1 & 6 & 134 & 7,5 \\
RIO COMPRIDO & 2025 & 2 & 7 & 349 & 5,7 \\
CIDADE NOVA & 2025 & 1 & 8 & 42 & 23,8 \\
ESTACIO & 2025 & 0 & 9 & 153 & 0,0 \\
SAO CRISTOVAO & 2025 & 0 & 10 & 290 & 0,0 \\
MANGUEIRA & 2025 & 1 & 11 & 199 & 5,0 \\
BENFICA & 2025 & 1 & 12 & 346 & 2,9 \\
PAQUETA & 2025 & 0 & 13 & 30 & 0,0 \\
SANTA TERESA & 2025 & 0 & 14 & 316 & 0,0 \\
FLAMENGO & 2025 & 1 & 15 & 274 & 3,6 \\
GLORIA & 2025 & 0 & 16 & 47 & 0,0 \\
LARANJEIRAS & 2025 & 0 & 17 & 206 & 0,0 \\
CATETE & 2025 & 0 & 18 & 145 & 0,0 \\
COSME VELHO & 2025 & 0 & 19 & 69 & 0,0 \\
BOTAFOGO & 2025 & 0 & 20 & 534 & 0,0 \\
HUMAITA & 2025 & 0 & 21 & 69 & 0,0 \\
URCA & 2025 & 0 & 22 & 57 & 0,0 \\
LEME & 2025 & 0 & 23 & 94 & 0,0 \\
COPACABANA & 2025 & 2 & 24 & 772 & 2,6 \\
IPANEMA & 2025 & 1 & 25 & 247 & 4,0 \\
LEBLON & 2025 & 0 & 26 & 275 & 0,0 \\
LAGOA & 2025 & 0 & 27 & 123 & 0,0 \\
JARDIM BOTANICO & 2025 & 0 & 28 & 86 & 0,0 \\
GAVEA & 2025 & 0 & 29 & 139 & 0,0 \\
VIDIGAL & 2025 & 0 & 30 & 145 & 0,0 \\
SAO CONRADO & 2025 & 0 & 31 & 85 & 0,0 \\
PRACA DA BANDEIRA & 2025 & 0 & 32 & 57 & 0,0 \\
TIJUCA & 2025 & 3 & 33 & 941 & 3,2 \\
ALTO DA BOA VISTA & 2025 & 0 & 34 & 67 & 0,0 \\
MARACANA & 2025 & 0 & 35 & 131 & 0,0 \\
VILA ISABEL & 2025 & 3 & 36 & 443 & 6,8 \\
ANDARAI & 2025 & 1 & 37 & 209 & 4,8 \\
GRAJAU & 2025 & 0 & 38 & 168 & 0,0 \\
MANGUINHOS & 2025 & 1 & 39 & 387 & 2,6 \\
BONSUCESSO & 2025 & 0 & 40 & 392 & 0,0 \\
RAMOS & 2025 & 1 & 41 & 479 & 2,1 \\
OLARIA & 2025 & 0 & 42 & 416 & 0,0 \\
PENHA & 2025 & 1 & 43 & 609 & 1,6 \\
PENHA CIRCULAR & 2025 & 1 & 44 & 356 & 2,8 \\
BRAS DE PINA & 2025 & 3 & 45 & 364 & 8,2 \\
CORDOVIL & 2025 & 3 & 46 & 316 & 9,5 \\
PARADA DE LUCAS & 2025 & 1 & 47 & 255 & 3,9 \\
VIGARIO GERAL & 2025 & 0 & 48 & 354 & 0,0 \\
JARDIM AMERICA & 2025 & 0 & 49 & 266 & 0,0 \\
HIGIENOPOLIS & 2025 & 0 & 50 & 118 & 0,0 \\
JACARE & 2025 & 0 & 51 & 166 & 0,0 \\
MARIA DA GRACA & 2025 & 0 & 52 & 100 & 0,0 \\
DEL CASTILHO & 2025 & 0 & 53 & 236 & 0,0 \\
INHAUMA & 2025 & 1 & 54 & 539 & 1,9 \\
ENGENHO DA RAINHA & 2025 & 1 & 55 & 217 & 4,6 \\
TOMAS COELHO & 2025 & 0 & 56 & 191 & 0,0 \\
SAO FRANCISCO XAVIER & 2025 & 0 & 57 & 55 & 0,0 \\
ROCHA & 2025 & 1 & 58 & 178 & 5,6 \\
RIACHUELO & 2025 & 1 & 59 & 66 & 15,2 \\
SAMPAIO & 2025 & 0 & 60 & 99 & 0,0 \\
ENGENHO NOVO & 2025 & 2 & 61 & 346 & 5,8 \\
LINS DE VASCONCELOS & 2025 & 0 & 62 & 304 & 0,0 \\
MEIER & 2025 & 2 & 63 & 247 & 8,1 \\
TODOS OS SANTOS & 2025 & 0 & 64 & 126 & 0,0 \\
CACHAMBI & 2025 & 2 & 65 & 274 & 7,3 \\
ENGENHO DE DENTRO & 2025 & 0 & 66 & 351 & 0,0 \\
AGUA SANTA & 2025 & 0 & 67 & 57 & 0,0 \\
ENCANTADO & 2025 & 0 & 68 & 81 & 0,0 \\
PIEDADE & 2025 & 0 & 69 & 328 & 0,0 \\
ABOLICAO & 2025 & 0 & 70 & 109 & 0,0 \\
PILARES & 2025 & 0 & 71 & 191 & 0,0 \\
VILA KOSMOS & 2025 & 0 & 72 & 109 & 0,0 \\
VICENTE DE CARVALHO & 2025 & 0 & 73 & 189 & 0,0 \\
VILA DA PENHA & 2025 & 0 & 74 & 148 & 0,0 \\
VISTA ALEGRE & 2025 & 0 & 75 & 50 & 0,0 \\
IRAJA & 2025 & 1 & 76 & 683 & 1,5 \\
COLEGIO & 2025 & 0 & 77 & 292 & 0,0 \\
CAMPINHO & 2025 & 1 & 78 & 95 & 10,5 \\
QUINTINO BOCAIUVA & 2025 & 0 & 79 & 202 & 0,0 \\
CAVALCANTE & 2025 & 1 & 80 & 145 & 6,9 \\
ENGENHEIRO LEAL & 2025 & 0 & 81 & 51 & 0,0 \\
CASCADURA & 2025 & 1 & 82 & 243 & 4,1 \\
MADUREIRA & 2025 & 2 & 83 & 378 & 5,3 \\
VAZ LOBO & 2025 & 1 & 84 & 139 & 7,2 \\
TURIACU & 2025 & 0 & 85 & 136 & 0,0 \\
ROCHA MIRANDA & 2025 & 0 & 86 & 324 & 0,0 \\
HONORIO GURGEL & 2025 & 1 & 87 & 183 & 5,5 \\
OSWALDO CRUZ & 2025 & 0 & 88 & 196 & 0,0 \\
BENTO RIBEIRO & 2025 & 1 & 89 & 300 & 3,3 \\
MARECHAL HERMES & 2025 & 2 & 90 & 334 & 6,0 \\
CACUIA & 2025 & 0 & 93 & 85 & 0,0 \\
PITANGUEIRAS & 2025 & 0 & 94 & 76 & 0,0 \\
PRAIA DA BANDEIRA & 2025 & 0 & 95 & 37 & 0,0 \\
COCOTA & 2025 & 0 & 96 & 53 & 0,0 \\
BANCARIOS & 2025 & 0 & 97 & 132 & 0,0 \\
FREGUESIA-ILHA & 2025 & 0 & 98 & 191 & 0,0 \\
JARDIM GUANABARA & 2025 & 0 & 99 & 170 & 0,0 \\
JARDIM CARIOCA & 2025 & 0 & 100 & 199 & 0,0 \\
TAUA & 2025 & 1 & 101 & 212 & 4,7 \\
MONERO & 2025 & 0 & 102 & 39 & 0,0 \\
PORTUGUESA & 2025 & 0 & 103 & 127 & 0,0 \\
GALEAO & 2025 & 0 & 104 & 248 & 0,0 \\
CIDADE UNIVERSITARIA & 2025 & 0 & 105 & 17 & 0,0 \\
GUADALUPE & 2025 & 0 & 106 & 355 & 0,0 \\
ANCHIETA & 2025 & 0 & 107 & 618 & 0,0 \\
PARQUE ANCHIETA & 2025 & 0 & 108 & 144 & 0,0 \\
RICARDO ALBUQUERQUE & 2025 & 0 & 109 & 262 & 0,0 \\
COELHO NETO & 2025 & 0 & 110 & 284 & 0,0 \\
ACARI & 2025 & 2 & 111 & 349 & 5,7 \\
BARROS FILHO & 2025 & 0 & 112 & 164 & 0,0 \\
COSTA BARROS & 2025 & 0 & 113 & 401 & 0,0 \\
PAVUNA & 2025 & 0 & 114 & 888 & 0,0 \\
JACAREPAGUA & 2025 & 8 & 115 & 1.878 & 4,3 \\
ANIL & 2025 & 0 & 116 & 377 & 0,0 \\
GARDENIA AZUL & 2025 & 1 & 117 & 400 & 2,5 \\
CIDADE DE DEUS & 2025 & 2 & 118 & 549 & 3,6 \\
CURICICA & 2025 & 1 & 119 & 369 & 2,7 \\
FREGUESIA-JPA & 2025 & 1 & 120 & 553 & 1,8 \\
PECHINCHA & 2025 & 1 & 121 & 339 & 2,9 \\
TAQUARA & 2025 & 2 & 122 & 1.005 & 2,0 \\
TANQUE & 2025 & 3 & 123 & 378 & 7,9 \\
PRACA SECA & 2025 & 1 & 124 & 567 & 1,8 \\
VILA VALQUEIRE & 2025 & 0 & 125 & 273 & 0,0 \\
ITANHANGA & 2025 & 2 & 127 & 791 & 2,5 \\
BARRA DA TIJUCA & 2025 & 2 & 128 & 1.115 & 1,8 \\
CAMORIM & 2025 & 0 & 129 & 35 & 0,0 \\
VARGEM PEQUENA & 2025 & 1 & 130 & 325 & 3,1 \\
VARGEM GRANDE & 2025 & 0 & 131 & 255 & 0,0 \\
RECREIO BANDEIRANTES & 2025 & 1 & 132 & 1.245 & 0,8 \\
DEODORO & 2025 & 0 & 134 & 124 & 0,0 \\
VILA MILITAR & 2025 & 1 & 135 & 108 & 9,3 \\
CAMPO DOS AFONSOS & 2025 & 0 & 136 & 36 & 0,0 \\
JARDIM SULACAP & 2025 & 0 & 137 & 132 & 0,0 \\
MAGALHAES BASTOS & 2025 & 2 & 138 & 208 & 9,6 \\
REALENGO & 2025 & 5 & 139 & 1.463 & 3,4 \\
PADRE MIGUEL & 2025 & 1 & 140 & 522 & 1,9 \\
BANGU & 2025 & 3 & 141 & 1.977 & 1,5 \\
SENADOR CAMARA & 2025 & 3 & 142 & 995 & 3,0 \\
SANTISSIMO & 2025 & 2 & 143 & 471 & 4,2 \\
CAMPO GRANDE & 2025 & 12 & 144 & 3.504 & 3,4 \\
SENADOR VASCONCELOS & 2025 & 0 & 145 & 298 & 0,0 \\
INHOAIBA & 2025 & 1 & 146 & 752 & 1,3 \\
COSMOS & 2025 & 3 & 147 & 954 & 3,1 \\
PACIENCIA & 2025 & 3 & 148 & 1.157 & 2,6 \\
SANTA CRUZ & 2025 & 12 & 149 & 3.186 & 3,8 \\
SEPETIBA & 2025 & 2 & 150 & 753 & 2,7 \\
GUARATIBA & 2025 & 3 & 151 & 1.867 & 1,6 \\
BARRA DE GUARATIBA & 2025 & 1 & 152 & 117 & 8,5 \\
PEDRA DE GUARATIBA & 2025 & 0 & 153 & 228 & 0,0 \\
ROCINHA & 2025 & 1 & 154 & 696 & 1,4 \\
JACAREZINHO & 2025 & 1 & 155 & 293 & 3,4 \\
COMPLEXO DO ALEMAO & 2025 & 1 & 156 & 361 & 2,8 \\
COMPLEXO DA MARE & 2025 & 5 & 157 & 1.494 & 3,3 \\
PARQUE COLUMBIA & 2025 & 1 & 158 & 83 & 12,0 \\
VASCO DA GAMA & 2025 & 0 & 159 & 131 & 0,0 \\
LAPA & 2025 & 0 & 161 & 39 & 0,0 \\
VILA KENNEDY & 2025 & 1 & 162 & 194 & 5,2 \\
BARRA OLIMPICA & 2025 & 1 & 165 & 503 & 2,0 \\
ILHA DO GOVERNADOR - IGNORADO & 2025 & 0 & 998 & 12 & 0,0 \\
IGNORADO BAIRRO-MRJ & 2025 & 0 & 999 & 19 & 0,0 \\
\end{longtable}}
\vspace{-6pt}\fonte{Elaboração IPP com dados de \citeonline{sms_rio_sim}; \citeonline{sms_rio_sinasc}.}


{\small\setlength{\tabcolsep}{5pt}
\begin{longtable}{@{}rr@{}}
\caption{Razão de mortalidade materna (durante a gravidez) (obitos gravidez por ano)}\label{tab:obitos-gravidez-por-ano}\\
\toprule \textbf{Ano} & \textbf{Óbitos-gravidez} \\ \midrule \endfirsthead
\multicolumn{2}{@{}l}{\footnotesize\itshape (continuação)}\\ \toprule \textbf{Ano} & \textbf{Óbitos-gravidez} \\ \midrule \endhead
\midrule \multicolumn{2}{r@{}}{\footnotesize\itshape (continua)}\\ \endfoot
\bottomrule \endlastfoot
2006 & 78 \\
2007 & 69 \\
2008 & 57 \\
2009 & 71 \\
2010 & 33 \\
2011 & 19 \\
2012 & 24 \\
2013 & 37 \\
2014 & 40 \\
2015 & 21 \\
2016 & 23 \\
2017 & 24 \\
2018 & 18 \\
2019 & 32 \\
2020 & 32 \\
2021 & 35 \\
2022 & 15 \\
2023 & 20 \\
2024 & 4 \\
2025 & 15 \\
\end{longtable}}
\vspace{-6pt}\fonte{Elaboração IPP com dados de \citeonline{sms_rio_sim}.}


{\small\setlength{\tabcolsep}{5pt}
\begin{longtable}{@{}lrrr@{}}
\caption{Razão de mortalidade materna (durante a gravidez) (obitos gravidez 2025)}\label{tab:tabela-mapa-obitos-gravidez-2025}\\
\toprule \textbf{Bairro} & \textbf{Ano} & \textbf{Óbitos-gravidez} & \textbf{Codigo} \\ \midrule \endfirsthead
\multicolumn{4}{@{}l}{\footnotesize\itshape (continuação)}\\ \toprule \textbf{Bairro} & \textbf{Ano} & \textbf{Óbitos-gravidez} & \textbf{Codigo} \\ \midrule \endhead
\midrule \multicolumn{4}{r@{}}{\footnotesize\itshape (continua)}\\ \endfoot
\bottomrule \endlastfoot
SAUDE & 2025 & 0 & 1 \\
GAMBOA & 2025 & 0 & 2 \\
CAJU & 2025 & 0 & 4 \\
CENTRO & 2025 & 0 & 5 \\
CATUMBI & 2025 & 0 & 6 \\
RIO COMPRIDO & 2025 & 0 & 7 \\
CIDADE NOVA & 2025 & 0 & 8 \\
ESTACIO & 2025 & 0 & 9 \\
SAO CRISTOVAO & 2025 & 0 & 10 \\
MANGUEIRA & 2025 & 0 & 11 \\
BENFICA & 2025 & 0 & 12 \\
SANTA TERESA & 2025 & 0 & 14 \\
FLAMENGO & 2025 & 0 & 15 \\
BOTAFOGO & 2025 & 0 & 20 \\
COPACABANA & 2025 & 1 & 24 \\
IPANEMA & 2025 & 0 & 25 \\
GAVEA & 2025 & 0 & 29 \\
VIDIGAL & 2025 & 0 & 30 \\
PRACA DA BANDEIRA & 2025 & 0 & 32 \\
TIJUCA & 2025 & 0 & 33 \\
MARACANA & 2025 & 0 & 35 \\
VILA ISABEL & 2025 & 0 & 36 \\
ANDARAI & 2025 & 0 & 37 \\
GRAJAU & 2025 & 0 & 38 \\
MANGUINHOS & 2025 & 0 & 39 \\
BONSUCESSO & 2025 & 0 & 40 \\
RAMOS & 2025 & 0 & 41 \\
OLARIA & 2025 & 0 & 42 \\
PENHA & 2025 & 0 & 43 \\
PENHA CIRCULAR & 2025 & 0 & 44 \\
BRAS DE PINA & 2025 & 0 & 45 \\
CORDOVIL & 2025 & 0 & 46 \\
PARADA DE LUCAS & 2025 & 0 & 47 \\
VIGARIO GERAL & 2025 & 2 & 48 \\
JARDIM AMERICA & 2025 & 0 & 49 \\
HIGIENOPOLIS & 2025 & 0 & 50 \\
JACARE & 2025 & 0 & 51 \\
DEL CASTILHO & 2025 & 0 & 53 \\
INHAUMA & 2025 & 0 & 54 \\
ENGENHO DA RAINHA & 2025 & 0 & 55 \\
TOMAS COELHO & 2025 & 0 & 56 \\
ROCHA & 2025 & 0 & 58 \\
ENGENHO NOVO & 2025 & 0 & 61 \\
LINS DE VASCONCELOS & 2025 & 0 & 62 \\
TODOS OS SANTOS & 2025 & 0 & 64 \\
CACHAMBI & 2025 & 0 & 65 \\
ENGENHO DE DENTRO & 2025 & 0 & 66 \\
PIEDADE & 2025 & 0 & 69 \\
PILARES & 2025 & 0 & 71 \\
VILA KOSMOS & 2025 & 0 & 72 \\
VICENTE DE CARVALHO & 2025 & 0 & 73 \\
VILA DA PENHA & 2025 & 0 & 74 \\
VISTA ALEGRE & 2025 & 0 & 75 \\
IRAJA & 2025 & 0 & 76 \\
COLEGIO & 2025 & 0 & 77 \\
QUINTINO BOCAIUVA & 2025 & 0 & 79 \\
CAVALCANTE & 2025 & 0 & 80 \\
ENGENHEIRO LEAL & 2025 & 0 & 81 \\
MADUREIRA & 2025 & 0 & 83 \\
ROCHA MIRANDA & 2025 & 0 & 86 \\
OSWALDO CRUZ & 2025 & 0 & 88 \\
BENTO RIBEIRO & 2025 & 0 & 89 \\
MARECHAL HERMES & 2025 & 0 & 90 \\
CACUIA & 2025 & 0 & 93 \\
PITANGUEIRAS & 2025 & 0 & 94 \\
FREGUESIA-ILHA & 2025 & 0 & 98 \\
JARDIM GUANABARA & 2025 & 0 & 99 \\
JARDIM CARIOCA & 2025 & 0 & 100 \\
TAUA & 2025 & 0 & 101 \\
MONERO & 2025 & 0 & 102 \\
PORTUGUESA & 2025 & 0 & 103 \\
GALEAO & 2025 & 0 & 104 \\
GUADALUPE & 2025 & 0 & 106 \\
ANCHIETA & 2025 & 0 & 107 \\
RICARDO ALBUQUERQUE & 2025 & 0 & 109 \\
COELHO NETO & 2025 & 0 & 110 \\
ACARI & 2025 & 0 & 111 \\
BARROS FILHO & 2025 & 0 & 112 \\
COSTA BARROS & 2025 & 1 & 113 \\
PAVUNA & 2025 & 1 & 114 \\
JACAREPAGUA & 2025 & 0 & 115 \\
ANIL & 2025 & 0 & 116 \\
GARDENIA AZUL & 2025 & 0 & 117 \\
CIDADE DE DEUS & 2025 & 0 & 118 \\
CURICICA & 2025 & 0 & 119 \\
FREGUESIA-JPA & 2025 & 0 & 120 \\
TAQUARA & 2025 & 0 & 122 \\
TANQUE & 2025 & 0 & 123 \\
PRACA SECA & 2025 & 0 & 124 \\
VILA VALQUEIRE & 2025 & 0 & 125 \\
ITANHANGA & 2025 & 0 & 127 \\
BARRA DA TIJUCA & 2025 & 0 & 128 \\
VARGEM PEQUENA & 2025 & 0 & 130 \\
RECREIO BANDEIRANTES & 2025 & 0 & 132 \\
VILA MILITAR & 2025 & 0 & 135 \\
MAGALHAES BASTOS & 2025 & 0 & 138 \\
REALENGO & 2025 & 1 & 139 \\
PADRE MIGUEL & 2025 & 0 & 140 \\
BANGU & 2025 & 1 & 141 \\
SENADOR CAMARA & 2025 & 1 & 142 \\
SANTISSIMO & 2025 & 0 & 143 \\
CAMPO GRANDE & 2025 & 0 & 144 \\
SENADOR VASCONCELOS & 2025 & 0 & 145 \\
INHOAIBA & 2025 & 0 & 146 \\
COSMOS & 2025 & 0 & 147 \\
PACIENCIA & 2025 & 0 & 148 \\
SANTA CRUZ & 2025 & 1 & 149 \\
SEPETIBA & 2025 & 1 & 150 \\
GUARATIBA & 2025 & 1 & 151 \\
BARRA DE GUARATIBA & 2025 & 0 & 152 \\
PEDRA DE GUARATIBA & 2025 & 0 & 153 \\
ROCINHA & 2025 & 2 & 154 \\
JACAREZINHO & 2025 & 0 & 155 \\
COMPLEXO DO ALEMAO & 2025 & 0 & 156 \\
COMPLEXO DA MARE & 2025 & 0 & 157 \\
PARQUE COLUMBIA & 2025 & 0 & 158 \\
VASCO DA GAMA & 2025 & 1 & 159 \\
ILHA DO GOVERNADOR - IGNORADO & 2025 & 0 & 998 \\
IGNORADO BAIRRO-MRJ & 2025 & 0 & 999 \\
\end{longtable}}
\vspace{-6pt}\fonte{Elaboração IPP com dados de \citeonline{sms_rio_sim}.}


{\small\setlength{\tabcolsep}{5pt}
\begin{longtable}{@{}rr@{}}
\caption{Razão de mortalidade materna (durante puerpério) (obitos puerperio por ano)}\label{tab:obitos-puerperio-por-ano}\\
\toprule \textbf{Ano} & \textbf{Óbitos-puerpério} \\ \midrule \endfirsthead
\multicolumn{2}{@{}l}{\footnotesize\itshape (continuação)}\\ \toprule \textbf{Ano} & \textbf{Óbitos-puerpério} \\ \midrule \endhead
\midrule \multicolumn{2}{r@{}}{\footnotesize\itshape (continua)}\\ \endfoot
\bottomrule \endlastfoot
2006 & 67 \\
2007 & 72 \\
2008 & 51 \\
2009 & 65 \\
2010 & 64 \\
2011 & 56 \\
2012 & 57 \\
2013 & 59 \\
2014 & 56 \\
2015 & 66 \\
2016 & 59 \\
2017 & 67 \\
2018 & 56 \\
2019 & 54 \\
2020 & 74 \\
2021 & 109 \\
2022 & 42 \\
2023 & 40 \\
2024 & 33 \\
2025 & 35 \\
\end{longtable}}
\vspace{-6pt}\fonte{Elaboração IPP com dados de \citeonline{sms_rio_sim}.}


{\small\setlength{\tabcolsep}{5pt}
\begin{longtable}{@{}lrrr@{}}
\caption{Razão de mortalidade materna (durante puerpério) (obitos puerperio 2025)}\label{tab:tabela-mapa-obitos-puerperio-2025}\\
\toprule \textbf{Bairro} & \textbf{Ano} & \textbf{Óbitos-puerpério} & \textbf{Codigo} \\ \midrule \endfirsthead
\multicolumn{4}{@{}l}{\footnotesize\itshape (continuação)}\\ \toprule \textbf{Bairro} & \textbf{Ano} & \textbf{Óbitos-puerpério} & \textbf{Codigo} \\ \midrule \endhead
\midrule \multicolumn{4}{r@{}}{\footnotesize\itshape (continua)}\\ \endfoot
\bottomrule \endlastfoot
GAMBOA & 2025 & 0 & 2 \\
SANTO CRISTO & 2025 & 0 & 3 \\
CAJU & 2025 & 0 & 4 \\
CENTRO & 2025 & 0 & 5 \\
CATUMBI & 2025 & 0 & 6 \\
RIO COMPRIDO & 2025 & 0 & 7 \\
CIDADE NOVA & 2025 & 0 & 8 \\
ESTACIO & 2025 & 1 & 9 \\
SAO CRISTOVAO & 2025 & 0 & 10 \\
MANGUEIRA & 2025 & 1 & 11 \\
BENFICA & 2025 & 0 & 12 \\
SANTA TERESA & 2025 & 0 & 14 \\
FLAMENGO & 2025 & 0 & 15 \\
LARANJEIRAS & 2025 & 0 & 17 \\
CATETE & 2025 & 0 & 18 \\
COSME VELHO & 2025 & 0 & 19 \\
BOTAFOGO & 2025 & 0 & 20 \\
URCA & 2025 & 0 & 22 \\
LEME & 2025 & 0 & 23 \\
COPACABANA & 2025 & 0 & 24 \\
JARDIM BOTANICO & 2025 & 0 & 28 \\
GAVEA & 2025 & 0 & 29 \\
VIDIGAL & 2025 & 0 & 30 \\
PRACA DA BANDEIRA & 2025 & 0 & 32 \\
TIJUCA & 2025 & 0 & 33 \\
ALTO DA BOA VISTA & 2025 & 1 & 34 \\
MARACANA & 2025 & 0 & 35 \\
VILA ISABEL & 2025 & 0 & 36 \\
GRAJAU & 2025 & 0 & 38 \\
MANGUINHOS & 2025 & 1 & 39 \\
BONSUCESSO & 2025 & 0 & 40 \\
RAMOS & 2025 & 0 & 41 \\
OLARIA & 2025 & 0 & 42 \\
PENHA & 2025 & 0 & 43 \\
PENHA CIRCULAR & 2025 & 0 & 44 \\
BRAS DE PINA & 2025 & 0 & 45 \\
CORDOVIL & 2025 & 0 & 46 \\
PARADA DE LUCAS & 2025 & 0 & 47 \\
VIGARIO GERAL & 2025 & 0 & 48 \\
JARDIM AMERICA & 2025 & 0 & 49 \\
HIGIENOPOLIS & 2025 & 0 & 50 \\
JACARE & 2025 & 0 & 51 \\
MARIA DA GRACA & 2025 & 0 & 52 \\
DEL CASTILHO & 2025 & 0 & 53 \\
INHAUMA & 2025 & 0 & 54 \\
ENGENHO DA RAINHA & 2025 & 0 & 55 \\
TOMAS COELHO & 2025 & 0 & 56 \\
SAO FRANCISCO XAVIER & 2025 & 0 & 57 \\
ROCHA & 2025 & 1 & 58 \\
RIACHUELO & 2025 & 0 & 59 \\
SAMPAIO & 2025 & 0 & 60 \\
ENGENHO NOVO & 2025 & 1 & 61 \\
LINS DE VASCONCELOS & 2025 & 0 & 62 \\
MEIER & 2025 & 0 & 63 \\
ENGENHO DE DENTRO & 2025 & 0 & 66 \\
ENCANTADO & 2025 & 0 & 68 \\
PIEDADE & 2025 & 0 & 69 \\
PILARES & 2025 & 0 & 71 \\
VILA KOSMOS & 2025 & 0 & 72 \\
VICENTE DE CARVALHO & 2025 & 0 & 73 \\
VILA DA PENHA & 2025 & 0 & 74 \\
VISTA ALEGRE & 2025 & 1 & 75 \\
IRAJA & 2025 & 1 & 76 \\
COLEGIO & 2025 & 0 & 77 \\
CAMPINHO & 2025 & 0 & 78 \\
QUINTINO BOCAIUVA & 2025 & 0 & 79 \\
CAVALCANTE & 2025 & 0 & 80 \\
ENGENHEIRO LEAL & 2025 & 0 & 81 \\
CASCADURA & 2025 & 0 & 82 \\
MADUREIRA & 2025 & 0 & 83 \\
VAZ LOBO & 2025 & 0 & 84 \\
TURIACU & 2025 & 0 & 85 \\
ROCHA MIRANDA & 2025 & 0 & 86 \\
HONORIO GURGEL & 2025 & 0 & 87 \\
OSWALDO CRUZ & 2025 & 0 & 88 \\
BENTO RIBEIRO & 2025 & 0 & 89 \\
MARECHAL HERMES & 2025 & 0 & 90 \\
ZUMBI & 2025 & 0 & 92 \\
CACUIA & 2025 & 0 & 93 \\
BANCARIOS & 2025 & 0 & 97 \\
FREGUESIA-ILHA & 2025 & 0 & 98 \\
JARDIM GUANABARA & 2025 & 0 & 99 \\
JARDIM CARIOCA & 2025 & 0 & 100 \\
TAUA & 2025 & 0 & 101 \\
GUADALUPE & 2025 & 0 & 106 \\
ANCHIETA & 2025 & 0 & 107 \\
PARQUE ANCHIETA & 2025 & 0 & 108 \\
RICARDO ALBUQUERQUE & 2025 & 0 & 109 \\
COELHO NETO & 2025 & 0 & 110 \\
ACARI & 2025 & 0 & 111 \\
BARROS FILHO & 2025 & 0 & 112 \\
COSTA BARROS & 2025 & 0 & 113 \\
PAVUNA & 2025 & 2 & 114 \\
JACAREPAGUA & 2025 & 2 & 115 \\
ANIL & 2025 & 1 & 116 \\
GARDENIA AZUL & 2025 & 1 & 117 \\
CIDADE DE DEUS & 2025 & 1 & 118 \\
CURICICA & 2025 & 0 & 119 \\
FREGUESIA-JPA & 2025 & 0 & 120 \\
PECHINCHA & 2025 & 0 & 121 \\
TAQUARA & 2025 & 0 & 122 \\
TANQUE & 2025 & 0 & 123 \\
PRACA SECA & 2025 & 0 & 124 \\
VILA VALQUEIRE & 2025 & 0 & 125 \\
ITANHANGA & 2025 & 1 & 127 \\
BARRA DA TIJUCA & 2025 & 0 & 128 \\
VARGEM PEQUENA & 2025 & 0 & 130 \\
VARGEM GRANDE & 2025 & 0 & 131 \\
RECREIO BANDEIRANTES & 2025 & 0 & 132 \\
DEODORO & 2025 & 0 & 134 \\
VILA MILITAR & 2025 & 0 & 135 \\
MAGALHAES BASTOS & 2025 & 0 & 138 \\
REALENGO & 2025 & 1 & 139 \\
PADRE MIGUEL & 2025 & 0 & 140 \\
BANGU & 2025 & 2 & 141 \\
SENADOR CAMARA & 2025 & 3 & 142 \\
SANTISSIMO & 2025 & 0 & 143 \\
CAMPO GRANDE & 2025 & 0 & 144 \\
SENADOR VASCONCELOS & 2025 & 0 & 145 \\
INHOAIBA & 2025 & 0 & 146 \\
COSMOS & 2025 & 1 & 147 \\
PACIENCIA & 2025 & 0 & 148 \\
SANTA CRUZ & 2025 & 1 & 149 \\
SEPETIBA & 2025 & 0 & 150 \\
GUARATIBA & 2025 & 2 & 151 \\
BARRA DE GUARATIBA & 2025 & 0 & 152 \\
PEDRA DE GUARATIBA & 2025 & 1 & 153 \\
ROCINHA & 2025 & 0 & 154 \\
JACAREZINHO & 2025 & 1 & 155 \\
COMPLEXO DO ALEMAO & 2025 & 2 & 156 \\
COMPLEXO DA MARE & 2025 & 1 & 157 \\
VASCO DA GAMA & 2025 & 0 & 159 \\
LAPA & 2025 & 0 & 161 \\
VILA KENNEDY & 2025 & 0 & 162 \\
IGNORADO BAIRRO-MRJ & 2025 & 0 & 999 \\
\end{longtable}}
\vspace{-6pt}\fonte{Elaboração IPP com dados de \citeonline{sms_rio_sim}.}


\begin{landscape}
{\scriptsize\setlength{\tabcolsep}{3pt}
\begin{longtable}{@{}rrrrrrrrrrrrrrrrrrr@{}}
\caption{Mortalidade infantil por raça/cor (menores de 1 ano) (mortalidade raca municipio ano)}\label{tab:mortalidade-raca-municipio-ano}\\
\toprule \textbf{Ano} & \textbf{Obitos amarela} & \textbf{Obitos branca} & \textbf{Obitos indigena} & \textbf{Obitos parda} & \textbf{Obitos preta} & \textbf{Obitos nao informado} & \textbf{Nascidos amarela} & \textbf{Nascidos branca} & \textbf{Nascidos indigena} & \textbf{Nascidos parda} & \textbf{Nascidos preta} & \textbf{Nascidos nao informado} & \textbf{Amarela (\%)} & \textbf{Branca (\%)} & \textbf{Indigena (\%)} & \textbf{Parda (\%)} & \textbf{Preta (\%)} & \textbf{Nao informado (\%)} \\ \midrule \endfirsthead
\multicolumn{19}{@{}l}{\footnotesize\itshape (continuação)}\\ \toprule \textbf{Ano} & \textbf{Obitos amarela} & \textbf{Obitos branca} & \textbf{Obitos indigena} & \textbf{Obitos parda} & \textbf{Obitos preta} & \textbf{Obitos nao informado} & \textbf{Nascidos amarela} & \textbf{Nascidos branca} & \textbf{Nascidos indigena} & \textbf{Nascidos parda} & \textbf{Nascidos preta} & \textbf{Nascidos nao informado} & \textbf{Amarela (\%)} & \textbf{Branca (\%)} & \textbf{Indigena (\%)} & \textbf{Parda (\%)} & \textbf{Preta (\%)} & \textbf{Nao informado (\%)} \\ \midrule \endhead
\midrule \multicolumn{19}{r@{}}{\footnotesize\itshape (continua)}\\ \endfoot
\bottomrule \endlastfoot
2006 & 0 & 468 & 0 & 373 & 89 & 167 & -- & -- & -- & -- & -- & -- & -- & -- & -- & -- & -- & -- \\
2007 & 0 & 441 & 1 & 425 & 74 & 145 & -- & -- & -- & -- & -- & -- & -- & -- & -- & -- & -- & -- \\
2008 & 0 & 440 & 0 & 449 & 85 & 137 & -- & -- & -- & -- & -- & -- & -- & -- & -- & -- & -- & -- \\
2009 & 0 & 459 & 0 & 456 & 109 & 121 & -- & -- & -- & -- & -- & -- & -- & -- & -- & -- & -- & -- \\
2010 & 3 & 438 & 0 & 477 & 78 & 86 & -- & -- & -- & -- & -- & -- & -- & -- & -- & -- & -- & -- \\
2011 & 1 & 443 & 1 & 456 & 86 & 84 & 187 & 25.203 & 36 & 23.058 & 6.026 & 31.215 & 0,5 & 1,8 & 2,8 & 2,0 & 1,4 & 0,3 \\
2012 & 0 & 481 & 0 & 469 & 60 & 86 & 253 & 35.276 & 49 & 39.358 & 9.710 & 1.375 & 0,0 & 1,4 & 0,0 & 1,2 & 0,6 & 6,2 \\
2013 & 0 & 422 & 0 & 535 & 61 & 91 & 350 & 34.911 & 53 & 41.993 & 9.227 & 889 & 0,0 & 1,2 & 0,0 & 1,3 & 0,7 & 10,2 \\
2014 & 0 & 403 & 0 & 485 & 55 & 74 & 367 & 33.847 & 53 & 46.182 & 8.337 & 974 & 0,0 & 1,2 & 0,0 & 1,1 & 0,7 & 7,6 \\
2015 & 3 & 419 & 0 & 529 & 49 & 97 & 289 & 33.108 & 44 & 47.308 & 8.003 & 1.551 & 1,0 & 1,3 & 0,0 & 1,1 & 0,6 & 6,2 \\
2016 & 1 & 412 & 2 & 502 & 58 & 85 & 214 & 30.312 & 49 & 43.417 & 7.779 & 1.083 & 0,5 & 1,4 & 4,1 & 1,2 & 0,8 & 7,8 \\
2017 & 0 & 340 & 1 & 449 & 79 & 81 & 168 & 30.215 & 43 & 43.963 & 8.283 & 1.667 & 0,0 & 1,1 & 2,3 & 1,0 & 0,9 & 4,9 \\
2018 & 0 & 332 & 0 & 472 & 55 & 106 & 229 & 29.414 & 55 & 42.342 & 9.054 & 1.390 & 0,0 & 1,1 & 0,0 & 1,1 & 0,6 & 7,6 \\
2019 & 2 & 367 & 0 & 427 & 51 & 84 & 321 & 26.822 & 55 & 38.600 & 9.079 & 1.600 & 0,6 & 1,4 & 0,0 & 1,1 & 0,6 & 5,2 \\
2020 & 1 & 322 & 1 & 434 & 51 & 73 & 450 & 23.945 & 73 & 35.764 & 9.912 & 2.270 & 0,2 & 1,3 & 1,4 & 1,2 & 0,5 & 3,2 \\
2021 & 1 & 287 & 0 & 440 & 49 & 62 & 409 & 22.156 & 62 & 33.633 & 9.736 & 648 & 0,2 & 1,3 & 0,0 & 1,3 & 0,5 & 9,6 \\
2022 & 2 & 297 & 1 & 375 & 77 & 53 & 407 & 21.530 & 72 & 30.188 & 10.075 & 2.418 & 0,5 & 1,4 & 1,4 & 1,2 & 0,8 & 2,2 \\
2023 & 2 & 288 & 0 & 425 & 67 & 36 & 454 & 21.743 & 72 & 29.550 & 10.543 & 425 & 0,4 & 1,3 & 0,0 & 1,4 & 0,6 & 8,5 \\
2024 & 2 & 284 & 0 & 360 & 57 & 48 & 285 & 20.311 & 82 & 26.659 & 9.742 & 553 & 0,7 & 1,4 & 0,0 & 1,4 & 0,6 & 8,7 \\
2025 & 0 & 272 & 1 & 396 & 57 & 47 & 234 & 20.448 & 69 & 27.610 & 10.281 & 529 & 0,0 & 1,3 & 1,4 & 1,4 & 0,6 & 8,9 \\
\end{longtable}}
\vspace{-6pt}\fonte{Elaboração IPP com dados de \citeonline{sms_rio_sim}.}
\end{landscape}


\begin{landscape}
{\scriptsize\setlength{\tabcolsep}{3pt}
\begin{longtable}{@{}rlrrrrrrrrrrrrrrrrrrrrrr@{}}
\caption{Mortalidade infantil por raça/cor (menores de 1 ano) (obitos raca total 2025)}\label{tab:tabela-mapa-obitos-raca-total-2025}\\
\toprule \textbf{Codigo} & \textbf{Bairro} & \textbf{Ano} & \textbf{Obitos amarela} & \textbf{Obitos branca} & \textbf{Obitos indigena} & \textbf{Obitos parda} & \textbf{Obitos preta} & \textbf{Obitos nao informado} & \textbf{Nascidos amarela} & \textbf{Nascidos branca} & \textbf{Nascidos indigena} & \textbf{Nascidos parda} & \textbf{Nascidos preta} & \textbf{Nascidos nao informado} & \textbf{Amarela (\%)} & \textbf{Branca (\%)} & \textbf{Indigena (\%)} & \textbf{Parda (\%)} & \textbf{Preta (\%)} & \textbf{Nao informado (\%)} & \textbf{Obitos total} & \textbf{Nascidos total} & \textbf{Total (\%)} \\ \midrule \endfirsthead
\multicolumn{24}{@{}l}{\footnotesize\itshape (continuação)}\\ \toprule \textbf{Codigo} & \textbf{Bairro} & \textbf{Ano} & \textbf{Obitos amarela} & \textbf{Obitos branca} & \textbf{Obitos indigena} & \textbf{Obitos parda} & \textbf{Obitos preta} & \textbf{Obitos nao informado} & \textbf{Nascidos amarela} & \textbf{Nascidos branca} & \textbf{Nascidos indigena} & \textbf{Nascidos parda} & \textbf{Nascidos preta} & \textbf{Nascidos nao informado} & \textbf{Amarela (\%)} & \textbf{Branca (\%)} & \textbf{Indigena (\%)} & \textbf{Parda (\%)} & \textbf{Preta (\%)} & \textbf{Nao informado (\%)} & \textbf{Obitos total} & \textbf{Nascidos total} & \textbf{Total (\%)} \\ \midrule \endhead
\midrule \multicolumn{24}{r@{}}{\footnotesize\itshape (continua)}\\ \endfoot
\bottomrule \endlastfoot
70 & ABOLICAO & 2025 & 0 & 1 & 0 & 1 & 0 & 0 & 1 & 44 & 0 & 44 & 18 & 2 & 0,0 & 2,3 & -- & 2,3 & 0,0 & 0,0 & 2 & 109 & 1,8 \\
111 & ACARI & 2025 & 0 & 2 & 0 & 2 & 1 & 1 & 0 & 72 & 0 & 208 & 67 & 2 & -- & 2,8 & -- & 1,0 & 1,5 & 50,0 & 6 & 349 & 1,7 \\
67 & AGUA SANTA & 2025 & 0 & 0 & 0 & 1 & 0 & 0 & 1 & 20 & 0 & 29 & 7 & 0 & 0,0 & 0,0 & -- & 3,5 & 0,0 & -- & 1 & 57 & 1,8 \\
34 & ALTO DA BOA VISTA & 2025 & 0 & 1 & 0 & 0 & 0 & 1 & 0 & 17 & 0 & 36 & 13 & 1 & -- & 5,9 & -- & 0,0 & 0,0 & 100,0 & 2 & 67 & 3,0 \\
107 & ANCHIETA & 2025 & 0 & 0 & 0 & 6 & 0 & 1 & 1 & 161 & 1 & 318 & 134 & 3 & 0,0 & 0,0 & 0,0 & 1,9 & 0,0 & 33,3 & 7 & 618 & 1,1 \\
37 & ANDARAI & 2025 & 0 & 2 & 0 & 0 & 0 & 0 & 3 & 113 & 1 & 62 & 30 & 0 & 0,0 & 1,8 & 0,0 & 0,0 & 0,0 & -- & 2 & 209 & 1,0 \\
116 & ANIL & 2025 & 0 & 3 & 0 & 3 & 1 & 0 & 2 & 134 & 0 & 201 & 37 & 3 & 0,0 & 2,2 & -- & 1,5 & 2,7 & 0,0 & 7 & 377 & 1,9 \\
97 & BANCARIOS & 2025 & 0 & 0 & 0 & 1 & 0 & 0 & 0 & 50 & 0 & 65 & 15 & 2 & -- & 0,0 & -- & 1,5 & 0,0 & 0,0 & 1 & 132 & 0,8 \\
141 & BANGU & 2025 & 0 & 6 & 0 & 12 & 2 & 0 & 14 & 556 & 6 & 969 & 425 & 7 & 0,0 & 1,1 & 0,0 & 1,2 & 0,5 & 0,0 & 20 & 1.977 & 1,0 \\
128 & BARRA DA TIJUCA & 2025 & 0 & 6 & 0 & 0 & 0 & 0 & 0 & 917 & 0 & 166 & 22 & 10 & -- & 0,7 & -- & 0,0 & 0,0 & 0,0 & 6 & 1.115 & 0,5 \\
152 & BARRA DE GUARATIBA & 2025 & 0 & 1 & 0 & 1 & 0 & 0 & 1 & 38 & 0 & 64 & 13 & 1 & 0,0 & 2,6 & -- & 1,6 & 0,0 & 0,0 & 2 & 117 & 1,7 \\
165 & BARRA OLIMPICA & 2025 & 0 & 2 & 0 & 0 & 0 & 1 & 0 & 366 & 0 & 114 & 22 & 1 & -- & 0,6 & -- & 0,0 & 0,0 & 100,0 & 3 & 503 & 0,6 \\
112 & BARROS FILHO & 2025 & 0 & 2 & 0 & 0 & 0 & 0 & 0 & 32 & 0 & 94 & 38 & 0 & -- & 6,2 & -- & 0,0 & 0,0 & -- & 2 & 164 & 1,2 \\
12 & BENFICA & 2025 & 0 & 1 & 0 & 5 & 1 & 1 & 3 & 104 & 1 & 165 & 73 & 0 & 0,0 & 1,0 & 0,0 & 3,0 & 1,4 & -- & 8 & 346 & 2,3 \\
89 & BENTO RIBEIRO & 2025 & 0 & 2 & 0 & 2 & 1 & 0 & 0 & 92 & 0 & 133 & 72 & 3 & -- & 2,2 & -- & 1,5 & 1,4 & 0,0 & 5 & 300 & 1,7 \\
40 & BONSUCESSO & 2025 & 0 & 1 & 0 & 2 & 0 & 0 & 1 & 127 & 0 & 198 & 65 & 1 & 0,0 & 0,8 & -- & 1,0 & 0,0 & 0,0 & 3 & 392 & 0,8 \\
20 & BOTAFOGO & 2025 & 0 & 2 & 0 & 0 & 0 & 0 & 3 & 377 & 0 & 106 & 42 & 6 & 0,0 & 0,5 & -- & 0,0 & 0,0 & 0,0 & 2 & 534 & 0,4 \\
45 & BRAS DE PINA & 2025 & 0 & 5 & 0 & 3 & 0 & 1 & 1 & 111 & 1 & 175 & 74 & 2 & 0,0 & 4,5 & 0,0 & 1,7 & 0,0 & 50,0 & 9 & 364 & 2,5 \\
65 & CACHAMBI & 2025 & 0 & 3 & 0 & 0 & 0 & 0 & 0 & 134 & 0 & 104 & 34 & 2 & -- & 2,2 & -- & 0,0 & 0,0 & 0,0 & 3 & 274 & 1,1 \\
93 & CACUIA & 2025 & 0 & 0 & 0 & 0 & 0 & 0 & 0 & 30 & 0 & 43 & 11 & 1 & -- & 0,0 & -- & 0,0 & 0,0 & 0,0 & 0 & 85 & 0,0 \\
4 & CAJU & 2025 & 0 & 4 & 0 & 3 & 1 & 1 & 1 & 78 & 0 & 173 & 52 & 1 & 0,0 & 5,1 & -- & 1,7 & 1,9 & 100,0 & 9 & 305 & 3,0 \\
129 & CAMORIM & 2025 & 0 & 1 & 0 & 0 & 0 & 0 & 0 & 10 & 0 & 21 & 4 & 0 & -- & 10,0 & -- & 0,0 & 0,0 & -- & 1 & 35 & 2,9 \\
78 & CAMPINHO & 2025 & 0 & 0 & 0 & 2 & 0 & 0 & 1 & 23 & 0 & 50 & 20 & 1 & 0,0 & 0,0 & -- & 4,0 & 0,0 & 0,0 & 2 & 95 & 2,1 \\
136 & CAMPO DOS AFONSOS & 2025 & 0 & 2 & 0 & 0 & 0 & 0 & 0 & 17 & 0 & 13 & 6 & 0 & -- & 11,8 & -- & 0,0 & 0,0 & -- & 2 & 36 & 5,6 \\
144 & CAMPO GRANDE & 2025 & 0 & 17 & 0 & 21 & 1 & 1 & 24 & 1.220 & 3 & 1.639 & 568 & 50 & 0,0 & 1,4 & 0,0 & 1,3 & 0,2 & 2,0 & 40 & 3.504 & 1,1 \\
82 & CASCADURA & 2025 & 0 & 2 & 0 & 1 & 0 & 1 & 0 & 79 & 0 & 108 & 50 & 6 & -- & 2,5 & -- & 0,9 & 0,0 & 16,7 & 4 & 243 & 1,6 \\
18 & CATETE & 2025 & 0 & 0 & 0 & 0 & 0 & 0 & 0 & 63 & 0 & 62 & 20 & 0 & -- & 0,0 & -- & 0,0 & 0,0 & -- & 0 & 145 & 0,0 \\
6 & CATUMBI & 2025 & 0 & 2 & 0 & 2 & 1 & 0 & 0 & 22 & 0 & 71 & 40 & 1 & -- & 9,1 & -- & 2,8 & 2,5 & 0,0 & 5 & 134 & 3,7 \\
80 & CAVALCANTE & 2025 & 0 & 1 & 0 & 1 & 0 & 0 & 0 & 37 & 0 & 74 & 31 & 3 & -- & 2,7 & -- & 1,4 & 0,0 & 0,0 & 2 & 145 & 1,4 \\
5 & CENTRO & 2025 & 0 & 2 & 0 & 2 & 0 & 0 & 3 & 75 & 3 & 76 & 37 & 0 & 0,0 & 2,7 & 0,0 & 2,6 & 0,0 & -- & 4 & 194 & 2,1 \\
118 & CIDADE DE DEUS & 2025 & 0 & 1 & 0 & 2 & 1 & 2 & 0 & 75 & 0 & 353 & 119 & 2 & -- & 1,3 & -- & 0,6 & 0,8 & 100,0 & 6 & 549 & 1,1 \\
8 & CIDADE NOVA & 2025 & 0 & 0 & 0 & 3 & 0 & 0 & 0 & 10 & 0 & 22 & 10 & 0 & -- & 0,0 & -- & 13,6 & 0,0 & -- & 3 & 42 & 7,1 \\
105 & CIDADE UNIVERSITARIA & 2025 & 0 & 1 & 0 & 0 & 0 & 0 & 0 & 6 & 1 & 8 & 2 & 0 & -- & 16,7 & 0,0 & 0,0 & 0,0 & -- & 1 & 17 & 5,9 \\
96 & COCOTA & 2025 & 0 & 0 & 0 & 0 & 0 & 0 & 1 & 26 & 0 & 20 & 5 & 1 & 0,0 & 0,0 & -- & 0,0 & 0,0 & 0,0 & 0 & 53 & 0,0 \\
110 & COELHO NETO & 2025 & 0 & 0 & 0 & 0 & 0 & 0 & 0 & 66 & 0 & 145 & 72 & 1 & -- & 0,0 & -- & 0,0 & 0,0 & 0,0 & 0 & 284 & 0,0 \\
77 & COLEGIO & 2025 & 0 & 4 & 0 & 2 & 0 & 0 & 1 & 89 & 1 & 137 & 64 & 0 & 0,0 & 4,5 & 0,0 & 1,5 & 0,0 & -- & 6 & 292 & 2,0 \\
157 & COMPLEXO DA MARE & 2025 & 0 & 2 & 0 & 9 & 2 & 2 & 8 & 389 & 4 & 854 & 231 & 8 & 0,0 & 0,5 & 0,0 & 1,1 & 0,9 & 25,0 & 15 & 1.494 & 1,0 \\
156 & COMPLEXO DO ALEMAO & 2025 & 0 & 1 & 0 & 5 & 0 & 0 & 0 & 86 & 0 & 188 & 86 & 1 & -- & 1,2 & -- & 2,7 & 0,0 & 0,0 & 6 & 361 & 1,7 \\
24 & COPACABANA & 2025 & 0 & 5 & 0 & 2 & 0 & 0 & 3 & 475 & 1 & 185 & 95 & 13 & 0,0 & 1,1 & 0,0 & 1,1 & 0,0 & 0,0 & 7 & 772 & 0,9 \\
46 & CORDOVIL & 2025 & 0 & 2 & 0 & 7 & 1 & 0 & 2 & 61 & 1 & 167 & 84 & 1 & 0,0 & 3,3 & 0,0 & 4,2 & 1,2 & 0,0 & 10 & 316 & 3,2 \\
19 & COSME VELHO & 2025 & 0 & 1 & 0 & 1 & 0 & 0 & 0 & 16 & 0 & 26 & 26 & 1 & -- & 6,2 & -- & 3,9 & 0,0 & 0,0 & 2 & 69 & 2,9 \\
147 & COSMOS & 2025 & 0 & 5 & 0 & 16 & 0 & 0 & 8 & 205 & 0 & 510 & 214 & 17 & 0,0 & 2,4 & -- & 3,1 & 0,0 & 0,0 & 21 & 954 & 2,2 \\
113 & COSTA BARROS & 2025 & 0 & 2 & 0 & 4 & 1 & 0 & 0 & 80 & 0 & 213 & 105 & 3 & -- & 2,5 & -- & 1,9 & 0,9 & 0,0 & 7 & 401 & 1,8 \\
119 & CURICICA & 2025 & 0 & 1 & 0 & 4 & 0 & 0 & 1 & 107 & 0 & 206 & 54 & 1 & 0,0 & 0,9 & -- & 1,9 & 0,0 & 0,0 & 5 & 369 & 1,4 \\
53 & DEL CASTILHO & 2025 & 0 & 2 & 0 & 0 & 0 & 1 & 1 & 107 & 0 & 91 & 35 & 2 & 0,0 & 1,9 & -- & 0,0 & 0,0 & 50,0 & 3 & 236 & 1,3 \\
134 & DEODORO & 2025 & 0 & 0 & 0 & 2 & 0 & 0 & 0 & 47 & 0 & 53 & 23 & 1 & -- & 0,0 & -- & 3,8 & 0,0 & 0,0 & 2 & 124 & 1,6 \\
68 & ENCANTADO & 2025 & 0 & 0 & 0 & 1 & 0 & 0 & 0 & 32 & 1 & 39 & 9 & 0 & -- & 0,0 & 0,0 & 2,6 & 0,0 & -- & 1 & 81 & 1,2 \\
81 & ENGENHEIRO LEAL & 2025 & 0 & 0 & 0 & 1 & 0 & 0 & 0 & 10 & 0 & 31 & 9 & 1 & -- & 0,0 & -- & 3,2 & 0,0 & 0,0 & 1 & 51 & 2,0 \\
55 & ENGENHO DA RAINHA & 2025 & 0 & 0 & 0 & 2 & 1 & 0 & 0 & 63 & 0 & 88 & 66 & 0 & -- & 0,0 & -- & 2,3 & 1,5 & -- & 3 & 217 & 1,4 \\
66 & ENGENHO DE DENTRO & 2025 & 0 & 2 & 0 & 1 & 0 & 0 & 0 & 148 & 0 & 142 & 59 & 2 & -- & 1,4 & -- & 0,7 & 0,0 & 0,0 & 3 & 351 & 0,8 \\
61 & ENGENHO NOVO & 2025 & 0 & 1 & 0 & 2 & 0 & 0 & 0 & 123 & 0 & 132 & 91 & 0 & -- & 0,8 & -- & 1,5 & 0,0 & -- & 3 & 346 & 0,9 \\
9 & ESTACIO & 2025 & 0 & 0 & 0 & 3 & 1 & 0 & 0 & 47 & 1 & 68 & 37 & 0 & -- & 0,0 & 0,0 & 4,4 & 2,7 & -- & 4 & 153 & 2,6 \\
15 & FLAMENGO & 2025 & 0 & 1 & 0 & 2 & 1 & 0 & 1 & 204 & 0 & 53 & 14 & 2 & 0,0 & 0,5 & -- & 3,8 & 7,1 & 0,0 & 4 & 274 & 1,5 \\
98 & FREGUESIA-ILHA & 2025 & 0 & 1 & 0 & 1 & 0 & 0 & 0 & 72 & 0 & 90 & 27 & 2 & -- & 1,4 & -- & 1,1 & 0,0 & 0,0 & 2 & 191 & 1,1 \\
120 & FREGUESIA-JPA & 2025 & 0 & 0 & 0 & 3 & 1 & 0 & 0 & 281 & 2 & 209 & 60 & 1 & -- & 0,0 & 0,0 & 1,4 & 1,7 & 0,0 & 4 & 553 & 0,7 \\
104 & GALEAO & 2025 & 0 & 0 & 0 & 0 & 1 & 0 & 1 & 84 & 1 & 133 & 27 & 2 & 0,0 & 0,0 & 0,0 & 0,0 & 3,7 & 0,0 & 1 & 248 & 0,4 \\
2 & GAMBOA & 2025 & 0 & 0 & 0 & 2 & 0 & 0 & 1 & 46 & 0 & 71 & 31 & 1 & 0,0 & 0,0 & -- & 2,8 & 0,0 & 0,0 & 2 & 150 & 1,3 \\
117 & GARDENIA AZUL & 2025 & 0 & 1 & 0 & 3 & 0 & 0 & 3 & 94 & 0 & 264 & 37 & 2 & 0,0 & 1,1 & -- & 1,1 & 0,0 & 0,0 & 4 & 400 & 1,0 \\
29 & GAVEA & 2025 & 0 & 1 & 0 & 1 & 0 & 0 & 0 & 88 & 0 & 40 & 9 & 2 & -- & 1,1 & -- & 2,5 & 0,0 & 0,0 & 2 & 139 & 1,4 \\
160 & GERICINo & 2025 & 0 & 0 & 0 & 1 & 0 & 0 & 0 & 4 & 0 & 6 & 4 & 0 & -- & 0,0 & -- & 16,7 & 0,0 & -- & 1 & 14 & 7,1 \\
16 & GLORIA & 2025 & 0 & 0 & 0 & 0 & 0 & 0 & 0 & 29 & 0 & 12 & 6 & 0 & -- & 0,0 & -- & 0,0 & 0,0 & -- & 0 & 47 & 0,0 \\
38 & GRAJAU & 2025 & 0 & 0 & 0 & 0 & 0 & 0 & 1 & 98 & 0 & 52 & 16 & 1 & 0,0 & 0,0 & -- & 0,0 & 0,0 & 0,0 & 0 & 168 & 0,0 \\
133 & GRUMARI & 2025 & 0 & 0 & 0 & 0 & 0 & 0 & 0 & 0 & 0 & 0 & 0 & 0 & -- & -- & -- & -- & -- & -- & 0 & 0 & -- \\
106 & GUADALUPE & 2025 & 0 & 2 & 0 & 2 & 0 & 1 & 1 & 102 & 0 & 187 & 64 & 1 & 0,0 & 2,0 & -- & 1,1 & 0,0 & 100,0 & 5 & 355 & 1,4 \\
151 & GUARATIBA & 2025 & 0 & 4 & 0 & 7 & 2 & 0 & 7 & 522 & 2 & 1.013 & 292 & 31 & 0,0 & 0,8 & 0,0 & 0,7 & 0,7 & 0,0 & 13 & 1.867 & 0,7 \\
50 & HIGIENOPOLIS & 2025 & 0 & 0 & 0 & 0 & 0 & 0 & 1 & 43 & 0 & 49 & 24 & 1 & 0,0 & 0,0 & -- & 0,0 & 0,0 & 0,0 & 0 & 118 & 0,0 \\
87 & HONORIO GURGEL & 2025 & 0 & 4 & 0 & 0 & 1 & 1 & 1 & 51 & 0 & 89 & 42 & 0 & 0,0 & 7,8 & -- & 0,0 & 2,4 & -- & 6 & 183 & 3,3 \\
21 & HUMAITA & 2025 & 0 & 0 & 0 & 1 & 0 & 0 & 0 & 51 & 0 & 13 & 5 & 0 & -- & 0,0 & -- & 7,7 & 0,0 & -- & 1 & 69 & 1,4 \\
999 & IGNORADO BAIRRO-MRJ & 2025 & 0 & 0 & 0 & 0 & 0 & 0 & 1 & 5 & 0 & 10 & 3 & 0 & 0,0 & 0,0 & -- & 0,0 & 0,0 & -- & 0 & 19 & 0,0 \\
164 & ILHA DE GUARATIBA & 2025 & 0 & 0 & 0 & 1 & 0 & 0 & 0 & 29 & 0 & 48 & 3 & 1 & -- & 0,0 & -- & 2,1 & 0,0 & 0,0 & 1 & 81 & 1,2 \\
998 & ILHA DO GOVERNADOR - IGNORADO & 2025 & 0 & 0 & 0 & 0 & 0 & 0 & 0 & 8 & 0 & 3 & 1 & 0 & -- & 0,0 & -- & 0,0 & 0,0 & -- & 0 & 12 & 0,0 \\
54 & INHAUMA & 2025 & 0 & 2 & 0 & 4 & 0 & 0 & 0 & 157 & 1 & 263 & 115 & 3 & -- & 1,3 & 0,0 & 1,5 & 0,0 & 0,0 & 6 & 539 & 1,1 \\
146 & INHOAIBA & 2025 & 0 & 4 & 0 & 3 & 0 & 0 & 3 & 194 & 0 & 403 & 147 & 5 & 0,0 & 2,1 & -- & 0,7 & 0,0 & 0,0 & 7 & 752 & 0,9 \\
25 & IPANEMA & 2025 & 0 & 5 & 0 & 0 & 0 & 1 & 1 & 197 & 1 & 33 & 10 & 5 & 0,0 & 2,5 & 0,0 & 0,0 & 0,0 & 20,0 & 6 & 247 & 2,4 \\
76 & IRAJA & 2025 & 0 & 5 & 0 & 3 & 0 & 1 & 1 & 259 & 1 & 297 & 116 & 9 & 0,0 & 1,9 & 0,0 & 1,0 & 0,0 & 11,1 & 9 & 683 & 1,3 \\
127 & ITANHANGA & 2025 & 0 & 6 & 0 & 4 & 0 & 1 & 8 & 251 & 2 & 470 & 55 & 5 & 0,0 & 2,4 & 0,0 & 0,8 & 0,0 & 20,0 & 11 & 791 & 1,4 \\
163 & JABOUR & 2025 & 0 & 0 & 0 & 0 & 0 & 0 & 0 & 7 & 0 & 9 & 1 & 0 & -- & 0,0 & -- & 0,0 & 0,0 & -- & 0 & 17 & 0,0 \\
51 & JACARE & 2025 & 0 & 0 & 0 & 0 & 0 & 0 & 1 & 35 & 0 & 86 & 44 & 0 & 0,0 & 0,0 & -- & 0,0 & 0,0 & -- & 0 & 166 & 0,0 \\
115 & JACAREPAGUA & 2025 & 0 & 13 & 0 & 16 & 1 & 2 & 4 & 578 & 1 & 1.077 & 208 & 10 & 0,0 & 2,2 & 0,0 & 1,5 & 0,5 & 20,0 & 32 & 1.878 & 1,7 \\
155 & JACAREZINHO & 2025 & 0 & 0 & 0 & 2 & 2 & 0 & 1 & 56 & 0 & 144 & 92 & 0 & 0,0 & 0,0 & -- & 1,4 & 2,2 & -- & 4 & 293 & 1,4 \\
49 & JARDIM AMERICA & 2025 & 0 & 1 & 0 & 1 & 0 & 0 & 1 & 84 & 0 & 135 & 44 & 2 & 0,0 & 1,2 & -- & 0,7 & 0,0 & 0,0 & 2 & 266 & 0,8 \\
28 & JARDIM BOTANICO & 2025 & 0 & 0 & 0 & 0 & 0 & 0 & 0 & 66 & 0 & 13 & 5 & 2 & -- & 0,0 & -- & 0,0 & 0,0 & 0,0 & 0 & 86 & 0,0 \\
100 & JARDIM CARIOCA & 2025 & 0 & 1 & 0 & 1 & 0 & 1 & 0 & 78 & 0 & 96 & 24 & 1 & -- & 1,3 & -- & 1,0 & 0,0 & 100,0 & 3 & 199 & 1,5 \\
99 & JARDIM GUANABARA & 2025 & 0 & 0 & 0 & 0 & 0 & 0 & 0 & 123 & 0 & 40 & 6 & 1 & -- & 0,0 & -- & 0,0 & 0,0 & 0,0 & 0 & 170 & 0,0 \\
137 & JARDIM SULACAP & 2025 & 0 & 1 & 0 & 0 & 0 & 0 & 0 & 65 & 0 & 46 & 19 & 2 & -- & 1,5 & -- & 0,0 & 0,0 & 0,0 & 1 & 132 & 0,8 \\
126 & JOA & 2025 & 0 & 0 & 0 & 0 & 0 & 0 & 0 & 8 & 0 & 0 & 0 & 0 & -- & 0,0 & -- & -- & -- & -- & 0 & 8 & 0,0 \\
27 & LAGOA & 2025 & 0 & 1 & 0 & 2 & 0 & 0 & 1 & 103 & 0 & 14 & 1 & 4 & 0,0 & 1,0 & -- & 14,3 & 0,0 & 0,0 & 3 & 123 & 2,4 \\
161 & LAPA & 2025 & 0 & 0 & 0 & 0 & 0 & 0 & 0 & 21 & 0 & 10 & 8 & 0 & -- & 0,0 & -- & 0,0 & 0,0 & -- & 0 & 39 & 0,0 \\
17 & LARANJEIRAS & 2025 & 0 & 0 & 0 & 1 & 0 & 0 & 1 & 156 & 0 & 35 & 14 & 0 & 0,0 & 0,0 & -- & 2,9 & 0,0 & -- & 1 & 206 & 0,5 \\
26 & LEBLON & 2025 & 0 & 1 & 0 & 1 & 0 & 0 & 2 & 226 & 1 & 31 & 7 & 8 & 0,0 & 0,4 & 0,0 & 3,2 & 0,0 & 0,0 & 2 & 275 & 0,7 \\
23 & LEME & 2025 & 0 & 1 & 0 & 0 & 0 & 0 & 2 & 39 & 0 & 40 & 11 & 2 & 0,0 & 2,6 & -- & 0,0 & 0,0 & 0,0 & 1 & 94 & 1,1 \\
62 & LINS DE VASCONCELOS & 2025 & 0 & 2 & 0 & 0 & 1 & 0 & 0 & 92 & 0 & 122 & 89 & 1 & -- & 2,2 & -- & 0,0 & 1,1 & 0,0 & 3 & 304 & 1,0 \\
83 & MADUREIRA & 2025 & 0 & 1 & 0 & 2 & 0 & 2 & 1 & 88 & 0 & 195 & 92 & 2 & 0,0 & 1,1 & -- & 1,0 & 0,0 & 100,0 & 5 & 378 & 1,3 \\
138 & MAGALHAES BASTOS & 2025 & 0 & 3 & 0 & 1 & 0 & 2 & 1 & 59 & 0 & 103 & 43 & 2 & 0,0 & 5,1 & -- & 1,0 & 0,0 & 100,0 & 6 & 208 & 2,9 \\
11 & MANGUEIRA & 2025 & 0 & 0 & 0 & 2 & 0 & 1 & 0 & 33 & 1 & 92 & 72 & 1 & -- & 0,0 & 0,0 & 2,2 & 0,0 & 100,0 & 3 & 199 & 1,5 \\
39 & MANGUINHOS & 2025 & 0 & 2 & 0 & 4 & 0 & 2 & 2 & 78 & 0 & 201 & 106 & 0 & 0,0 & 2,6 & -- & 2,0 & 0,0 & -- & 8 & 387 & 2,1 \\
35 & MARACANA & 2025 & 0 & 0 & 0 & 0 & 0 & 0 & 0 & 82 & 0 & 35 & 12 & 2 & -- & 0,0 & -- & 0,0 & 0,0 & 0,0 & 0 & 131 & 0,0 \\
90 & MARECHAL HERMES & 2025 & 0 & 1 & 0 & 4 & 1 & 1 & 3 & 112 & 0 & 142 & 77 & 0 & 0,0 & 0,9 & -- & 2,8 & 1,3 & -- & 7 & 334 & 2,1 \\
52 & MARIA DA GRACA & 2025 & 0 & 1 & 0 & 0 & 0 & 0 & 0 & 35 & 0 & 46 & 17 & 2 & -- & 2,9 & -- & 0,0 & 0,0 & 0,0 & 1 & 100 & 1,0 \\
63 & MEIER & 2025 & 0 & 3 & 0 & 2 & 1 & 0 & 0 & 149 & 0 & 66 & 30 & 2 & -- & 2,0 & -- & 3,0 & 3,3 & 0,0 & 6 & 247 & 2,4 \\
102 & MONERO & 2025 & 0 & 0 & 0 & 0 & 0 & 0 & 0 & 27 & 0 & 9 & 3 & 0 & -- & 0,0 & -- & 0,0 & 0,0 & -- & 0 & 39 & 0,0 \\
42 & OLARIA & 2025 & 0 & 5 & 0 & 0 & 0 & 1 & 3 & 132 & 0 & 182 & 97 & 2 & 0,0 & 3,8 & -- & 0,0 & 0,0 & 50,0 & 6 & 416 & 1,4 \\
88 & OSWALDO CRUZ & 2025 & 0 & 1 & 0 & 1 & 0 & 0 & 1 & 50 & 0 & 98 & 47 & 0 & 0,0 & 2,0 & -- & 1,0 & 0,0 & -- & 2 & 196 & 1,0 \\
148 & PACIENCIA & 2025 & 0 & 4 & 0 & 8 & 2 & 0 & 2 & 243 & 0 & 618 & 270 & 24 & 0,0 & 1,6 & -- & 1,3 & 0,7 & 0,0 & 14 & 1.157 & 1,2 \\
140 & PADRE MIGUEL & 2025 & 0 & 2 & 0 & 7 & 1 & 0 & 3 & 125 & 2 & 293 & 98 & 1 & 0,0 & 1,6 & 0,0 & 2,4 & 1,0 & 0,0 & 10 & 522 & 1,9 \\
13 & PAQUETA & 2025 & 0 & 0 & 0 & 0 & 0 & 0 & 0 & 8 & 0 & 16 & 6 & 0 & -- & 0,0 & -- & 0,0 & 0,0 & -- & 0 & 30 & 0,0 \\
47 & PARADA DE LUCAS & 2025 & 0 & 0 & 0 & 3 & 0 & 1 & 0 & 87 & 0 & 114 & 53 & 1 & -- & 0,0 & -- & 2,6 & 0,0 & 100,0 & 4 & 255 & 1,6 \\
108 & PARQUE ANCHIETA & 2025 & 0 & 1 & 0 & 1 & 0 & 0 & 0 & 42 & 0 & 63 & 36 & 3 & -- & 2,4 & -- & 1,6 & 0,0 & 0,0 & 2 & 144 & 1,4 \\
158 & PARQUE COLUMBIA & 2025 & 0 & 0 & 0 & 2 & 0 & 1 & 0 & 19 & 0 & 45 & 19 & 0 & -- & 0,0 & -- & 4,4 & 0,0 & -- & 3 & 83 & 3,6 \\
114 & PAVUNA & 2025 & 0 & 4 & 0 & 15 & 1 & 0 & 1 & 168 & 2 & 475 & 230 & 12 & 0,0 & 2,4 & 0,0 & 3,2 & 0,4 & 0,0 & 20 & 888 & 2,2 \\
121 & PECHINCHA & 2025 & 0 & 3 & 0 & 0 & 0 & 0 & 1 & 192 & 1 & 118 & 27 & 0 & 0,0 & 1,6 & 0,0 & 0,0 & 0,0 & -- & 3 & 339 & 0,9 \\
153 & PEDRA DE GUARATIBA & 2025 & 0 & 1 & 0 & 0 & 0 & 0 & 1 & 62 & 0 & 135 & 27 & 3 & 0,0 & 1,6 & -- & 0,0 & 0,0 & 0,0 & 1 & 228 & 0,4 \\
43 & PENHA & 2025 & 0 & 3 & 0 & 2 & 1 & 1 & 2 & 152 & 2 & 307 & 146 & 0 & 0,0 & 2,0 & 0,0 & 0,7 & 0,7 & -- & 7 & 609 & 1,1 \\
44 & PENHA CIRCULAR & 2025 & 0 & 5 & 0 & 2 & 1 & 0 & 0 & 121 & 1 & 172 & 62 & 0 & -- & 4,1 & 0,0 & 1,2 & 1,6 & -- & 8 & 356 & 2,2 \\
69 & PIEDADE & 2025 & 0 & 0 & 0 & 3 & 0 & 0 & 0 & 125 & 0 & 142 & 59 & 2 & -- & 0,0 & -- & 2,1 & 0,0 & 0,0 & 3 & 328 & 0,9 \\
71 & PILARES & 2025 & 0 & 1 & 0 & 0 & 1 & 0 & 1 & 67 & 0 & 78 & 45 & 0 & 0,0 & 1,5 & -- & 0,0 & 2,2 & -- & 2 & 191 & 1,1 \\
94 & PITANGUEIRAS & 2025 & 0 & 0 & 0 & 2 & 0 & 0 & 0 & 26 & 0 & 39 & 10 & 1 & -- & 0,0 & -- & 5,1 & 0,0 & 0,0 & 2 & 76 & 2,6 \\
103 & PORTUGUESA & 2025 & 0 & 0 & 0 & 1 & 0 & 1 & 1 & 45 & 1 & 58 & 20 & 2 & 0,0 & 0,0 & 0,0 & 1,7 & 0,0 & 50,0 & 2 & 127 & 1,6 \\
32 & PRACA DA BANDEIRA & 2025 & 0 & 0 & 0 & 0 & 0 & 0 & 0 & 28 & 0 & 23 & 5 & 1 & -- & 0,0 & -- & 0,0 & 0,0 & 0,0 & 0 & 57 & 0,0 \\
124 & PRACA SECA & 2025 & 0 & 3 & 0 & 5 & 2 & 0 & 2 & 180 & 0 & 281 & 96 & 8 & 0,0 & 1,7 & -- & 1,8 & 2,1 & 0,0 & 10 & 567 & 1,8 \\
95 & PRAIA DA BANDEIRA & 2025 & 0 & 0 & 0 & 0 & 0 & 0 & 2 & 19 & 0 & 12 & 4 & 0 & 0,0 & 0,0 & -- & 0,0 & 0,0 & -- & 0 & 37 & 0,0 \\
79 & QUINTINO BOCAIUVA & 2025 & 0 & 0 & 0 & 0 & 1 & 0 & 0 & 60 & 0 & 96 & 44 & 2 & -- & 0,0 & -- & 0,0 & 2,3 & 0,0 & 1 & 202 & 0,5 \\
41 & RAMOS & 2025 & 0 & 2 & 0 & 2 & 0 & 0 & 4 & 159 & 0 & 214 & 99 & 3 & 0,0 & 1,3 & -- & 0,9 & 0,0 & 0,0 & 4 & 479 & 0,8 \\
139 & REALENGO & 2025 & 0 & 4 & 0 & 13 & 0 & 0 & 5 & 385 & 4 & 750 & 311 & 8 & 0,0 & 1,0 & 0,0 & 1,7 & 0,0 & 0,0 & 17 & 1.463 & 1,2 \\
132 & RECREIO BANDEIRANTES & 2025 & 0 & 7 & 0 & 2 & 0 & 0 & 7 & 733 & 0 & 435 & 63 & 7 & 0,0 & 0,9 & -- & 0,5 & 0,0 & 0,0 & 9 & 1.245 & 0,7 \\
59 & RIACHUELO & 2025 & 0 & 1 & 0 & 0 & 1 & 0 & 1 & 28 & 0 & 29 & 8 & 0 & 0,0 & 3,6 & -- & 0,0 & 12,5 & -- & 2 & 66 & 3,0 \\
91 & RIBEIRA & 2025 & 0 & 0 & 0 & 1 & 0 & 0 & 0 & 11 & 0 & 8 & 3 & 1 & -- & 0,0 & -- & 12,5 & 0,0 & 0,0 & 1 & 23 & 4,3 \\
109 & RICARDO ALBUQUERQUE & 2025 & 0 & 1 & 0 & 1 & 0 & 0 & 1 & 56 & 0 & 152 & 52 & 1 & 0,0 & 1,8 & -- & 0,7 & 0,0 & 0,0 & 2 & 262 & 0,8 \\
7 & RIO COMPRIDO & 2025 & 0 & 2 & 0 & 3 & 0 & 0 & 3 & 129 & 2 & 142 & 71 & 2 & 0,0 & 1,6 & 0,0 & 2,1 & 0,0 & 0,0 & 5 & 349 & 1,4 \\
58 & ROCHA & 2025 & 0 & 0 & 0 & 4 & 0 & 1 & 0 & 54 & 0 & 79 & 45 & 0 & -- & 0,0 & -- & 5,1 & 0,0 & -- & 5 & 178 & 2,8 \\
86 & ROCHA MIRANDA & 2025 & 0 & 1 & 0 & 2 & 1 & 1 & 0 & 90 & 0 & 131 & 102 & 1 & -- & 1,1 & -- & 1,5 & 1,0 & 100,0 & 5 & 324 & 1,5 \\
154 & ROCINHA & 2025 & 0 & 4 & 0 & 4 & 0 & 1 & 1 & 197 & 2 & 409 & 80 & 7 & 0,0 & 2,0 & 0,0 & 1,0 & 0,0 & 14,3 & 9 & 696 & 1,3 \\
60 & SAMPAIO & 2025 & 0 & 0 & 0 & 2 & 1 & 0 & 0 & 29 & 0 & 50 & 20 & 0 & -- & 0,0 & -- & 4,0 & 5,0 & -- & 3 & 99 & 3,0 \\
149 & SANTA CRUZ & 2025 & 0 & 14 & 0 & 36 & 2 & 1 & 12 & 678 & 0 & 1.704 & 709 & 83 & 0,0 & 2,1 & -- & 2,1 & 0,3 & 1,2 & 53 & 3.186 & 1,7 \\
14 & SANTA TERESA & 2025 & 0 & 1 & 0 & 3 & 0 & 0 & 2 & 104 & 1 & 138 & 68 & 3 & 0,0 & 1,0 & 0,0 & 2,2 & 0,0 & 0,0 & 4 & 316 & 1,3 \\
143 & SANTISSIMO & 2025 & 0 & 0 & 0 & 8 & 0 & 0 & 2 & 125 & 2 & 255 & 83 & 4 & 0,0 & 0,0 & 0,0 & 3,1 & 0,0 & 0,0 & 8 & 471 & 1,7 \\
3 & SANTO CRISTO & 2025 & 0 & 0 & 0 & 2 & 0 & 1 & 2 & 32 & 1 & 66 & 28 & 0 & 0,0 & 0,0 & 0,0 & 3,0 & 0,0 & -- & 3 & 129 & 2,3 \\
31 & SAO CONRADO & 2025 & 0 & 0 & 0 & 0 & 0 & 0 & 0 & 59 & 0 & 20 & 6 & 0 & -- & 0,0 & -- & 0,0 & 0,0 & -- & 0 & 85 & 0,0 \\
10 & SAO CRISTOVAO & 2025 & 0 & 1 & 0 & 1 & 0 & 1 & 0 & 98 & 1 & 130 & 60 & 1 & -- & 1,0 & 0,0 & 0,8 & 0,0 & 100,0 & 3 & 290 & 1,0 \\
57 & SAO FRANCISCO XAVIER & 2025 & 0 & 0 & 0 & 1 & 0 & 0 & 0 & 16 & 0 & 24 & 15 & 0 & -- & 0,0 & -- & 4,2 & 0,0 & -- & 1 & 55 & 1,8 \\
1 & SAUDE & 2025 & 0 & 0 & 0 & 0 & 0 & 0 & 0 & 8 & 1 & 21 & 6 & 0 & -- & 0,0 & 0,0 & 0,0 & 0,0 & -- & 0 & 36 & 0,0 \\
142 & SENADOR CAMARA & 2025 & 0 & 3 & 0 & 5 & 0 & 0 & 7 & 240 & 0 & 529 & 212 & 7 & 0,0 & 1,2 & -- & 0,9 & 0,0 & 0,0 & 8 & 995 & 0,8 \\
145 & SENADOR VASCONCELOS & 2025 & 0 & 2 & 0 & 5 & 1 & 0 & 2 & 78 & 0 & 161 & 55 & 2 & 0,0 & 2,6 & -- & 3,1 & 1,8 & 0,0 & 8 & 298 & 2,7 \\
150 & SEPETIBA & 2025 & 0 & 1 & 0 & 5 & 3 & 0 & 3 & 193 & 0 & 397 & 148 & 12 & 0,0 & 0,5 & -- & 1,3 & 2,0 & 0,0 & 9 & 753 & 1,2 \\
123 & TANQUE & 2025 & 0 & 2 & 0 & 2 & 1 & 0 & 2 & 98 & 0 & 220 & 57 & 1 & 0,0 & 2,0 & -- & 0,9 & 1,8 & 0,0 & 5 & 378 & 1,3 \\
122 & TAQUARA & 2025 & 0 & 5 & 0 & 8 & 1 & 0 & 2 & 401 & 0 & 474 & 126 & 2 & 0,0 & 1,2 & -- & 1,7 & 0,8 & 0,0 & 14 & 1.005 & 1,4 \\
101 & TAUA & 2025 & 0 & 1 & 0 & 1 & 0 & 0 & 4 & 59 & 0 & 113 & 35 & 1 & 0,0 & 1,7 & -- & 0,9 & 0,0 & 0,0 & 2 & 212 & 0,9 \\
33 & TIJUCA & 2025 & 0 & 3 & 0 & 2 & 1 & 2 & 6 & 541 & 3 & 252 & 133 & 6 & 0,0 & 0,6 & 0,0 & 0,8 & 0,8 & 33,3 & 8 & 941 & 0,8 \\
64 & TODOS OS SANTOS & 2025 & 0 & 1 & 0 & 0 & 0 & 0 & 0 & 78 & 0 & 34 & 12 & 2 & -- & 1,3 & -- & 0,0 & 0,0 & 0,0 & 1 & 126 & 0,8 \\
56 & TOMAS COELHO & 2025 & 0 & 2 & 0 & 0 & 2 & 2 & 0 & 60 & 0 & 90 & 41 & 0 & -- & 3,3 & -- & 0,0 & 4,9 & -- & 6 & 191 & 3,1 \\
85 & TURIACU & 2025 & 0 & 1 & 0 & 0 & 0 & 0 & 1 & 30 & 0 & 62 & 43 & 0 & 0,0 & 3,3 & -- & 0,0 & 0,0 & -- & 1 & 136 & 0,7 \\
22 & URCA & 2025 & 0 & 0 & 0 & 0 & 0 & 0 & 0 & 40 & 0 & 12 & 1 & 4 & -- & 0,0 & -- & 0,0 & 0,0 & 0,0 & 0 & 57 & 0,0 \\
131 & VARGEM GRANDE & 2025 & 0 & 3 & 0 & 0 & 0 & 0 & 1 & 90 & 0 & 137 & 25 & 2 & 0,0 & 3,3 & -- & 0,0 & 0,0 & 0,0 & 3 & 255 & 1,2 \\
130 & VARGEM PEQUENA & 2025 & 0 & 0 & 0 & 3 & 3 & 1 & 0 & 133 & 0 & 153 & 35 & 4 & -- & 0,0 & -- & 2,0 & 8,6 & 25,0 & 7 & 325 & 2,1 \\
159 & VASCO DA GAMA & 2025 & 0 & 0 & 0 & 0 & 0 & 0 & 1 & 53 & 0 & 52 & 24 & 1 & 0,0 & 0,0 & -- & 0,0 & 0,0 & 0,0 & 0 & 131 & 0,0 \\
84 & VAZ LOBO & 2025 & 0 & 2 & 0 & 1 & 0 & 0 & 0 & 30 & 0 & 70 & 38 & 1 & -- & 6,7 & -- & 1,4 & 0,0 & 0,0 & 3 & 139 & 2,2 \\
73 & VICENTE DE CARVALHO & 2025 & 0 & 1 & 0 & 0 & 0 & 0 & 0 & 57 & 0 & 97 & 34 & 1 & -- & 1,8 & -- & 0,0 & 0,0 & 0,0 & 1 & 189 & 0,5 \\
30 & VIDIGAL & 2025 & 0 & 1 & 1 & 0 & 0 & 0 & 1 & 38 & 1 & 75 & 26 & 4 & 0,0 & 2,6 & 100,0 & 0,0 & 0,0 & 0,0 & 2 & 145 & 1,4 \\
48 & VIGARIO GERAL & 2025 & 0 & 0 & 0 & 3 & 1 & 0 & 3 & 87 & 0 & 173 & 90 & 1 & 0,0 & 0,0 & -- & 1,7 & 1,1 & 0,0 & 4 & 354 & 1,1 \\
74 & VILA DA PENHA & 2025 & 0 & 0 & 0 & 0 & 0 & 0 & 0 & 66 & 1 & 60 & 17 & 4 & -- & 0,0 & 0,0 & 0,0 & 0,0 & 0,0 & 0 & 148 & 0,0 \\
36 & VILA ISABEL & 2025 & 0 & 2 & 0 & 3 & 0 & 0 & 4 & 209 & 0 & 128 & 97 & 5 & 0,0 & 1,0 & -- & 2,3 & 0,0 & 0,0 & 5 & 443 & 1,1 \\
162 & VILA KENNEDY & 2025 & 0 & 0 & 0 & 4 & 0 & 0 & 0 & 37 & 1 & 109 & 46 & 1 & -- & 0,0 & 0,0 & 3,7 & 0,0 & 0,0 & 4 & 194 & 2,1 \\
72 & VILA KOSMOS & 2025 & 0 & 1 & 0 & 0 & 0 & 0 & 0 & 33 & 0 & 58 & 17 & 1 & -- & 3,0 & -- & 0,0 & 0,0 & 0,0 & 1 & 109 & 0,9 \\
135 & VILA MILITAR & 2025 & 0 & 1 & 0 & 1 & 1 & 0 & 0 & 42 & 0 & 45 & 20 & 1 & -- & 2,4 & -- & 2,2 & 5,0 & 0,0 & 3 & 108 & 2,8 \\
125 & VILA VALQUEIRE & 2025 & 0 & 0 & 0 & 1 & 0 & 1 & 1 & 124 & 1 & 102 & 44 & 1 & 0,0 & 0,0 & 0,0 & 1,0 & 0,0 & 100,0 & 2 & 273 & 0,7 \\
75 & VISTA ALEGRE & 2025 & 0 & 0 & 0 & 0 & 0 & 0 & 0 & 21 & 0 & 24 & 5 & 0 & -- & 0,0 & -- & 0,0 & 0,0 & -- & 0 & 50 & 0,0 \\
92 & ZUMBI & 2025 & 0 & 0 & 0 & 0 & 0 & 0 & 0 & 12 & 0 & 7 & 1 & 0 & -- & 0,0 & -- & 0,0 & 0,0 & -- & 0 & 20 & 0,0 \\
\end{longtable}}
\vspace{-6pt}\fonte{Elaboração IPP com dados de \citeonline{sms_rio_sim}.}
\end{landscape}


\begin{landscape}
{\scriptsize\setlength{\tabcolsep}{3pt}
\begin{longtable}{@{}rrlrrrrr@{}}
\caption{Mortalidade infantil por causas evitáveis (0 a 6, 7 a 27, 28 a 364 dias) (mortalidade evitaveis cap 2025)}\label{tab:mortalidade-evitaveis-cap-2025}\\
\toprule \textbf{Cod ap sms} & \textbf{Ano} & \textbf{Faixa etaria} & \textbf{Evitaveis} & \textbf{Mal definidas} & \textbf{Demais} & \textbf{Total} & \textbf{Evitaveis (\%)} \\ \midrule \endfirsthead
\multicolumn{8}{@{}l}{\footnotesize\itshape (continuação)}\\ \toprule \textbf{Cod ap sms} & \textbf{Ano} & \textbf{Faixa etaria} & \textbf{Evitaveis} & \textbf{Mal definidas} & \textbf{Demais} & \textbf{Total} & \textbf{Evitaveis (\%)} \\ \midrule \endhead
\midrule \multicolumn{8}{r@{}}{\footnotesize\itshape (continua)}\\ \endfoot
\bottomrule \endlastfoot
1,0 & 2025 & de 1 a 4 anos & 1 & 0 & 4 & 5 & 20,0 \\
1,0 & 2025 & menores de 1 ano & 27 & 0 & 11 & 38 & 71,0 \\
1,0 & 2025 & menores de 5 anos & 28 & 0 & 15 & 43 & 65,1 \\
2,1 & 2025 & de 1 a 4 anos & 2 & 0 & 1 & 3 & 66,7 \\
2,1 & 2025 & menores de 1 ano & 12 & 4 & 16 & 32 & 37,5 \\
2,1 & 2025 & menores de 5 anos & 14 & 4 & 17 & 35 & 40,0 \\
2,2 & 2025 & de 1 a 4 anos & 6 & 0 & 0 & 6 & 100,0 \\
2,2 & 2025 & menores de 1 ano & 9 & 0 & 5 & 14 & 64,3 \\
2,2 & 2025 & menores de 5 anos & 15 & 0 & 5 & 20 & 75,0 \\
3,1 & 2025 & de 1 a 4 anos & 8 & 1 & 9 & 18 & 44,4 \\
3,1 & 2025 & menores de 1 ano & 64 & 3 & 30 & 97 & 66,0 \\
3,1 & 2025 & menores de 5 anos & 72 & 4 & 39 & 115 & 62,6 \\
3,2 & 2025 & de 1 a 4 anos & 4 & 2 & 9 & 15 & 26,7 \\
3,2 & 2025 & menores de 1 ano & 34 & 0 & 18 & 52 & 65,4 \\
3,2 & 2025 & menores de 5 anos & 38 & 2 & 27 & 67 & 56,7 \\
3,3 & 2025 & de 1 a 4 anos & 12 & 1 & 7 & 20 & 60,0 \\
3,3 & 2025 & menores de 1 ano & 84 & 2 & 25 & 111 & 75,7 \\
3,3 & 2025 & menores de 5 anos & 96 & 3 & 32 & 131 & 73,3 \\
4,0 & 2025 & de 1 a 4 anos & 7 & 2 & 9 & 18 & 38,9 \\
4,0 & 2025 & menores de 1 ano & 89 & 4 & 39 & 132 & 67,4 \\
4,0 & 2025 & menores de 5 anos & 96 & 6 & 48 & 150 & 64,0 \\
5,1 & 2025 & de 1 a 4 anos & 9 & 1 & 6 & 16 & 56,2 \\
5,1 & 2025 & menores de 1 ano & 48 & 1 & 23 & 72 & 66,7 \\
5,1 & 2025 & menores de 5 anos & 57 & 2 & 29 & 88 & 64,8 \\
5,2 & 2025 & de 1 a 4 anos & 7 & 0 & 8 & 15 & 46,7 \\
5,2 & 2025 & menores de 1 ano & 67 & 0 & 33 & 100 & 67,0 \\
5,2 & 2025 & menores de 5 anos & 74 & 0 & 41 & 115 & 64,3 \\
5,3 & 2025 & de 1 a 4 anos & 10 & 0 & 6 & 16 & 62,5 \\
5,3 & 2025 & menores de 1 ano & 57 & 1 & 18 & 76 & 75,0 \\
5,3 & 2025 & menores de 5 anos & 67 & 1 & 24 & 92 & 72,8 \\
\end{longtable}}
\vspace{-6pt}\fonte{Elaboração IPP com dados de \citeonline{sms_rio_sim}.}
\end{landscape}


{\small\setlength{\tabcolsep}{5pt}
\begin{longtable}{@{}rrrr@{}}
\caption{Mortalidade infantil por causas evitáveis (0 a 6, 7 a 27, 28 a 364 dias) (taxa mortalidade evitaveis menores 5 municipio ano)}\label{tab:taxa-mortalidade-evitaveis-menores-5-municipio-ano}\\
\toprule \textbf{Ano} & \textbf{Obitos} & \textbf{Nascidos vivos} & \textbf{Taxa por mil} \\ \midrule \endfirsthead
\multicolumn{4}{@{}l}{\footnotesize\itshape (continuação)}\\ \toprule \textbf{Ano} & \textbf{Obitos} & \textbf{Nascidos vivos} & \textbf{Taxa por mil} \\ \midrule \endhead
\midrule \multicolumn{4}{r@{}}{\footnotesize\itshape (continua)}\\ \endfoot
\bottomrule \endlastfoot
2006 & 1.353 & 82.101 & 16,5 \\
2007 & 1.286 & 82.019 & 15,7 \\
2008 & 1.310 & 82.341 & 15,9 \\
2009 & 1.349 & 84.415 & 16,0 \\
2010 & 1.304 & 83.257 & 15,7 \\
2011 & 1.249 & 85.939 & 14,5 \\
2012 & 1.278 & 86.377 & 14,8 \\
2013 & 1.279 & 87.474 & 14,6 \\
2014 & 1.178 & 89.923 & 13,1 \\
2015 & 1.252 & 90.539 & 13,8 \\
2016 & 1.241 & 83.057 & 14,9 \\
2017 & 1.135 & 84.486 & 13,4 \\
2018 & 1.119 & 82.485 & 13,6 \\
2019 & 1.099 & 76.574 & 14,4 \\
2020 & 1.008 & 72.820 & 13,8 \\
2021 & 972 & 68.611 & 14,2 \\
2022 & 923 & 64.635 & 14,3 \\
2023 & 938 & 62.645 & 15,0 \\
2024 & 820 & 57.443 & 14,3 \\
2025 & 856 & 58.694 & 14,6 \\
\end{longtable}}
\vspace{-6pt}\fonte{Elaboração IPP com dados de \citeonline{sms_rio_sim}.}


{\small\setlength{\tabcolsep}{5pt}
\begin{longtable}{@{}rrrrrr@{}}
\caption{Taxa de mortalidade na primeira infância (menores de 1, 1-4, 0-5 anos) (mortalidade infantil pos neonatal total por ano)}\label{tab:mortalidade-infantil-pos-neonatal-total-por-ano}\\
\toprule \textbf{Ano} & \textbf{Obitos 0 364} & \textbf{Obitos 28 364} & \textbf{Nascidos vivos} & \textbf{Taxa mortalidade infantil} & \textbf{Taxa mortalidade pos neonatal} \\ \midrule \endfirsthead
\multicolumn{6}{@{}l}{\footnotesize\itshape (continuação)}\\ \toprule \textbf{Ano} & \textbf{Obitos 0 364} & \textbf{Obitos 28 364} & \textbf{Nascidos vivos} & \textbf{Taxa mortalidade infantil} & \textbf{Taxa mortalidade pos neonatal} \\ \midrule \endhead
\midrule \multicolumn{6}{r@{}}{\footnotesize\itshape (continua)}\\ \endfoot
\bottomrule \endlastfoot
2006 & 1.097 & 376 & 82.068 & 13,4 & 4,6 \\
2007 & 1.086 & 402 & 81.911 & 13,3 & 4,9 \\
2008 & 1.111 & 409 & 82.182 & 13,5 & 5,0 \\
2009 & 1.145 & 421 & 84.242 & 13,6 & 5,0 \\
2010 & 1.082 & 387 & 83.092 & 13,0 & 4,7 \\
2011 & 1.071 & 419 & 85.725 & 12,5 & 4,9 \\
2012 & 1.096 & 396 & 86.021 & 12,7 & 4,6 \\
2013 & 1.109 & 377 & 87.423 & 12,7 & 4,3 \\
2014 & 1.017 & 346 & 89.760 & 11,3 & 3,9 \\
2015 & 1.097 & 377 & 90.303 & 12,1 & 4,2 \\
2016 & 1.060 & 372 & 82.854 & 12,8 & 4,5 \\
2017 & 950 & 316 & 84.339 & 11,3 & 3,7 \\
2018 & 965 & 323 & 82.484 & 11,7 & 3,9 \\
2019 & 931 & 319 & 76.477 & 12,2 & 4,2 \\
2020 & 882 & 268 & 72.414 & 12,2 & 3,7 \\
2021 & 839 & 272 & 68.412 & 12,3 & 4,0 \\
2022 & 805 & 292 & 64.690 & 12,4 & 4,5 \\
2023 & 818 & 289 & 62.787 & 13,0 & 4,6 \\
2024 & 751 & 268 & 57.632 & 13,0 & 4,7 \\
2025 & 773 & 265 & 59.171 & 13,1 & 4,5 \\
\end{longtable}}
\vspace{-6pt}\fonte{Elaboração IPP com dados de \citeonline{sms_rio_sim}.}


\begin{landscape}
{\scriptsize\setlength{\tabcolsep}{3pt}
\begin{longtable}{@{}rlrrrrrrrr@{}}
\caption{Taxa de mortalidade na primeira infância (menores de 1, 1-4, 0-5 anos) (mortalidade infantil 2025)}\label{tab:tabela-mapa-mortalidade-infantil-2025}\\
\toprule \textbf{Codigo} & \textbf{Bairro} & \textbf{Ano} & \textbf{Obitos 0 364} & \textbf{Obitos 0 6} & \textbf{Obitos 7 27} & \textbf{Nascidos vivos} & \textbf{Obitos 28 364} & \textbf{Taxa mortalidade infantil} & \textbf{Taxa mortalidade pos neonatal} \\ \midrule \endfirsthead
\multicolumn{10}{@{}l}{\footnotesize\itshape (continuação)}\\ \toprule \textbf{Codigo} & \textbf{Bairro} & \textbf{Ano} & \textbf{Obitos 0 364} & \textbf{Obitos 0 6} & \textbf{Obitos 7 27} & \textbf{Nascidos vivos} & \textbf{Obitos 28 364} & \textbf{Taxa mortalidade infantil} & \textbf{Taxa mortalidade pos neonatal} \\ \midrule \endhead
\midrule \multicolumn{10}{r@{}}{\footnotesize\itshape (continua)}\\ \endfoot
\bottomrule \endlastfoot
1 & SAUDE & 2025 & 0 & 0 & 0 & 36 & 0 & 0,0 & 0,0 \\
2 & GAMBOA & 2025 & 2 & 2 & 0 & 150 & 0 & 13,3 & 0,0 \\
3 & SANTO CRISTO & 2025 & 3 & 1 & 1 & 129 & 1 & 23,3 & 7,8 \\
4 & CAJU & 2025 & 9 & 3 & 3 & 305 & 3 & 29,5 & 9,8 \\
5 & CENTRO & 2025 & 4 & 1 & 2 & 194 & 1 & 20,6 & 5,2 \\
6 & CATUMBI & 2025 & 5 & 4 & 1 & 134 & 0 & 37,3 & 0,0 \\
7 & RIO COMPRIDO & 2025 & 5 & 2 & 2 & 349 & 1 & 14,3 & 2,9 \\
8 & CIDADE NOVA & 2025 & 3 & 0 & 1 & 42 & 2 & 71,4 & 47,6 \\
9 & ESTACIO & 2025 & 4 & 3 & 0 & 153 & 1 & 26,1 & 6,5 \\
10 & SAO CRISTOVAO & 2025 & 3 & 0 & 0 & 290 & 3 & 10,3 & 10,3 \\
11 & MANGUEIRA & 2025 & 3 & 2 & 1 & 199 & 0 & 15,1 & 0,0 \\
12 & BENFICA & 2025 & 8 & 3 & 1 & 346 & 4 & 23,1 & 11,6 \\
13 & PAQUETA & 2025 & 0 & 0 & 0 & 30 & 0 & 0,0 & 0,0 \\
14 & SANTA TERESA & 2025 & 4 & 2 & 0 & 316 & 2 & 12,7 & 6,3 \\
15 & FLAMENGO & 2025 & 4 & 3 & 1 & 274 & 0 & 14,6 & 0,0 \\
16 & GLORIA & 2025 & 0 & 0 & 0 & 47 & 0 & 0,0 & 0,0 \\
17 & LARANJEIRAS & 2025 & 1 & 0 & 0 & 206 & 1 & 4,9 & 4,9 \\
18 & CATETE & 2025 & 0 & 0 & 0 & 145 & 0 & 0,0 & 0,0 \\
19 & COSME VELHO & 2025 & 2 & 1 & 0 & 69 & 1 & 29,0 & 14,5 \\
20 & BOTAFOGO & 2025 & 2 & 0 & 0 & 534 & 2 & 3,7 & 3,7 \\
21 & HUMAITA & 2025 & 1 & 0 & 0 & 69 & 1 & 14,5 & 14,5 \\
22 & URCA & 2025 & 0 & 0 & 0 & 57 & 0 & 0,0 & 0,0 \\
23 & LEME & 2025 & 1 & 1 & 0 & 94 & 0 & 10,6 & 0,0 \\
24 & COPACABANA & 2025 & 7 & 3 & 2 & 772 & 2 & 9,1 & 2,6 \\
25 & IPANEMA & 2025 & 6 & 5 & 1 & 247 & 0 & 24,3 & 0,0 \\
26 & LEBLON & 2025 & 2 & 0 & 0 & 275 & 2 & 7,3 & 7,3 \\
27 & LAGOA & 2025 & 3 & 1 & 0 & 123 & 2 & 24,4 & 16,3 \\
28 & JARDIM BOTANICO & 2025 & 0 & 0 & 0 & 86 & 0 & 0,0 & 0,0 \\
29 & GAVEA & 2025 & 2 & 1 & 0 & 139 & 1 & 14,4 & 7,2 \\
30 & VIDIGAL & 2025 & 2 & 1 & 0 & 145 & 1 & 13,8 & 6,9 \\
31 & SAO CONRADO & 2025 & 0 & 0 & 0 & 85 & 0 & 0,0 & 0,0 \\
32 & PRACA DA BANDEIRA & 2025 & 0 & 0 & 0 & 57 & 0 & 0,0 & 0,0 \\
33 & TIJUCA & 2025 & 8 & 3 & 3 & 941 & 2 & 8,5 & 2,1 \\
34 & ALTO DA BOA VISTA & 2025 & 2 & 1 & 0 & 67 & 1 & 29,9 & 14,9 \\
35 & MARACANA & 2025 & 0 & 0 & 0 & 131 & 0 & 0,0 & 0,0 \\
36 & VILA ISABEL & 2025 & 5 & 2 & 3 & 443 & 0 & 11,3 & 0,0 \\
37 & ANDARAI & 2025 & 2 & 1 & 1 & 209 & 0 & 9,6 & 0,0 \\
38 & GRAJAU & 2025 & 0 & 0 & 0 & 168 & 0 & 0,0 & 0,0 \\
39 & MANGUINHOS & 2025 & 8 & 4 & 1 & 387 & 3 & 20,7 & 7,8 \\
40 & BONSUCESSO & 2025 & 3 & 0 & 0 & 392 & 3 & 7,7 & 7,7 \\
41 & RAMOS & 2025 & 4 & 3 & 1 & 479 & 0 & 8,4 & 0,0 \\
42 & OLARIA & 2025 & 6 & 3 & 0 & 416 & 3 & 14,4 & 7,2 \\
43 & PENHA & 2025 & 7 & 4 & 1 & 609 & 2 & 11,5 & 3,3 \\
44 & PENHA CIRCULAR & 2025 & 8 & 4 & 1 & 356 & 3 & 22,5 & 8,4 \\
45 & BRAS DE PINA & 2025 & 9 & 3 & 3 & 364 & 3 & 24,7 & 8,2 \\
46 & CORDOVIL & 2025 & 10 & 5 & 3 & 316 & 2 & 31,6 & 6,3 \\
47 & PARADA DE LUCAS & 2025 & 4 & 1 & 1 & 255 & 2 & 15,7 & 7,8 \\
48 & VIGARIO GERAL & 2025 & 4 & 3 & 0 & 354 & 1 & 11,3 & 2,8 \\
49 & JARDIM AMERICA & 2025 & 2 & 1 & 0 & 266 & 1 & 7,5 & 3,8 \\
50 & HIGIENOPOLIS & 2025 & 0 & 0 & 0 & 118 & 0 & 0,0 & 0,0 \\
51 & JACARE & 2025 & 0 & 0 & 0 & 166 & 0 & 0,0 & 0,0 \\
52 & MARIA DA GRACA & 2025 & 1 & 0 & 0 & 100 & 1 & 10,0 & 10,0 \\
53 & DEL CASTILHO & 2025 & 3 & 2 & 0 & 236 & 1 & 12,7 & 4,2 \\
54 & INHAUMA & 2025 & 6 & 3 & 1 & 539 & 2 & 11,1 & 3,7 \\
55 & ENGENHO DA RAINHA & 2025 & 3 & 2 & 1 & 217 & 0 & 13,8 & 0,0 \\
56 & TOMAS COELHO & 2025 & 6 & 5 & 0 & 191 & 1 & 31,4 & 5,2 \\
57 & SAO FRANCISCO XAVIER & 2025 & 1 & 1 & 0 & 55 & 0 & 18,2 & 0,0 \\
58 & ROCHA & 2025 & 5 & 3 & 1 & 178 & 1 & 28,1 & 5,6 \\
59 & RIACHUELO & 2025 & 2 & 0 & 1 & 66 & 1 & 30,3 & 15,2 \\
60 & SAMPAIO & 2025 & 3 & 2 & 0 & 99 & 1 & 30,3 & 10,1 \\
61 & ENGENHO NOVO & 2025 & 3 & 1 & 2 & 346 & 0 & 8,7 & 0,0 \\
62 & LINS DE VASCONCELOS & 2025 & 3 & 2 & 0 & 304 & 1 & 9,9 & 3,3 \\
63 & MEIER & 2025 & 6 & 2 & 2 & 247 & 2 & 24,3 & 8,1 \\
64 & TODOS OS SANTOS & 2025 & 1 & 1 & 0 & 126 & 0 & 7,9 & 0,0 \\
65 & CACHAMBI & 2025 & 3 & 0 & 2 & 274 & 1 & 10,9 & 3,6 \\
66 & ENGENHO DE DENTRO & 2025 & 3 & 0 & 0 & 351 & 3 & 8,5 & 8,5 \\
67 & AGUA SANTA & 2025 & 1 & 1 & 0 & 57 & 0 & 17,5 & 0,0 \\
68 & ENCANTADO & 2025 & 1 & 1 & 0 & 81 & 0 & 12,3 & 0,0 \\
69 & PIEDADE & 2025 & 3 & 2 & 0 & 328 & 1 & 9,1 & 3,0 \\
70 & ABOLICAO & 2025 & 2 & 1 & 0 & 109 & 1 & 18,3 & 9,2 \\
71 & PILARES & 2025 & 2 & 1 & 0 & 191 & 1 & 10,5 & 5,2 \\
72 & VILA KOSMOS & 2025 & 1 & 0 & 0 & 109 & 1 & 9,2 & 9,2 \\
73 & VICENTE DE CARVALHO & 2025 & 1 & 1 & 0 & 189 & 0 & 5,3 & 0,0 \\
74 & VILA DA PENHA & 2025 & 0 & 0 & 0 & 148 & 0 & 0,0 & 0,0 \\
75 & VISTA ALEGRE & 2025 & 0 & 0 & 0 & 50 & 0 & 0,0 & 0,0 \\
76 & IRAJA & 2025 & 9 & 3 & 1 & 683 & 5 & 13,2 & 7,3 \\
77 & COLEGIO & 2025 & 6 & 4 & 0 & 292 & 2 & 20,5 & 6,8 \\
78 & CAMPINHO & 2025 & 2 & 1 & 1 & 95 & 0 & 21,1 & 0,0 \\
79 & QUINTINO BOCAIUVA & 2025 & 1 & 0 & 0 & 202 & 1 & 5,0 & 5,0 \\
80 & CAVALCANTE & 2025 & 2 & 1 & 1 & 145 & 0 & 13,8 & 0,0 \\
81 & ENGENHEIRO LEAL & 2025 & 1 & 1 & 0 & 51 & 0 & 19,6 & 0,0 \\
82 & CASCADURA & 2025 & 4 & 3 & 1 & 243 & 0 & 16,5 & 0,0 \\
83 & MADUREIRA & 2025 & 5 & 1 & 2 & 378 & 2 & 13,2 & 5,3 \\
84 & VAZ LOBO & 2025 & 3 & 0 & 1 & 139 & 2 & 21,6 & 14,4 \\
85 & TURIACU & 2025 & 1 & 0 & 0 & 136 & 1 & 7,4 & 7,4 \\
86 & ROCHA MIRANDA & 2025 & 5 & 3 & 0 & 324 & 2 & 15,4 & 6,2 \\
87 & HONORIO GURGEL & 2025 & 6 & 3 & 1 & 183 & 2 & 32,8 & 10,9 \\
88 & OSWALDO CRUZ & 2025 & 2 & 0 & 0 & 196 & 2 & 10,2 & 10,2 \\
89 & BENTO RIBEIRO & 2025 & 5 & 4 & 1 & 300 & 0 & 16,7 & 0,0 \\
90 & MARECHAL HERMES & 2025 & 7 & 4 & 2 & 334 & 1 & 21,0 & 3,0 \\
91 & RIBEIRA & 2025 & 1 & 1 & 0 & 23 & 0 & 43,5 & 0,0 \\
92 & ZUMBI & 2025 & 0 & 0 & 0 & 20 & 0 & 0,0 & 0,0 \\
93 & CACUIA & 2025 & 0 & 0 & 0 & 85 & 0 & 0,0 & 0,0 \\
94 & PITANGUEIRAS & 2025 & 2 & 0 & 0 & 76 & 2 & 26,3 & 26,3 \\
95 & PRAIA DA BANDEIRA & 2025 & 0 & 0 & 0 & 37 & 0 & 0,0 & 0,0 \\
96 & COCOTA & 2025 & 0 & 0 & 0 & 53 & 0 & 0,0 & 0,0 \\
97 & BANCARIOS & 2025 & 1 & 0 & 0 & 132 & 1 & 7,6 & 7,6 \\
98 & FREGUESIA-ILHA & 2025 & 2 & 0 & 0 & 191 & 2 & 10,5 & 10,5 \\
99 & JARDIM GUANABARA & 2025 & 0 & 0 & 0 & 170 & 0 & 0,0 & 0,0 \\
100 & JARDIM CARIOCA & 2025 & 3 & 2 & 0 & 199 & 1 & 15,1 & 5,0 \\
101 & TAUA & 2025 & 2 & 0 & 1 & 212 & 1 & 9,4 & 4,7 \\
102 & MONERO & 2025 & 0 & 0 & 0 & 39 & 0 & 0,0 & 0,0 \\
103 & PORTUGUESA & 2025 & 2 & 2 & 0 & 127 & 0 & 15,7 & 0,0 \\
104 & GALEAO & 2025 & 1 & 0 & 0 & 248 & 1 & 4,0 & 4,0 \\
105 & CIDADE UNIVERSITARIA & 2025 & 1 & 1 & 0 & 17 & 0 & 58,8 & 0,0 \\
106 & GUADALUPE & 2025 & 5 & 2 & 0 & 355 & 3 & 14,1 & 8,5 \\
107 & ANCHIETA & 2025 & 7 & 5 & 0 & 618 & 2 & 11,3 & 3,2 \\
108 & PARQUE ANCHIETA & 2025 & 2 & 1 & 0 & 144 & 1 & 13,9 & 6,9 \\
109 & RICARDO ALBUQUERQUE & 2025 & 2 & 1 & 0 & 262 & 1 & 7,6 & 3,8 \\
110 & COELHO NETO & 2025 & 0 & 0 & 0 & 284 & 0 & 0,0 & 0,0 \\
111 & ACARI & 2025 & 6 & 2 & 2 & 349 & 2 & 17,2 & 5,7 \\
112 & BARROS FILHO & 2025 & 2 & 1 & 0 & 164 & 1 & 12,2 & 6,1 \\
113 & COSTA BARROS & 2025 & 7 & 4 & 0 & 401 & 3 & 17,5 & 7,5 \\
114 & PAVUNA & 2025 & 20 & 14 & 0 & 888 & 6 & 22,5 & 6,8 \\
115 & JACAREPAGUA & 2025 & 32 & 11 & 8 & 1.878 & 13 & 17,0 & 6,9 \\
116 & ANIL & 2025 & 7 & 5 & 0 & 377 & 2 & 18,6 & 5,3 \\
117 & GARDENIA AZUL & 2025 & 4 & 1 & 1 & 400 & 2 & 10,0 & 5,0 \\
118 & CIDADE DE DEUS & 2025 & 6 & 2 & 2 & 549 & 2 & 10,9 & 3,6 \\
119 & CURICICA & 2025 & 5 & 4 & 1 & 369 & 0 & 13,6 & 0,0 \\
120 & FREGUESIA-JPA & 2025 & 4 & 2 & 1 & 553 & 1 & 7,2 & 1,8 \\
121 & PECHINCHA & 2025 & 3 & 2 & 1 & 339 & 0 & 8,8 & 0,0 \\
122 & TAQUARA & 2025 & 14 & 8 & 2 & 1.005 & 4 & 13,9 & 4,0 \\
123 & TANQUE & 2025 & 5 & 1 & 3 & 378 & 1 & 13,2 & 2,6 \\
124 & PRACA SECA & 2025 & 10 & 4 & 1 & 567 & 5 & 17,6 & 8,8 \\
125 & VILA VALQUEIRE & 2025 & 2 & 2 & 0 & 273 & 0 & 7,3 & 0,0 \\
126 & JOA & 2025 & 0 & 0 & 0 & 8 & 0 & 0,0 & 0,0 \\
127 & ITANHANGA & 2025 & 11 & 4 & 2 & 791 & 5 & 13,9 & 6,3 \\
128 & BARRA DA TIJUCA & 2025 & 6 & 4 & 2 & 1.115 & 0 & 5,4 & 0,0 \\
129 & CAMORIM & 2025 & 1 & 0 & 0 & 35 & 1 & 28,6 & 28,6 \\
130 & VARGEM PEQUENA & 2025 & 7 & 4 & 1 & 325 & 2 & 21,5 & 6,2 \\
131 & VARGEM GRANDE & 2025 & 3 & 1 & 0 & 255 & 2 & 11,8 & 7,8 \\
132 & RECREIO BANDEIRANTES & 2025 & 9 & 3 & 1 & 1.245 & 5 & 7,2 & 4,0 \\
133 & GRUMARI & 2025 & 0 & 0 & 0 & 0 & 0 & -- & -- \\
134 & DEODORO & 2025 & 2 & 2 & 0 & 124 & 0 & 16,1 & 0,0 \\
135 & VILA MILITAR & 2025 & 3 & 2 & 1 & 108 & 0 & 27,8 & 0,0 \\
136 & CAMPO DOS AFONSOS & 2025 & 2 & 2 & 0 & 36 & 0 & 55,6 & 0,0 \\
137 & JARDIM SULACAP & 2025 & 1 & 1 & 0 & 132 & 0 & 7,6 & 0,0 \\
138 & MAGALHAES BASTOS & 2025 & 6 & 2 & 2 & 208 & 2 & 28,8 & 9,6 \\
139 & REALENGO & 2025 & 17 & 5 & 5 & 1.463 & 7 & 11,6 & 4,8 \\
140 & PADRE MIGUEL & 2025 & 10 & 5 & 1 & 522 & 4 & 19,2 & 7,7 \\
141 & BANGU & 2025 & 20 & 7 & 3 & 1.977 & 10 & 10,1 & 5,1 \\
142 & SENADOR CAMARA & 2025 & 8 & 2 & 3 & 995 & 3 & 8,0 & 3,0 \\
143 & SANTISSIMO & 2025 & 8 & 5 & 2 & 471 & 1 & 17,0 & 2,1 \\
144 & CAMPO GRANDE & 2025 & 40 & 16 & 12 & 3.504 & 12 & 11,4 & 3,4 \\
145 & SENADOR VASCONCELOS & 2025 & 8 & 5 & 0 & 298 & 3 & 26,8 & 10,1 \\
146 & INHOAIBA & 2025 & 7 & 3 & 1 & 752 & 3 & 9,3 & 4,0 \\
147 & COSMOS & 2025 & 21 & 10 & 3 & 954 & 8 & 22,0 & 8,4 \\
148 & PACIENCIA & 2025 & 14 & 9 & 3 & 1.157 & 2 & 12,1 & 1,7 \\
149 & SANTA CRUZ & 2025 & 53 & 28 & 12 & 3.186 & 13 & 16,6 & 4,1 \\
150 & SEPETIBA & 2025 & 9 & 2 & 2 & 753 & 5 & 12,0 & 6,6 \\
151 & GUARATIBA & 2025 & 13 & 6 & 3 & 1.867 & 4 & 7,0 & 2,1 \\
152 & BARRA DE GUARATIBA & 2025 & 2 & 0 & 1 & 117 & 1 & 17,1 & 8,5 \\
153 & PEDRA DE GUARATIBA & 2025 & 1 & 0 & 0 & 228 & 1 & 4,4 & 4,4 \\
154 & ROCINHA & 2025 & 9 & 3 & 1 & 696 & 5 & 12,9 & 7,2 \\
155 & JACAREZINHO & 2025 & 4 & 3 & 1 & 293 & 0 & 13,7 & 0,0 \\
156 & COMPLEXO DO ALEMAO & 2025 & 6 & 0 & 1 & 361 & 5 & 16,6 & 13,9 \\
157 & COMPLEXO DA MARE & 2025 & 15 & 6 & 5 & 1.494 & 4 & 10,0 & 2,7 \\
158 & PARQUE COLUMBIA & 2025 & 3 & 0 & 1 & 83 & 2 & 36,1 & 24,1 \\
159 & VASCO DA GAMA & 2025 & 0 & 0 & 0 & 131 & 0 & 0,0 & 0,0 \\
160 & GERICINo & 2025 & 1 & 1 & 0 & 14 & 0 & 71,4 & 0,0 \\
161 & LAPA & 2025 & 0 & 0 & 0 & 39 & 0 & 0,0 & 0,0 \\
162 & VILA KENNEDY & 2025 & 4 & 2 & 1 & 194 & 1 & 20,6 & 5,2 \\
163 & JABOUR & 2025 & 0 & 0 & 0 & 17 & 0 & 0,0 & 0,0 \\
164 & ILHA DE GUARATIBA & 2025 & 1 & 1 & 0 & 81 & 0 & 12,3 & 0,0 \\
165 & BARRA OLIMPICA & 2025 & 3 & 1 & 1 & 503 & 1 & 6,0 & 2,0 \\
998 & ILHA DO GOVERNADOR - IGNORADO & 2025 & 0 & 0 & 0 & 12 & 0 & 0,0 & 0,0 \\
999 & IGNORADO BAIRRO-MRJ & 2025 & 0 & 0 & 0 & 19 & 0 & 0,0 & 0,0 \\
\end{longtable}}
\vspace{-6pt}\fonte{Elaboração IPP com dados de \citeonline{sms_rio_sim}.}
\end{landscape}


{\small\setlength{\tabcolsep}{5pt}
\begin{longtable}{@{}rrrr@{}}
\caption{Mortalidade infantil por causas evitáveis, por tipo de causa (grupo/subgrupo CID-10) (mortalidade causas evitaveis grupo ano)}\label{tab:mortalidade-causas-evitaveis-grupo-ano}\\
\toprule \textbf{Ano} & \textbf{1. Causas evitáveis} & \textbf{2. Causas mal definidas} & \textbf{3. Demais causas (não claramente evitáveis)} \\ \midrule \endfirsthead
\multicolumn{4}{@{}l}{\footnotesize\itshape (continuação)}\\ \toprule \textbf{Ano} & \textbf{1. Causas evitáveis} & \textbf{2. Causas mal definidas} & \textbf{3. Demais causas (não claramente evitáveis)} \\ \midrule \endhead
\midrule \multicolumn{4}{r@{}}{\footnotesize\itshape (continua)}\\ \endfoot
\bottomrule \endlastfoot
1996 & 1.599 & 163 & 376 \\
1997 & 1.542 & 131 & 379 \\
1998 & 1.295 & 153 & 339 \\
1999 & 1.295 & 134 & 395 \\
2000 & 1.171 & 146 & 340 \\
2001 & 992 & 87 & 332 \\
2002 & 923 & 100 & 323 \\
2003 & 910 & 110 & 362 \\
2004 & 827 & 107 & 405 \\
2005 & 790 & 79 & 301 \\
2006 & 779 & 65 & 284 \\
2007 & 728 & 49 & 301 \\
2008 & 772 & 44 & 315 \\
2009 & 835 & 42 & 293 \\
2010 & 727 & 54 & 316 \\
2011 & 711 & 48 & 322 \\
2012 & 768 & 42 & 310 \\
2013 & 782 & 39 & 294 \\
2014 & 685 & 34 & 303 \\
2015 & 763 & 38 & 295 \\
2016 & 733 & 33 & 301 \\
2017 & 640 & 30 & 278 \\
2018 & 665 & 25 & 278 \\
2019 & 640 & 25 & 268 \\
2020 & 603 & 27 & 255 \\
2021 & 592 & 26 & 222 \\
2022 & 518 & 14 & 235 \\
2023 & 527 & 10 & 234 \\
2024 & 469 & 16 & 224 \\
2025 & 502 & 15 & 206 \\
\end{longtable}}
\vspace{-6pt}\fonte{Elaboração IPP com dados de \citeonline{sms_rio_sim}.}


\begin{landscape}
{\scriptsize\setlength{\tabcolsep}{3pt}
\begin{longtable}{@{}rrrrrrrr@{}}
\caption{Mortalidade infantil por causas evitáveis, por tipo de causa (grupo/subgrupo CID-10) (mortalidade causas evitaveis subgrupo ano)}\label{tab:mortalidade-causas-evitaveis-subgrupo-ano}\\
\toprule \textbf{Ano} & \textbf{1.1. Reduzível pelas ações de imunização} & \textbf{1.2.1 Reduzíveis atenção à mulher na gestação} & \textbf{1.2.2 Reduz por adequada atenção à mulher no parto} & \textbf{1.2.3 Reduzíveis adequada atenção ao recém-nascido} & \textbf{1.3. Reduz ações diagnóstico e tratamento adequado} & \textbf{1.4. Reduz. ações promoção vinc. ações de atenção} & \textbf{2. Causas mal definidas} \\ \midrule \endfirsthead
\multicolumn{8}{@{}l}{\footnotesize\itshape (continuação)}\\ \toprule \textbf{Ano} & \textbf{1.1. Reduzível pelas ações de imunização} & \textbf{1.2.1 Reduzíveis atenção à mulher na gestação} & \textbf{1.2.2 Reduz por adequada atenção à mulher no parto} & \textbf{1.2.3 Reduzíveis adequada atenção ao recém-nascido} & \textbf{1.3. Reduz ações diagnóstico e tratamento adequado} & \textbf{1.4. Reduz. ações promoção vinc. ações de atenção} & \textbf{2. Causas mal definidas} \\ \midrule \endhead
\midrule \multicolumn{8}{r@{}}{\footnotesize\itshape (continua)}\\ \endfoot
\bottomrule \endlastfoot
1996 & 10 & 438 & 164 & 543 & 301 & 143 & 163 \\
1997 & 10 & 441 & 169 & 568 & 263 & 91 & 131 \\
1998 & 7 & 428 & 176 & 378 & 213 & 93 & 153 \\
1999 & 4 & 441 & 163 & 383 & 230 & 74 & 134 \\
2000 & 2 & 443 & 164 & 318 & 166 & 78 & 146 \\
2001 & 2 & 333 & 135 & 314 & 142 & 66 & 87 \\
2002 & 1 & 384 & 113 & 236 & 132 & 57 & 100 \\
2003 & 2 & 311 & 122 & 279 & 138 & 58 & 110 \\
2004 & 4 & 336 & 89 & 222 & 117 & 59 & 107 \\
2005 & 0 & 323 & 87 & 218 & 108 & 54 & 79 \\
2006 & 2 & 321 & 85 & 201 & 108 & 62 & 65 \\
2007 & 3 & 287 & 111 & 149 & 126 & 52 & 49 \\
2008 & 3 & 311 & 112 & 174 & 99 & 73 & 44 \\
2009 & 1 & 371 & 110 & 168 & 114 & 71 & 42 \\
2010 & 3 & 352 & 87 & 136 & 82 & 67 & 54 \\
2011 & 4 & 331 & 83 & 117 & 110 & 66 & 48 \\
2012 & 6 & 373 & 84 & 125 & 101 & 79 & 42 \\
2013 & 0 & 391 & 82 & 142 & 108 & 59 & 39 \\
2014 & 3 & 328 & 78 & 143 & 73 & 60 & 34 \\
2015 & 5 & 348 & 94 & 161 & 86 & 69 & 38 \\
2016 & 1 & 334 & 76 & 154 & 93 & 75 & 33 \\
2017 & 2 & 288 & 85 & 131 & 73 & 61 & 30 \\
2018 & 0 & 310 & 96 & 121 & 82 & 56 & 25 \\
2019 & 2 & 273 & 104 & 127 & 78 & 56 & 25 \\
2020 & 1 & 305 & 90 & 108 & 42 & 57 & 27 \\
2021 & 1 & 254 & 87 & 132 & 49 & 69 & 26 \\
2022 & 2 & 220 & 70 & 104 & 63 & 59 & 14 \\
2023 & 0 & 230 & 74 & 94 & 78 & 51 & 10 \\
2024 & 0 & 241 & 53 & 53 & 62 & 60 & 16 \\
2025 & 1 & 263 & 68 & 67 & 46 & 57 & 15 \\
\end{longtable}}
\vspace{-6pt}\fonte{Elaboração IPP com dados de \citeonline{sms_rio_sim}.}
\end{landscape}


{\small\setlength{\tabcolsep}{5pt}
\begin{longtable}{@{}rrrr@{}}
\caption{Mortalidade infantil por causas evitáveis, por tipo de causa (grupo/subgrupo CID-10) (mortalidade causas evitaveis grupo 0 a 6 dias ano)}\label{tab:mortalidade-causas-evitaveis-grupo-0-a-6-dias-ano}\\
\toprule \textbf{Ano} & \textbf{1. Causas evitáveis} & \textbf{2. Causas mal definidas} & \textbf{3. Demais causas (não claramente evitáveis)} \\ \midrule \endfirsthead
\multicolumn{4}{@{}l}{\footnotesize\itshape (continuação)}\\ \toprule \textbf{Ano} & \textbf{1. Causas evitáveis} & \textbf{2. Causas mal definidas} & \textbf{3. Demais causas (não claramente evitáveis)} \\ \midrule \endhead
\midrule \multicolumn{4}{r@{}}{\footnotesize\itshape (continua)}\\ \endfoot
\bottomrule \endlastfoot
1996 & 889 & 36 & 160 \\
1997 & 849 & 26 & 151 \\
1998 & 714 & 23 & 130 \\
1999 & 708 & 17 & 148 \\
2000 & 653 & 25 & 127 \\
2001 & 558 & 16 & 148 \\
2002 & 557 & 9 & 137 \\
2003 & 504 & 15 & 156 \\
2004 & 449 & 13 & 190 \\
2005 & 404 & 8 & 122 \\
2006 & 444 & 8 & 110 \\
2007 & 377 & 5 & 115 \\
2008 & 379 & 7 & 118 \\
2009 & 455 & 6 & 93 \\
2010 & 417 & 6 & 114 \\
2011 & 375 & 1 & 117 \\
2012 & 400 & 7 & 107 \\
2013 & 420 & 6 & 112 \\
2014 & 377 & 2 & 113 \\
2015 & 408 & 6 & 104 \\
2016 & 393 & 9 & 117 \\
2017 & 337 & 2 & 110 \\
2018 & 366 & 1 & 85 \\
2019 & 339 & 1 & 92 \\
2020 & 325 & 1 & 86 \\
2021 & 301 & 4 & 86 \\
2022 & 285 & 1 & 74 \\
2023 & 276 & 1 & 69 \\
2024 & 246 & 5 & 71 \\
2025 & 269 & 2 & 62 \\
\end{longtable}}
\vspace{-6pt}\fonte{Elaboração IPP com dados de \citeonline{sms_rio_sim}.}


\begin{landscape}
{\scriptsize\setlength{\tabcolsep}{3pt}
\begin{longtable}{@{}rrrrrrrr@{}}
\caption{Mortalidade infantil por causas evitáveis, por tipo de causa (grupo/subgrupo CID-10) (mortalidade causas evitaveis subgrupo 0 a 6 dias ano)}\label{tab:mortalidade-causas-evitaveis-subgrupo-0-a-6-dias-ano}\\
\toprule \textbf{Ano} & \textbf{1.1. Reduzível pelas ações de imunização} & \textbf{1.2.1 Reduzíveis atenção à mulher na gestação} & \textbf{1.2.2 Reduz por adequada atenção à mulher no parto} & \textbf{1.2.3 Reduzíveis adequada atenção ao recém-nascido} & \textbf{1.3. Reduz ações diagnóstico e tratamento adequado} & \textbf{1.4. Reduz. ações promoção vinc. ações de atenção} & \textbf{2. Causas mal definidas} \\ \midrule \endfirsthead
\multicolumn{8}{@{}l}{\footnotesize\itshape (continuação)}\\ \toprule \textbf{Ano} & \textbf{1.1. Reduzível pelas ações de imunização} & \textbf{1.2.1 Reduzíveis atenção à mulher na gestação} & \textbf{1.2.2 Reduz por adequada atenção à mulher no parto} & \textbf{1.2.3 Reduzíveis adequada atenção ao recém-nascido} & \textbf{1.3. Reduz ações diagnóstico e tratamento adequado} & \textbf{1.4. Reduz. ações promoção vinc. ações de atenção} & \textbf{2. Causas mal definidas} \\ \midrule \endhead
\midrule \multicolumn{8}{r@{}}{\footnotesize\itshape (continua)}\\ \endfoot
\bottomrule \endlastfoot
1996 & 1 & 355 & 152 & 370 & 4 & 7 & 36 \\
1997 & 1 & 342 & 153 & 343 & 3 & 7 & 26 \\
1998 & 0 & 316 & 154 & 236 & 2 & 6 & 23 \\
1999 & 0 & 336 & 132 & 233 & 4 & 3 & 17 \\
2000 & 0 & 322 & 132 & 196 & 1 & 2 & 25 \\
2001 & 0 & 262 & 107 & 186 & 2 & 1 & 16 \\
2002 & 0 & 303 & 93 & 158 & 1 & 2 & 9 \\
2003 & 0 & 243 & 94 & 166 & 0 & 1 & 15 \\
2004 & 0 & 255 & 70 & 123 & 1 & 0 & 13 \\
2005 & 0 & 227 & 58 & 116 & 0 & 3 & 8 \\
2006 & 0 & 255 & 68 & 117 & 3 & 1 & 8 \\
2007 & 0 & 211 & 88 & 76 & 1 & 1 & 5 \\
2008 & 0 & 214 & 80 & 85 & 0 & 0 & 7 \\
2009 & 0 & 269 & 89 & 96 & 0 & 1 & 6 \\
2010 & 0 & 260 & 74 & 81 & 0 & 2 & 6 \\
2011 & 0 & 241 & 68 & 65 & 1 & 0 & 1 \\
2012 & 0 & 271 & 57 & 67 & 2 & 3 & 7 \\
2013 & 0 & 281 & 53 & 86 & 0 & 0 & 6 \\
2014 & 0 & 238 & 53 & 85 & 1 & 0 & 2 \\
2015 & 0 & 242 & 70 & 96 & 0 & 0 & 6 \\
2016 & 0 & 235 & 58 & 100 & 0 & 0 & 9 \\
2017 & 0 & 213 & 62 & 58 & 2 & 2 & 2 \\
2018 & 0 & 229 & 66 & 70 & 1 & 0 & 1 \\
2019 & 0 & 191 & 84 & 63 & 1 & 0 & 1 \\
2020 & 0 & 199 & 63 & 62 & 0 & 1 & 1 \\
2021 & 0 & 172 & 60 & 68 & 0 & 1 & 4 \\
2022 & 0 & 167 & 56 & 59 & 2 & 1 & 1 \\
2023 & 0 & 168 & 56 & 50 & 1 & 1 & 1 \\
2024 & 0 & 174 & 42 & 30 & 0 & 0 & 5 \\
2025 & 0 & 183 & 50 & 36 & 0 & 0 & 2 \\
\end{longtable}}
\vspace{-6pt}\fonte{Elaboração IPP com dados de \citeonline{sms_rio_sim}.}
\end{landscape}


{\small\setlength{\tabcolsep}{5pt}
\begin{longtable}{@{}rrrr@{}}
\caption{Mortalidade infantil por causas evitáveis, por tipo de causa (grupo/subgrupo CID-10) (mortalidade causas evitaveis grupo 7 a 27 dias ano)}\label{tab:mortalidade-causas-evitaveis-grupo-7-a-27-dias-ano}\\
\toprule \textbf{Ano} & \textbf{1. Causas evitáveis} & \textbf{2. Causas mal definidas} & \textbf{3. Demais causas (não claramente evitáveis)} \\ \midrule \endfirsthead
\multicolumn{4}{@{}l}{\footnotesize\itshape (continuação)}\\ \toprule \textbf{Ano} & \textbf{1. Causas evitáveis} & \textbf{2. Causas mal definidas} & \textbf{3. Demais causas (não claramente evitáveis)} \\ \midrule \endhead
\midrule \multicolumn{4}{r@{}}{\footnotesize\itshape (continua)}\\ \endfoot
\bottomrule \endlastfoot
1996 & 249 & 19 & 57 \\
1997 & 284 & 12 & 54 \\
1998 & 254 & 14 & 46 \\
1999 & 270 & 10 & 63 \\
2000 & 232 & 13 & 61 \\
2001 & 179 & 8 & 54 \\
2002 & 161 & 14 & 45 \\
2003 & 179 & 15 & 60 \\
2004 & 166 & 7 & 66 \\
2005 & 189 & 8 & 54 \\
2006 & 134 & 7 & 38 \\
2007 & 130 & 6 & 42 \\
2008 & 155 & 8 & 52 \\
2009 & 134 & 7 & 47 \\
2010 & 110 & 6 & 56 \\
2011 & 112 & 3 & 52 \\
2012 & 147 & 4 & 50 \\
2013 & 151 & 3 & 45 \\
2014 & 125 & 4 & 55 \\
2015 & 146 & 7 & 47 \\
2016 & 128 & 3 & 40 \\
2017 & 140 & 5 & 38 \\
2018 & 128 & 2 & 62 \\
2019 & 133 & 1 & 48 \\
2020 & 139 & 10 & 54 \\
2021 & 134 & 2 & 42 \\
2022 & 90 & 0 & 42 \\
2023 & 118 & 2 & 40 \\
2024 & 86 & 3 & 48 \\
2025 & 99 & 0 & 42 \\
\end{longtable}}
\vspace{-6pt}\fonte{Elaboração IPP com dados de \citeonline{sms_rio_sim}.}


\begin{landscape}
{\scriptsize\setlength{\tabcolsep}{3pt}
\begin{longtable}{@{}rrrrrrrr@{}}
\caption{Mortalidade infantil por causas evitáveis, por tipo de causa (grupo/subgrupo CID-10) (mortalidade causas evitaveis subgrupo 7 a 27 dias ano)}\label{tab:mortalidade-causas-evitaveis-subgrupo-7-a-27-dias-ano}\\
\toprule \textbf{Ano} & \textbf{1.1. Reduzível pelas ações de imunização} & \textbf{1.2.1 Reduzíveis atenção à mulher na gestação} & \textbf{1.2.2 Reduz por adequada atenção à mulher no parto} & \textbf{1.2.3 Reduzíveis adequada atenção ao recém-nascido} & \textbf{1.3. Reduz ações diagnóstico e tratamento adequado} & \textbf{1.4. Reduz. ações promoção vinc. ações de atenção} & \textbf{2. Causas mal definidas} \\ \midrule \endfirsthead
\multicolumn{8}{@{}l}{\footnotesize\itshape (continuação)}\\ \toprule \textbf{Ano} & \textbf{1.1. Reduzível pelas ações de imunização} & \textbf{1.2.1 Reduzíveis atenção à mulher na gestação} & \textbf{1.2.2 Reduz por adequada atenção à mulher no parto} & \textbf{1.2.3 Reduzíveis adequada atenção ao recém-nascido} & \textbf{1.3. Reduz ações diagnóstico e tratamento adequado} & \textbf{1.4. Reduz. ações promoção vinc. ações de atenção} & \textbf{2. Causas mal definidas} \\ \midrule \endhead
\midrule \multicolumn{8}{r@{}}{\footnotesize\itshape (continua)}\\ \endfoot
\bottomrule \endlastfoot
1996 & 2 & 59 & 10 & 157 & 14 & 7 & 19 \\
1997 & 2 & 64 & 11 & 192 & 7 & 8 & 12 \\
1998 & 1 & 91 & 17 & 127 & 12 & 6 & 14 \\
1999 & 0 & 78 & 30 & 141 & 18 & 3 & 10 \\
2000 & 0 & 93 & 19 & 104 & 14 & 2 & 13 \\
2001 & 0 & 56 & 16 & 100 & 5 & 2 & 8 \\
2002 & 0 & 67 & 13 & 68 & 10 & 3 & 14 \\
2003 & 0 & 52 & 19 & 94 & 12 & 2 & 15 \\
2004 & 0 & 61 & 16 & 85 & 3 & 1 & 7 \\
2005 & 0 & 75 & 18 & 87 & 7 & 2 & 8 \\
2006 & 0 & 44 & 15 & 67 & 3 & 5 & 7 \\
2007 & 0 & 55 & 14 & 56 & 4 & 1 & 6 \\
2008 & 0 & 69 & 15 & 63 & 3 & 5 & 8 \\
2009 & 0 & 68 & 9 & 53 & 0 & 4 & 7 \\
2010 & 1 & 53 & 11 & 40 & 3 & 2 & 6 \\
2011 & 0 & 62 & 11 & 33 & 6 & 0 & 3 \\
2012 & 3 & 75 & 20 & 37 & 6 & 6 & 4 \\
2013 & 0 & 74 & 22 & 47 & 7 & 1 & 3 \\
2014 & 0 & 62 & 20 & 40 & 2 & 1 & 4 \\
2015 & 0 & 73 & 18 & 51 & 4 & 0 & 7 \\
2016 & 0 & 71 & 8 & 36 & 4 & 9 & 3 \\
2017 & 0 & 60 & 18 & 52 & 4 & 6 & 5 \\
2018 & 0 & 66 & 24 & 34 & 1 & 3 & 2 \\
2019 & 0 & 59 & 14 & 48 & 6 & 6 & 1 \\
2020 & 0 & 79 & 19 & 32 & 4 & 5 & 10 \\
2021 & 0 & 61 & 17 & 46 & 3 & 7 & 2 \\
2022 & 0 & 42 & 10 & 33 & 2 & 3 & 0 \\
2023 & 0 & 54 & 16 & 40 & 2 & 6 & 2 \\
2024 & 0 & 53 & 10 & 15 & 2 & 6 & 3 \\
2025 & 0 & 64 & 11 & 20 & 3 & 1 & 0 \\
\end{longtable}}
\vspace{-6pt}\fonte{Elaboração IPP com dados de \citeonline{sms_rio_sim}.}
\end{landscape}


{\small\setlength{\tabcolsep}{5pt}
\begin{longtable}{@{}rrrr@{}}
\caption{Mortalidade infantil por causas evitáveis, por tipo de causa (grupo/subgrupo CID-10) (mortalidade causas evitaveis grupo 28 a 364 dias ano)}\label{tab:mortalidade-causas-evitaveis-grupo-28-a-364-dias-ano}\\
\toprule \textbf{Ano} & \textbf{1. Causas evitáveis} & \textbf{2. Causas mal definidas} & \textbf{3. Demais causas (não claramente evitáveis)} \\ \midrule \endfirsthead
\multicolumn{4}{@{}l}{\footnotesize\itshape (continuação)}\\ \toprule \textbf{Ano} & \textbf{1. Causas evitáveis} & \textbf{2. Causas mal definidas} & \textbf{3. Demais causas (não claramente evitáveis)} \\ \midrule \endhead
\midrule \multicolumn{4}{r@{}}{\footnotesize\itshape (continua)}\\ \endfoot
\bottomrule \endlastfoot
1996 & 461 & 108 & 159 \\
1997 & 409 & 93 & 174 \\
1998 & 327 & 116 & 163 \\
1999 & 317 & 107 & 184 \\
2000 & 286 & 108 & 152 \\
2001 & 255 & 63 & 130 \\
2002 & 205 & 77 & 141 \\
2003 & 227 & 80 & 146 \\
2004 & 212 & 87 & 149 \\
2005 & 197 & 63 & 125 \\
2006 & 201 & 50 & 136 \\
2007 & 221 & 38 & 144 \\
2008 & 238 & 29 & 145 \\
2009 & 246 & 29 & 153 \\
2010 & 200 & 42 & 146 \\
2011 & 224 & 44 & 153 \\
2012 & 221 & 31 & 153 \\
2013 & 211 & 30 & 137 \\
2014 & 183 & 28 & 135 \\
2015 & 209 & 25 & 144 \\
2016 & 212 & 21 & 144 \\
2017 & 163 & 23 & 130 \\
2018 & 171 & 22 & 131 \\
2019 & 168 & 23 & 128 \\
2020 & 139 & 16 & 115 \\
2021 & 157 & 20 & 94 \\
2022 & 143 & 13 & 119 \\
2023 & 133 & 7 & 125 \\
2024 & 137 & 8 & 105 \\
2025 & 134 & 13 & 102 \\
\end{longtable}}
\vspace{-6pt}\fonte{Elaboração IPP com dados de \citeonline{sms_rio_sim}.}


\begin{landscape}
{\scriptsize\setlength{\tabcolsep}{3pt}
\begin{longtable}{@{}rrrrrrrr@{}}
\caption{Mortalidade infantil por causas evitáveis, por tipo de causa (grupo/subgrupo CID-10) (mortalidade causas evitaveis subgrupo 28 a 364 dias ano)}\label{tab:mortalidade-causas-evitaveis-subgrupo-28-a-364-dias-ano}\\
\toprule \textbf{Ano} & \textbf{1.1. Reduzível pelas ações de imunização} & \textbf{1.2.1 Reduzíveis atenção à mulher na gestação} & \textbf{1.2.2 Reduz por adequada atenção à mulher no parto} & \textbf{1.2.3 Reduzíveis adequada atenção ao recém-nascido} & \textbf{1.3. Reduz ações diagnóstico e tratamento adequado} & \textbf{1.4. Reduz. ações promoção vinc. ações de atenção} & \textbf{2. Causas mal definidas} \\ \midrule \endfirsthead
\multicolumn{8}{@{}l}{\footnotesize\itshape (continuação)}\\ \toprule \textbf{Ano} & \textbf{1.1. Reduzível pelas ações de imunização} & \textbf{1.2.1 Reduzíveis atenção à mulher na gestação} & \textbf{1.2.2 Reduz por adequada atenção à mulher no parto} & \textbf{1.2.3 Reduzíveis adequada atenção ao recém-nascido} & \textbf{1.3. Reduz ações diagnóstico e tratamento adequado} & \textbf{1.4. Reduz. ações promoção vinc. ações de atenção} & \textbf{2. Causas mal definidas} \\ \midrule \endhead
\midrule \multicolumn{8}{r@{}}{\footnotesize\itshape (continua)}\\ \endfoot
\bottomrule \endlastfoot
1996 & 7 & 24 & 2 & 16 & 283 & 129 & 108 \\
1997 & 7 & 35 & 5 & 33 & 253 & 76 & 93 \\
1998 & 6 & 21 & 5 & 15 & 199 & 81 & 116 \\
1999 & 4 & 27 & 1 & 9 & 208 & 68 & 107 \\
2000 & 2 & 28 & 13 & 18 & 151 & 74 & 108 \\
2001 & 2 & 15 & 12 & 28 & 135 & 63 & 63 \\
2002 & 1 & 14 & 7 & 10 & 121 & 52 & 77 \\
2003 & 2 & 16 & 9 & 19 & 126 & 55 & 80 \\
2004 & 4 & 20 & 3 & 14 & 113 & 58 & 87 \\
2005 & 0 & 21 & 11 & 15 & 101 & 49 & 63 \\
2006 & 2 & 22 & 2 & 17 & 102 & 56 & 50 \\
2007 & 3 & 21 & 9 & 17 & 121 & 50 & 38 \\
2008 & 3 & 28 & 17 & 26 & 96 & 68 & 29 \\
2009 & 1 & 34 & 12 & 19 & 114 & 66 & 29 \\
2010 & 2 & 39 & 2 & 15 & 79 & 63 & 42 \\
2011 & 4 & 28 & 4 & 19 & 103 & 66 & 44 \\
2012 & 3 & 27 & 7 & 21 & 93 & 70 & 31 \\
2013 & 0 & 36 & 7 & 9 & 101 & 58 & 30 \\
2014 & 3 & 28 & 5 & 18 & 70 & 59 & 28 \\
2015 & 5 & 33 & 6 & 14 & 82 & 69 & 25 \\
2016 & 1 & 28 & 10 & 18 & 89 & 66 & 21 \\
2017 & 2 & 15 & 5 & 21 & 67 & 53 & 23 \\
2018 & 0 & 15 & 6 & 17 & 80 & 53 & 22 \\
2019 & 2 & 23 & 6 & 16 & 71 & 50 & 23 \\
2020 & 1 & 27 & 8 & 14 & 38 & 51 & 16 \\
2021 & 1 & 21 & 10 & 18 & 46 & 61 & 20 \\
2022 & 2 & 11 & 4 & 12 & 59 & 55 & 13 \\
2023 & 0 & 8 & 2 & 4 & 75 & 44 & 7 \\
2024 & 0 & 14 & 1 & 8 & 60 & 54 & 8 \\
2025 & 1 & 16 & 7 & 11 & 43 & 56 & 13 \\
\end{longtable}}
\vspace{-6pt}\fonte{Elaboração IPP com dados de \citeonline{sms_rio_sim}.}
\end{landscape}


{\small\setlength{\tabcolsep}{5pt}
\begin{longtable}{@{}rlrl@{}}
\caption{Mortalidade infantil por causas evitáveis, por tipo de causa (grupo/subgrupo CID-10) (mortalidade causas evitaveis subgrupo faixa 2025)}\label{tab:mortalidade-causas-evitaveis-subgrupo-faixa-2025}\\
\toprule \textbf{Ano} & \textbf{Subgrupo} & \textbf{Obitos} & \textbf{Faixa etaria} \\ \midrule \endfirsthead
\multicolumn{4}{@{}l}{\footnotesize\itshape (continuação)}\\ \toprule \textbf{Ano} & \textbf{Subgrupo} & \textbf{Obitos} & \textbf{Faixa etaria} \\ \midrule \endhead
\midrule \multicolumn{4}{r@{}}{\footnotesize\itshape (continua)}\\ \endfoot
\bottomrule \endlastfoot
2025 & 1.1. Reduzível pelas ações de imunização & 0 & 0-6 dias \\
2025 & 1.2.1 Reduzíveis atenção à mulher na gestação & 183 & 0-6 dias \\
2025 & 1.2.2 Reduz por adequada atenção à mulher no parto & 50 & 0-6 dias \\
2025 & 1.2.3 Reduzíveis adequada atenção ao recém-nascido & 36 & 0-6 dias \\
2025 & 1.3. Reduz ações diagnóstico e tratamento adequado & 0 & 0-6 dias \\
2025 & 1.4. Reduz. ações promoção vinc. ações de atenção & 0 & 0-6 dias \\
2025 & 2. Causas mal definidas & 2 & 0-6 dias \\
2025 & 1.1. Reduzível pelas ações de imunização & 0 & 7-27 dias \\
2025 & 1.2.1 Reduzíveis atenção à mulher na gestação & 64 & 7-27 dias \\
2025 & 1.2.2 Reduz por adequada atenção à mulher no parto & 11 & 7-27 dias \\
2025 & 1.2.3 Reduzíveis adequada atenção ao recém-nascido & 20 & 7-27 dias \\
2025 & 1.3. Reduz ações diagnóstico e tratamento adequado & 3 & 7-27 dias \\
2025 & 1.4. Reduz. ações promoção vinc. ações de atenção & 1 & 7-27 dias \\
2025 & 2. Causas mal definidas & 0 & 7-27 dias \\
2025 & 1.1. Reduzível pelas ações de imunização & 1 & 28-364 dias \\
2025 & 1.2.1 Reduzíveis atenção à mulher na gestação & 16 & 28-364 dias \\
2025 & 1.2.2 Reduz por adequada atenção à mulher no parto & 7 & 28-364 dias \\
2025 & 1.2.3 Reduzíveis adequada atenção ao recém-nascido & 11 & 28-364 dias \\
2025 & 1.3. Reduz ações diagnóstico e tratamento adequado & 43 & 28-364 dias \\
2025 & 1.4. Reduz. ações promoção vinc. ações de atenção & 56 & 28-364 dias \\
2025 & 2. Causas mal definidas & 13 & 28-364 dias \\
\end{longtable}}
\vspace{-6pt}\fonte{Elaboração IPP com dados de \citeonline{sms_rio_sim}.}


{\small\setlength{\tabcolsep}{5pt}
\begin{longtable}{@{}lrr@{}}
\caption{Mortalidade infantil por causas evitáveis, por tipo de causa (grupo/subgrupo CID-10) (obitos evitaveis menores 5 subgrupo municipio ano)}\label{tab:obitos-evitaveis-menores-5-subgrupo-municipio-ano}\\
\toprule \textbf{Subgrupo} & \textbf{Ano} & \textbf{Obitos} \\ \midrule \endfirsthead
\multicolumn{3}{@{}l}{\footnotesize\itshape (continuação)}\\ \toprule \textbf{Subgrupo} & \textbf{Ano} & \textbf{Obitos} \\ \midrule \endhead
\midrule \multicolumn{3}{r@{}}{\footnotesize\itshape (continua)}\\ \endfoot
\bottomrule \endlastfoot
1.1. Reduzível pelas ações de imunização & 2006 & 2 \\
1.2.1. Red por at à mulher na gestação & 2006 & 325 \\
1.2.2. Red por at à mulher no parto & 2006 & 85 \\
1.2.3. Red por at ao recém-nascido & 2006 & 204 \\
1.3. Red por ações de diag e trat adequado & 2006 & 164 \\
1.4. Red por ações promoção vinc a atenção & 2006 & 99 \\
2. Causas mal definidas & 2006 & 99 \\
3. Demais causas (não claramente evitáveis) & 2006 & 375 \\
1.1. Reduzível pelas ações de imunização & 2007 & 3 \\
1.2.1. Red por at à mulher na gestação & 2007 & 300 \\
1.2.2. Red por at à mulher no parto & 2007 & 112 \\
1.2.3. Red por at ao recém-nascido & 2007 & 149 \\
1.3. Red por ações de diag e trat adequado & 2007 & 161 \\
1.4. Red por ações promoção vinc a atenção & 2007 & 106 \\
2. Causas mal definidas & 2007 & 67 \\
3. Demais causas (não claramente evitáveis) & 2007 & 388 \\
1.1. Reduzível pelas ações de imunização & 2008 & 3 \\
1.2.1. Red por at à mulher na gestação & 2008 & 308 \\
1.2.2. Red por at à mulher no parto & 2008 & 112 \\
1.2.3. Red por at ao recém-nascido & 2008 & 174 \\
1.3. Red por ações de diag e trat adequado & 2008 & 148 \\
1.4. Red por ações promoção vinc a atenção & 2008 & 112 \\
2. Causas mal definidas & 2008 & 61 \\
3. Demais causas (não claramente evitáveis) & 2008 & 392 \\
1.1. Reduzível pelas ações de imunização & 2009 & 1 \\
1.2.1. Red por at à mulher na gestação & 2009 & 370 \\
1.2.2. Red por at à mulher no parto & 2009 & 109 \\
1.2.3. Red por at ao recém-nascido & 2009 & 162 \\
1.3. Red por ações de diag e trat adequado & 2009 & 154 \\
1.4. Red por ações promoção vinc a atenção & 2009 & 118 \\
2. Causas mal definidas & 2009 & 48 \\
3. Demais causas (não claramente evitáveis) & 2009 & 387 \\
1.1. Reduzível pelas ações de imunização & 2010 & 4 \\
1.2.1. Red por at à mulher na gestação & 2010 & 351 \\
1.2.2. Red por at à mulher no parto & 2010 & 87 \\
1.2.3. Red por at ao recém-nascido & 2010 & 134 \\
1.3. Red por ações de diag e trat adequado & 2010 & 142 \\
1.4. Red por ações promoção vinc a atenção & 2010 & 119 \\
2. Causas mal definidas & 2010 & 72 \\
3. Demais causas (não claramente evitáveis) & 2010 & 395 \\
1.1. Reduzível pelas ações de imunização & 2011 & 4 \\
1.2.1. Red por at à mulher na gestação & 2011 & 334 \\
1.2.2. Red por at à mulher no parto & 2011 & 84 \\
1.2.3. Red por at ao recém-nascido & 2011 & 120 \\
1.3. Red por ações de diag e trat adequado & 2011 & 149 \\
1.4. Red por ações promoção vinc a atenção & 2011 & 101 \\
2. Causas mal definidas & 2011 & 65 \\
3. Demais causas (não claramente evitáveis) & 2011 & 392 \\
1.1. Reduzível pelas ações de imunização & 2012 & 5 \\
1.2.1. Red por at à mulher na gestação & 2012 & 366 \\
1.2.2. Red por at à mulher no parto & 2012 & 83 \\
1.2.3. Red por at ao recém-nascido & 2012 & 126 \\
1.3. Red por ações de diag e trat adequado & 2012 & 145 \\
1.4. Red por ações promoção vinc a atenção & 2012 & 118 \\
2. Causas mal definidas & 2012 & 54 \\
3. Demais causas (não claramente evitáveis) & 2012 & 381 \\
1.1. Reduzível pelas ações de imunização & 2013 & 0 \\
1.2.1. Red por at à mulher na gestação & 2013 & 390 \\
1.2.2. Red por at à mulher no parto & 2013 & 81 \\
1.2.3. Red por at ao recém-nascido & 2013 & 141 \\
1.3. Red por ações de diag e trat adequado & 2013 & 150 \\
1.4. Red por ações promoção vinc a atenção & 2013 & 88 \\
2. Causas mal definidas & 2013 & 52 \\
3. Demais causas (não claramente evitáveis) & 2013 & 377 \\
1.1. Reduzível pelas ações de imunização & 2014 & 5 \\
1.2.1. Red por at à mulher na gestação & 2014 & 325 \\
1.2.2. Red por at à mulher no parto & 2014 & 77 \\
1.2.3. Red por at ao recém-nascido & 2014 & 143 \\
1.3. Red por ações de diag e trat adequado & 2014 & 106 \\
1.4. Red por ações promoção vinc a atenção & 2014 & 103 \\
2. Causas mal definidas & 2014 & 43 \\
3. Demais causas (não claramente evitáveis) & 2014 & 376 \\
1.1. Reduzível pelas ações de imunização & 2015 & 6 \\
1.2.1. Red por at à mulher na gestação & 2015 & 352 \\
1.2.2. Red por at à mulher no parto & 2015 & 94 \\
1.2.3. Red por at ao recém-nascido & 2015 & 163 \\
1.3. Red por ações de diag e trat adequado & 2015 & 127 \\
1.4. Red por ações promoção vinc a atenção & 2015 & 110 \\
2. Causas mal definidas & 2015 & 46 \\
3. Demais causas (não claramente evitáveis) & 2015 & 354 \\
1.1. Reduzível pelas ações de imunização & 2016 & 1 \\
1.2.1. Red por at à mulher na gestação & 2016 & 339 \\
1.2.2. Red por at à mulher no parto & 2016 & 76 \\
1.2.3. Red por at ao recém-nascido & 2016 & 158 \\
1.3. Red por ações de diag e trat adequado & 2016 & 141 \\
1.4. Red por ações promoção vinc a atenção & 2016 & 100 \\
2. Causas mal definidas & 2016 & 44 \\
3. Demais causas (não claramente evitáveis) & 2016 & 382 \\
1.1. Reduzível pelas ações de imunização & 2017 & 3 \\
1.2.1. Red por at à mulher na gestação & 2017 & 289 \\
1.2.2. Red por at à mulher no parto & 2017 & 86 \\
1.2.3. Red por at ao recém-nascido & 2017 & 132 \\
1.3. Red por ações de diag e trat adequado & 2017 & 113 \\
1.4. Red por ações promoção vinc a atenção & 2017 & 105 \\
2. Causas mal definidas & 2017 & 40 \\
3. Demais causas (não claramente evitáveis) & 2017 & 367 \\
1.1. Reduzível pelas ações de imunização & 2018 & 1 \\
1.2.1. Red por at à mulher na gestação & 2018 & 312 \\
1.2.2. Red por at à mulher no parto & 2018 & 101 \\
1.2.3. Red por at ao recém-nascido & 2018 & 117 \\
1.3. Red por ações de diag e trat adequado & 2018 & 118 \\
1.4. Red por ações promoção vinc a atenção & 2018 & 89 \\
2. Causas mal definidas & 2018 & 30 \\
3. Demais causas (não claramente evitáveis) & 2018 & 351 \\
1.1. Reduzível pelas ações de imunização & 2019 & 2 \\
1.2.1. Red por at à mulher na gestação & 2019 & 273 \\
1.2.2. Red por at à mulher no parto & 2019 & 103 \\
1.2.3. Red por at ao recém-nascido & 2019 & 127 \\
1.3. Red por ações de diag e trat adequado & 2019 & 114 \\
1.4. Red por ações promoção vinc a atenção & 2019 & 85 \\
2. Causas mal definidas & 2019 & 36 \\
3. Demais causas (não claramente evitáveis) & 2019 & 359 \\
1.1. Reduzível pelas ações de imunização & 2020 & 1 \\
1.2.1. Red por at à mulher na gestação & 2020 & 304 \\
1.2.2. Red por at à mulher no parto & 2020 & 91 \\
1.2.3. Red por at ao recém-nascido & 2020 & 108 \\
1.3. Red por ações de diag e trat adequado & 2020 & 64 \\
1.4. Red por ações promoção vinc a atenção & 2020 & 89 \\
2. Causas mal definidas & 2020 & 33 \\
3. Demais causas (não claramente evitáveis) & 2020 & 318 \\
1.1. Reduzível pelas ações de imunização & 2021 & 1 \\
1.2.1. Red por at à mulher na gestação & 2021 & 254 \\
1.2.2. Red por at à mulher no parto & 2021 & 87 \\
1.2.3. Red por at ao recém-nascido & 2021 & 129 \\
1.3. Red por ações de diag e trat adequado & 2021 & 77 \\
1.4. Red por ações promoção vinc a atenção & 2021 & 107 \\
2. Causas mal definidas & 2021 & 34 \\
3. Demais causas (não claramente evitáveis) & 2021 & 283 \\
1.1. Reduzível pelas ações de imunização & 2022 & 4 \\
1.2.1. Red por at à mulher na gestação & 2022 & 219 \\
1.2.2. Red por at à mulher no parto & 2022 & 70 \\
1.2.3. Red por at ao recém-nascido & 2022 & 104 \\
1.3. Red por ações de diag e trat adequado & 2022 & 97 \\
1.4. Red por ações promoção vinc a atenção & 2022 & 84 \\
2. Causas mal definidas & 2022 & 30 \\
3. Demais causas (não claramente evitáveis) & 2022 & 315 \\
1.1. Reduzível pelas ações de imunização & 2023 & 1 \\
1.2.1. Red por at à mulher na gestação & 2023 & 228 \\
1.2.2. Red por at à mulher no parto & 2023 & 74 \\
1.2.3. Red por at ao recém-nascido & 2023 & 95 \\
1.3. Red por ações de diag e trat adequado & 2023 & 129 \\
1.4. Red por ações promoção vinc a atenção & 2023 & 96 \\
2. Causas mal definidas & 2023 & 15 \\
3. Demais causas (não claramente evitáveis) & 2023 & 300 \\
1.1. Reduzível pelas ações de imunização & 2024 & 0 \\
1.2.1. Red por at à mulher na gestação & 2024 & 241 \\
1.2.2. Red por at à mulher no parto & 2024 & 53 \\
1.2.3. Red por at ao recém-nascido & 2024 & 54 \\
1.3. Red por ações de diag e trat adequado & 2024 & 91 \\
1.4. Red por ações promoção vinc a atenção & 2024 & 85 \\
2. Causas mal definidas & 2024 & 19 \\
3. Demais causas (não claramente evitáveis) & 2024 & 277 \\
1.1. Reduzível pelas ações de imunização & 2025 & 2 \\
1.2.1. Red por at à mulher na gestação & 2025 & 266 \\
1.2.2. Red por at à mulher no parto & 2025 & 67 \\
1.2.3. Red por at ao recém-nascido & 2025 & 55 \\
1.3. Red por ações de diag e trat adequado & 2025 & 75 \\
1.4. Red por ações promoção vinc a atenção & 2025 & 92 \\
2. Causas mal definidas & 2025 & 22 \\
3. Demais causas (não claramente evitáveis) & 2025 & 277 \\
\end{longtable}}
\vspace{-6pt}\fonte{Elaboração IPP com dados de \citeonline{sms_rio_sim}.}


\begin{landscape}
{\scriptsize\setlength{\tabcolsep}{3pt}
\begin{longtable}{@{}rlrrrrrrrr@{}}
\caption{Mortalidade infantil por causas evitáveis, por tipo de causa (grupo/subgrupo CID-10) (mortalidade evitaveis subgrupo cap 2025)}\label{tab:mortalidade-evitaveis-subgrupo-cap-2025}\\
\toprule \textbf{Cod ap sms} & \textbf{Faixa etaria} & \textbf{1.1. Reduzível pelas ações de imunização} & \textbf{1.2.1. Red por at à mulher na gestação} & \textbf{1.2.2. Red por at à mulher no parto} & \textbf{1.2.3. Red por at ao recém-nascido} & \textbf{1.3. Red por ações de diag e trat adequado} & \textbf{1.4. Red por ações promoção vinc a atenção} & \textbf{2. Causas mal definidas} & \textbf{3. Demais causas (não claramente evitáveis)} \\ \midrule \endfirsthead
\multicolumn{10}{@{}l}{\footnotesize\itshape (continuação)}\\ \toprule \textbf{Cod ap sms} & \textbf{Faixa etaria} & \textbf{1.1. Reduzível pelas ações de imunização} & \textbf{1.2.1. Red por at à mulher na gestação} & \textbf{1.2.2. Red por at à mulher no parto} & \textbf{1.2.3. Red por at ao recém-nascido} & \textbf{1.3. Red por ações de diag e trat adequado} & \textbf{1.4. Red por ações promoção vinc a atenção} & \textbf{2. Causas mal definidas} & \textbf{3. Demais causas (não claramente evitáveis)} \\ \midrule \endhead
\midrule \multicolumn{10}{r@{}}{\footnotesize\itshape (continua)}\\ \endfoot
\bottomrule \endlastfoot
1,0 & de 1 a 4 anos & 0 & 0 & 0 & 0 & 0 & 1 & 0 & 4 \\
1,0 & menores de 1 ano & 0 & 13 & 8 & 1 & 2 & 3 & 0 & 11 \\
1,0 & menores de 5 anos & 0 & 13 & 8 & 1 & 2 & 4 & 0 & 15 \\
2,1 & de 1 a 4 anos & 0 & 0 & 0 & 0 & 2 & 0 & 0 & 1 \\
2,1 & menores de 1 ano & 0 & 7 & 1 & 1 & 3 & 0 & 4 & 16 \\
2,1 & menores de 5 anos & 0 & 7 & 1 & 1 & 5 & 0 & 4 & 17 \\
2,2 & de 1 a 4 anos & 0 & 0 & 0 & 0 & 2 & 4 & 0 & 0 \\
2,2 & menores de 1 ano & 0 & 3 & 3 & 2 & 0 & 1 & 0 & 5 \\
2,2 & menores de 5 anos & 0 & 3 & 3 & 2 & 2 & 5 & 0 & 5 \\
3,1 & de 1 a 4 anos & 0 & 0 & 0 & 0 & 6 & 2 & 1 & 9 \\
3,1 & menores de 1 ano & 0 & 28 & 8 & 8 & 10 & 10 & 3 & 30 \\
3,1 & menores de 5 anos & 0 & 28 & 8 & 8 & 16 & 12 & 4 & 39 \\
3,2 & de 1 a 4 anos & 0 & 0 & 0 & 0 & 1 & 3 & 2 & 9 \\
3,2 & menores de 1 ano & 0 & 24 & 6 & 2 & 0 & 2 & 0 & 18 \\
3,2 & menores de 5 anos & 0 & 24 & 6 & 2 & 1 & 5 & 2 & 27 \\
3,3 & de 1 a 4 anos & 1 & 1 & 0 & 0 & 6 & 4 & 1 & 7 \\
3,3 & menores de 1 ano & 0 & 39 & 10 & 16 & 12 & 7 & 2 & 25 \\
3,3 & menores de 5 anos & 1 & 40 & 10 & 16 & 18 & 11 & 3 & 32 \\
4,0 & de 1 a 4 anos & 0 & 0 & 0 & 0 & 2 & 5 & 2 & 9 \\
4,0 & menores de 1 ano & 0 & 47 & 10 & 11 & 7 & 14 & 4 & 39 \\
4,0 & menores de 5 anos & 0 & 47 & 10 & 11 & 9 & 19 & 6 & 48 \\
5,1 & de 1 a 4 anos & 0 & 0 & 0 & 0 & 3 & 6 & 1 & 6 \\
5,1 & menores de 1 ano & 0 & 31 & 4 & 2 & 6 & 5 & 1 & 23 \\
5,1 & menores de 5 anos & 0 & 31 & 4 & 2 & 9 & 11 & 2 & 29 \\
5,2 & de 1 a 4 anos & 0 & 1 & 0 & 0 & 3 & 3 & 0 & 8 \\
5,2 & menores de 1 ano & 0 & 40 & 7 & 6 & 4 & 10 & 0 & 33 \\
5,2 & menores de 5 anos & 0 & 41 & 7 & 6 & 7 & 13 & 0 & 41 \\
5,3 & de 1 a 4 anos & 0 & 0 & 0 & 0 & 4 & 6 & 0 & 6 \\
5,3 & menores de 1 ano & 1 & 32 & 10 & 6 & 2 & 6 & 1 & 18 \\
5,3 & menores de 5 anos & 1 & 32 & 10 & 6 & 6 & 12 & 1 & 24 \\
\end{longtable}}
\vspace{-6pt}\fonte{Elaboração IPP com dados de \citeonline{sms_rio_sim}.}
\end{landscape}


\section*{Tabelas disponíveis em formato digital}

As tabelas abaixo têm mais de 200 linhas e não são impressas; estão em \texttt{tabelas\_finais/} no repositório do projeto.

\begin{itemize}\item \texttt{mortalidade\_neonatal\_precoce\_bairro\_ano.csv} --- Taxa de mortalidade neonatal precoce (0 a 6 dias) (3.300 linhas)\item \texttt{mortalidade\_neonatal\_tardia\_bairro\_ano.csv} --- Taxa de mortalidade neonatal tardia (7 a 27 dias) (3.220 linhas)\item \texttt{obitos\_gravidez\_bairro\_ano.csv} --- Razão de mortalidade materna (durante a gravidez) (2.400 linhas)\item \texttt{obitos\_puerperio\_bairro\_ano.csv} --- Razão de mortalidade materna (durante puerpério) (2.720 linhas)\item \texttt{mortalidade\_raca\_bairro\_ano.csv} --- Mortalidade infantil por raça/cor (menores de 1 ano) (3.330 linhas)\item \texttt{mortalidade\_evitaveis\_cap\_faixa\_ano.csv} --- Mortalidade infantil por causas evitáveis (0 a 6, 7 a 27, 28 a 364 dias) (4.800 linhas)\item \texttt{mortalidade\_infantil\_pos\_neonatal\_total\_bairro\_ano.csv} --- Taxa de mortalidade na primeira infância (menores de 1, 1-4, 0-5 anos) (3.340 linhas)\item \texttt{mortalidade\_evitaveis\_grupo\_cap\_faixa\_ano.csv} --- Mortalidade infantil por causas evitáveis, por tipo de causa (grupo/subgrupo CID-10) (600 linhas)\end{itemize}

\chapter{Tabelas do eixo Inclusão}\label{ap:eixo-2}

{\small\setlength{\tabcolsep}{5pt}
\begin{longtable}{@{}lrrr@{}}
\caption{Crianças até 6 anos, por sexo (censo sidra populacao 0 6 sexo 2022)}\label{tab:censo-sidra-populacao-0-6-sexo-2022}\\
\toprule \textbf{Idade} & \textbf{Homens} & \textbf{Mulheres} & \textbf{Total} \\ \midrule \endfirsthead
\multicolumn{4}{@{}l}{\footnotesize\itshape (continuação)}\\ \toprule \textbf{Idade} & \textbf{Homens} & \textbf{Mulheres} & \textbf{Total} \\ \midrule \endhead
\midrule \multicolumn{4}{r@{}}{\footnotesize\itshape (continua)}\\ \endfoot
\bottomrule \endlastfoot
1 ano & 28.155 & 27.841 & 55.996 \\
2 anos & 31.707 & 31.303 & 63.010 \\
3 anos & 33.798 & 33.003 & 66.801 \\
4 anos & 35.648 & 34.856 & 70.504 \\
5 anos & 34.820 & 34.141 & 68.961 \\
6 anos & 37.794 & 35.963 & 73.757 \\
Menos de 1 ano & 27.633 & 26.704 & 54.337 \\
\textbf{Total} & \textbf{2.882.579} & \textbf{3.328.644} & \textbf{6.211.223} \\
\end{longtable}}
\vspace{-6pt}\fonte{Elaboração IPP com dados de \citeonline{ibge_censo2022}.}


{\small\setlength{\tabcolsep}{5pt}
\begin{longtable}{@{}lrrrrrr@{}}
\caption{Crianças até 6 anos, por raça/cor (censo sidra populacao 0 6 raca 2022)}\label{tab:censo-sidra-populacao-0-6-raca-2022}\\
\toprule \textbf{Idade} & \textbf{Amarela} & \textbf{Branca} & \textbf{Indígena} & \textbf{Parda} & \textbf{Preta} & \textbf{Total} \\ \midrule \endfirsthead
\multicolumn{7}{@{}l}{\footnotesize\itshape (continuação)}\\ \toprule \textbf{Idade} & \textbf{Amarela} & \textbf{Branca} & \textbf{Indígena} & \textbf{Parda} & \textbf{Preta} & \textbf{Total} \\ \midrule \endhead
\midrule \multicolumn{7}{r@{}}{\footnotesize\itshape (continua)}\\ \endfoot
\bottomrule \endlastfoot
1 ano & 67 & 26.742 & 25 & 22.615 & 6.542 & 55.996 \\
2 anos & 81 & 29.501 & 19 & 25.892 & 7.517 & 63.010 \\
3 anos & 81 & 30.184 & 43 & 28.147 & 8.342 & 66.801 \\
4 anos & 68 & 31.496 & 33 & 29.716 & 9.186 & 70.504 \\
5 anos & 87 & 29.746 & 38 & 29.837 & 9.252 & 68.961 \\
6 anos & 97 & 31.345 & 36 & 32.416 & 9.863 & 73.757 \\
Menos de 1 ano & 66 & 26.909 & 23 & 21.576 & 5.762 & 54.337 \\
\textbf{Total} & \textbf{10.514} & \textbf{2.821.619} & \textbf{6.531} & \textbf{2.403.895} & \textbf{968.428} & \textbf{6.211.223} \\
\end{longtable}}
\vspace{-6pt}\fonte{Elaboração IPP com dados de \citeonline{ibge_censo2022}.}


{\small\setlength{\tabcolsep}{5pt}
\begin{longtable}{@{}lrrrrrr@{}}
\caption{Crianças até 6 anos frequentando escola/creche, por raça/cor (sidra frequencia escola 0 5 raca 2022)}\label{tab:sidra-frequencia-escola-0-5-raca-2022}\\
\toprule \textbf{Idade} & \textbf{Amarela} & \textbf{Branca} & \textbf{Indígena} & \textbf{Parda} & \textbf{Preta} & \textbf{Total} \\ \midrule \endfirsthead
\multicolumn{7}{@{}l}{\footnotesize\itshape (continuação)}\\ \toprule \textbf{Idade} & \textbf{Amarela} & \textbf{Branca} & \textbf{Indígena} & \textbf{Parda} & \textbf{Preta} & \textbf{Total} \\ \midrule \endhead
\midrule \multicolumn{7}{r@{}}{\footnotesize\itshape (continua)}\\ \endfoot
\bottomrule \endlastfoot
0 ano & 0 & 1.835 & 0 & 1.883 & 641 & 4.358 \\
1 ano & 0 & 8.285 & 24 & 6.692 & 2.553 & 17.597 \\
2 anos & 63 & 16.579 & 36 & 13.916 & 4.313 & 34.907 \\
3 anos & 19 & 22.588 & 20 & 19.718 & 6.897 & 49.278 \\
4 anos & 42 & 27.262 & 23 & 25.676 & 8.203 & 61.206 \\
5 anos & 93 & 28.432 & 20 & 28.467 & 9.151 & 66.163 \\
\textbf{Total} & \textbf{217} & \textbf{104.981} & \textbf{122} & \textbf{96.352} & \textbf{31.757} & \textbf{233.509} \\
\end{longtable}}
\vspace{-6pt}\fonte{Elaboração IPP com dados de \citeonline{ibge_censo2022}.}


{\small\setlength{\tabcolsep}{5pt}
\begin{longtable}{@{}lrrrrrr@{}}
\caption{Crianças até 6 anos frequentando escola/creche, por raça/cor (sidra taxa frequencia 0 6 raca 2022)}\label{tab:sidra-taxa-frequencia-0-6-raca-2022}\\
\toprule \textbf{Idade} & \textbf{Amarela} & \textbf{Branca} & \textbf{Indígena} & \textbf{Parda} & \textbf{Preta} & \textbf{Total} \\ \midrule \endfirsthead
\multicolumn{7}{@{}l}{\footnotesize\itshape (continuação)}\\ \toprule \textbf{Idade} & \textbf{Amarela} & \textbf{Branca} & \textbf{Indígena} & \textbf{Parda} & \textbf{Preta} & \textbf{Total} \\ \midrule \endhead
\midrule \multicolumn{7}{r@{}}{\footnotesize\itshape (continua)}\\ \endfoot
\bottomrule \endlastfoot
0 ano & 0,0 & 6,7 & 0,0 & 8,9 & 11,3 & 8,0 \\
1 ano & 0,0 & 29,4 & 100,0 & 29,9 & 33,8 & 30,3 \\
2 anos & 48,8 & 57,4 & 100,0 & 51,7 & 59,3 & 55,2 \\
3 anos & 100,0 & 74,3 & 33,4 & 69,5 & 72,0 & 72,0 \\
4 anos & 100,0 & 84,0 & 100,0 & 82,5 & 80,8 & 82,9 \\
5 anos & 100,0 & 90,7 & 77,8 & 90,8 & 89,2 & 90,6 \\
6 anos & 100,0 & 97,2 & 100,0 & 96,9 & 96,8 & 97,0 \\
\textbf{Total} & \textbf{20,9} & \textbf{24,4} & \textbf{20,1} & \textbf{26,8} & \textbf{24,5} & \textbf{25,3} \\
\end{longtable}}
\vspace{-6pt}\fonte{Elaboração IPP com dados de \citeonline{ibge_censo2022}.}


{\small\setlength{\tabcolsep}{5pt}
\begin{longtable}{@{}lrrr@{}}
\caption{Crianças até 6 anos frequentando escola/creche, por sexo (sidra frequencia escola 0 5 sexo 2022)}\label{tab:sidra-frequencia-escola-0-5-sexo-2022}\\
\toprule \textbf{Idade} & \textbf{Homens} & \textbf{Mulheres} & \textbf{Total} \\ \midrule \endfirsthead
\multicolumn{4}{@{}l}{\footnotesize\itshape (continuação)}\\ \toprule \textbf{Idade} & \textbf{Homens} & \textbf{Mulheres} & \textbf{Total} \\ \midrule \endhead
\midrule \multicolumn{4}{r@{}}{\footnotesize\itshape (continua)}\\ \endfoot
\bottomrule \endlastfoot
0 ano & 2.334 & 2.025 & 4.358 \\
1 ano & 9.056 & 8.541 & 17.597 \\
2 anos & 17.291 & 17.616 & 34.907 \\
3 anos & 25.425 & 23.852 & 49.278 \\
4 anos & 31.299 & 29.907 & 61.206 \\
5 anos & 34.898 & 31.264 & 66.163 \\
\textbf{Total} & \textbf{120.304} & \textbf{113.205} & \textbf{233.509} \\
\end{longtable}}
\vspace{-6pt}\fonte{Elaboração IPP com dados de \citeonline{ibge_censo2022}.}


{\small\setlength{\tabcolsep}{5pt}
\begin{longtable}{@{}lrrr@{}}
\caption{Crianças até 6 anos frequentando escola/creche, por sexo (sidra taxa frequencia 0 6 sexo 2022)}\label{tab:sidra-taxa-frequencia-0-6-sexo-2022}\\
\toprule \textbf{Idade} & \textbf{Homens} & \textbf{Mulheres} & \textbf{Total} \\ \midrule \endfirsthead
\multicolumn{4}{@{}l}{\footnotesize\itshape (continuação)}\\ \toprule \textbf{Idade} & \textbf{Homens} & \textbf{Mulheres} & \textbf{Total} \\ \midrule \endhead
\midrule \multicolumn{4}{r@{}}{\footnotesize\itshape (continua)}\\ \endfoot
\bottomrule \endlastfoot
0 ano & 8,5 & 7,6 & 8,0 \\
1 ano & 31,1 & 29,5 & 30,3 \\
2 anos & 56,1 & 54,3 & 55,2 \\
3 anos & 72,5 & 71,5 & 72,0 \\
4 anos & 82,5 & 83,3 & 82,9 \\
5 anos & 91,6 & 89,5 & 90,6 \\
6 anos & 97,1 & 97,0 & 97,0 \\
\textbf{Total} & \textbf{26,8} & \textbf{24,1} & \textbf{25,3} \\
\end{longtable}}
\vspace{-6pt}\fonte{Elaboração IPP com dados de \citeonline{ibge_censo2022}.}


{\small\setlength{\tabcolsep}{5pt}
\begin{longtable}{@{}lrrrrr@{}}
\caption{Crianças de 0 a 5 anos no CadÚnico em relação à população do município (cadunico razao populacao 0 a 5 2026)}\label{tab:cadunico-razao-populacao-0-a-5-2026}\\
\toprule \textbf{Data particao} & \textbf{Criancas cadunico 0 a 5} & \textbf{Familias cadunico} & \textbf{Ano populacao} & \textbf{Populacao ripsa 0 a 5} & \textbf{Razao (\%)} \\ \midrule \endfirsthead
\multicolumn{6}{@{}l}{\footnotesize\itshape (continuação)}\\ \toprule \textbf{Data particao} & \textbf{Criancas cadunico 0 a 5} & \textbf{Familias cadunico} & \textbf{Ano populacao} & \textbf{Populacao ripsa 0 a 5} & \textbf{Razao (\%)} \\ \midrule \endhead
\midrule \multicolumn{6}{r@{}}{\footnotesize\itshape (continua)}\\ \endfoot
\bottomrule \endlastfoot
2026-06-12 & 194.138 & 173.768 & 2.025 & 393.073 & 49,4 \\
\end{longtable}}
\vspace{-6pt}\fonte{Elaboração IPP com dados de \citeonline{ms_ripsa_populacao}; \citeonline{mds_cadunico}.}


{\small\setlength{\tabcolsep}{5pt}
\begin{longtable}{@{}llrrrrr@{}}
\caption{Famílias no CadÚnico com crianças até 6 anos, por sexo (cadunico por sexo 2026)}\label{tab:cadunico-por-sexo-2026}\\
\toprule \textbf{Recorte} & \textbf{Categoria} & \textbf{Crianças} & \textbf{Famílias com ao menos uma} & \textbf{\% das crianças} & \textbf{Famílias} & \textbf{\% das famílias} \\ \midrule \endfirsthead
\multicolumn{7}{@{}l}{\footnotesize\itshape (continuação)}\\ \toprule \textbf{Recorte} & \textbf{Categoria} & \textbf{Crianças} & \textbf{Famílias com ao menos uma} & \textbf{\% das crianças} & \textbf{Famílias} & \textbf{\% das famílias} \\ \midrule \endhead
\midrule \multicolumn{7}{r@{}}{\footnotesize\itshape (continua)}\\ \endfoot
\bottomrule \endlastfoot
Crianças por sexo & Feminino & 94.778 & 89.463 & 48,8 & -- & -- \\
Crianças por sexo & Masculino & 99.360 & 93.769 & 51,2 & -- & -- \\
Crianças por sexo & Total (famílias não somam) & 194.138 & 173.768 & 100,0 & -- & -- \\
Famílias por sexo das crianças & Só meninas & 84.879 & -- & -- & 79.999 & 46,0 \\
Famílias por sexo das crianças & Só meninos & 89.490 & -- & -- & 84.305 & 48,5 \\
Famílias por sexo das crianças & Meninas e meninos & 19.769 & -- & -- & 9.464 & 5,4 \\
Famílias por sexo das crianças & Total & 194.138 & -- & -- & 173.768 & 100,0 \\
\end{longtable}}
\vspace{-6pt}\fonte{Elaboração IPP com dados de \citeonline{mds_cadunico}.}


\begin{landscape}
{\scriptsize\setlength{\tabcolsep}{3pt}
\begin{longtable}{@{}lrrrrrrrrrr@{}}
\caption{Famílias no CadÚnico com crianças até 6 anos, por sexo (cadunico recortes bairro 2026)}\label{tab:tabela-mapa-cadunico-recortes-bairro-2026}\\
\toprule \textbf{Bairro} & \textbf{Crianças} & \textbf{Meninas} & \textbf{Crianças negras} & \textbf{Famílias} & \textbf{Famílias com uma adulta} & \textbf{Codbairro} & \textbf{\% meninas} & \textbf{\% crianças negras} & \textbf{\% famílias com uma adulta} & \textbf{Suprimido} \\ \midrule \endfirsthead
\multicolumn{11}{@{}l}{\footnotesize\itshape (continuação)}\\ \toprule \textbf{Bairro} & \textbf{Crianças} & \textbf{Meninas} & \textbf{Crianças negras} & \textbf{Famílias} & \textbf{Famílias com uma adulta} & \textbf{Codbairro} & \textbf{\% meninas} & \textbf{\% crianças negras} & \textbf{\% famílias com uma adulta} & \textbf{Suprimido} \\ \midrule \endhead
\midrule \multicolumn{11}{r@{}}{\footnotesize\itshape (continua)}\\ \endfoot
\bottomrule \endlastfoot
Abolição & 177 & 93 & 96 & 161 & 135 & 70 & 52,5 & 54,2 & 83,9 & 0 \\
Acari & 1.596 & 811 & 928 & 1.396 & 1.205 & 111 & 50,8 & 58,1 & 86,3 & 0 \\
Alto da Boa Vista & 187 & 91 & 120 & 159 & 122 & 34 & 48,7 & 64,2 & 76,7 & 0 \\
Anchieta & 2.317 & 1.126 & 1.531 & 2.050 & 1.637 & 107 & 48,6 & 66,1 & 79,9 & 0 \\
Andaraí & 451 & 231 & 293 & 411 & 331 & 37 & 51,2 & 65,0 & 80,5 & 0 \\
Anil & 3.285 & 1.587 & 1.679 & 2.953 & 2.470 & 116 & 48,3 & 51,1 & 83,6 & 0 \\
Argentino & -- & -- & -- & -- & -- & 166 & -- & -- & -- & 1 \\
Bancários & 438 & 220 & 223 & 382 & 289 & 97 & 50,2 & 50,9 & 75,7 & 0 \\
Bangu & 7.705 & 3.763 & 5.239 & 6.942 & 5.694 & 141 & 48,8 & 68,0 & 82,0 & 0 \\
Barra Olímpica & 141 & 72 & 97 & 126 & 96 & 165 & 51,1 & 68,8 & 76,2 & 0 \\
Barra da Tijuca & 333 & 166 & 191 & 305 & 238 & 128 & 49,8 & 57,4 & 78,0 & 0 \\
Barra de Guaratiba & 127 & 71 & 74 & 116 & 87 & 152 & 55,9 & 58,3 & 75,0 & 0 \\
Barros Filho & 688 & 358 & 467 & 614 & 516 & 112 & 52,0 & 67,9 & 84,0 & 0 \\
Benfica & 1.144 & 549 & 792 & 1.046 & 879 & 12 & 48,0 & 69,2 & 84,0 & 0 \\
Bento Ribeiro & 895 & 440 & 546 & 798 & 648 & 89 & 49,2 & 61,0 & 81,2 & 0 \\
Bonsucesso & 2.936 & 1.377 & 1.887 & 2.598 & 2.174 & 40 & 46,9 & 64,3 & 83,7 & 0 \\
Botafogo & 312 & 149 & 149 & 283 & 186 & 20 & 47,8 & 47,8 & 65,7 & 0 \\
Brás de Pina & 1.229 & 607 & 789 & 1.097 & 877 & 45 & 49,4 & 64,2 & 79,9 & 0 \\
Cachambi & 401 & 188 & 252 & 358 & 267 & 65 & 46,9 & 62,8 & 74,6 & 0 \\
Cacuia & 219 & 121 & 116 & 199 & 145 & 93 & 55,3 & 53,0 & 72,9 & 0 \\
Caju & 1.160 & 577 & 753 & 1.035 & 859 & 4 & 49,7 & 64,9 & 83,0 & 0 \\
Camorim & 556 & 274 & 369 & 504 & 402 & 129 & 49,3 & 66,4 & 79,8 & 0 \\
Campinho & 254 & 118 & 157 & 231 & 182 & 78 & 46,5 & 61,8 & 78,8 & 0 \\
Campo Grande & 10.459 & 5.110 & 7.170 & 9.354 & 7.355 & 144 & 48,9 & 68,6 & 78,6 & 0 \\
Campo dos Afonsos & -- & -- & -- & -- & -- & 136 & -- & -- & -- & 1 \\
Cascadura & 910 & 462 & 591 & 816 & 682 & 82 & 50,8 & 64,9 & 83,6 & 0 \\
Catete & 175 & 90 & 64 & 171 & 126 & 18 & 51,4 & 36,6 & 73,7 & 0 \\
Catumbi & 507 & 252 & 322 & 460 & 388 & 6 & 49,7 & 63,5 & 84,3 & 0 \\
Cavalcanti & 595 & 274 & 402 & 528 & 436 & 80 & 46,1 & 67,6 & 82,6 & 0 \\
Centro & 535 & 253 & 312 & 480 & 367 & 5 & 47,3 & 58,3 & 76,5 & 0 \\
Cidade Nova & 112 & 67 & 56 & 105 & 85 & 8 & 59,8 & 50,0 & 81,0 & 0 \\
Cidade Universitária & 31 & 16 & 22 & 26 & 22 & 105 & 51,6 & 71,0 & 84,6 & 0 \\
Cidade de Deus & 2.639 & 1.272 & 2.042 & 2.317 & 1.954 & 118 & 48,2 & 77,4 & 84,3 & 0 \\
Cocotá & 88 & 43 & 41 & 83 & 64 & 96 & 48,9 & 46,6 & 77,1 & 0 \\
Coelho Neto & 1.414 & 649 & 949 & 1.253 & 1.019 & 110 & 45,9 & 67,1 & 81,3 & 0 \\
Colégio & 984 & 492 & 630 & 896 & 744 & 77 & 50,0 & 64,0 & 83,0 & 0 \\
Complexo do Alemão & 1.363 & 662 & 922 & 1.216 & 1.024 & 156 & 48,6 & 67,6 & 84,2 & 0 \\
Copacabana & 1.227 & 615 & 657 & 1.127 & 877 & 24 & 50,1 & 53,5 & 77,8 & 0 \\
Cordovil & 1.279 & 650 & 829 & 1.147 & 935 & 46 & 50,8 & 64,8 & 81,5 & 0 \\
Cosme Velho & 180 & 78 & 106 & 160 & 127 & 19 & 43,3 & 58,9 & 79,4 & 0 \\
Cosmos & 3.228 & 1.527 & 2.223 & 2.867 & 2.290 & 147 & 47,3 & 68,9 & 79,9 & 0 \\
Costa Barros & 1.781 & 846 & 1.286 & 1.580 & 1.352 & 113 & 47,5 & 72,2 & 85,6 & 0 \\
Curicica & 1.141 & 556 & 744 & 1.020 & 815 & 119 & 48,7 & 65,2 & 79,9 & 0 \\
Del Castilho & 610 & 291 & 387 & 557 & 447 & 53 & 47,7 & 63,4 & 80,3 & 0 \\
Deodoro & 521 & 262 & 328 & 468 & 386 & 134 & 50,3 & 63,0 & 82,5 & 0 \\
Encantado & 274 & 126 & 166 & 251 & 182 & 68 & 46,0 & 60,6 & 72,5 & 0 \\
Engenheiro Leal & 254 & 135 & 172 & 232 & 187 & 81 & 53,1 & 67,7 & 80,6 & 0 \\
Engenho Novo & 1.317 & 660 & 839 & 1.177 & 924 & 61 & 50,1 & 63,7 & 78,5 & 0 \\
Engenho da Rainha & 738 & 365 & 517 & 661 & 521 & 55 & 49,5 & 70,1 & 78,8 & 0 \\
Engenho de Dentro & 1.096 & 558 & 723 & 967 & 743 & 66 & 50,9 & 66,0 & 76,8 & 0 \\
Estácio & 853 & 408 & 473 & 749 & 604 & 9 & 47,8 & 55,5 & 80,6 & 0 \\
Flamengo & 99 & 54 & 48 & 91 & 71 & 15 & 54,5 & 48,5 & 78,0 & 0 \\
Freguesia (Ilha) & 525 & 258 & 304 & 472 & 367 & 98 & 49,1 & 57,9 & 77,8 & 0 \\
Freguesia (Jacarepaguá) & 939 & 449 & 644 & 850 & 685 & 120 & 47,8 & 68,6 & 80,6 & 0 \\
Galeão & 734 & 355 & 369 & 650 & 512 & 104 & 48,4 & 50,3 & 78,8 & 0 \\
Gamboa & 502 & 252 & 322 & 454 & 398 & 2 & 50,2 & 64,1 & 87,7 & 0 \\
Gardênia Azul & 1.654 & 807 & 959 & 1.473 & 1.197 & 117 & 48,8 & 58,0 & 81,3 & 0 \\
Glória & 163 & 77 & 73 & 146 & 124 & 16 & 47,2 & 44,8 & 84,9 & 0 \\
Grajaú & 377 & 208 & 204 & 336 & 250 & 38 & 55,2 & 54,1 & 74,4 & 0 \\
Grumari & -- & -- & -- & -- & -- & 133 & -- & -- & -- & 1 \\
Guadalupe & 1.211 & 590 & 614 & 1.097 & 892 & 106 & 48,7 & 50,7 & 81,3 & 0 \\
Guaratiba & 6.960 & 3.406 & 4.446 & 6.244 & 4.953 & 151 & 48,9 & 63,9 & 79,3 & 0 \\
Gávea & 1.858 & 894 & 906 & 1.690 & 1.348 & 29 & 48,1 & 48,8 & 79,8 & 0 \\
Higienópolis & 215 & 107 & 132 & 198 & 142 & 50 & 49,8 & 61,4 & 71,7 & 0 \\
Honório Gurgel & 663 & 339 & 418 & 595 & 477 & 87 & 51,1 & 63,0 & 80,2 & 0 \\
Humaitá & -- & -- & -- & -- & -- & 21 & -- & -- & -- & 1 \\
Inhaúma & 1.874 & 923 & 1.239 & 1.692 & 1.386 & 54 & 49,3 & 66,1 & 81,9 & 0 \\
Inhoaíba & 2.895 & 1.442 & 1.959 & 2.572 & 2.072 & 146 & 49,8 & 67,7 & 80,6 & 0 \\
Ipanema & -- & -- & -- & -- & -- & 25 & -- & -- & -- & 1 \\
Irajá & 1.542 & 730 & 861 & 1.379 & 1.095 & 76 & 47,3 & 55,8 & 79,4 & 0 \\
Itanhangá & 1.917 & 937 & 913 & 1.706 & 1.412 & 127 & 48,9 & 47,6 & 82,8 & 0 \\
Jacarepaguá & 3.427 & 1.696 & 2.027 & 3.054 & 2.436 & 115 & 49,5 & 59,1 & 79,8 & 0 \\
Jacarezinho & 277 & 141 & 201 & 253 & 220 & 155 & 50,9 & 72,6 & 87,0 & 0 \\
Jacaré & 1.409 & 642 & 988 & 1.260 & 1.056 & 51 & 45,6 & 70,1 & 83,8 & 0 \\
Jardim América & 1.096 & 544 & 718 & 989 & 795 & 49 & 49,6 & 65,5 & 80,4 & 0 \\
Jardim Botânico & 38 & 19 & 19 & 37 & 21 & 28 & 50,0 & 50,0 & 56,8 & 0 \\
Jardim Carioca & 473 & 228 & 235 & 426 & 324 & 100 & 48,2 & 49,7 & 76,1 & 0 \\
Jardim Guanabara & 42 & 27 & 18 & 37 & 25 & 99 & 64,3 & 42,9 & 67,6 & 0 \\
Jardim Sulacap & 206 & 103 & 110 & 186 & 142 & 137 & 50,0 & 53,4 & 76,3 & 0 \\
Joá & -- & -- & -- & -- & -- & 126 & -- & -- & -- & 1 \\
Lagoa & -- & -- & -- & -- & -- & 27 & -- & -- & -- & 1 \\
Laranjeiras & 112 & 51 & 49 & 102 & 74 & 17 & 45,5 & 43,8 & 72,5 & 0 \\
Leblon & 49 & 19 & 35 & 47 & 22 & 26 & 38,8 & 71,4 & 46,8 & 0 \\
Leme & 172 & 86 & 104 & 159 & 119 & 23 & 50,0 & 60,5 & 74,8 & 0 \\
Lins de Vasconcelos & 947 & 446 & 637 & 840 & 690 & 62 & 47,1 & 67,3 & 82,1 & 0 \\
Madureira & 1.435 & 724 & 961 & 1.283 & 1.079 & 83 & 50,5 & 67,0 & 84,1 & 0 \\
Magalhães Bastos & 571 & 293 & 398 & 518 & 422 & 138 & 51,3 & 69,7 & 81,5 & 0 \\
Mangueira & 821 & 423 & 591 & 738 & 626 & 11 & 51,5 & 72,0 & 84,8 & 0 \\
Manguinhos & 1.927 & 965 & 1.268 & 1.689 & 1.440 & 39 & 50,1 & 65,8 & 85,3 & 0 \\
Maracanã & 67 & 33 & 41 & 56 & 43 & 35 & 49,3 & 61,2 & 76,8 & 0 \\
Marechal Hermes & 1.084 & 535 & 565 & 977 & 800 & 90 & 49,4 & 52,1 & 81,9 & 0 \\
Maria da Graça & 85 & 36 & 48 & 79 & 60 & 52 & 42,4 & 56,5 & 75,9 & 0 \\
Maré & 3.405 & 1.633 & 2.261 & 3.042 & 2.604 & 157 & 48,0 & 66,4 & 85,6 & 0 \\
Moneró & -- & -- & -- & -- & -- & 102 & -- & -- & -- & 1 \\
Méier & 398 & 193 & 236 & 362 & 279 & 63 & 48,5 & 59,3 & 77,1 & 0 \\
Olaria & 1.231 & 580 & 766 & 1.118 & 903 & 42 & 47,1 & 62,2 & 80,8 & 0 \\
Osvaldo Cruz & 688 & 340 & 461 & 614 & 488 & 88 & 49,4 & 67,0 & 79,5 & 0 \\
Paciência & 4.702 & 2.324 & 3.216 & 4.202 & 3.487 & 148 & 49,4 & 68,4 & 83,0 & 0 \\
Padre Miguel & 1.525 & 712 & 1.001 & 1.377 & 1.116 & 140 & 46,7 & 65,6 & 81,0 & 0 \\
Paquetá & 101 & 44 & 53 & 85 & 72 & 13 & 43,6 & 52,5 & 84,7 & 0 \\
Parada de Lucas & 594 & 294 & 391 & 529 & 442 & 47 & 49,5 & 65,8 & 83,6 & 0 \\
Parque Anchieta & 448 & 219 & 260 & 401 & 311 & 108 & 48,9 & 58,0 & 77,6 & 0 \\
Parque Colúmbia & 123 & 61 & 77 & 109 & 91 & 159 & 49,6 & 62,6 & 83,5 & 0 \\
Pavuna & 2.786 & 1.306 & 1.868 & 2.508 & 2.031 & 114 & 46,9 & 67,0 & 81,0 & 0 \\
Pechincha & 366 & 168 & 209 & 334 & 257 & 121 & 45,9 & 57,1 & 76,9 & 0 \\
Pedra de Guaratiba & 742 & 355 & 439 & 676 & 542 & 153 & 47,8 & 59,2 & 80,2 & 0 \\
Penha & 2.007 & 980 & 1.376 & 1.799 & 1.518 & 43 & 48,8 & 68,6 & 84,4 & 0 \\
Penha Circular & 1.275 & 610 & 795 & 1.172 & 947 & 44 & 47,8 & 62,4 & 80,8 & 0 \\
Piedade & 1.206 & 598 & 786 & 1.078 & 854 & 69 & 49,6 & 65,2 & 79,2 & 0 \\
Pilares & 748 & 364 & 515 & 664 & 514 & 71 & 48,7 & 68,9 & 77,4 & 0 \\
Pitangueiras & 135 & 58 & 70 & 123 & 85 & 94 & 43,0 & 51,9 & 69,1 & 0 \\
Portuguesa & 262 & 129 & 142 & 241 & 179 & 103 & 49,2 & 54,2 & 74,3 & 0 \\
Praia da Bandeira & 21 & 7 & 6 & 20 & 7 & 95 & 33,3 & 28,6 & 35,0 & 0 \\
Praça Seca & 1.833 & 905 & 1.202 & 1.634 & 1.316 & 124 & 49,4 & 65,6 & 80,5 & 0 \\
Praça da Bandeira & 69 & 40 & 40 & 64 & 47 & 32 & 58,0 & 58,0 & 73,4 & 0 \\
Quintino Bocaiúva & 708 & 355 & 430 & 631 & 532 & 79 & 50,1 & 60,7 & 84,3 & 0 \\
Ramos & 2.661 & 1.303 & 1.630 & 2.393 & 1.969 & 41 & 49,0 & 61,3 & 82,3 & 0 \\
Realengo & 5.498 & 2.684 & 3.668 & 4.941 & 3.913 & 139 & 48,8 & 66,7 & 79,2 & 0 \\
Recreio dos Bandeirantes & 1.028 & 489 & 574 & 931 & 730 & 132 & 47,6 & 55,8 & 78,4 & 0 \\
Riachuelo & 126 & 61 & 65 & 109 & 80 & 59 & 48,4 & 51,6 & 73,4 & 0 \\
Ribeira & 26 & 13 & 11 & 25 & 14 & 91 & 50,0 & 42,3 & 56,0 & 0 \\
Ricardo de Albuquerque & 1.191 & 591 & 662 & 1.034 & 860 & 109 & 49,6 & 55,6 & 83,2 & 0 \\
Rio Comprido & 947 & 432 & 539 & 855 & 695 & 7 & 45,6 & 56,9 & 81,3 & 0 \\
Rocha & 433 & 221 & 293 & 398 & 326 & 58 & 51,0 & 67,7 & 81,9 & 0 \\
Rocha Miranda & 1.244 & 604 & 856 & 1.092 & 901 & 86 & 48,6 & 68,8 & 82,5 & 0 \\
Rocinha & 1.237 & 592 & 622 & 1.125 & 915 & 154 & 47,9 & 50,3 & 81,3 & 0 \\
Sampaio & 525 & 260 & 337 & 479 & 400 & 60 & 49,5 & 64,2 & 83,5 & 0 \\
Santa Cruz & 11.321 & 5.560 & 7.354 & 10.099 & 8.265 & 149 & 49,1 & 65,0 & 81,8 & 0 \\
Santa Teresa & 889 & 450 & 470 & 825 & 663 & 14 & 50,6 & 52,9 & 80,4 & 0 \\
Santo Cristo & 557 & 282 & 348 & 499 & 414 & 3 & 50,6 & 62,5 & 83,0 & 0 \\
Santíssimo & 1.981 & 957 & 1.303 & 1.785 & 1.436 & 143 & 48,3 & 65,8 & 80,4 & 0 \\
Saúde & 53 & 28 & 33 & 48 & 34 & 1 & 52,8 & 62,3 & 70,8 & 0 \\
Senador Camará & 3.520 & 1.739 & 2.343 & 3.123 & 2.562 & 142 & 49,4 & 66,6 & 82,0 & 0 \\
Senador Vasconcelos & 1.154 & 560 & 836 & 1.011 & 808 & 145 & 48,5 & 72,4 & 79,9 & 0 \\
Sepetiba & 2.937 & 1.436 & 1.745 & 2.608 & 2.109 & 150 & 48,9 & 59,4 & 80,9 & 0 \\
São Conrado & 76 & 37 & 38 & 68 & 53 & 31 & 48,7 & 50,0 & 77,9 & 0 \\
Imperial de São Cristóvão & 1.561 & 771 & 1.060 & 1.407 & 1.165 & 10 & 49,4 & 67,9 & 82,8 & 0 \\
São Francisco Xavier & 127 & 57 & 84 & 113 & 88 & 57 & 44,9 & 66,1 & 77,9 & 0 \\
Tanque & 1.501 & 736 & 1.014 & 1.314 & 1.067 & 123 & 49,0 & 67,6 & 81,2 & 0 \\
Taquara & 2.502 & 1.189 & 1.656 & 2.212 & 1.736 & 122 & 47,5 & 66,2 & 78,5 & 0 \\
Tauá & 615 & 301 & 355 & 556 & 429 & 101 & 48,9 & 57,7 & 77,2 & 0 \\
Tijuca & 1.313 & 661 & 879 & 1.166 & 904 & 33 & 50,3 & 66,9 & 77,5 & 0 \\
Todos os Santos & 91 & 48 & 56 & 85 & 60 & 64 & 52,7 & 61,5 & 70,6 & 0 \\
Tomás Coelho & 725 & 356 & 471 & 654 & 529 & 56 & 49,1 & 65,0 & 80,9 & 0 \\
Turiaçú & 566 & 271 & 394 & 495 & 415 & 85 & 47,9 & 69,6 & 83,8 & 0 \\
Urca & -- & -- & -- & -- & -- & 22 & -- & -- & -- & 1 \\
Vargem Grande & 764 & 392 & 488 & 686 & 522 & 131 & 51,3 & 63,9 & 76,1 & 0 \\
Vargem Pequena & 625 & 306 & 367 & 572 & 445 & 130 & 49,0 & 58,7 & 77,8 & 0 \\
Vasco da Gama & 56 & 29 & 30 & 49 & 37 & 158 & 51,8 & 53,6 & 75,5 & 0 \\
Vaz Lobo & 580 & 293 & 391 & 487 & 403 & 84 & 50,5 & 67,4 & 82,8 & 0 \\
Vicente de Carvalho & 760 & 364 & 474 & 679 & 547 & 73 & 47,9 & 62,4 & 80,6 & 0 \\
Vidigal & 481 & 223 & 253 & 440 & 341 & 30 & 46,4 & 52,6 & 77,5 & 0 \\
Vigário Geral & 1.264 & 615 & 893 & 1.134 & 946 & 48 & 48,7 & 70,6 & 83,4 & 0 \\
Vila Isabel & 960 & 447 & 615 & 864 & 708 & 36 & 46,6 & 64,1 & 81,9 & 0 \\
Vila Kosmos & 314 & 133 & 178 & 275 & 224 & 72 & 42,4 & 56,7 & 81,5 & 0 \\
Vila Militar & 106 & 50 & 76 & 99 & 77 & 135 & 47,2 & 71,7 & 77,8 & 0 \\
Vila Valqueire & 496 & 236 & 311 & 438 & 352 & 125 & 47,6 & 62,7 & 80,4 & 0 \\
Vila da Penha & 171 & 80 & 89 & 158 & 116 & 74 & 46,8 & 52,0 & 73,4 & 0 \\
Vista Alegre & 128 & 63 & 59 & 117 & 74 & 75 & 49,2 & 46,1 & 63,2 & 0 \\
Zumbi & -- & -- & -- & -- & -- & 92 & -- & -- & -- & 1 \\
Água Santa & 154 & 63 & 104 & 142 & 104 & 67 & 40,9 & 67,5 & 73,2 & 0 \\
\end{longtable}}
\vspace{-6pt}\fonte{Elaboração IPP com dados de \citeonline{mds_cadunico}.}
\end{landscape}


{\small\setlength{\tabcolsep}{5pt}
\begin{longtable}{@{}lrrrl@{}}
\caption{Famílias no CadÚnico com crianças até 6 anos, por raça/cor (cadunico por raca cor 2026)}\label{tab:cadunico-por-raca-cor-2026}\\
\toprule \textbf{Raça/cor da criança} & \textbf{Crianças} & \textbf{Famílias com ao menos uma} & \textbf{\% das crianças} & \textbf{Nota} \\ \midrule \endfirsthead
\multicolumn{5}{@{}l}{\footnotesize\itshape (continuação)}\\ \toprule \textbf{Raça/cor da criança} & \textbf{Crianças} & \textbf{Famílias com ao menos uma} & \textbf{\% das crianças} & \textbf{Nota} \\ \midrule \endhead
\midrule \multicolumn{5}{r@{}}{\footnotesize\itshape (continua)}\\ \endfoot
\bottomrule \endlastfoot
Parda & 103.991 & 96.016 & 53,6 & famílias não exclusivas entre categorias; não somar \\
Branca & 68.200 & 63.845 & 35,1 & famílias não exclusivas entre categorias; não somar \\
Preta & 20.670 & 19.452 & 10,6 & famílias não exclusivas entre categorias; não somar \\
Amarela & 1.208 & 1.193 & 0,6 & famílias não exclusivas entre categorias; não somar \\
Indígena & 69 & 64 & 0,0 & famílias não exclusivas entre categorias; não somar \\
Negra (preta + parda) & 124.661 & 113.754 & 64,2 & famílias não exclusivas entre categorias; não somar \\
\textbf{Total (famílias não somam)} & \textbf{194.138} & \textbf{173.768} & \textbf{100,0} & \textbf{famílias não exclusivas entre categorias; não somar} \\
\end{longtable}}
\vspace{-6pt}\fonte{Elaboração IPP com dados de \citeonline{mds_cadunico}.}


{\small\setlength{\tabcolsep}{5pt}
\begin{longtable}{@{}lrrr@{}}
\caption{Famílias no CadÚnico com crianças até 6 anos, por renda e arranjo familiar (cadunico familias por arranjo 2026)}\label{tab:cadunico-familias-por-arranjo-2026}\\
\toprule \textbf{Arranjo familiar} & \textbf{Famílias} & \textbf{Crianças} & \textbf{\% das famílias} \\ \midrule \endfirsthead
\multicolumn{4}{@{}l}{\footnotesize\itshape (continuação)}\\ \toprule \textbf{Arranjo familiar} & \textbf{Famílias} & \textbf{Crianças} & \textbf{\% das famílias} \\ \midrule \endhead
\midrule \multicolumn{4}{r@{}}{\footnotesize\itshape (continua)}\\ \endfoot
\bottomrule \endlastfoot
Uma adulta (mulher) & 141.051 & 158.437 & 81,2 \\
Dois adultos (homem e mulher) & 18.168 & 20.042 & 10,5 \\
Dois adultos (outra composição) & 6.377 & 6.932 & 3,7 \\
Um adulto (homem) & 3.944 & 4.147 & 2,3 \\
Três ou mais adultos & 3.515 & 3.834 & 2,0 \\
Sem adulto (18+) & 713 & 746 & 0,4 \\
\textbf{Total} & \textbf{173.768} & \textbf{194.138} & \textbf{100,0} \\
\end{longtable}}
\vspace{-6pt}\fonte{Elaboração IPP com dados de \citeonline{mds_cadunico}.}


{\small\setlength{\tabcolsep}{5pt}
\begin{longtable}{@{}lrrlr@{}}
\caption{Famílias no CadÚnico com crianças até 6 anos, por renda e arranjo familiar (cadunico familias arranjo renda 2026)}\label{tab:cadunico-familias-arranjo-renda-2026}\\
\toprule \textbf{Arranjo} & \textbf{Famílias} & \textbf{\% no arranjo} & \textbf{Faixa de renda per capita} & \textbf{Suprimido} \\ \midrule \endfirsthead
\multicolumn{5}{@{}l}{\footnotesize\itshape (continuação)}\\ \toprule \textbf{Arranjo} & \textbf{Famílias} & \textbf{\% no arranjo} & \textbf{Faixa de renda per capita} & \textbf{Suprimido} \\ \midrule \endhead
\midrule \multicolumn{5}{r@{}}{\footnotesize\itshape (continua)}\\ \endfoot
\bottomrule \endlastfoot
Uma adulta (mulher) & 112.380 & 79,7 & Extrema pobreza (até R\$ 218) & 0 \\
Uma adulta (mulher) & 21.653 & 15,4 & Pobreza/baixa renda (R\$ 218 a 810) & 0 \\
Uma adulta (mulher) & 7.018 & 5,0 & Acima de 1/2 SM (mais de R\$ 810) & 0 \\
Dois adultos (homem e mulher) & 9.830 & 54,1 & Extrema pobreza (até R\$ 218) & 0 \\
Dois adultos (homem e mulher) & 5.538 & 30,5 & Pobreza/baixa renda (R\$ 218 a 810) & 0 \\
Dois adultos (homem e mulher) & 2.800 & 15,4 & Acima de 1/2 SM (mais de R\$ 810) & 0 \\
Dois adultos (outra composição) & 4.262 & 66,8 & Extrema pobreza (até R\$ 218) & 0 \\
Dois adultos (outra composição) & 1.718 & 26,9 & Pobreza/baixa renda (R\$ 218 a 810) & 0 \\
Dois adultos (outra composição) & 397 & 6,2 & Acima de 1/2 SM (mais de R\$ 810) & 0 \\
Um adulto (homem) & 2.889 & 73,3 & Extrema pobreza (até R\$ 218) & 0 \\
Um adulto (homem) & 707 & 17,9 & Pobreza/baixa renda (R\$ 218 a 810) & 0 \\
Um adulto (homem) & 348 & 8,8 & Acima de 1/2 SM (mais de R\$ 810) & 0 \\
Três ou mais adultos & 1.761 & 50,1 & Extrema pobreza (até R\$ 218) & 0 \\
Três ou mais adultos & 1.312 & 37,3 & Pobreza/baixa renda (R\$ 218 a 810) & 0 \\
Três ou mais adultos & 442 & 12,6 & Acima de 1/2 SM (mais de R\$ 810) & 0 \\
Sem adulto (18+) & 686 & 96,2 & Extrema pobreza (até R\$ 218) & 0 \\
Sem adulto (18+) & 20 & 2,8 & Pobreza/baixa renda (R\$ 218 a 810) & 0 \\
Sem adulto (18+) & -- & -- & Acima de 1/2 SM (mais de R\$ 810) & 1 \\
\end{longtable}}
\vspace{-6pt}\fonte{Elaboração IPP com dados de \citeonline{mds_cadunico}.}


\chapter{Tabelas do eixo Família e Cuidados}\label{ap:eixo-3}

{\small\setlength{\tabcolsep}{5pt}
\begin{longtable}{@{}rrr@{}}
\caption{Crianças até 6 anos no Cadastro Único (número) (cadunico por idade 2026)}\label{tab:cadunico-por-idade-2026}\\
\toprule \textbf{Idade} & \textbf{Crianças} & \textbf{Famílias} \\ \midrule \endfirsthead
\multicolumn{3}{@{}l}{\footnotesize\itshape (continuação)}\\ \toprule \textbf{Idade} & \textbf{Crianças} & \textbf{Famílias} \\ \midrule \endhead
\midrule \multicolumn{3}{r@{}}{\footnotesize\itshape (continua)}\\ \endfoot
\bottomrule \endlastfoot
0 & 11.328 & 11.155 \\
1 & 27.357 & 26.998 \\
2 & 33.116 & 32.689 \\
3 & 37.838 & 37.410 \\
4 & 41.312 & 40.864 \\
5 & 43.187 & 42.708 \\
\end{longtable}}
\vspace{-6pt}\fonte{Elaboração IPP com dados de \citeonline{mds_cadunico}.}


{\small\setlength{\tabcolsep}{5pt}
\begin{longtable}{@{}lrrr@{}}
\caption{Crianças até 6 anos no Cadastro Único (número) (cadunico por bairro 2026)}\label{tab:cadunico-por-bairro-2026}\\
\toprule \textbf{Bairro} & \textbf{Crianças} & \textbf{Famílias} & \textbf{Suprimido} \\ \midrule \endfirsthead
\multicolumn{4}{@{}l}{\footnotesize\itshape (continuação)}\\ \toprule \textbf{Bairro} & \textbf{Crianças} & \textbf{Famílias} & \textbf{Suprimido} \\ \midrule \endhead
\midrule \multicolumn{4}{r@{}}{\footnotesize\itshape (continua)}\\ \endfoot
\bottomrule \endlastfoot
Abolição & 177 & 161 & 0 \\
Acari & 1.596 & 1.396 & 0 \\
Alto da Boa Vista & 187 & 159 & 0 \\
Anchieta & 2.317 & 2.050 & 0 \\
Andaraí & 451 & 411 & 0 \\
Anil & 3.285 & 2.953 & 0 \\
Argentino & -- & -- & 1 \\
Bancários & 438 & 382 & 0 \\
Bangu & 7.705 & 6.942 & 0 \\
Barra Olímpica & 141 & 126 & 0 \\
Barra da Tijuca & 333 & 305 & 0 \\
Barra de Guaratiba & 127 & 116 & 0 \\
Barros Filho & 688 & 614 & 0 \\
Benfica & 1.144 & 1.046 & 0 \\
Bento Ribeiro & 895 & 798 & 0 \\
Bonsucesso & 2.936 & 2.598 & 0 \\
Botafogo & 312 & 283 & 0 \\
Brás de Pina & 1.229 & 1.097 & 0 \\
Cachambi & 401 & 358 & 0 \\
Cacuia & 219 & 199 & 0 \\
Caju & 1.160 & 1.035 & 0 \\
Camorim & 556 & 504 & 0 \\
Campinho & 254 & 231 & 0 \\
Campo Grande & 10.459 & 9.354 & 0 \\
Campo dos Afonsos & -- & -- & 1 \\
Cascadura & 910 & 816 & 0 \\
Catete & 175 & 171 & 0 \\
Catumbi & 507 & 460 & 0 \\
Cavalcanti & 595 & 528 & 0 \\
Centro & 535 & 480 & 0 \\
Cidade Nova & 112 & 105 & 0 \\
Cidade Universitária & 31 & 26 & 0 \\
Cidade de Deus & 2.639 & 2.317 & 0 \\
Cocotá & 88 & 83 & 0 \\
Coelho Neto & 1.414 & 1.253 & 0 \\
Colégio & 984 & 896 & 0 \\
Complexo do Alemão & 1.363 & 1.216 & 0 \\
Copacabana & 1.227 & 1.127 & 0 \\
Cordovil & 1.279 & 1.147 & 0 \\
Cosme Velho & 180 & 160 & 0 \\
Cosmos & 3.228 & 2.867 & 0 \\
Costa Barros & 1.781 & 1.580 & 0 \\
Curicica & 1.141 & 1.020 & 0 \\
Del Castilho & 610 & 557 & 0 \\
Dendê & 64 & 61 & 0 \\
Deodoro & 521 & 468 & 0 \\
Dumas & 98 & 90 & 0 \\
Encantado & 274 & 251 & 0 \\
Engenheiro Leal & 254 & 232 & 0 \\
Engenho Novo & 1.317 & 1.177 & 0 \\
Engenho da Rainha & 738 & 661 & 0 \\
Engenho de Dentro & 1.096 & 967 & 0 \\
Estácio & 853 & 749 & 0 \\
Flamengo & 99 & 91 & 0 \\
Freguesia (Ilha do Governador) & 525 & 472 & 0 \\
Freguesia (Jacarepaguá) & 939 & 850 & 0 \\
Galeão & 734 & 650 & 0 \\
Gamboa & 502 & 454 & 0 \\
Gardênia Azul & 1.654 & 1.473 & 0 \\
Glória & 163 & 146 & 0 \\
Grajaú & 377 & 336 & 0 \\
Grumari & -- & -- & 1 \\
Guadalupe & 1.211 & 1.097 & 0 \\
Guarabu & 39 & 35 & 0 \\
Guaratiba & 6.960 & 6.244 & 0 \\
Gávea & 1.858 & 1.690 & 0 \\
Higienópolis & 215 & 198 & 0 \\
Honório Gurgel & 663 & 595 & 0 \\
Humaitá & -- & -- & 1 \\
Inhaúma & 1.874 & 1.692 & 0 \\
Inhoaíba & 2.895 & 2.572 & 0 \\
Ipanema & -- & -- & 1 \\
Irajá & 1.542 & 1.379 & 0 \\
Itacolomi & 62 & 59 & 0 \\
Itanhangá & 1.917 & 1.706 & 0 \\
Jacarepaguá & 3.427 & 3.054 & 0 \\
Jacarezinho & 277 & 253 & 0 \\
Jacaré & 1.409 & 1.260 & 0 \\
Jardim América & 1.096 & 989 & 0 \\
Jardim Botânico & 38 & 37 & 0 \\
Jardim Carioca & 473 & 426 & 0 \\
Jardim Guanabara & 42 & 37 & 0 \\
Jardim Sulacap & 206 & 186 & 0 \\
Joá & -- & -- & 1 \\
Lagoa & -- & -- & 1 \\
Laranjeiras & 112 & 102 & 0 \\
Leblon & 49 & 47 & 0 \\
Leme & 172 & 159 & 0 \\
Lins de Vasconcelos & 947 & 840 & 0 \\
Madureira & 1.435 & 1.283 & 0 \\
Magalhães Bastos & 571 & 518 & 0 \\
Mangueira & 821 & 738 & 0 \\
Manguinhos & 1.927 & 1.689 & 0 \\
Maracanã & 67 & 56 & 0 \\
Marechal Hermes & 1.084 & 977 & 0 \\
Maria da Graça & 85 & 79 & 0 \\
Maré & 3.405 & 3.042 & 0 \\
Moneró & -- & -- & 1 \\
Méier & 398 & 362 & 0 \\
Nossa Senhora das Graças & 25 & 24 & 0 \\
Olaria & 1.231 & 1.118 & 0 \\
Oswaldo Cruz & 688 & 614 & 0 \\
Paciência & 4.702 & 4.202 & 0 \\
Padre Miguel & 1.525 & 1.377 & 0 \\
Paquetá & 101 & 85 & 0 \\
Parada de Lucas & 594 & 529 & 0 \\
Parque Anchieta & 448 & 401 & 0 \\
Parque Colúmbia & 123 & 109 & 0 \\
Pavuna & 2.786 & 2.508 & 0 \\
Pechincha & 366 & 334 & 0 \\
Pedra de Guaratiba & 742 & 676 & 0 \\
Penha & 2.007 & 1.799 & 0 \\
Penha Circular & 1.275 & 1.172 & 0 \\
Piedade & 1.206 & 1.078 & 0 \\
Pilares & 748 & 664 & 0 \\
Pitangueiras & 135 & 123 & 0 \\
Portuguesa & 262 & 241 & 0 \\
Praia da Bandeira & 21 & 20 & 0 \\
Praça Seca & 1.833 & 1.634 & 0 \\
Praça da Bandeira & 69 & 64 & 0 \\
Quintino Bocaiúva & 708 & 631 & 0 \\
Ramos & 2.661 & 2.393 & 0 \\
Realengo & 5.498 & 4.941 & 0 \\
Recreio dos Bandeirantes & 1.028 & 931 & 0 \\
Riachuelo & 126 & 109 & 0 \\
Ribeira & 26 & 25 & 0 \\
Ricardo de Albuquerque & 1.191 & 1.034 & 0 \\
Rio Comprido & 947 & 855 & 0 \\
Rocha & 433 & 398 & 0 \\
Rocha Miranda & 1.244 & 1.092 & 0 \\
Rocinha & 1.237 & 1.125 & 0 \\
Sampaio & 525 & 479 & 0 \\
Santa Cruz & 11.321 & 10.099 & 0 \\
Santa Teresa & 889 & 825 & 0 \\
Santo Cristo & 557 & 499 & 0 \\
Santíssimo & 1.981 & 1.785 & 0 \\
Saúde & 53 & 48 & 0 \\
Senador Camará & 3.520 & 3.123 & 0 \\
Senador Vasconcelos & 1.154 & 1.011 & 0 \\
Sepetiba & 2.937 & 2.608 & 0 \\
São Conrado & 76 & 68 & 0 \\
São Cristóvão & 1.561 & 1.407 & 0 \\
São Francisco Xavier & 127 & 113 & 0 \\
Tanque & 1.501 & 1.314 & 0 \\
Taquara & 2.502 & 2.212 & 0 \\
Tauá & 615 & 556 & 0 \\
Tijuca & 1.313 & 1.166 & 0 \\
Todos os Santos & 91 & 85 & 0 \\
Tomás Coelho & 725 & 654 & 0 \\
Tubiacanga & 92 & 83 & 0 \\
Turiaçu & 566 & 495 & 0 \\
Urca & -- & -- & 1 \\
Vargem Grande & 764 & 686 & 0 \\
Vargem Pequena & 625 & 572 & 0 \\
Vasco da Gama & 56 & 49 & 0 \\
Vaz Lobo & 580 & 487 & 0 \\
Vicente de Carvalho & 760 & 679 & 0 \\
Vidigal & 481 & 440 & 0 \\
Vigário Geral & 1.264 & 1.134 & 0 \\
Vila Isabel & 960 & 864 & 0 \\
Vila Kosmos & 314 & 275 & 0 \\
Vila Militar & 106 & 99 & 0 \\
Vila Valqueire & 496 & 438 & 0 \\
Vila da Penha & 171 & 158 & 0 \\
Vista Alegre & 128 & 117 & 0 \\
Zumbi & -- & -- & 1 \\
Água Santa & 154 & 142 & 0 \\
Sem bairro identificado (CEP fora da lista) & 15.809 & 14.159 & 0 \\
\textbf{Total} & \textbf{194.138} & \textbf{173.768} & \textbf{0} \\
\end{longtable}}
\vspace{-6pt}\fonte{Elaboração IPP com dados de \citeonline{mds_cadunico}.}


{\small\setlength{\tabcolsep}{5pt}
\begin{longtable}{@{}lrrr@{}}
\caption{Crianças até 6 anos no Cadastro Único (número) (cadunico por bairro ate 4 2026)}\label{tab:cadunico-por-bairro-ate-4-2026}\\
\toprule \textbf{Bairro} & \textbf{Crianças} & \textbf{Famílias} & \textbf{Suprimido} \\ \midrule \endfirsthead
\multicolumn{4}{@{}l}{\footnotesize\itshape (continuação)}\\ \toprule \textbf{Bairro} & \textbf{Crianças} & \textbf{Famílias} & \textbf{Suprimido} \\ \midrule \endhead
\midrule \multicolumn{4}{r@{}}{\footnotesize\itshape (continua)}\\ \endfoot
\bottomrule \endlastfoot
Abolição & 128 & 119 & 0 \\
Acari & 1.268 & 1.162 & 0 \\
Alto da Boa Vista & 144 & 125 & 0 \\
Anchieta & 1.802 & 1.631 & 0 \\
Andaraí & 348 & 326 & 0 \\
Anil & 2.542 & 2.361 & 0 \\
Argentino & -- & -- & 1 \\
Bancários & 357 & 319 & 0 \\
Bangu & 5.975 & 5.546 & 0 \\
Barra Olímpica & 107 & 100 & 0 \\
Barra da Tijuca & 274 & 257 & 0 \\
Barra de Guaratiba & 102 & 94 & 0 \\
Barros Filho & 555 & 503 & 0 \\
Benfica & 899 & 842 & 0 \\
Bento Ribeiro & 721 & 660 & 0 \\
Bonsucesso & 2.293 & 2.094 & 0 \\
Botafogo & 232 & 210 & 0 \\
Brás de Pina & 974 & 900 & 0 \\
Cachambi & 310 & 280 & 0 \\
Cacuia & 174 & 160 & 0 \\
Caju & 922 & 846 & 0 \\
Camorim & 423 & 398 & 0 \\
Campinho & 197 & 185 & 0 \\
Campo Grande & 8.100 & 7.469 & 0 \\
Campo dos Afonsos & -- & -- & 1 \\
Cascadura & 706 & 643 & 0 \\
Catete & 138 & 136 & 0 \\
Catumbi & 373 & 350 & 0 \\
Cavalcanti & 459 & 420 & 0 \\
Centro & 417 & 377 & 0 \\
Cidade Nova & 87 & 83 & 0 \\
Cidade Universitária & 23 & 20 & 0 \\
Cidade de Deus & 2.072 & 1.893 & 0 \\
Cocotá & 65 & 62 & 0 \\
Coelho Neto & 1.105 & 1.006 & 0 \\
Colégio & 764 & 720 & 0 \\
Complexo do Alemão & 1.060 & 980 & 0 \\
Copacabana & 925 & 877 & 0 \\
Cordovil & 995 & 916 & 0 \\
Cosme Velho & 135 & 126 & 0 \\
Cosmos & 2.512 & 2.309 & 0 \\
Costa Barros & 1.397 & 1.272 & 0 \\
Curicica & 900 & 820 & 0 \\
Del Castilho & 463 & 429 & 0 \\
Dendê & 47 & 47 & 0 \\
Deodoro & 402 & 371 & 0 \\
Dumas & 70 & 67 & 0 \\
Encantado & 203 & 189 & 0 \\
Engenheiro Leal & 194 & 178 & 0 \\
Engenho Novo & 1.002 & 922 & 0 \\
Engenho da Rainha & 576 & 536 & 0 \\
Engenho de Dentro & 826 & 751 & 0 \\
Estácio & 646 & 589 & 0 \\
Flamengo & 80 & 75 & 0 \\
Freguesia (Ilha do Governador) & 389 & 355 & 0 \\
Freguesia (Jacarepaguá) & 731 & 675 & 0 \\
Galeão & 556 & 508 & 0 \\
Gamboa & 404 & 376 & 0 \\
Gardênia Azul & 1.305 & 1.194 & 0 \\
Glória & 119 & 109 & 0 \\
Grajaú & 281 & 262 & 0 \\
Grumari & -- & -- & 1 \\
Guadalupe & 929 & 863 & 0 \\
Guarabu & 30 & 28 & 0 \\
Guaratiba & 5.408 & 5.013 & 0 \\
Gávea & 1.437 & 1.329 & 0 \\
Higienópolis & 167 & 158 & 0 \\
Honório Gurgel & 523 & 481 & 0 \\
Humaitá & -- & -- & 1 \\
Inhaúma & 1.473 & 1.368 & 0 \\
Inhoaíba & 2.273 & 2.083 & 0 \\
Ipanema & -- & -- & 1 \\
Irajá & 1.198 & 1.102 & 0 \\
Itacolomi & 50 & 47 & 0 \\
Itanhangá & 1.539 & 1.413 & 0 \\
Jacarepaguá & 2.621 & 2.428 & 0 \\
Jacarezinho & 211 & 197 & 0 \\
Jacaré & 1.071 & 994 & 0 \\
Jardim América & 854 & 798 & 0 \\
Jardim Botânico & 34 & 33 & 0 \\
Jardim Carioca & 368 & 344 & 0 \\
Jardim Guanabara & 36 & 32 & 0 \\
Jardim Sulacap & 151 & 137 & 0 \\
Joá & -- & -- & 1 \\
Lagoa & -- & -- & 1 \\
Laranjeiras & 91 & 85 & 0 \\
Leblon & 32 & 31 & 0 \\
Leme & 134 & 124 & 0 \\
Lins de Vasconcelos & 740 & 681 & 0 \\
Madureira & 1.145 & 1.052 & 0 \\
Magalhães Bastos & 441 & 408 & 0 \\
Mangueira & 644 & 598 & 0 \\
Manguinhos & 1.494 & 1.361 & 0 \\
Maracanã & 44 & 41 & 0 \\
Marechal Hermes & 832 & 771 & 0 \\
Maria da Graça & 63 & 59 & 0 \\
Maré & 2.691 & 2.475 & 0 \\
Moneró & -- & -- & 1 \\
Méier & 304 & 280 & 0 \\
Nossa Senhora das Graças & 20 & 20 & 0 \\
Olaria & 966 & 896 & 0 \\
Oswaldo Cruz & 534 & 488 & 0 \\
Paciência & 3.683 & 3.389 & 0 \\
Padre Miguel & 1.179 & 1.097 & 0 \\
Paquetá & 67 & 62 & 0 \\
Parada de Lucas & 451 & 413 & 0 \\
Parque Anchieta & 336 & 308 & 0 \\
Parque Colúmbia & 88 & 82 & 0 \\
Pavuna & 2.161 & 1.996 & 0 \\
Pechincha & 290 & 273 & 0 \\
Pedra de Guaratiba & 566 & 525 & 0 \\
Penha & 1.568 & 1.445 & 0 \\
Penha Circular & 961 & 907 & 0 \\
Piedade & 951 & 874 & 0 \\
Pilares & 579 & 528 & 0 \\
Pitangueiras & 99 & 94 & 0 \\
Portuguesa & 198 & 188 & 0 \\
Praia da Bandeira & -- & -- & 1 \\
Praça Seca & 1.421 & 1.305 & 0 \\
Praça da Bandeira & 53 & 51 & 0 \\
Quintino Bocaiúva & 537 & 497 & 0 \\
Ramos & 2.052 & 1.901 & 0 \\
Realengo & 4.223 & 3.920 & 0 \\
Recreio dos Bandeirantes & 795 & 741 & 0 \\
Riachuelo & 91 & 83 & 0 \\
Ribeira & 20 & 20 & 0 \\
Ricardo de Albuquerque & 920 & 833 & 0 \\
Rio Comprido & 740 & 688 & 0 \\
Rocha & 333 & 317 & 0 \\
Rocha Miranda & 962 & 871 & 0 \\
Rocinha & 989 & 918 & 0 \\
Sampaio & 412 & 391 & 0 \\
Santa Cruz & 8.838 & 8.167 & 0 \\
Santa Teresa & 704 & 663 & 0 \\
Santo Cristo & 432 & 400 & 0 \\
Santíssimo & 1.542 & 1.433 & 0 \\
Saúde & 38 & 36 & 0 \\
Senador Camará & 2.741 & 2.502 & 0 \\
Senador Vasconcelos & 900 & 820 & 0 \\
Sepetiba & 2.301 & 2.129 & 0 \\
São Conrado & 61 & 58 & 0 \\
São Cristóvão & 1.203 & 1.102 & 0 \\
São Francisco Xavier & 100 & 90 & 0 \\
Tanque & 1.187 & 1.073 & 0 \\
Taquara & 1.930 & 1.758 & 0 \\
Tauá & 473 & 441 & 0 \\
Tijuca & 1.009 & 919 & 0 \\
Todos os Santos & 65 & 62 & 0 \\
Tomás Coelho & 561 & 521 & 0 \\
Tubiacanga & 72 & 66 & 0 \\
Turiaçu & 444 & 392 & 0 \\
Urca & -- & -- & 1 \\
Vargem Grande & 591 & 539 & 0 \\
Vargem Pequena & 488 & 457 & 0 \\
Vasco da Gama & 46 & 40 & 0 \\
Vaz Lobo & 449 & 392 & 0 \\
Vicente de Carvalho & 580 & 535 & 0 \\
Vidigal & 359 & 341 & 0 \\
Vigário Geral & 993 & 921 & 0 \\
Vila Isabel & 733 & 680 & 0 \\
Vila Kosmos & 253 & 224 & 0 \\
Vila Militar & 85 & 80 & 0 \\
Vila Valqueire & 389 & 354 & 0 \\
Vila da Penha & 140 & 132 & 0 \\
Vista Alegre & 106 & 97 & 0 \\
Zumbi & -- & -- & 1 \\
Água Santa & 120 & 112 & 0 \\
\textbf{Total} & \textbf{138.558} & \textbf{127.774} & \textbf{0} \\
\end{longtable}}
\vspace{-6pt}\fonte{Elaboração IPP com dados de \citeonline{mds_cadunico}.}


{\small\setlength{\tabcolsep}{5pt}
\begin{longtable}{@{}lrrrr@{}}
\caption{Crianças até 6 anos no Cadastro Único (número) (cadunico criancas 2026)}\label{tab:tabela-mapa-cadunico-criancas-2026}\\
\toprule \textbf{Bairro} & \textbf{Crianças} & \textbf{Famílias} & \textbf{Codbairro} & \textbf{Suprimido} \\ \midrule \endfirsthead
\multicolumn{5}{@{}l}{\footnotesize\itshape (continuação)}\\ \toprule \textbf{Bairro} & \textbf{Crianças} & \textbf{Famílias} & \textbf{Codbairro} & \textbf{Suprimido} \\ \midrule \endhead
\midrule \multicolumn{5}{r@{}}{\footnotesize\itshape (continua)}\\ \endfoot
\bottomrule \endlastfoot
Abolição & 177 & 161 & 70 & 0 \\
Acari & 1.596 & 1.396 & 111 & 0 \\
Alto da Boa Vista & 187 & 159 & 34 & 0 \\
Anchieta & 2.317 & 2.050 & 107 & 0 \\
Andaraí & 451 & 411 & 37 & 0 \\
Anil & 3.285 & 2.953 & 116 & 0 \\
Argentino & -- & -- & 166 & 1 \\
Bancários & 438 & 382 & 97 & 0 \\
Bangu & 7.705 & 6.942 & 141 & 0 \\
Barra Olímpica & 141 & 126 & 165 & 0 \\
Barra da Tijuca & 333 & 305 & 128 & 0 \\
Barra de Guaratiba & 127 & 116 & 152 & 0 \\
Barros Filho & 688 & 614 & 112 & 0 \\
Benfica & 1.144 & 1.046 & 12 & 0 \\
Bento Ribeiro & 895 & 798 & 89 & 0 \\
Bonsucesso & 2.936 & 2.598 & 40 & 0 \\
Botafogo & 312 & 283 & 20 & 0 \\
Brás de Pina & 1.229 & 1.097 & 45 & 0 \\
Cachambi & 401 & 358 & 65 & 0 \\
Cacuia & 219 & 199 & 93 & 0 \\
Caju & 1.160 & 1.035 & 4 & 0 \\
Camorim & 556 & 504 & 129 & 0 \\
Campinho & 254 & 231 & 78 & 0 \\
Campo Grande & 10.459 & 9.354 & 144 & 0 \\
Campo dos Afonsos & -- & -- & 136 & 1 \\
Cascadura & 910 & 816 & 82 & 0 \\
Catete & 175 & 171 & 18 & 0 \\
Catumbi & 507 & 460 & 6 & 0 \\
Cavalcanti & 595 & 528 & 80 & 0 \\
Centro & 535 & 480 & 5 & 0 \\
Cidade Nova & 112 & 105 & 8 & 0 \\
Cidade Universitária & 31 & 26 & 105 & 0 \\
Cidade de Deus & 2.639 & 2.317 & 118 & 0 \\
Cocotá & 88 & 83 & 96 & 0 \\
Coelho Neto & 1.414 & 1.253 & 110 & 0 \\
Colégio & 984 & 896 & 77 & 0 \\
Complexo do Alemão & 1.363 & 1.216 & 156 & 0 \\
Copacabana & 1.227 & 1.127 & 24 & 0 \\
Cordovil & 1.279 & 1.147 & 46 & 0 \\
Cosme Velho & 180 & 160 & 19 & 0 \\
Cosmos & 3.228 & 2.867 & 147 & 0 \\
Costa Barros & 1.781 & 1.580 & 113 & 0 \\
Curicica & 1.141 & 1.020 & 119 & 0 \\
Del Castilho & 610 & 557 & 53 & 0 \\
Deodoro & 521 & 468 & 134 & 0 \\
Encantado & 274 & 251 & 68 & 0 \\
Engenheiro Leal & 254 & 232 & 81 & 0 \\
Engenho Novo & 1.317 & 1.177 & 61 & 0 \\
Engenho da Rainha & 738 & 661 & 55 & 0 \\
Engenho de Dentro & 1.096 & 967 & 66 & 0 \\
Estácio & 853 & 749 & 9 & 0 \\
Flamengo & 99 & 91 & 15 & 0 \\
Freguesia (Ilha) & 525 & 472 & 98 & 0 \\
Freguesia (Jacarepaguá) & 939 & 850 & 120 & 0 \\
Galeão & 734 & 650 & 104 & 0 \\
Gamboa & 502 & 454 & 2 & 0 \\
Gardênia Azul & 1.654 & 1.473 & 117 & 0 \\
Glória & 163 & 146 & 16 & 0 \\
Grajaú & 377 & 336 & 38 & 0 \\
Grumari & -- & -- & 133 & 1 \\
Guadalupe & 1.211 & 1.097 & 106 & 0 \\
Guaratiba & 6.960 & 6.244 & 151 & 0 \\
Gávea & 1.858 & 1.690 & 29 & 0 \\
Higienópolis & 215 & 198 & 50 & 0 \\
Honório Gurgel & 663 & 595 & 87 & 0 \\
Humaitá & -- & -- & 21 & 1 \\
Inhaúma & 1.874 & 1.692 & 54 & 0 \\
Inhoaíba & 2.895 & 2.572 & 146 & 0 \\
Ipanema & -- & -- & 25 & 1 \\
Irajá & 1.542 & 1.379 & 76 & 0 \\
Itanhangá & 1.917 & 1.706 & 127 & 0 \\
Jacarepaguá & 3.427 & 3.054 & 115 & 0 \\
Jacarezinho & 277 & 253 & 155 & 0 \\
Jacaré & 1.409 & 1.260 & 51 & 0 \\
Jardim América & 1.096 & 989 & 49 & 0 \\
Jardim Botânico & 38 & 37 & 28 & 0 \\
Jardim Carioca & 473 & 426 & 100 & 0 \\
Jardim Guanabara & 42 & 37 & 99 & 0 \\
Jardim Sulacap & 206 & 186 & 137 & 0 \\
Joá & -- & -- & 126 & 1 \\
Lagoa & -- & -- & 27 & 1 \\
Laranjeiras & 112 & 102 & 17 & 0 \\
Leblon & 49 & 47 & 26 & 0 \\
Leme & 172 & 159 & 23 & 0 \\
Lins de Vasconcelos & 947 & 840 & 62 & 0 \\
Madureira & 1.435 & 1.283 & 83 & 0 \\
Magalhães Bastos & 571 & 518 & 138 & 0 \\
Mangueira & 821 & 738 & 11 & 0 \\
Manguinhos & 1.927 & 1.689 & 39 & 0 \\
Maracanã & 67 & 56 & 35 & 0 \\
Marechal Hermes & 1.084 & 977 & 90 & 0 \\
Maria da Graça & 85 & 79 & 52 & 0 \\
Maré & 3.405 & 3.042 & 157 & 0 \\
Moneró & -- & -- & 102 & 1 \\
Méier & 398 & 362 & 63 & 0 \\
Olaria & 1.231 & 1.118 & 42 & 0 \\
Osvaldo Cruz & 688 & 614 & 88 & 0 \\
Paciência & 4.702 & 4.202 & 148 & 0 \\
Padre Miguel & 1.525 & 1.377 & 140 & 0 \\
Paquetá & 101 & 85 & 13 & 0 \\
Parada de Lucas & 594 & 529 & 47 & 0 \\
Parque Anchieta & 448 & 401 & 108 & 0 \\
Parque Colúmbia & 123 & 109 & 159 & 0 \\
Pavuna & 2.786 & 2.508 & 114 & 0 \\
Pechincha & 366 & 334 & 121 & 0 \\
Pedra de Guaratiba & 742 & 676 & 153 & 0 \\
Penha & 2.007 & 1.799 & 43 & 0 \\
Penha Circular & 1.275 & 1.172 & 44 & 0 \\
Piedade & 1.206 & 1.078 & 69 & 0 \\
Pilares & 748 & 664 & 71 & 0 \\
Pitangueiras & 135 & 123 & 94 & 0 \\
Portuguesa & 262 & 241 & 103 & 0 \\
Praia da Bandeira & 21 & 20 & 95 & 0 \\
Praça Seca & 1.833 & 1.634 & 124 & 0 \\
Praça da Bandeira & 69 & 64 & 32 & 0 \\
Quintino Bocaiúva & 708 & 631 & 79 & 0 \\
Ramos & 2.661 & 2.393 & 41 & 0 \\
Realengo & 5.498 & 4.941 & 139 & 0 \\
Recreio dos Bandeirantes & 1.028 & 931 & 132 & 0 \\
Riachuelo & 126 & 109 & 59 & 0 \\
Ribeira & 26 & 25 & 91 & 0 \\
Ricardo de Albuquerque & 1.191 & 1.034 & 109 & 0 \\
Rio Comprido & 947 & 855 & 7 & 0 \\
Rocha & 433 & 398 & 58 & 0 \\
Rocha Miranda & 1.244 & 1.092 & 86 & 0 \\
Rocinha & 1.237 & 1.125 & 154 & 0 \\
Sampaio & 525 & 479 & 60 & 0 \\
Santa Cruz & 11.321 & 10.099 & 149 & 0 \\
Santa Teresa & 889 & 825 & 14 & 0 \\
Santo Cristo & 557 & 499 & 3 & 0 \\
Santíssimo & 1.981 & 1.785 & 143 & 0 \\
Saúde & 53 & 48 & 1 & 0 \\
Senador Camará & 3.520 & 3.123 & 142 & 0 \\
Senador Vasconcelos & 1.154 & 1.011 & 145 & 0 \\
Sepetiba & 2.937 & 2.608 & 150 & 0 \\
São Conrado & 76 & 68 & 31 & 0 \\
Imperial de São Cristóvão & 1.561 & 1.407 & 10 & 0 \\
São Francisco Xavier & 127 & 113 & 57 & 0 \\
Tanque & 1.501 & 1.314 & 123 & 0 \\
Taquara & 2.502 & 2.212 & 122 & 0 \\
Tauá & 615 & 556 & 101 & 0 \\
Tijuca & 1.313 & 1.166 & 33 & 0 \\
Todos os Santos & 91 & 85 & 64 & 0 \\
Tomás Coelho & 725 & 654 & 56 & 0 \\
Turiaçú & 566 & 495 & 85 & 0 \\
Urca & -- & -- & 22 & 1 \\
Vargem Grande & 764 & 686 & 131 & 0 \\
Vargem Pequena & 625 & 572 & 130 & 0 \\
Vasco da Gama & 56 & 49 & 158 & 0 \\
Vaz Lobo & 580 & 487 & 84 & 0 \\
Vicente de Carvalho & 760 & 679 & 73 & 0 \\
Vidigal & 481 & 440 & 30 & 0 \\
Vigário Geral & 1.264 & 1.134 & 48 & 0 \\
Vila Isabel & 960 & 864 & 36 & 0 \\
Vila Kosmos & 314 & 275 & 72 & 0 \\
Vila Militar & 106 & 99 & 135 & 0 \\
Vila Valqueire & 496 & 438 & 125 & 0 \\
Vila da Penha & 171 & 158 & 74 & 0 \\
Vista Alegre & 128 & 117 & 75 & 0 \\
Zumbi & -- & -- & 92 & 1 \\
Água Santa & 154 & 142 & 67 & 0 \\
\end{longtable}}
\vspace{-6pt}\fonte{Elaboração IPP com dados de \citeonline{mds_cadunico}.}


\begin{landscape}
{\scriptsize\setlength{\tabcolsep}{3pt}
\begin{longtable}{@{}lrrrrrrr@{}}
\caption{Crianças até 6 anos no Cadastro Único (número) (cadunico criancas 0 a 4 2026)}\label{tab:tabela-mapa-cadunico-criancas-0-a-4-2026}\\
\toprule \textbf{Bairro} & \textbf{Crianças} & \textbf{Famílias} & \textbf{Codbairro} & \textbf{0 a 4 anos} & \textbf{Primeira Inf. Cadúnico} & \textbf{Percentual Primeira Inf. Cadúnico} & \textbf{Suprimido} \\ \midrule \endfirsthead
\multicolumn{8}{@{}l}{\footnotesize\itshape (continuação)}\\ \toprule \textbf{Bairro} & \textbf{Crianças} & \textbf{Famílias} & \textbf{Codbairro} & \textbf{0 a 4 anos} & \textbf{Primeira Inf. Cadúnico} & \textbf{Percentual Primeira Inf. Cadúnico} & \textbf{Suprimido} \\ \midrule \endhead
\midrule \multicolumn{8}{r@{}}{\footnotesize\itshape (continua)}\\ \endfoot
\bottomrule \endlastfoot
Saúde & 38 & 36 & 1 & 54 & 0,7 & 70,4 & 0 \\
Gamboa & 404 & 376 & 2 & 669 & 0,6 & 60,4 & 0 \\
Santo Cristo & 432 & 400 & 3 & 534 & 0,8 & 80,9 & 0 \\
Caju & 922 & 846 & 4 & 1.529 & 0,6 & 60,3 & 0 \\
Centro & 417 & 377 & 5 & 622 & 0,7 & 67,0 & 0 \\
Catumbi & 373 & 350 & 6 & 662 & 0,6 & 56,3 & 0 \\
Rio Comprido & 740 & 688 & 7 & 1.908 & 0,4 & 38,8 & 0 \\
Cidade Nova & 87 & 83 & 8 & 211 & 0,4 & 41,2 & 0 \\
Estácio & 646 & 589 & 9 & 1.280 & 0,5 & 50,5 & 0 \\
Imperial de São Cristóvão & 1.203 & 1.102 & 10 & 1.257 & 1,0 & 95,7 & 0 \\
Mangueira & 644 & 598 & 11 & 1.228 & 0,5 & 52,4 & 0 \\
Benfica & 899 & 842 & 12 & 1.055 & 0,9 & 85,2 & 0 \\
Paquetá & 67 & 62 & 13 & 145 & 0,5 & 46,2 & 0 \\
Santa Teresa & 704 & 663 & 14 & 1.594 & 0,4 & 44,2 & 0 \\
Flamengo & 80 & 75 & 15 & 1.210 & 0,1 & 6,6 & 0 \\
Glória & 119 & 109 & 16 & 169 & 0,7 & 70,4 & 0 \\
Laranjeiras & 91 & 85 & 17 & 1.389 & 0,1 & 6,6 & 0 \\
Catete & 138 & 136 & 18 & 830 & 0,2 & 16,6 & 0 \\
Cosme Velho & 135 & 126 & 19 & 325 & 0,4 & 41,5 & 0 \\
Botafogo & 232 & 210 & 20 & 2.761 & 0,1 & 8,4 & 0 \\
Humaitá & -- & -- & 21 & 454 & -- & -- & 1 \\
Urca & -- & -- & 22 & 252 & -- & -- & 1 \\
Leme & 134 & 124 & 23 & 422 & 0,3 & 31,8 & 0 \\
Copacabana & 925 & 877 & 24 & 3.528 & 0,3 & 26,2 & 0 \\
Ipanema & -- & -- & 25 & 1.314 & -- & -- & 1 \\
Leblon & 32 & 31 & 26 & 1.199 & 0,0 & 2,7 & 0 \\
Lagoa & -- & -- & 27 & 772 & -- & -- & 1 \\
Jardim Botânico & 34 & 33 & 28 & 587 & 0,1 & 5,8 & 0 \\
Gávea & 1.437 & 1.329 & 29 & 459 & 3,1 & 313,1 & 0 \\
Vidigal & 359 & 341 & 30 & 853 & 0,4 & 42,1 & 0 \\
São Conrado & 61 & 58 & 31 & 337 & 0,2 & 18,1 & 0 \\
Praça da Bandeira & 53 & 51 & 32 & 261 & 0,2 & 20,3 & 0 \\
Tijuca & 1.009 & 919 & 33 & 4.764 & 0,2 & 21,2 & 0 \\
Alto da Boa Vista & 144 & 125 & 34 & 418 & 0,3 & 34,4 & 0 \\
Maracanã & 44 & 41 & 35 & 837 & 0,1 & 5,3 & 0 \\
Vila Isabel & 733 & 680 & 36 & 2.358 & 0,3 & 31,1 & 0 \\
Andaraí & 348 & 326 & 37 & 1.336 & 0,3 & 26,0 & 0 \\
Grajaú & 281 & 262 & 38 & 1.238 & 0,2 & 22,7 & 0 \\
Manguinhos & 1.494 & 1.361 & 39 & 1.915 & 0,8 & 78,0 & 0 \\
Bonsucesso & 2.293 & 2.094 & 40 & 672 & 3,4 & 341,2 & 0 \\
Ramos & 2.052 & 1.901 & 41 & 1.527 & 1,3 & 134,4 & 0 \\
Olaria & 966 & 896 & 42 & 2.133 & 0,5 & 45,3 & 0 \\
Penha & 1.568 & 1.445 & 43 & 2.958 & 0,5 & 53,0 & 0 \\
Penha Circular & 961 & 907 & 44 & 1.639 & 0,6 & 58,6 & 0 \\
Brás de Pina & 974 & 900 & 45 & 1.760 & 0,6 & 55,3 & 0 \\
Cordovil & 995 & 916 & 46 & 1.641 & 0,6 & 60,6 & 0 \\
Parada de Lucas & 451 & 413 & 47 & 1.171 & 0,4 & 38,5 & 0 \\
Vigário Geral & 993 & 921 & 48 & 2.050 & 0,5 & 48,4 & 0 \\
Jardim América & 854 & 798 & 49 & 908 & 0,9 & 94,1 & 0 \\
Higienópolis & 167 & 158 & 50 & 503 & 0,3 & 33,2 & 0 \\
Jacaré & 1.071 & 994 & 51 & 443 & 2,4 & 241,8 & 0 \\
Maria da Graça & 63 & 59 & 52 & 277 & 0,2 & 22,7 & 0 \\
Del Castilho & 463 & 429 & 53 & 991 & 0,5 & 46,7 & 0 \\
Inhaúma & 1.473 & 1.368 & 54 & 1.718 & 0,9 & 85,7 & 0 \\
Engenho da Rainha & 576 & 536 & 55 & 1.130 & 0,5 & 51,0 & 0 \\
Tomás Coelho & 561 & 521 & 56 & 1.135 & 0,5 & 49,4 & 0 \\
São Francisco Xavier & 100 & 90 & 57 & 255 & 0,4 & 39,2 & 0 \\
Rocha & 333 & 317 & 58 & 885 & 0,4 & 37,6 & 0 \\
Riachuelo & 91 & 83 & 59 & 327 & 0,3 & 27,8 & 0 \\
Sampaio & 412 & 391 & 60 & 453 & 0,9 & 90,9 & 0 \\
Engenho Novo & 1.002 & 922 & 61 & 1.905 & 0,5 & 52,6 & 0 \\
Lins de Vasconcelos & 740 & 681 & 62 & 1.502 & 0,5 & 49,3 & 0 \\
Méier & 304 & 280 & 63 & 1.080 & 0,3 & 28,1 & 0 \\
Todos os Santos & 65 & 62 & 64 & 746 & 0,1 & 8,7 & 0 \\
Cachambi & 310 & 280 & 65 & 1.412 & 0,2 & 22,0 & 0 \\
Engenho de Dentro & 826 & 751 & 66 & 1.772 & 0,5 & 46,6 & 0 \\
Água Santa & 120 & 112 & 67 & 280 & 0,4 & 42,9 & 0 \\
Encantado & 203 & 189 & 68 & 442 & 0,5 & 45,9 & 0 \\
Piedade & 951 & 874 & 69 & 1.616 & 0,6 & 58,8 & 0 \\
Abolição & 128 & 119 & 70 & 313 & 0,4 & 40,9 & 0 \\
Pilares & 579 & 528 & 71 & 1.036 & 0,6 & 55,9 & 0 \\
Vila Kosmos & 253 & 224 & 72 & 544 & 0,5 & 46,5 & 0 \\
Vicente de Carvalho & 580 & 535 & 73 & 1.042 & 0,6 & 55,7 & 0 \\
Vila da Penha & 140 & 132 & 74 & 754 & 0,2 & 18,6 & 0 \\
Vista Alegre & 106 & 97 & 75 & 645 & 0,2 & 16,4 & 0 \\
Irajá & 1.198 & 1.102 & 76 & 3.596 & 0,3 & 33,3 & 0 \\
Colégio & 764 & 720 & 77 & 1.496 & 0,5 & 51,1 & 0 \\
Campinho & 197 & 185 & 78 & 392 & 0,5 & 50,3 & 0 \\
Quintino Bocaiúva & 537 & 497 & 79 & 1.021 & 0,5 & 52,6 & 0 \\
Cavalcanti & 459 & 420 & 80 & 597 & 0,8 & 76,9 & 0 \\
Engenheiro Leal & 194 & 178 & 81 & 245 & 0,8 & 79,2 & 0 \\
Cascadura & 706 & 643 & 82 & 1.184 & 0,6 & 59,6 & 0 \\
Madureira & 1.145 & 1.052 & 83 & 1.766 & 0,6 & 64,8 & 0 \\
Vaz Lobo & 449 & 392 & 84 & 582 & 0,8 & 77,1 & 0 \\
Turiaçú & 444 & 392 & 85 & 649 & 0,7 & 68,4 & 0 \\
Rocha Miranda & 962 & 871 & 86 & 1.724 & 0,6 & 55,8 & 0 \\
Honório Gurgel & 523 & 481 & 87 & 933 & 0,6 & 56,1 & 0 \\
Osvaldo Cruz & 534 & 488 & 88 & 1.289 & 0,4 & 41,4 & 0 \\
Bento Ribeiro & 721 & 660 & 89 & 1.541 & 0,5 & 46,8 & 0 \\
Marechal Hermes & 832 & 771 & 90 & 1.963 & 0,4 & 42,4 & 0 \\
Ribeira & 20 & 20 & 91 & 129 & 0,2 & 15,5 & 0 \\
Zumbi & -- & -- & 92 & 78 & -- & -- & 1 \\
Cacuia & 174 & 160 & 93 & 439 & 0,4 & 39,6 & 0 \\
Pitangueiras & 99 & 94 & 94 & 503 & 0,2 & 19,7 & 0 \\
Praia da Bandeira & -- & -- & 95 & 205 & -- & -- & 1 \\
Cocotá & 65 & 62 & 96 & 174 & 0,4 & 37,4 & 0 \\
Bancários & 357 & 319 & 97 & 600 & 0,6 & 59,5 & 0 \\
Freguesia (Ilha) & 389 & 355 & 98 & 828 & 0,5 & 47,0 & 0 \\
Jardim Guanabara & 36 & 32 & 99 & 967 & 0,0 & 3,7 & 0 \\
Jardim Carioca & 368 & 344 & 100 & 790 & 0,5 & 46,6 & 0 \\
Tauá & 473 & 441 & 101 & 1.230 & 0,4 & 38,5 & 0 \\
Moneró & -- & -- & 102 & 181 & -- & -- & 1 \\
Portuguesa & 198 & 188 & 103 & 839 & 0,2 & 23,6 & 0 \\
Galeão & 556 & 508 & 104 & 1.448 & 0,4 & 38,4 & 0 \\
Cidade Universitária & 23 & 20 & 105 & 65 & 0,4 & 35,4 & 0 \\
Guadalupe & 929 & 863 & 106 & 2.073 & 0,4 & 44,8 & 0 \\
Anchieta & 1.802 & 1.631 & 107 & 2.866 & 0,6 & 62,9 & 0 \\
Parque Anchieta & 336 & 308 & 108 & 962 & 0,3 & 34,9 & 0 \\
Ricardo de Albuquerque & 920 & 833 & 109 & 1.345 & 0,7 & 68,4 & 0 \\
Coelho Neto & 1.105 & 1.006 & 110 & 1.190 & 0,9 & 92,9 & 0 \\
Acari & 1.268 & 1.162 & 111 & 2.112 & 0,6 & 60,0 & 0 \\
Barros Filho & 555 & 503 & 112 & 1.400 & 0,4 & 39,6 & 0 \\
Costa Barros & 1.397 & 1.272 & 113 & 1.952 & 0,7 & 71,6 & 0 \\
Pavuna & 2.161 & 1.996 & 114 & 5.892 & 0,4 & 36,7 & 0 \\
Jacarepaguá & 2.621 & 2.428 & 115 & 11.499 & 0,2 & 22,8 & 0 \\
Anil & 2.542 & 2.361 & 116 & 1.504 & 1,7 & 169,0 & 0 \\
Gardênia Azul & 1.305 & 1.194 & 117 & 1.227 & 1,1 & 106,4 & 0 \\
Cidade de Deus & 2.072 & 1.893 & 118 & 1.968 & 1,1 & 105,3 & 0 \\
Curicica & 900 & 820 & 119 & 1.485 & 0,6 & 60,6 & 0 \\
Freguesia (Jacarepaguá) & 731 & 675 & 120 & 3.266 & 0,2 & 22,4 & 0 \\
Pechincha & 290 & 273 & 121 & 1.656 & 0,2 & 17,5 & 0 \\
Taquara & 1.930 & 1.758 & 122 & 5.078 & 0,4 & 38,0 & 0 \\
Tanque & 1.187 & 1.073 & 123 & 1.818 & 0,7 & 65,3 & 0 \\
Praça Seca & 1.421 & 1.305 & 124 & 3.399 & 0,4 & 41,8 & 0 \\
Vila Valqueire & 389 & 354 & 125 & 1.382 & 0,3 & 28,1 & 0 \\
Joá & -- & -- & 126 & 22 & -- & -- & 1 \\
Itanhangá & 1.539 & 1.413 & 127 & 4.220 & 0,4 & 36,5 & 0 \\
Barra da Tijuca & 274 & 257 & 128 & 5.310 & 0,1 & 5,2 & 0 \\
Camorim & 423 & 398 & 129 & 83 & 5,1 & 509,6 & 0 \\
Vargem Pequena & 488 & 457 & 130 & 1.827 & 0,3 & 26,7 & 0 \\
Vargem Grande & 591 & 539 & 131 & 1.191 & 0,5 & 49,6 & 0 \\
Recreio dos Bandeirantes & 795 & 741 & 132 & 6.450 & 0,1 & 12,3 & 0 \\
Grumari & -- & -- & 133 & 15 & -- & -- & 1 \\
Deodoro & 402 & 371 & 134 & 742 & 0,5 & 54,2 & 0 \\
Vila Militar & 85 & 80 & 135 & 798 & 0,1 & 10,7 & 0 \\
Campo dos Afonsos & -- & -- & 136 & 178 & -- & -- & 1 \\
Jardim Sulacap & 151 & 137 & 137 & 696 & 0,2 & 21,7 & 0 \\
Magalhães Bastos & 441 & 408 & 138 & 968 & 0,5 & 45,6 & 0 \\
Realengo & 4.223 & 3.920 & 139 & 8.532 & 0,5 & 49,5 & 0 \\
Padre Miguel & 1.179 & 1.097 & 140 & 3.101 & 0,4 & 38,0 & 0 \\
Bangu & 5.975 & 5.546 & 141 & 11.418 & 0,5 & 52,3 & 0 \\
Senador Camará & 2.741 & 2.502 & 142 & 5.702 & 0,5 & 48,1 & 0 \\
Santíssimo & 1.542 & 1.433 & 143 & 2.482 & 0,6 & 62,1 & 0 \\
Campo Grande & 8.100 & 7.469 & 144 & 18.505 & 0,4 & 43,8 & 0 \\
Senador Vasconcelos & 900 & 820 & 145 & 1.638 & 0,5 & 54,9 & 0 \\
Inhoaíba & 2.273 & 2.083 & 146 & 4.495 & 0,5 & 50,6 & 0 \\
Cosmos & 2.512 & 2.309 & 147 & 5.662 & 0,4 & 44,4 & 0 \\
Paciência & 3.683 & 3.389 & 148 & 6.333 & 0,6 & 58,2 & 0 \\
Santa Cruz & 8.838 & 8.167 & 149 & 16.093 & 0,5 & 54,9 & 0 \\
Sepetiba & 2.301 & 2.129 & 150 & 3.682 & 0,6 & 62,5 & 0 \\
Guaratiba & 5.408 & 5.013 & 151 & 10.088 & 0,5 & 53,6 & 0 \\
Barra de Guaratiba & 102 & 94 & 152 & 208 & 0,5 & 49,0 & 0 \\
Pedra de Guaratiba & 566 & 525 & 153 & 599 & 0,9 & 94,5 & 0 \\
Rocinha & 989 & 918 & 154 & 4.208 & 0,2 & 23,5 & 0 \\
Jacarezinho & 211 & 197 & 155 & 2.375 & 0,1 & 8,9 & 0 \\
Complexo do Alemão & 1.060 & 980 & 156 & 3.477 & 0,3 & 30,5 & 0 \\
Maré & 2.691 & 2.475 & 157 & 8.553 & 0,3 & 31,5 & 0 \\
Vasco da Gama & 46 & 40 & 158 & 745 & 0,1 & 6,2 & 0 \\
Parque Colúmbia & 88 & 82 & 159 & 483 & 0,2 & 18,2 & 0 \\
Gericinó & -- & -- & 160 & 178 & -- & -- & 0 \\
Lapa & -- & -- & 161 & 206 & -- & -- & 0 \\
Vila Kennedy & -- & -- & 162 & 1.199 & -- & -- & 0 \\
Jabour & -- & -- & 163 & 173 & -- & -- & 0 \\
Ilha de Guaratiba & -- & -- & 164 & 309 & -- & -- & 0 \\
Barra Olímpica & 107 & 100 & 165 & 3.251 & 0,0 & 3,3 & 0 \\
Argentino & -- & -- & 166 & 33 & -- & -- & 1 \\
\end{longtable}}
\vspace{-6pt}\fonte{Elaboração IPP com dados de \citeonline{mds_cadunico}.}
\end{landscape}


{\small\setlength{\tabcolsep}{5pt}
\begin{longtable}{@{}lrrl@{}}
\caption{Famílias com crianças até 6 anos no Cadastro Único, por renda (cadunico por faixa renda 2026)}\label{tab:cadunico-por-faixa-renda-2026}\\
\toprule \textbf{Faixa de renda} & \textbf{Crianças} & \textbf{Famílias} & \textbf{Faixa de renda (descrição)} \\ \midrule \endfirsthead
\multicolumn{4}{@{}l}{\footnotesize\itshape (continuação)}\\ \toprule \textbf{Faixa de renda} & \textbf{Crianças} & \textbf{Famílias} & \textbf{Faixa de renda (descrição)} \\ \midrule \endhead
\midrule \multicolumn{4}{r@{}}{\footnotesize\itshape (continua)}\\ \endfoot
\bottomrule \endlastfoot
0-218 & 149.426 & 131.808 & Extrema pobreza (até R\$ 218) \\
219-810 & 33.435 & 30.948 & Pobreza/baixa renda (R\$ 218 a 810) \\
811-1621 & 9.742 & 9.511 & 1/2 a 1 SM (R\$ 810 a 1.621) \\
1621-3242 & 1.374 & 1.347 & 1 a 2 SM (R\$ 1.621 a 3.242) \\
3242+ & 161 & 154 & Acima de 2 SM (mais de R\$ 3.242) \\
\textbf{Total} & \textbf{194.138} & \textbf{173.768} & \textbf{Total} \\
\end{longtable}}
\vspace{-6pt}\fonte{Elaboração IPP com dados de \citeonline{mds_cadunico}.}


{\small\setlength{\tabcolsep}{5pt}
\begin{longtable}{@{}lrll@{}}
\caption{Taxa bruta de frequência escolar da população até 6 anos (frequencia escolar pnad por idade)}\label{tab:frequencia-escolar-pnad-por-idade}\\
\toprule \textbf{Idade} & \textbf{Total} & \textbf{Homens} & \textbf{Mulheres} \\ \midrule \endfirsthead
\multicolumn{4}{@{}l}{\footnotesize\itshape (continuação)}\\ \toprule \textbf{Idade} & \textbf{Total} & \textbf{Homens} & \textbf{Mulheres} \\ \midrule \endhead
\midrule \multicolumn{4}{r@{}}{\footnotesize\itshape (continua)}\\ \endfoot
\bottomrule \endlastfoot
0 ano & 0,1 & 8,5 & 7,57 \\
1 ano & 0,3 & 31,1 & 29,46 \\
2 anos & 0,6 & 56,09 & 54,33 \\
3 anos & 0,7 & 72,49 & 71,47 \\
4 anos & 0,8 & 82,51 & 83,35 \\
5 anos & 0,9 & 91,56 & 89,49 \\
\end{longtable}}
\vspace{-6pt}\fonte{Elaboração IPP com dados de \citeonline{ibge_pnadc}.}


\begin{landscape}
{\scriptsize\setlength{\tabcolsep}{3pt}
\begin{longtable}{@{}rrrrrrrrrrrr@{}}
\caption{Cobertura vacinal de rotina em crianças até 2 anos (cobertura vacinal epi por ano)}\label{tab:cobertura-vacinal-epi-por-ano}\\
\toprule \textbf{Ano} & \textbf{BCG} & \textbf{DTP 1ºREF} & \textbf{FEBRE AMARELA} & \textbf{HEPATITE A} & \textbf{MENINGO C} & \textbf{PNEUMO 10} & \textbf{POLIOMIELITE} & \textbf{ROTAVÍRUS} & \textbf{TRÍPLICE VIRAL D1} & \textbf{TRÍPLICE VIRAL D2} & \textbf{VARICELA} \\ \midrule \endfirsthead
\multicolumn{12}{@{}l}{\footnotesize\itshape (continuação)}\\ \toprule \textbf{Ano} & \textbf{BCG} & \textbf{DTP 1ºREF} & \textbf{FEBRE AMARELA} & \textbf{HEPATITE A} & \textbf{MENINGO C} & \textbf{PNEUMO 10} & \textbf{POLIOMIELITE} & \textbf{ROTAVÍRUS} & \textbf{TRÍPLICE VIRAL D1} & \textbf{TRÍPLICE VIRAL D2} & \textbf{VARICELA} \\ \midrule \endhead
\midrule \multicolumn{12}{r@{}}{\footnotesize\itshape (continua)}\\ \endfoot
\bottomrule \endlastfoot
2016 & 110,3 & 79,5 & 0,0 & 73,3 & 103,0 & 119,8 & 101,2 & 99,2 & 105,5 & 79,4 & 79,4 \\
2017 & 128,1 & 85,1 & 0,0 & 92,3 & 97,3 & 107,8 & 103,8 & 96,2 & 102,9 & 80,1 & 89,7 \\
2018 & 114,2 & 61,5 & 52,5 & 91,2 & 97,3 & 101,3 & 98,7 & 99,4 & 105,6 & 76,1 & 79,1 \\
2019 & 78,8 & 47,4 & 69,0 & 83,4 & 80,6 & 82,7 & 78,7 & 78,4 & 100,4 & 81,3 & 83,5 \\
2020 & 81,0 & 77,4 & 60,1 & 75,0 & 76,3 & 80,0 & 76,4 & 76,9 & 81,5 & 54,0 & 93,8 \\
2021 & 70,7 & 71,1 & 65,3 & 74,1 & 74,7 & 77,5 & 75,5 & 81,9 & 79,0 & 54,2 & 95,1 \\
2022 & 103,4 & 59,9 & 54,6 & 69,1 & 67,2 & 68,6 & 66,8 & 65,8 & 79,0 & 70,7 & 76,5 \\
2023 & 121,3 & 77,2 & 72,7 & 92,7 & 92,2 & 92,4 & 91,5 & 91,0 & 96,7 & 66,7 & 90,2 \\
2024 & 101,4 & 88,4 & 76,9 & 93,7 & 94,7 & 95,7 & 95,3 & 93,8 & 94,8 & 80,4 & 81,5 \\
2025 & 108,6 & 88,1 & 80,2 & 83,9 & 94,7 & 96,3 & 93,6 & 95,4 & 102,3 & 85,7 & 87,4 \\
2026 & 99,5 & 73,7 & 74,6 & 77,8 & 84,8 & 86,2 & 84,8 & 85,7 & 89,1 & 74,7 & 78,6 \\
\end{longtable}}
\vspace{-6pt}\fonte{Elaboração IPP com dados de \citeonline{sms_rio_epi_vacinal}.}
\end{landscape}


{\small\setlength{\tabcolsep}{5pt}
\begin{longtable}{@{}rlr@{}}
\caption{Cobertura vacinal de rotina em crianças até 2 anos (cobertura vacinal epi comparativo anos)}\label{tab:cobertura-vacinal-epi-comparativo-anos}\\
\toprule \textbf{Ano} & \textbf{Imunobiologico} & \textbf{Cobertura} \\ \midrule \endfirsthead
\multicolumn{3}{@{}l}{\footnotesize\itshape (continuação)}\\ \toprule \textbf{Ano} & \textbf{Imunobiologico} & \textbf{Cobertura} \\ \midrule \endhead
\midrule \multicolumn{3}{r@{}}{\footnotesize\itshape (continua)}\\ \endfoot
\bottomrule \endlastfoot
2016 & BCG & 110,3 \\
2016 & DTP 1ºREF & 79,5 \\
2016 & FEBRE AMARELA & 0,0 \\
2016 & HEPATITE A & 73,3 \\
2016 & MENINGO C & 103,0 \\
2016 & PNEUMO 10 & 119,8 \\
2016 & POLIOMIELITE & 101,2 \\
2016 & ROTAVÍRUS & 99,2 \\
2016 & TRÍPLICE VIRAL D1 & 105,5 \\
2016 & TRÍPLICE VIRAL D2 & 79,4 \\
2016 & VARICELA & 79,4 \\
2019 & BCG & 78,8 \\
2019 & DTP 1ºREF & 47,4 \\
2019 & FEBRE AMARELA & 69,0 \\
2019 & HEPATITE A & 83,4 \\
2019 & MENINGO C & 80,6 \\
2019 & PNEUMO 10 & 82,7 \\
2019 & POLIOMIELITE & 78,7 \\
2019 & ROTAVÍRUS & 78,4 \\
2019 & TRÍPLICE VIRAL D1 & 100,4 \\
2019 & TRÍPLICE VIRAL D2 & 81,3 \\
2019 & VARICELA & 83,5 \\
2022 & BCG & 103,4 \\
2022 & DTP 1ºREF & 59,9 \\
2022 & FEBRE AMARELA & 54,6 \\
2022 & HEPATITE A & 69,1 \\
2022 & MENINGO C & 67,2 \\
2022 & PNEUMO 10 & 68,6 \\
2022 & POLIOMIELITE & 66,8 \\
2022 & ROTAVÍRUS & 65,8 \\
2022 & TRÍPLICE VIRAL D1 & 79,0 \\
2022 & TRÍPLICE VIRAL D2 & 70,7 \\
2022 & VARICELA & 76,5 \\
2025 & BCG & 108,6 \\
2025 & DTP 1ºREF & 88,1 \\
2025 & FEBRE AMARELA & 80,2 \\
2025 & HEPATITE A & 83,9 \\
2025 & MENINGO C & 94,7 \\
2025 & PNEUMO 10 & 96,3 \\
2025 & POLIOMIELITE & 93,6 \\
2025 & ROTAVÍRUS & 95,4 \\
2025 & TRÍPLICE VIRAL D1 & 102,3 \\
2025 & TRÍPLICE VIRAL D2 & 85,7 \\
2025 & VARICELA & 87,4 \\
\end{longtable}}
\vspace{-6pt}\fonte{Elaboração IPP com dados de \citeonline{sms_rio_epi_vacinal}.}


{\small\setlength{\tabcolsep}{5pt}
\begin{longtable}{@{}lr@{}}
\caption{Crianças até 6 anos frequentando escola/creche (geral) (sidra frequencia escola 0 5 total 2022)}\label{tab:sidra-frequencia-escola-0-5-total-2022}\\
\toprule \textbf{Idade} & \textbf{Crianças} \\ \midrule \endfirsthead
\multicolumn{2}{@{}l}{\footnotesize\itshape (continuação)}\\ \toprule \textbf{Idade} & \textbf{Crianças} \\ \midrule \endhead
\midrule \multicolumn{2}{r@{}}{\footnotesize\itshape (continua)}\\ \endfoot
\bottomrule \endlastfoot
0 ano & 4.358 \\
1 ano & 17.597 \\
2 anos & 34.907 \\
3 anos & 49.278 \\
4 anos & 61.206 \\
5 anos & 66.163 \\
\end{longtable}}
\vspace{-6pt}\fonte{Elaboração IPP com dados de \citeonline{ibge_censo2022}.}


\begin{landscape}
{\scriptsize\setlength{\tabcolsep}{3pt}
\begin{longtable}{@{}rrrrrrrrrrrr@{}}
\caption{Matrículas na educação básica de crianças de 0 a 5 anos (matriculas 0 a 5 por ano)}\label{tab:matriculas-0-a-5-por-ano}\\
\toprule \textbf{Ano} & \textbf{Matriculas} & \textbf{Matriculas 0 a 3} & \textbf{Matriculas 4 a 5} & \textbf{Matriculas publica} & \textbf{Matriculas privada} & \textbf{Populacao 0 a 3} & \textbf{Populacao 4 a 5} & \textbf{Populacao 0 a 5} & \textbf{Taxa atendimento 0 a 3} & \textbf{Taxa atendimento 4 a 5} & \textbf{Taxa atendimento 0 a 5} \\ \midrule \endfirsthead
\multicolumn{12}{@{}l}{\footnotesize\itshape (continuação)}\\ \toprule \textbf{Ano} & \textbf{Matriculas} & \textbf{Matriculas 0 a 3} & \textbf{Matriculas 4 a 5} & \textbf{Matriculas publica} & \textbf{Matriculas privada} & \textbf{Populacao 0 a 3} & \textbf{Populacao 4 a 5} & \textbf{Populacao 0 a 5} & \textbf{Taxa atendimento 0 a 3} & \textbf{Taxa atendimento 4 a 5} & \textbf{Taxa atendimento 0 a 5} \\ \midrule \endhead
\midrule \multicolumn{12}{r@{}}{\footnotesize\itshape (continua)}\\ \endfoot
\bottomrule \endlastfoot
2007 & 170.650 & 60.997 & 109.653 & 107.599 & 63.051 & 323.157 & 173.017 & 496.174 & 18,9 & 63,4 & 34,4 \\
2008 & 196.904 & 74.455 & 122.449 & 104.919 & 91.985 & 321.319 & 168.111 & 489.430 & 23,2 & 72,8 & 40,2 \\
2009 & 193.124 & 72.318 & 120.806 & 104.076 & 89.048 & 320.426 & 165.119 & 485.545 & 22,6 & 73,2 & 39,8 \\
2010 & 202.850 & 80.325 & 122.525 & 105.600 & 97.250 & 319.474 & 163.196 & 482.670 & 25,1 & 75,1 & 42,0 \\
2011 & 222.432 & 95.636 & 126.796 & 116.243 & 106.189 & 319.284 & 161.364 & 480.648 & 30,0 & 78,6 & 46,3 \\
2012 & 228.523 & 100.517 & 128.006 & 120.029 & 108.494 & 320.584 & 159.958 & 480.542 & 31,4 & 80,0 & 47,6 \\
2013 & 241.730 & 108.859 & 132.871 & 125.616 & 116.114 & 322.541 & 159.293 & 481.834 & 33,8 & 83,4 & 50,2 \\
2014 & 247.280 & 111.591 & 135.689 & 127.060 & 120.220 & 326.749 & 158.702 & 485.451 & 34,2 & 85,5 & 50,9 \\
2015 & 243.176 & 109.873 & 133.303 & 127.291 & 115.885 & 332.546 & 159.135 & 491.681 & 33,0 & 83,8 & 49,5 \\
2016 & 247.022 & 110.358 & 136.664 & 130.762 & 116.260 & 334.015 & 161.298 & 495.313 & 33,0 & 84,7 & 49,9 \\
2017 & 253.641 & 114.609 & 139.032 & 139.031 & 114.610 & 331.642 & 163.003 & 494.645 & 34,6 & 85,3 & 51,3 \\
2018 & 253.518 & 113.789 & 139.729 & 139.935 & 113.583 & 327.060 & 164.808 & 491.868 & 34,8 & 84,8 & 51,5 \\
2019 & 258.868 & 115.723 & 143.145 & 144.314 & 114.554 & 317.445 & 168.360 & 485.805 & 36,5 & 85,0 & 53,3 \\
2020 & 247.133 & 107.422 & 139.711 & 143.909 & 103.224 & 307.006 & 167.316 & 474.322 & 35,0 & 83,5 & 52,1 \\
2021 & 230.260 & 101.421 & 128.839 & 140.722 & 89.538 & 296.299 & 160.981 & 457.280 & 34,2 & 80,0 & 50,4 \\
2022 & 246.522 & 114.822 & 131.700 & 136.691 & 109.831 & 282.173 & 157.734 & 439.907 & 40,7 & 83,5 & 56,0 \\
2023 & 246.119 & 113.935 & 132.184 & 131.935 & 114.184 & 268.654 & 155.085 & 423.739 & 42,4 & 85,2 & 58,1 \\
2024 & 242.575 & 114.934 & 127.641 & 127.851 & 114.724 & 259.038 & 148.499 & 407.537 & 44,4 & 86,0 & 59,5 \\
2025 & 230.284 & 110.097 & 120.187 & 122.236 & 108.048 & 251.936 & 141.137 & 393.073 & 43,7 & 85,2 & 58,6 \\
\end{longtable}}
\vspace{-6pt}\fonte{Elaboração IPP com dados de \citeonline{ms_ripsa_populacao}; \citeonline{inep_censo_escolar}.}
\end{landscape}


\chapter{Tabelas do eixo Proteção}\label{ap:eixo-4}

{\small\setlength{\tabcolsep}{5pt}
\begin{longtable}{@{}rlrrr@{}}
\caption{Violência territorial (violencia territorial ra 2024)}\label{tab:tabela-mapa-violencia-territorial-ra-2024}\\
\toprule \textbf{Codra} & \textbf{Regiao adm} & \textbf{Taxa homicidios} & \textbf{Homicidios acao policial} & \textbf{Homicidios jovens negros} \\ \midrule \endfirsthead
\multicolumn{5}{@{}l}{\footnotesize\itshape (continuação)}\\ \toprule \textbf{Codra} & \textbf{Regiao adm} & \textbf{Taxa homicidios} & \textbf{Homicidios acao policial} & \textbf{Homicidios jovens negros} \\ \midrule \endhead
\midrule \multicolumn{5}{r@{}}{\footnotesize\itshape (continua)}\\ \endfoot
\bottomrule \endlastfoot
1 & PORTUARIA & 19,0 & 8,4 & 4,2 \\
2 & CENTRO & 76,3 & 9,2 & 6,1 \\
3 & RIO COMPRIDO & 18,7 & 6,2 & 3,7 \\
4 & BOTAFOGO & 4,2 & 0,5 & 0,5 \\
5 & COPACABANA & 5,7 & 0,0 & 0,7 \\
6 & LAGOA & 4,7 & 1,3 & 0,7 \\
7 & SAO CRISTOVAO & 19,5 & 2,4 & 4,9 \\
8 & TIJUCA & 4,4 & 1,3 & 0,0 \\
9 & VILA ISABEL & 8,3 & 0,6 & 0,6 \\
10 & RAMOS & 16,7 & 7,6 & 3,0 \\
11 & PENHA & 27,8 & 15,0 & 5,7 \\
12 & INHAUMA & 23,4 & 4,8 & 4,0 \\
13 & MEIER & 12,2 & 9,0 & 3,1 \\
14 & IRAJA & 22,5 & 12,9 & 5,6 \\
15 & MADUREIRA & 33,5 & 11,5 & 9,2 \\
16 & JACAREPAGUA & 24,2 & 11,4 & 4,4 \\
17 & BANGU & 15,4 & 11,1 & 3,5 \\
18 & CAMPO GRANDE & 6,7 & 0,5 & 1,3 \\
19 & SANTA CRUZ & 17,9 & 0,7 & 2,4 \\
20 & ILHA DO GOVERNADOR & 11,4 & 1,1 & 3,2 \\
22 & ANCHIETA & 20,3 & 2,8 & 3,5 \\
23 & SANTA TEREZA & 5,5 & 0,0 & 0,0 \\
24 & BARRA DA TIJUCA & 20,5 & 5,6 & 5,8 \\
25 & PAVUNA & 29,5 & 10,2 & 6,8 \\
26 & GUARATIBA & 16,4 & 1,1 & 4,0 \\
27 & ROCINHA & 2,8 & 0,0 & 1,4 \\
28 & JACAREZINHO & 8,5 & 2,8 & 2,8 \\
29 & COMPLEXO DO ALEMÃO & 3,7 & 0,0 & 1,8 \\
30 & COMPLEXO DA MARE & 4,8 & 2,4 & 0,8 \\
31 & VIGARIO GERAL & 12,5 & 6,3 & 0,9 \\
33 & REALENGO & 23,4 & 7,5 & 5,3 \\
34 & CIDADE DE DEUS & 22,9 & 19,6 & 6,5 \\
\end{longtable}}
\vspace{-6pt}\fonte{Elaboração IPP com dados de \citeonline{ipp_ips2024}.}


\begin{landscape}
{\scriptsize\setlength{\tabcolsep}{3pt}
\begin{longtable}{@{}rrrrrrrrr@{}}
\caption{Violência familiar (0 a 5 anos, por vínculo do provável autor) (violencia familiar por vinculo ano)}\label{tab:violencia-familiar-por-vinculo-ano}\\
\toprule \textbf{Ano} & \textbf{Mae} & \textbf{Pai} & \textbf{Padrasto} & \textbf{Irmao} & \textbf{Conjuge} & \textbf{Exconjuge} & \textbf{Filho} & \textbf{Outros} \\ \midrule \endfirsthead
\multicolumn{9}{@{}l}{\footnotesize\itshape (continuação)}\\ \toprule \textbf{Ano} & \textbf{Mae} & \textbf{Pai} & \textbf{Padrasto} & \textbf{Irmao} & \textbf{Conjuge} & \textbf{Exconjuge} & \textbf{Filho} & \textbf{Outros} \\ \midrule \endhead
\midrule \multicolumn{9}{r@{}}{\footnotesize\itshape (continua)}\\ \endfoot
\bottomrule \endlastfoot
2011 & 324 & 135 & 16 & 6 & 0 & 0 & 0 & 22 \\
2012 & 464 & 188 & 19 & 7 & 0 & 0 & 0 & 26 \\
2013 & 429 & 222 & 33 & 15 & 0 & 0 & 1 & 49 \\
2014 & 564 & 246 & 24 & 8 & 0 & 0 & 0 & 32 \\
2015 & 452 & 234 & 20 & 17 & 0 & 0 & 0 & 37 \\
2016 & 600 & 371 & 37 & 11 & 0 & 0 & 0 & 48 \\
2017 & 1.514 & 1.261 & 46 & 26 & 0 & 0 & 0 & 72 \\
2018 & 1.265 & 954 & 75 & 25 & 9 & 2 & 3 & 114 \\
2019 & 868 & 522 & 48 & 20 & 2 & 3 & 2 & 75 \\
2020 & 734 & 464 & 34 & 11 & 4 & 1 & 2 & 52 \\
2021 & 805 & 581 & 60 & 14 & 5 & 2 & 5 & 86 \\
2022 & 1.359 & 1.003 & 53 & 27 & 0 & 0 & 0 & 80 \\
2023 & 1.469 & 1.068 & 54 & 28 & 0 & 0 & 0 & 82 \\
2024 & 1.515 & 1.105 & 63 & 27 & 0 & 0 & 0 & 90 \\
2025 & 1.756 & 1.404 & 73 & 26 & 5 & 3 & 3 & 110 \\
\end{longtable}}
\vspace{-6pt}\fonte{Elaboração IPP com dados de \citeonline{sms_rio_sinan}.}
\end{landscape}


{\small\setlength{\tabcolsep}{5pt}
\begin{longtable}{@{}rlr@{}}
\caption{Violência familiar (0 a 5 anos, por vínculo do provável autor) (violencia familiar mae 2025)}\label{tab:tabela-mapa-violencia-familiar-mae-2025}\\
\toprule \textbf{Codbairro} & \textbf{Bairro} & \textbf{Mae} \\ \midrule \endfirsthead
\multicolumn{3}{@{}l}{\footnotesize\itshape (continuação)}\\ \toprule \textbf{Codbairro} & \textbf{Bairro} & \textbf{Mae} \\ \midrule \endhead
\midrule \multicolumn{3}{r@{}}{\footnotesize\itshape (continua)}\\ \endfoot
\bottomrule \endlastfoot
1 & Saúde & 1 \\
2 & Gamboa & 2 \\
3 & Santo Cristo & 3 \\
4 & Caju & 1 \\
5 & Centro & 3 \\
6 & Catumbi & 1 \\
7 & Rio Comprido & 5 \\
8 & Cidade Nova & 0 \\
9 & Estácio & 2 \\
10 & Imperial de São Cristóvão & 4 \\
11 & Mangueira & 4 \\
12 & Benfica & 12 \\
13 & Paquetá & 0 \\
14 & Santa Teresa & 3 \\
15 & Flamengo & 2 \\
16 & Glória & 0 \\
17 & Laranjeiras & 0 \\
18 & Catete & 2 \\
19 & Cosme Velho & 0 \\
20 & Botafogo & 5 \\
21 & Humaitá & 0 \\
22 & Urca & 0 \\
23 & Leme & 4 \\
24 & Copacabana & 3 \\
25 & Ipanema & 1 \\
26 & Leblon & 2 \\
27 & Lagoa & 0 \\
28 & Jardim Botânico & 1 \\
29 & Gávea & 1 \\
30 & Vidigal & 3 \\
31 & São Conrado & 0 \\
32 & Praça da Bandeira & 0 \\
33 & Tijuca & 16 \\
34 & Alto da Boa Vista & 0 \\
35 & Maracanã & 0 \\
36 & Vila Isabel & 0 \\
37 & Andaraí & 0 \\
38 & Grajaú & 3 \\
39 & Manguinhos & 8 \\
40 & Bonsucesso & 5 \\
41 & Ramos & 14 \\
42 & Olaria & 17 \\
43 & Penha & 23 \\
44 & Penha Circular & 6 \\
45 & Brás de Pina & 34 \\
46 & Cordovil & 10 \\
47 & Parada de Lucas & 14 \\
48 & Vigário Geral & 14 \\
49 & Jardim América & 7 \\
50 & Higienópolis & 15 \\
51 & Jacaré & 4 \\
52 & Maria da Graça & 0 \\
53 & Del Castilho & 14 \\
54 & Inhaúma & 4 \\
55 & Engenho da Rainha & 6 \\
56 & Tomás Coelho & 12 \\
57 & São Francisco Xavier & 5 \\
58 & Rocha & 4 \\
59 & Riachuelo & 0 \\
60 & Sampaio & 2 \\
61 & Engenho Novo & 3 \\
62 & Lins de Vasconcelos & 3 \\
63 & Méier & 21 \\
64 & Todos os Santos & 6 \\
65 & Cachambi & 13 \\
66 & Engenho de Dentro & 2 \\
67 & Água Santa & 6 \\
68 & Encantado & 13 \\
69 & Piedade & 1 \\
70 & Abolição & 2 \\
71 & Pilares & 4 \\
72 & Vila Kosmos & 1 \\
73 & Vicente de Carvalho & 11 \\
74 & Vila da Penha & 11 \\
75 & Vista Alegre & 7 \\
76 & Irajá & 6 \\
77 & Colégio & 2 \\
78 & Campinho & 24 \\
79 & Quintino Bocaiúva & 14 \\
80 & Cavalcanti & 3 \\
81 & Engenheiro Leal & 13 \\
82 & Cascadura & 0 \\
83 & Madureira & 0 \\
84 & Vaz Lobo & 10 \\
85 & Turiaçú & 13 \\
86 & Rocha Miranda & 5 \\
87 & Honório Gurgel & 5 \\
88 & Osvaldo Cruz & 7 \\
89 & Bento Ribeiro & 7 \\
90 & Marechal Hermes & 2 \\
91 & Ribeira & 8 \\
92 & Zumbi & 25 \\
93 & Cacuia & 0 \\
94 & Pitangueiras & 0 \\
95 & Praia da Bandeira & 3 \\
96 & Cocotá & 1 \\
97 & Bancários & 0 \\
98 & Freguesia (Ilha) & 0 \\
99 & Jardim Guanabara & 2 \\
100 & Jardim Carioca & 2 \\
101 & Tauá & 2 \\
102 & Moneró & 3 \\
103 & Portuguesa & 4 \\
104 & Galeão & 0 \\
105 & Cidade Universitária & 3 \\
106 & Guadalupe & 9 \\
107 & Anchieta & 0 \\
108 & Parque Anchieta & 22 \\
109 & Ricardo de Albuquerque & 38 \\
110 & Coelho Neto & 9 \\
111 & Acari & 17 \\
112 & Barros Filho & 19 \\
113 & Costa Barros & 27 \\
114 & Pavuna & 7 \\
115 & Jacarepaguá & 16 \\
116 & Anil & 27 \\
117 & Gardênia Azul & 33 \\
118 & Cidade de Deus & 10 \\
119 & Curicica & 12 \\
120 & Freguesia (Jacarepaguá) & 15 \\
121 & Pechincha & 5 \\
122 & Taquara & 1 \\
123 & Tanque & 1 \\
124 & Praça Seca & 10 \\
125 & Vila Valqueire & 2 \\
126 & Joá & 11 \\
127 & Itanhangá & 4 \\
128 & Barra da Tijuca & 0 \\
129 & Camorim & 10 \\
130 & Vargem Pequena & 5 \\
131 & Vargem Grande & 0 \\
132 & Recreio dos Bandeirantes & 4 \\
133 & Grumari & 0 \\
134 & Deodoro & 6 \\
135 & Vila Militar & 0 \\
136 & Campo dos Afonsos & 4 \\
137 & Jardim Sulacap & 1 \\
138 & Magalhães Bastos & 1 \\
139 & Realengo & 9 \\
140 & Padre Miguel & 13 \\
141 & Bangu & 127 \\
142 & Senador Camará & 60 \\
143 & Santíssimo & 165 \\
144 & Campo Grande & 49 \\
145 & Senador Vasconcelos & 23 \\
146 & Inhoaíba & 81 \\
147 & Cosmos & 4 \\
148 & Paciência & 13 \\
149 & Santa Cruz & 23 \\
150 & Sepetiba & 27 \\
151 & Guaratiba & 116 \\
152 & Barra de Guaratiba & 43 \\
153 & Pedra de Guaratiba & 50 \\
154 & Rocinha & 4 \\
155 & Jacarezinho & 2 \\
156 & Complexo do Alemão & 27 \\
157 & Maré & 6 \\
158 & Vasco da Gama & 2 \\
159 & Parque Colúmbia & 16 \\
160 & Gericinó & 0 \\
161 & Lapa & 0 \\
162 & Vila Kennedy & 0 \\
163 & Jabour & 0 \\
164 & Ilha de Guaratiba & 0 \\
165 & Barra Olímpica & 0 \\
166 & Argentino & 0 \\
\end{longtable}}
\vspace{-6pt}\fonte{Elaboração IPP com dados de \citeonline{sms_rio_sinan}.}


{\small\setlength{\tabcolsep}{5pt}
\begin{longtable}{@{}rlr@{}}
\caption{Violência familiar (0 a 5 anos, por vínculo do provável autor) (violencia familiar pai 2025)}\label{tab:tabela-mapa-violencia-familiar-pai-2025}\\
\toprule \textbf{Codbairro} & \textbf{Bairro} & \textbf{Pai} \\ \midrule \endfirsthead
\multicolumn{3}{@{}l}{\footnotesize\itshape (continuação)}\\ \toprule \textbf{Codbairro} & \textbf{Bairro} & \textbf{Pai} \\ \midrule \endhead
\midrule \multicolumn{3}{r@{}}{\footnotesize\itshape (continua)}\\ \endfoot
\bottomrule \endlastfoot
1 & Saúde & 0 \\
2 & Gamboa & 0 \\
3 & Santo Cristo & 0 \\
4 & Caju & 3 \\
5 & Centro & 2 \\
6 & Catumbi & 0 \\
7 & Rio Comprido & 3 \\
8 & Cidade Nova & 0 \\
9 & Estácio & 3 \\
10 & Imperial de São Cristóvão & 1 \\
11 & Mangueira & 2 \\
12 & Benfica & 10 \\
13 & Paquetá & 0 \\
14 & Santa Teresa & 4 \\
15 & Flamengo & 2 \\
16 & Glória & 0 \\
17 & Laranjeiras & 0 \\
18 & Catete & 1 \\
19 & Cosme Velho & 0 \\
20 & Botafogo & 4 \\
21 & Humaitá & 0 \\
22 & Urca & 1 \\
23 & Leme & 3 \\
24 & Copacabana & 3 \\
25 & Ipanema & 1 \\
26 & Leblon & 1 \\
27 & Lagoa & 0 \\
28 & Jardim Botânico & 1 \\
29 & Gávea & 2 \\
30 & Vidigal & 5 \\
31 & São Conrado & 0 \\
32 & Praça da Bandeira & 0 \\
33 & Tijuca & 9 \\
34 & Alto da Boa Vista & 0 \\
35 & Maracanã & 0 \\
36 & Vila Isabel & 0 \\
37 & Andaraí & 0 \\
38 & Grajaú & 6 \\
39 & Manguinhos & 3 \\
40 & Bonsucesso & 4 \\
41 & Ramos & 10 \\
42 & Olaria & 10 \\
43 & Penha & 11 \\
44 & Penha Circular & 1 \\
45 & Brás de Pina & 19 \\
46 & Cordovil & 6 \\
47 & Parada de Lucas & 10 \\
48 & Vigário Geral & 10 \\
49 & Jardim América & 6 \\
50 & Higienópolis & 8 \\
51 & Jacaré & 3 \\
52 & Maria da Graça & 0 \\
53 & Del Castilho & 8 \\
54 & Inhaúma & 3 \\
55 & Engenho da Rainha & 3 \\
56 & Tomás Coelho & 9 \\
57 & São Francisco Xavier & 3 \\
58 & Rocha & 4 \\
59 & Riachuelo & 0 \\
60 & Sampaio & 2 \\
61 & Engenho Novo & 1 \\
62 & Lins de Vasconcelos & 2 \\
63 & Méier & 18 \\
64 & Todos os Santos & 4 \\
65 & Cachambi & 9 \\
66 & Engenho de Dentro & 2 \\
67 & Água Santa & 6 \\
68 & Encantado & 12 \\
69 & Piedade & 0 \\
70 & Abolição & 2 \\
71 & Pilares & 3 \\
72 & Vila Kosmos & 1 \\
73 & Vicente de Carvalho & 7 \\
74 & Vila da Penha & 3 \\
75 & Vista Alegre & 6 \\
76 & Irajá & 4 \\
77 & Colégio & 1 \\
78 & Campinho & 14 \\
79 & Quintino Bocaiúva & 9 \\
80 & Cavalcanti & 2 \\
81 & Engenheiro Leal & 7 \\
82 & Cascadura & 0 \\
83 & Madureira & 0 \\
84 & Vaz Lobo & 8 \\
85 & Turiaçú & 8 \\
86 & Rocha Miranda & 3 \\
87 & Honório Gurgel & 1 \\
88 & Osvaldo Cruz & 4 \\
89 & Bento Ribeiro & 6 \\
90 & Marechal Hermes & 2 \\
91 & Ribeira & 4 \\
92 & Zumbi & 18 \\
93 & Cacuia & 0 \\
94 & Pitangueiras & 0 \\
95 & Praia da Bandeira & 2 \\
96 & Cocotá & 3 \\
97 & Bancários & 0 \\
98 & Freguesia (Ilha) & 1 \\
99 & Jardim Guanabara & 2 \\
100 & Jardim Carioca & 1 \\
101 & Tauá & 1 \\
102 & Moneró & 4 \\
103 & Portuguesa & 0 \\
104 & Galeão & 0 \\
105 & Cidade Universitária & 1 \\
106 & Guadalupe & 2 \\
107 & Anchieta & 0 \\
108 & Parque Anchieta & 20 \\
109 & Ricardo de Albuquerque & 35 \\
110 & Coelho Neto & 8 \\
111 & Acari & 16 \\
112 & Barros Filho & 17 \\
113 & Costa Barros & 15 \\
114 & Pavuna & 6 \\
115 & Jacarepaguá & 9 \\
116 & Anil & 19 \\
117 & Gardênia Azul & 25 \\
118 & Cidade de Deus & 7 \\
119 & Curicica & 9 \\
120 & Freguesia (Jacarepaguá) & 10 \\
121 & Pechincha & 6 \\
122 & Taquara & 1 \\
123 & Tanque & 3 \\
124 & Praça Seca & 10 \\
125 & Vila Valqueire & 1 \\
126 & Joá & 10 \\
127 & Itanhangá & 4 \\
128 & Barra da Tijuca & 0 \\
129 & Camorim & 5 \\
130 & Vargem Pequena & 3 \\
131 & Vargem Grande & 0 \\
132 & Recreio dos Bandeirantes & 3 \\
133 & Grumari & 1 \\
134 & Deodoro & 7 \\
135 & Vila Militar & 1 \\
136 & Campo dos Afonsos & 6 \\
137 & Jardim Sulacap & 2 \\
138 & Magalhães Bastos & 1 \\
139 & Realengo & 9 \\
140 & Padre Miguel & 12 \\
141 & Bangu & 121 \\
142 & Senador Camará & 53 \\
143 & Santíssimo & 132 \\
144 & Campo Grande & 50 \\
145 & Senador Vasconcelos & 22 \\
146 & Inhoaíba & 92 \\
147 & Cosmos & 5 \\
148 & Paciência & 11 \\
149 & Santa Cruz & 20 \\
150 & Sepetiba & 23 \\
151 & Guaratiba & 100 \\
152 & Barra de Guaratiba & 34 \\
153 & Pedra de Guaratiba & 37 \\
154 & Rocinha & 4 \\
155 & Jacarezinho & 3 \\
156 & Complexo do Alemão & 14 \\
157 & Maré & 5 \\
158 & Vasco da Gama & 1 \\
159 & Parque Colúmbia & 12 \\
160 & Gericinó & 0 \\
161 & Lapa & 0 \\
162 & Vila Kennedy & 0 \\
163 & Jabour & 0 \\
164 & Ilha de Guaratiba & 0 \\
165 & Barra Olímpica & 0 \\
166 & Argentino & 0 \\
\end{longtable}}
\vspace{-6pt}\fonte{Elaboração IPP com dados de \citeonline{sms_rio_sinan}.}


{\small\setlength{\tabcolsep}{5pt}
\begin{longtable}{@{}rlr@{}}
\caption{Violência familiar (0 a 5 anos, por vínculo do provável autor) (violencia familiar outros 2021 2025)}\label{tab:tabela-mapa-violencia-familiar-outros-2021-2025}\\
\toprule \textbf{Codbairro} & \textbf{Bairro} & \textbf{Outros 2021 2025} \\ \midrule \endfirsthead
\multicolumn{3}{@{}l}{\footnotesize\itshape (continuação)}\\ \toprule \textbf{Codbairro} & \textbf{Bairro} & \textbf{Outros 2021 2025} \\ \midrule \endhead
\midrule \multicolumn{3}{r@{}}{\footnotesize\itshape (continua)}\\ \endfoot
\bottomrule \endlastfoot
1 & Saúde & 0 \\
2 & Gamboa & 0 \\
3 & Santo Cristo & 0 \\
4 & Caju & 2 \\
5 & Centro & 3 \\
6 & Catumbi & 3 \\
7 & Rio Comprido & 0 \\
8 & Cidade Nova & 0 \\
9 & Estácio & 1 \\
10 & Imperial de São Cristóvão & 2 \\
11 & Mangueira & 1 \\
12 & Benfica & 4 \\
13 & Paquetá & 0 \\
14 & Santa Teresa & 1 \\
15 & Flamengo & 0 \\
16 & Glória & 0 \\
17 & Laranjeiras & 0 \\
18 & Catete & 0 \\
19 & Cosme Velho & 0 \\
20 & Botafogo & 1 \\
21 & Humaitá & 0 \\
22 & Urca & 0 \\
23 & Leme & 1 \\
24 & Copacabana & 3 \\
25 & Ipanema & 0 \\
26 & Leblon & 0 \\
27 & Lagoa & 0 \\
28 & Jardim Botânico & 0 \\
29 & Gávea & 1 \\
30 & Vidigal & 2 \\
31 & São Conrado & 0 \\
32 & Praça da Bandeira & 0 \\
33 & Tijuca & 7 \\
34 & Alto da Boa Vista & 0 \\
35 & Maracanã & 0 \\
36 & Vila Isabel & 0 \\
37 & Andaraí & 0 \\
38 & Grajaú & 4 \\
39 & Manguinhos & 1 \\
40 & Bonsucesso & 2 \\
41 & Ramos & 2 \\
42 & Olaria & 8 \\
43 & Penha & 3 \\
44 & Penha Circular & 3 \\
45 & Brás de Pina & 5 \\
46 & Cordovil & 3 \\
47 & Parada de Lucas & 0 \\
48 & Vigário Geral & 0 \\
49 & Jardim América & 1 \\
50 & Higienópolis & 3 \\
51 & Jacaré & 2 \\
52 & Maria da Graça & 1 \\
53 & Del Castilho & 3 \\
54 & Inhaúma & 0 \\
55 & Engenho da Rainha & 0 \\
56 & Tomás Coelho & 2 \\
57 & São Francisco Xavier & 3 \\
58 & Rocha & 1 \\
59 & Riachuelo & 0 \\
60 & Sampaio & 1 \\
61 & Engenho Novo & 0 \\
62 & Lins de Vasconcelos & 4 \\
63 & Méier & 2 \\
64 & Todos os Santos & 1 \\
65 & Cachambi & 0 \\
66 & Engenho de Dentro & 0 \\
67 & Água Santa & 0 \\
68 & Encantado & 2 \\
69 & Piedade & 1 \\
70 & Abolição & 0 \\
71 & Pilares & 2 \\
72 & Vila Kosmos & 0 \\
73 & Vicente de Carvalho & 1 \\
74 & Vila da Penha & 2 \\
75 & Vista Alegre & 2 \\
76 & Irajá & 1 \\
77 & Colégio & 0 \\
78 & Campinho & 4 \\
79 & Quintino Bocaiúva & 3 \\
80 & Cavalcanti & 0 \\
81 & Engenheiro Leal & 2 \\
82 & Cascadura & 0 \\
83 & Madureira & 0 \\
84 & Vaz Lobo & 1 \\
85 & Turiaçú & 0 \\
86 & Rocha Miranda & 1 \\
87 & Honório Gurgel & 0 \\
88 & Osvaldo Cruz & 1 \\
89 & Bento Ribeiro & 1 \\
90 & Marechal Hermes & 1 \\
91 & Ribeira & 1 \\
92 & Zumbi & 0 \\
93 & Cacuia & 1 \\
94 & Pitangueiras & 0 \\
95 & Praia da Bandeira & 2 \\
96 & Cocotá & 1 \\
97 & Bancários & 0 \\
98 & Freguesia (Ilha) & 0 \\
99 & Jardim Guanabara & 0 \\
100 & Jardim Carioca & 2 \\
101 & Tauá & 0 \\
102 & Moneró & 0 \\
103 & Portuguesa & 2 \\
104 & Galeão & 0 \\
105 & Cidade Universitária & 0 \\
106 & Guadalupe & 3 \\
107 & Anchieta & 2 \\
108 & Parque Anchieta & 5 \\
109 & Ricardo de Albuquerque & 8 \\
110 & Coelho Neto & 6 \\
111 & Acari & 4 \\
112 & Barros Filho & 1 \\
113 & Costa Barros & 3 \\
114 & Pavuna & 6 \\
115 & Jacarepaguá & 6 \\
116 & Anil & 5 \\
117 & Gardênia Azul & 6 \\
118 & Cidade de Deus & 3 \\
119 & Curicica & 4 \\
120 & Freguesia (Jacarepaguá) & 6 \\
121 & Pechincha & 7 \\
122 & Taquara & 2 \\
123 & Tanque & 1 \\
124 & Praça Seca & 8 \\
125 & Vila Valqueire & 0 \\
126 & Joá & 2 \\
127 & Itanhangá & 0 \\
128 & Barra da Tijuca & 0 \\
129 & Camorim & 2 \\
130 & Vargem Pequena & 3 \\
131 & Vargem Grande & 0 \\
132 & Recreio dos Bandeirantes & 1 \\
133 & Grumari & 1 \\
134 & Deodoro & 2 \\
135 & Vila Militar & 0 \\
136 & Campo dos Afonsos & 0 \\
137 & Jardim Sulacap & 1 \\
138 & Magalhães Bastos & 0 \\
139 & Realengo & 0 \\
140 & Padre Miguel & 1 \\
141 & Bangu & 11 \\
142 & Senador Camará & 3 \\
143 & Santíssimo & 26 \\
144 & Campo Grande & 16 \\
145 & Senador Vasconcelos & 3 \\
146 & Inhoaíba & 26 \\
147 & Cosmos & 3 \\
148 & Paciência & 4 \\
149 & Santa Cruz & 8 \\
150 & Sepetiba & 15 \\
151 & Guaratiba & 50 \\
152 & Barra de Guaratiba & 18 \\
153 & Pedra de Guaratiba & 18 \\
154 & Rocinha & 2 \\
155 & Jacarezinho & 5 \\
156 & Complexo do Alemão & 5 \\
157 & Maré & 1 \\
158 & Vasco da Gama & 4 \\
159 & Parque Colúmbia & 10 \\
160 & Gericinó & 0 \\
161 & Lapa & 0 \\
162 & Vila Kennedy & 0 \\
163 & Jabour & 0 \\
164 & Ilha de Guaratiba & 0 \\
165 & Barra Olímpica & 0 \\
166 & Argentino & 0 \\
\end{longtable}}
\vspace{-6pt}\fonte{Elaboração IPP com dados de \citeonline{sms_rio_sinan}.}


{\small\setlength{\tabcolsep}{5pt}
\begin{longtable}{@{}rrrrrrr@{}}
\caption{Violência familiar — composição de "outros" vínculos (violencia familiar outros detalhe)}\label{tab:violencia-familiar-outros-detalhe}\\
\toprule \textbf{Ano} & \textbf{Padrasto} & \textbf{Irmao} & \textbf{Conjuge} & \textbf{Exconjuge} & \textbf{Filho} & \textbf{Outros} \\ \midrule \endfirsthead
\multicolumn{7}{@{}l}{\footnotesize\itshape (continuação)}\\ \toprule \textbf{Ano} & \textbf{Padrasto} & \textbf{Irmao} & \textbf{Conjuge} & \textbf{Exconjuge} & \textbf{Filho} & \textbf{Outros} \\ \midrule \endhead
\midrule \multicolumn{7}{r@{}}{\footnotesize\itshape (continua)}\\ \endfoot
\bottomrule \endlastfoot
2011 & 16 & 6 & 0 & 0 & 0 & 22 \\
2012 & 19 & 7 & 0 & 0 & 0 & 26 \\
2013 & 33 & 15 & 0 & 0 & 1 & 49 \\
2014 & 24 & 8 & 0 & 0 & 0 & 32 \\
2015 & 20 & 17 & 0 & 0 & 0 & 37 \\
2016 & 37 & 11 & 0 & 0 & 0 & 48 \\
2017 & 46 & 26 & 0 & 0 & 0 & 72 \\
2018 & 75 & 25 & 9 & 2 & 3 & 114 \\
2019 & 48 & 20 & 2 & 3 & 2 & 75 \\
2020 & 34 & 11 & 4 & 1 & 2 & 52 \\
2021 & 60 & 14 & 5 & 2 & 5 & 86 \\
2022 & 53 & 27 & 0 & 0 & 0 & 80 \\
2023 & 54 & 28 & 0 & 0 & 0 & 82 \\
2024 & 63 & 27 & 0 & 0 & 0 & 90 \\
2025 & 73 & 26 & 5 & 3 & 3 & 110 \\
\end{longtable}}
\vspace{-6pt}\fonte{Elaboração IPP com dados de \citeonline{sms_rio_sinan}.}


{\small\setlength{\tabcolsep}{5pt}
\begin{longtable}{@{}rllr@{}}
\caption{Violência familiar — bairros com mais notificações (2025) (violencia familiar top bairros 2025)}\label{tab:violencia-familiar-top-bairros-2025}\\
\toprule \textbf{Codbairro} & \textbf{Bairro} & \textbf{Vinculo} & \textbf{Notificações} \\ \midrule \endfirsthead
\multicolumn{4}{@{}l}{\footnotesize\itshape (continuação)}\\ \toprule \textbf{Codbairro} & \textbf{Bairro} & \textbf{Vinculo} & \textbf{Notificações} \\ \midrule \endhead
\midrule \multicolumn{4}{r@{}}{\footnotesize\itshape (continua)}\\ \endfoot
\bottomrule \endlastfoot
143 & Santíssimo & Mãe & 165 \\
141 & Bangu & Mãe & 127 \\
151 & Guaratiba & Mãe & 116 \\
146 & Inhoaíba & Mãe & 81 \\
142 & Senador Camará & Mãe & 60 \\
153 & Pedra de Guaratiba & Mãe & 50 \\
144 & Campo Grande & Mãe & 49 \\
152 & Barra de Guaratiba & Mãe & 43 \\
109 & Ricardo de Albuquerque & Mãe & 38 \\
45 & Brás de Pina & Mãe & 34 \\
143 & Santíssimo & Pai & 132 \\
141 & Bangu & Pai & 121 \\
151 & Guaratiba & Pai & 100 \\
146 & Inhoaíba & Pai & 92 \\
142 & Senador Camará & Pai & 53 \\
153 & Pedra de Guaratiba & Pai & 37 \\
144 & Campo Grande & Pai & 50 \\
152 & Barra de Guaratiba & Pai & 34 \\
109 & Ricardo de Albuquerque & Pai & 35 \\
45 & Brás de Pina & Pai & 19 \\
\end{longtable}}
\vspace{-6pt}\fonte{Elaboração IPP com dados de \citeonline{sms_rio_sinan}.}


\begin{landscape}
{\scriptsize\setlength{\tabcolsep}{3pt}
\begin{longtable}{@{}rrrrrrrrr@{}}
\caption{Violência familiar — por CAP (violencia familiar por cap)}\label{tab:violencia-familiar-por-cap}\\
\toprule \textbf{Cod ap sms} & \textbf{Ano} & \textbf{Mae} & \textbf{Pai} & \textbf{Outros} & \textbf{Pop 0 4} & \textbf{Taxa por mil mae} & \textbf{Taxa por mil pai} & \textbf{Taxa por mil outros} \\ \midrule \endfirsthead
\multicolumn{9}{@{}l}{\footnotesize\itshape (continuação)}\\ \toprule \textbf{Cod ap sms} & \textbf{Ano} & \textbf{Mae} & \textbf{Pai} & \textbf{Outros} & \textbf{Pop 0 4} & \textbf{Taxa por mil mae} & \textbf{Taxa por mil pai} & \textbf{Taxa por mil outros} \\ \midrule \endhead
\midrule \multicolumn{9}{r@{}}{\footnotesize\itshape (continua)}\\ \endfoot
\bottomrule \endlastfoot
1,0 & 2011 & 10 & 3 & 1 & 13.699 & 0,7 & 0,2 & 0,1 \\
2,1 & 2011 & 1 & 3 & 0 & 21.069 & 0,0 & 0,1 & 0,0 \\
2,2 & 2011 & 3 & 1 & 3 & 11.212 & 0,3 & 0,1 & 0,3 \\
3,1 & 2011 & 102 & 42 & 2 & 38.913 & 2,6 & 1,1 & 0,1 \\
3,2 & 2011 & 20 & 14 & 3 & 22.596 & 0,9 & 0,6 & 0,1 \\
3,3 & 2011 & 42 & 19 & 6 & 42.238 & 1,0 & 0,4 & 0,1 \\
4,0 & 2011 & 42 & 15 & 2 & 56.651 & 0,7 & 0,3 & 0,0 \\
5,1 & 2011 & 5 & 5 & 0 & 33.685 & 0,1 & 0,1 & 0,0 \\
5,2 & 2011 & 78 & 28 & 3 & 43.986 & 1,8 & 0,6 & 0,1 \\
5,3 & 2011 & 21 & 5 & 2 & 26.108 & 0,8 & 0,2 & 0,1 \\
1,0 & 2012 & 12 & 8 & 4 & 13.699 & 0,9 & 0,6 & 0,3 \\
2,1 & 2012 & 4 & 2 & 1 & 21.069 & 0,2 & 0,1 & 0,0 \\
2,2 & 2012 & 8 & 2 & 0 & 11.212 & 0,7 & 0,2 & 0,0 \\
3,1 & 2012 & 116 & 66 & 6 & 38.913 & 3,0 & 1,7 & 0,2 \\
3,2 & 2012 & 43 & 15 & 1 & 22.596 & 1,9 & 0,7 & 0,0 \\
3,3 & 2012 & 49 & 25 & 4 & 42.238 & 1,2 & 0,6 & 0,1 \\
4,0 & 2012 & 55 & 29 & 3 & 56.651 & 1,0 & 0,5 & 0,1 \\
5,1 & 2012 & 9 & 9 & 3 & 33.685 & 0,3 & 0,3 & 0,1 \\
5,2 & 2012 & 124 & 26 & 3 & 43.986 & 2,8 & 0,6 & 0,1 \\
5,3 & 2012 & 44 & 6 & 1 & 26.108 & 1,7 & 0,2 & 0,0 \\
1,0 & 2013 & 19 & 10 & 3 & 13.699 & 1,4 & 0,7 & 0,2 \\
2,1 & 2013 & 8 & 5 & 1 & 21.069 & 0,4 & 0,2 & 0,0 \\
2,2 & 2013 & 13 & 7 & 5 & 11.212 & 1,2 & 0,6 & 0,4 \\
3,1 & 2013 & 88 & 55 & 14 & 38.913 & 2,3 & 1,4 & 0,4 \\
3,2 & 2013 & 41 & 21 & 5 & 22.596 & 1,8 & 0,9 & 0,2 \\
3,3 & 2013 & 44 & 19 & 2 & 42.238 & 1,0 & 0,4 & 0,0 \\
4,0 & 2013 & 38 & 28 & 3 & 56.651 & 0,7 & 0,5 & 0,1 \\
5,1 & 2013 & 31 & 21 & 3 & 33.685 & 0,9 & 0,6 & 0,1 \\
5,2 & 2013 & 125 & 47 & 10 & 43.986 & 2,8 & 1,1 & 0,2 \\
5,3 & 2013 & 22 & 9 & 3 & 26.108 & 0,8 & 0,3 & 0,1 \\
1,0 & 2014 & 33 & 22 & 2 & 13.699 & 2,4 & 1,6 & 0,1 \\
2,1 & 2014 & 6 & 2 & 0 & 21.069 & 0,3 & 0,1 & 0,0 \\
2,2 & 2014 & 22 & 9 & 2 & 11.212 & 2,0 & 0,8 & 0,2 \\
3,1 & 2014 & 80 & 43 & 2 & 38.913 & 2,1 & 1,1 & 0,1 \\
3,2 & 2014 & 52 & 21 & 3 & 22.596 & 2,3 & 0,9 & 0,1 \\
3,3 & 2014 & 58 & 31 & 5 & 42.238 & 1,4 & 0,7 & 0,1 \\
4,0 & 2014 & 59 & 44 & 3 & 56.651 & 1,0 & 0,8 & 0,1 \\
5,1 & 2014 & 36 & 20 & 4 & 33.685 & 1,1 & 0,6 & 0,1 \\
5,2 & 2014 & 170 & 44 & 10 & 43.986 & 3,9 & 1,0 & 0,2 \\
5,3 & 2014 & 48 & 10 & 1 & 26.108 & 1,8 & 0,4 & 0,0 \\
1,0 & 2015 & 31 & 18 & 3 & 13.699 & 2,3 & 1,3 & 0,2 \\
2,1 & 2015 & 8 & 5 & 1 & 21.069 & 0,4 & 0,2 & 0,0 \\
2,2 & 2015 & 18 & 10 & 0 & 11.212 & 1,6 & 0,9 & 0,0 \\
3,1 & 2015 & 56 & 32 & 2 & 38.913 & 1,4 & 0,8 & 0,1 \\
3,2 & 2015 & 44 & 18 & 2 & 22.596 & 1,9 & 0,8 & 0,1 \\
3,3 & 2015 & 54 & 25 & 8 & 42.238 & 1,3 & 0,6 & 0,2 \\
4,0 & 2015 & 52 & 33 & 6 & 56.651 & 0,9 & 0,6 & 0,1 \\
5,1 & 2015 & 20 & 14 & 1 & 33.685 & 0,6 & 0,4 & 0,0 \\
5,2 & 2015 & 133 & 65 & 11 & 43.986 & 3,0 & 1,5 & 0,3 \\
5,3 & 2015 & 36 & 14 & 3 & 26.108 & 1,4 & 0,5 & 0,1 \\
1,0 & 2016 & 40 & 23 & 4 & 13.699 & 2,9 & 1,7 & 0,3 \\
2,1 & 2016 & 7 & 6 & 0 & 21.069 & 0,3 & 0,3 & 0,0 \\
2,2 & 2016 & 14 & 13 & 2 & 11.212 & 1,2 & 1,2 & 0,2 \\
3,1 & 2016 & 52 & 44 & 2 & 38.913 & 1,3 & 1,1 & 0,1 \\
3,2 & 2016 & 66 & 33 & 10 & 22.596 & 2,9 & 1,5 & 0,4 \\
3,3 & 2016 & 71 & 48 & 6 & 42.238 & 1,7 & 1,1 & 0,1 \\
4,0 & 2016 & 67 & 52 & 6 & 56.651 & 1,2 & 0,9 & 0,1 \\
5,1 & 2016 & 44 & 23 & 1 & 33.685 & 1,3 & 0,7 & 0,0 \\
5,2 & 2016 & 204 & 104 & 15 & 43.986 & 4,6 & 2,4 & 0,3 \\
5,3 & 2016 & 35 & 25 & 2 & 26.108 & 1,3 & 1,0 & 0,1 \\
1,0 & 2017 & 77 & 43 & 3 & 13.699 & 5,6 & 3,1 & 0,2 \\
2,1 & 2017 & 13 & 5 & 4 & 21.069 & 0,6 & 0,2 & 0,2 \\
2,2 & 2017 & 19 & 14 & 3 & 11.212 & 1,7 & 1,2 & 0,3 \\
3,1 & 2017 & 104 & 66 & 11 & 38.913 & 2,7 & 1,7 & 0,3 \\
3,2 & 2017 & 43 & 37 & 11 & 22.596 & 1,9 & 1,6 & 0,5 \\
3,3 & 2017 & 110 & 88 & 8 & 42.238 & 2,6 & 2,1 & 0,2 \\
4,0 & 2017 & 104 & 68 & 10 & 56.651 & 1,8 & 1,2 & 0,2 \\
5,1 & 2017 & 191 & 177 & 6 & 33.685 & 5,7 & 5,3 & 0,2 \\
5,2 & 2017 & 717 & 633 & 12 & 43.986 & 16,3 & 14,4 & 0,3 \\
5,3 & 2017 & 136 & 130 & 4 & 26.108 & 5,2 & 5,0 & 0,2 \\
1,0 & 2018 & 71 & 47 & 11 & 13.699 & 5,2 & 3,4 & 0,8 \\
2,1 & 2018 & 15 & 6 & 2 & 21.069 & 0,7 & 0,3 & 0,1 \\
2,2 & 2018 & 23 & 16 & 1 & 11.212 & 2,1 & 1,4 & 0,1 \\
3,1 & 2018 & 133 & 96 & 10 & 38.913 & 3,4 & 2,5 & 0,3 \\
3,2 & 2018 & 58 & 37 & 8 & 22.596 & 2,6 & 1,6 & 0,4 \\
3,3 & 2018 & 105 & 80 & 12 & 42.238 & 2,5 & 1,9 & 0,3 \\
4,0 & 2018 & 157 & 122 & 14 & 56.651 & 2,8 & 2,2 & 0,2 \\
5,1 & 2018 & 162 & 138 & 7 & 33.685 & 4,8 & 4,1 & 0,2 \\
5,2 & 2018 & 465 & 358 & 41 & 43.986 & 10,6 & 8,1 & 0,9 \\
5,3 & 2018 & 76 & 54 & 8 & 26.108 & 2,9 & 2,1 & 0,3 \\
1,0 & 2019 & 45 & 29 & 1 & 13.699 & 3,3 & 2,1 & 0,1 \\
2,1 & 2019 & 6 & 12 & 2 & 21.069 & 0,3 & 0,6 & 0,1 \\
2,2 & 2019 & 11 & 7 & 2 & 11.212 & 1,0 & 0,6 & 0,2 \\
3,1 & 2019 & 113 & 69 & 5 & 38.913 & 2,9 & 1,8 & 0,1 \\
3,2 & 2019 & 50 & 37 & 7 & 22.596 & 2,2 & 1,6 & 0,3 \\
3,3 & 2019 & 111 & 82 & 12 & 42.238 & 2,6 & 1,9 & 0,3 \\
4,0 & 2019 & 61 & 50 & 10 & 56.651 & 1,1 & 0,9 & 0,2 \\
5,1 & 2019 & 72 & 62 & 3 & 33.685 & 2,1 & 1,8 & 0,1 \\
5,2 & 2019 & 323 & 153 & 32 & 43.986 & 7,3 & 3,5 & 0,7 \\
5,3 & 2019 & 76 & 21 & 1 & 26.108 & 2,9 & 0,8 & 0,0 \\
1,0 & 2020 & 28 & 20 & 2 & 13.699 & 2,0 & 1,5 & 0,1 \\
2,1 & 2020 & 11 & 14 & 3 & 21.069 & 0,5 & 0,7 & 0,1 \\
2,2 & 2020 & 10 & 6 & 1 & 11.212 & 0,9 & 0,5 & 0,1 \\
3,1 & 2020 & 135 & 90 & 5 & 38.913 & 3,5 & 2,3 & 0,1 \\
3,2 & 2020 & 35 & 27 & 7 & 22.596 & 1,5 & 1,2 & 0,3 \\
3,3 & 2020 & 118 & 81 & 10 & 42.238 & 2,8 & 1,9 & 0,2 \\
4,0 & 2020 & 41 & 37 & 4 & 56.651 & 0,7 & 0,7 & 0,1 \\
5,1 & 2020 & 56 & 35 & 2 & 33.685 & 1,7 & 1,0 & 0,1 \\
5,2 & 2020 & 250 & 127 & 15 & 43.986 & 5,7 & 2,9 & 0,3 \\
5,3 & 2020 & 50 & 27 & 3 & 26.108 & 1,9 & 1,0 & 0,1 \\
1,0 & 2021 & 32 & 16 & 7 & 13.699 & 2,3 & 1,2 & 0,5 \\
2,1 & 2021 & 4 & 2 & 2 & 21.069 & 0,2 & 0,1 & 0,1 \\
2,2 & 2021 & 2 & 1 & 2 & 11.212 & 0,2 & 0,1 & 0,2 \\
3,1 & 2021 & 152 & 101 & 5 & 38.913 & 3,9 & 2,6 & 0,1 \\
3,2 & 2021 & 35 & 17 & 3 & 22.596 & 1,5 & 0,8 & 0,1 \\
3,3 & 2021 & 129 & 76 & 13 & 42.238 & 3,1 & 1,8 & 0,3 \\
4,0 & 2021 & 65 & 46 & 11 & 56.651 & 1,1 & 0,8 & 0,2 \\
5,1 & 2021 & 89 & 66 & 6 & 33.685 & 2,6 & 2,0 & 0,2 \\
5,2 & 2021 & 261 & 220 & 34 & 43.986 & 5,9 & 5,0 & 0,8 \\
5,3 & 2021 & 36 & 36 & 3 & 26.108 & 1,4 & 1,4 & 0,1 \\
1,0 & 2022 & 28 & 21 & 2 & 13.699 & 2,0 & 1,5 & 0,1 \\
2,1 & 2022 & 6 & 10 & 2 & 21.069 & 0,3 & 0,5 & 0,1 \\
2,2 & 2022 & 13 & 10 & 2 & 11.212 & 1,2 & 0,9 & 0,2 \\
3,1 & 2022 & 148 & 79 & 9 & 38.913 & 3,8 & 2,0 & 0,2 \\
3,2 & 2022 & 51 & 33 & 6 & 22.596 & 2,3 & 1,5 & 0,3 \\
3,3 & 2022 & 186 & 112 & 14 & 42.238 & 4,4 & 2,7 & 0,3 \\
4,0 & 2022 & 94 & 65 & 9 & 56.651 & 1,7 & 1,1 & 0,2 \\
5,1 & 2022 & 269 & 223 & 4 & 33.685 & 8,0 & 6,6 & 0,1 \\
5,2 & 2022 & 517 & 417 & 28 & 43.986 & 11,8 & 9,5 & 0,6 \\
5,3 & 2022 & 47 & 33 & 4 & 26.108 & 1,8 & 1,3 & 0,2 \\
1,0 & 2023 & 57 & 40 & 7 & 13.699 & 4,2 & 2,9 & 0,5 \\
2,1 & 2023 & 14 & 11 & 3 & 21.069 & 0,7 & 0,5 & 0,1 \\
2,2 & 2023 & 12 & 12 & 1 & 11.212 & 1,1 & 1,1 & 0,1 \\
3,1 & 2023 & 200 & 118 & 7 & 38.913 & 5,1 & 3,0 & 0,2 \\
3,2 & 2023 & 72 & 54 & 4 & 22.596 & 3,2 & 2,4 & 0,2 \\
3,3 & 2023 & 245 & 170 & 11 & 42.238 & 5,8 & 4,0 & 0,3 \\
4,0 & 2023 & 136 & 98 & 14 & 56.651 & 2,4 & 1,7 & 0,2 \\
5,1 & 2023 & 192 & 150 & 1 & 33.685 & 5,7 & 4,5 & 0,0 \\
5,2 & 2023 & 471 & 356 & 27 & 43.986 & 10,7 & 8,1 & 0,6 \\
5,3 & 2023 & 70 & 59 & 7 & 26.108 & 2,7 & 2,3 & 0,3 \\
1,0 & 2024 & 61 & 31 & 1 & 13.699 & 4,5 & 2,3 & 0,1 \\
2,1 & 2024 & 19 & 11 & 1 & 21.069 & 0,9 & 0,5 & 0,0 \\
2,2 & 2024 & 21 & 13 & 2 & 11.212 & 1,9 & 1,2 & 0,2 \\
3,1 & 2024 & 169 & 105 & 8 & 38.913 & 4,3 & 2,7 & 0,2 \\
3,2 & 2024 & 130 & 77 & 10 & 22.596 & 5,8 & 3,4 & 0,4 \\
3,3 & 2024 & 262 & 170 & 14 & 42.238 & 6,2 & 4,0 & 0,3 \\
4,0 & 2024 & 161 & 133 & 11 & 56.651 & 2,8 & 2,3 & 0,2 \\
5,1 & 2024 & 130 & 103 & 4 & 33.685 & 3,9 & 3,1 & 0,1 \\
5,2 & 2024 & 476 & 391 & 35 & 43.986 & 10,8 & 8,9 & 0,8 \\
5,3 & 2024 & 86 & 71 & 4 & 26.108 & 3,3 & 2,7 & 0,2 \\
1,0 & 2025 & 43 & 29 & 4 & 13.699 & 3,1 & 2,1 & 0,3 \\
2,1 & 2025 & 28 & 28 & 2 & 21.069 & 1,3 & 1,3 & 0,1 \\
2,2 & 2025 & 19 & 15 & 4 & 11.212 & 1,7 & 1,3 & 0,4 \\
3,1 & 2025 & 238 & 146 & 14 & 38.913 & 6,1 & 3,8 & 0,4 \\
3,2 & 2025 & 142 & 105 & 10 & 22.596 & 6,3 & 4,6 & 0,4 \\
3,3 & 2025 & 305 & 217 & 16 & 42.238 & 7,2 & 5,1 & 0,4 \\
4,0 & 2025 & 166 & 126 & 12 & 56.651 & 2,9 & 2,2 & 0,2 \\
5,1 & 2025 & 221 & 212 & 3 & 33.685 & 6,6 & 6,3 & 0,1 \\
5,2 & 2025 & 531 & 472 & 36 & 43.986 & 12,1 & 10,7 & 0,8 \\
5,3 & 2025 & 63 & 54 & 9 & 26.108 & 2,4 & 2,1 & 0,3 \\
\end{longtable}}
\vspace{-6pt}\fonte{Elaboração IPP com dados de \citeonline{ipp_datario_censo}; \citeonline{sms_rio_sinan}.}
\end{landscape}


{\small\setlength{\tabcolsep}{5pt}
\begin{longtable}{@{}rlr@{}}
\caption{Notificações de violência interpessoal/autoprovocada (menores de 1 ano, 1 a 5 anos) (notif autoprovocada 2026)}\label{tab:tabela-mapa-notif-autoprovocada-2026}\\
\toprule \textbf{Codbairro} & \textbf{Bairro} & \textbf{Casos} \\ \midrule \endfirsthead
\multicolumn{3}{@{}l}{\footnotesize\itshape (continuação)}\\ \toprule \textbf{Codbairro} & \textbf{Bairro} & \textbf{Casos} \\ \midrule \endhead
\midrule \multicolumn{3}{r@{}}{\footnotesize\itshape (continua)}\\ \endfoot
\bottomrule \endlastfoot
1 & Saúde & 0 \\
2 & Gamboa & 0 \\
3 & Santo Cristo & 1 \\
4 & Caju & 0 \\
5 & Centro & 0 \\
6 & Catumbi & 0 \\
7 & Rio Comprido & 0 \\
8 & Cidade Nova & 0 \\
9 & Estácio & 0 \\
10 & Imperial de São Cristóvão & 0 \\
11 & Mangueira & 1 \\
12 & Benfica & 0 \\
13 & Paquetá & 0 \\
14 & Santa Teresa & 0 \\
15 & Flamengo & 0 \\
16 & Glória & 0 \\
17 & Laranjeiras & 0 \\
18 & Catete & 0 \\
19 & Cosme Velho & 0 \\
20 & Botafogo & 0 \\
21 & Humaitá & 0 \\
22 & Urca & 0 \\
23 & Leme & 0 \\
24 & Copacabana & 0 \\
25 & Ipanema & 0 \\
26 & Leblon & 0 \\
27 & Lagoa & 0 \\
28 & Jardim Botânico & 0 \\
29 & Gávea & 0 \\
30 & Vidigal & 0 \\
31 & São Conrado & 0 \\
32 & Praça da Bandeira & 0 \\
33 & Tijuca & 0 \\
34 & Alto da Boa Vista & 0 \\
35 & Maracanã & 0 \\
36 & Vila Isabel & 0 \\
37 & Andaraí & 0 \\
38 & Grajaú & 0 \\
39 & Manguinhos & 0 \\
40 & Bonsucesso & 0 \\
41 & Ramos & 0 \\
42 & Olaria & 2 \\
43 & Penha & 2 \\
44 & Penha Circular & 0 \\
45 & Brás de Pina & 0 \\
46 & Cordovil & 0 \\
47 & Parada de Lucas & 0 \\
48 & Vigário Geral & 1 \\
49 & Jardim América & 0 \\
50 & Higienópolis & 0 \\
51 & Jacaré & 0 \\
52 & Maria da Graça & 0 \\
53 & Del Castilho & 1 \\
54 & Inhaúma & 0 \\
55 & Engenho da Rainha & 0 \\
56 & Tomás Coelho & 1 \\
57 & São Francisco Xavier & 0 \\
58 & Rocha & 0 \\
59 & Riachuelo & 0 \\
60 & Sampaio & 0 \\
61 & Engenho Novo & 0 \\
62 & Lins de Vasconcelos & 0 \\
63 & Méier & 1 \\
64 & Todos os Santos & 0 \\
65 & Cachambi & 0 \\
66 & Engenho de Dentro & 0 \\
67 & Água Santa & 0 \\
68 & Encantado & 0 \\
69 & Piedade & 0 \\
70 & Abolição & 0 \\
71 & Pilares & 0 \\
72 & Vila Kosmos & 0 \\
73 & Vicente de Carvalho & 0 \\
74 & Vila da Penha & 0 \\
75 & Vista Alegre & 3 \\
76 & Irajá & 0 \\
77 & Colégio & 1 \\
78 & Campinho & 0 \\
79 & Quintino Bocaiúva & 0 \\
80 & Cavalcanti & 0 \\
81 & Engenheiro Leal & 0 \\
82 & Cascadura & 0 \\
83 & Madureira & 0 \\
84 & Vaz Lobo & 0 \\
85 & Turiaçú & 0 \\
86 & Rocha Miranda & 0 \\
87 & Honório Gurgel & 0 \\
88 & Osvaldo Cruz & 0 \\
89 & Bento Ribeiro & 0 \\
90 & Marechal Hermes & 0 \\
91 & Ribeira & 0 \\
92 & Zumbi & 1 \\
93 & Cacuia & 0 \\
94 & Pitangueiras & 0 \\
95 & Praia da Bandeira & 1 \\
96 & Cocotá & 0 \\
97 & Bancários & 0 \\
98 & Freguesia (Ilha) & 0 \\
99 & Jardim Guanabara & 0 \\
100 & Jardim Carioca & 0 \\
101 & Tauá & 0 \\
102 & Moneró & 1 \\
103 & Portuguesa & 1 \\
104 & Galeão & 0 \\
105 & Cidade Universitária & 0 \\
106 & Guadalupe & 0 \\
107 & Anchieta & 0 \\
108 & Parque Anchieta & 0 \\
109 & Ricardo de Albuquerque & 0 \\
110 & Coelho Neto & 0 \\
111 & Acari & 0 \\
112 & Barros Filho & 0 \\
113 & Costa Barros & 0 \\
114 & Pavuna & 0 \\
115 & Jacarepaguá & 0 \\
116 & Anil & 0 \\
117 & Gardênia Azul & 0 \\
118 & Cidade de Deus & 0 \\
119 & Curicica & 0 \\
120 & Freguesia (Jacarepaguá) & 0 \\
121 & Pechincha & 0 \\
122 & Taquara & 0 \\
123 & Tanque & 0 \\
124 & Praça Seca & 0 \\
125 & Vila Valqueire & 0 \\
126 & Joá & 0 \\
127 & Itanhangá & 0 \\
128 & Barra da Tijuca & 0 \\
129 & Camorim & 0 \\
130 & Vargem Pequena & 0 \\
131 & Vargem Grande & 0 \\
132 & Recreio dos Bandeirantes & 0 \\
133 & Grumari & 0 \\
134 & Deodoro & 0 \\
135 & Vila Militar & 0 \\
136 & Campo dos Afonsos & 1 \\
137 & Jardim Sulacap & 0 \\
138 & Magalhães Bastos & 0 \\
139 & Realengo & 0 \\
140 & Padre Miguel & 0 \\
141 & Bangu & 0 \\
142 & Senador Camará & 0 \\
143 & Santíssimo & 3 \\
144 & Campo Grande & 1 \\
145 & Senador Vasconcelos & 0 \\
146 & Inhoaíba & 0 \\
147 & Cosmos & 0 \\
148 & Paciência & 0 \\
149 & Santa Cruz & 0 \\
150 & Sepetiba & 0 \\
151 & Guaratiba & 1 \\
152 & Barra de Guaratiba & 0 \\
153 & Pedra de Guaratiba & 2 \\
154 & Rocinha & 0 \\
155 & Jacarezinho & 0 \\
156 & Complexo do Alemão & 2 \\
157 & Maré & 0 \\
158 & Vasco da Gama & 0 \\
159 & Parque Colúmbia & 5 \\
160 & Gericinó & 0 \\
161 & Lapa & 0 \\
162 & Vila Kennedy & 0 \\
163 & Jabour & 0 \\
164 & Ilha de Guaratiba & 0 \\
165 & Barra Olímpica & 0 \\
166 & Argentino & 0 \\
\end{longtable}}
\vspace{-6pt}\fonte{Elaboração IPP com dados de \citeonline{sms_rio_sinan}.}


\begin{landscape}
{\scriptsize\setlength{\tabcolsep}{3pt}
\begin{longtable}{@{}rrrrrrrr@{}}
\caption{Taxa de notificações de violência (0 a 6 anos) (violencia familiar taxa municipio ano)}\label{tab:violencia-familiar-taxa-municipio-ano}\\
\toprule \textbf{Ano} & \textbf{Mae} & \textbf{Pai} & \textbf{Outros} & \textbf{Populacao 0 a 5} & \textbf{Taxa por mil mae} & \textbf{Taxa por mil pai} & \textbf{Taxa por mil outros} \\ \midrule \endfirsthead
\multicolumn{8}{@{}l}{\footnotesize\itshape (continuação)}\\ \toprule \textbf{Ano} & \textbf{Mae} & \textbf{Pai} & \textbf{Outros} & \textbf{Populacao 0 a 5} & \textbf{Taxa por mil mae} & \textbf{Taxa por mil pai} & \textbf{Taxa por mil outros} \\ \midrule \endhead
\midrule \multicolumn{8}{r@{}}{\footnotesize\itshape (continua)}\\ \endfoot
\bottomrule \endlastfoot
2011 & 324 & 135 & 22 & 480.648 & 0,7 & 0,3 & 0,0 \\
2012 & 464 & 188 & 26 & 480.542 & 1,0 & 0,4 & 0,1 \\
2013 & 429 & 222 & 49 & 481.834 & 0,9 & 0,5 & 0,1 \\
2014 & 564 & 246 & 32 & 485.451 & 1,2 & 0,5 & 0,1 \\
2015 & 452 & 234 & 37 & 491.681 & 0,9 & 0,5 & 0,1 \\
2016 & 600 & 371 & 48 & 495.313 & 1,2 & 0,7 & 0,1 \\
2017 & 1.514 & 1.261 & 72 & 494.645 & 3,1 & 2,5 & 0,1 \\
2018 & 1.265 & 954 & 114 & 491.868 & 2,6 & 1,9 & 0,2 \\
2019 & 868 & 522 & 75 & 485.805 & 1,8 & 1,1 & 0,2 \\
2020 & 734 & 464 & 52 & 474.322 & 1,5 & 1,0 & 0,1 \\
2021 & 805 & 581 & 86 & 457.280 & 1,8 & 1,3 & 0,2 \\
2022 & 1.359 & 1.003 & 80 & 439.907 & 3,1 & 2,3 & 0,2 \\
2023 & 1.469 & 1.068 & 82 & 423.739 & 3,5 & 2,5 & 0,2 \\
2024 & 1.515 & 1.105 & 90 & 407.537 & 3,7 & 2,7 & 0,2 \\
2025 & 1.756 & 1.404 & 110 & 393.073 & 4,5 & 3,6 & 0,3 \\
\end{longtable}}
\vspace{-6pt}\fonte{Elaboração IPP com dados de \citeonline{ipp_datario_censo}; \citeonline{ms_ripsa_populacao}; \citeonline{sms_rio_sinan}.}
\end{landscape}


\begin{landscape}
{\scriptsize\setlength{\tabcolsep}{3pt}
\begin{longtable}{@{}rlrrrrrrr@{}}
\caption{Taxa de notificações de violência (0 a 6 anos) (violencia familiar taxa por bairro)}\label{tab:violencia-familiar-taxa-por-bairro}\\
\toprule \textbf{Codbairro} & \textbf{Bairro} & \textbf{Casos mae 2025} & \textbf{Casos pai 2025} & \textbf{Casos outros 2021 2025} & \textbf{Pop 0 4} & \textbf{Taxa por mil mae 2025} & \textbf{Taxa por mil pai 2025} & \textbf{Taxa por mil outros 2021 2025} \\ \midrule \endfirsthead
\multicolumn{9}{@{}l}{\footnotesize\itshape (continuação)}\\ \toprule \textbf{Codbairro} & \textbf{Bairro} & \textbf{Casos mae 2025} & \textbf{Casos pai 2025} & \textbf{Casos outros 2021 2025} & \textbf{Pop 0 4} & \textbf{Taxa por mil mae 2025} & \textbf{Taxa por mil pai 2025} & \textbf{Taxa por mil outros 2021 2025} \\ \midrule \endhead
\midrule \multicolumn{9}{r@{}}{\footnotesize\itshape (continua)}\\ \endfoot
\bottomrule \endlastfoot
1 & Saúde & 1 & 0 & 0 & 54 & 18,5 & 0,0 & 0,0 \\
2 & Gamboa & 2 & 0 & 0 & 669 & 3,0 & 0,0 & 0,0 \\
3 & Santo Cristo & 3 & 0 & 0 & 534 & 5,6 & 0,0 & 0,0 \\
4 & Caju & 1 & 3 & 2 & 1.529 & 0,7 & 2,0 & 1,3 \\
5 & Centro & 3 & 2 & 3 & 622 & 4,8 & 3,2 & 4,8 \\
6 & Catumbi & 1 & 0 & 3 & 662 & 1,5 & 0,0 & 4,5 \\
7 & Rio Comprido & 5 & 3 & 0 & 1.908 & 2,6 & 1,6 & 0,0 \\
8 & Cidade Nova & 0 & 0 & 0 & 211 & 0,0 & 0,0 & 0,0 \\
9 & Estácio & 2 & 3 & 1 & 1.280 & 1,6 & 2,3 & 0,8 \\
10 & Imperial de São Cristóvão & 4 & 1 & 2 & 1.257 & 3,2 & 0,8 & 1,6 \\
11 & Mangueira & 4 & 2 & 1 & 1.228 & 3,3 & 1,6 & 0,8 \\
12 & Benfica & 12 & 10 & 4 & 1.055 & 11,4 & 9,5 & 3,8 \\
13 & Paquetá & 0 & 0 & 0 & 145 & 0,0 & 0,0 & 0,0 \\
14 & Santa Teresa & 3 & 4 & 1 & 1.594 & 1,9 & 2,5 & 0,6 \\
15 & Flamengo & 2 & 2 & 0 & 1.210 & 1,7 & 1,7 & 0,0 \\
16 & Glória & 0 & 0 & 0 & 169 & 0,0 & 0,0 & 0,0 \\
17 & Laranjeiras & 0 & 0 & 0 & 1.389 & 0,0 & 0,0 & 0,0 \\
18 & Catete & 2 & 1 & 0 & 830 & 2,4 & 1,2 & 0,0 \\
19 & Cosme Velho & 0 & 0 & 0 & 325 & 0,0 & 0,0 & 0,0 \\
20 & Botafogo & 5 & 4 & 1 & 2.761 & 1,8 & 1,4 & 0,4 \\
21 & Humaitá & 0 & 0 & 0 & 454 & 0,0 & 0,0 & 0,0 \\
22 & Urca & 0 & 1 & 0 & 252 & 0,0 & 4,0 & 0,0 \\
23 & Leme & 4 & 3 & 1 & 422 & 9,5 & 7,1 & 2,4 \\
24 & Copacabana & 3 & 3 & 3 & 3.528 & 0,9 & 0,9 & 0,9 \\
25 & Ipanema & 1 & 1 & 0 & 1.314 & 0,8 & 0,8 & 0,0 \\
26 & Leblon & 2 & 1 & 0 & 1.199 & 1,7 & 0,8 & 0,0 \\
27 & Lagoa & 0 & 0 & 0 & 772 & 0,0 & 0,0 & 0,0 \\
28 & Jardim Botânico & 1 & 1 & 0 & 587 & 1,7 & 1,7 & 0,0 \\
29 & Gávea & 1 & 2 & 1 & 459 & 2,2 & 4,4 & 2,2 \\
30 & Vidigal & 3 & 5 & 2 & 853 & 3,5 & 5,9 & 2,3 \\
31 & São Conrado & 0 & 0 & 0 & 337 & 0,0 & 0,0 & 0,0 \\
32 & Praça da Bandeira & 0 & 0 & 0 & 261 & 0,0 & 0,0 & 0,0 \\
33 & Tijuca & 16 & 9 & 7 & 4.764 & 3,4 & 1,9 & 1,5 \\
34 & Alto da Boa Vista & 0 & 0 & 0 & 418 & 0,0 & 0,0 & 0,0 \\
35 & Maracanã & 0 & 0 & 0 & 837 & 0,0 & 0,0 & 0,0 \\
36 & Vila Isabel & 0 & 0 & 0 & 2.358 & 0,0 & 0,0 & 0,0 \\
37 & Andaraí & 0 & 0 & 0 & 1.336 & 0,0 & 0,0 & 0,0 \\
38 & Grajaú & 3 & 6 & 4 & 1.238 & 2,4 & 4,8 & 3,2 \\
39 & Manguinhos & 8 & 3 & 1 & 1.915 & 4,2 & 1,6 & 0,5 \\
40 & Bonsucesso & 5 & 4 & 2 & 672 & 7,4 & 6,0 & 3,0 \\
41 & Ramos & 14 & 10 & 2 & 1.527 & 9,2 & 6,5 & 1,3 \\
42 & Olaria & 17 & 10 & 8 & 2.133 & 8,0 & 4,7 & 3,8 \\
43 & Penha & 23 & 11 & 3 & 2.958 & 7,8 & 3,7 & 1,0 \\
44 & Penha Circular & 6 & 1 & 3 & 1.639 & 3,7 & 0,6 & 1,8 \\
45 & Brás de Pina & 34 & 19 & 5 & 1.760 & 19,3 & 10,8 & 2,8 \\
46 & Cordovil & 10 & 6 & 3 & 1.641 & 6,1 & 3,7 & 1,8 \\
47 & Parada de Lucas & 14 & 10 & 0 & 1.171 & 12,0 & 8,5 & 0,0 \\
48 & Vigário Geral & 14 & 10 & 0 & 2.050 & 6,8 & 4,9 & 0,0 \\
49 & Jardim América & 7 & 6 & 1 & 908 & 7,7 & 6,6 & 1,1 \\
50 & Higienópolis & 15 & 8 & 3 & 503 & 29,8 & 15,9 & 6,0 \\
51 & Jacaré & 4 & 3 & 2 & 443 & 9,0 & 6,8 & 4,5 \\
52 & Maria da Graça & 0 & 0 & 1 & 277 & 0,0 & 0,0 & 3,6 \\
53 & Del Castilho & 14 & 8 & 3 & 991 & 14,1 & 8,1 & 3,0 \\
54 & Inhaúma & 4 & 3 & 0 & 1.718 & 2,3 & 1,7 & 0,0 \\
55 & Engenho da Rainha & 6 & 3 & 0 & 1.130 & 5,3 & 2,7 & 0,0 \\
56 & Tomás Coelho & 12 & 9 & 2 & 1.135 & 10,6 & 7,9 & 1,8 \\
57 & São Francisco Xavier & 5 & 3 & 3 & 255 & 19,6 & 11,8 & 11,8 \\
58 & Rocha & 4 & 4 & 1 & 885 & 4,5 & 4,5 & 1,1 \\
59 & Riachuelo & 0 & 0 & 0 & 327 & 0,0 & 0,0 & 0,0 \\
60 & Sampaio & 2 & 2 & 1 & 453 & 4,4 & 4,4 & 2,2 \\
61 & Engenho Novo & 3 & 1 & 0 & 1.905 & 1,6 & 0,5 & 0,0 \\
62 & Lins de Vasconcelos & 3 & 2 & 4 & 1.502 & 2,0 & 1,3 & 2,7 \\
63 & Méier & 21 & 18 & 2 & 1.080 & 19,4 & 16,7 & 1,9 \\
64 & Todos os Santos & 6 & 4 & 1 & 746 & 8,0 & 5,4 & 1,3 \\
65 & Cachambi & 13 & 9 & 0 & 1.412 & 9,2 & 6,4 & 0,0 \\
66 & Engenho de Dentro & 2 & 2 & 0 & 1.772 & 1,1 & 1,1 & 0,0 \\
67 & Água Santa & 6 & 6 & 0 & 280 & 21,4 & 21,4 & 0,0 \\
68 & Encantado & 13 & 12 & 2 & 442 & 29,4 & 27,1 & 4,5 \\
69 & Piedade & 1 & 0 & 1 & 1.616 & 0,6 & 0,0 & 0,6 \\
70 & Abolição & 2 & 2 & 0 & 313 & 6,4 & 6,4 & 0,0 \\
71 & Pilares & 4 & 3 & 2 & 1.036 & 3,9 & 2,9 & 1,9 \\
72 & Vila Kosmos & 1 & 1 & 0 & 544 & 1,8 & 1,8 & 0,0 \\
73 & Vicente de Carvalho & 11 & 7 & 1 & 1.042 & 10,6 & 6,7 & 1,0 \\
74 & Vila da Penha & 11 & 3 & 2 & 754 & 14,6 & 4,0 & 2,7 \\
75 & Vista Alegre & 7 & 6 & 2 & 645 & 10,9 & 9,3 & 3,1 \\
76 & Irajá & 6 & 4 & 1 & 3.596 & 1,7 & 1,1 & 0,3 \\
77 & Colégio & 2 & 1 & 0 & 1.496 & 1,3 & 0,7 & 0,0 \\
78 & Campinho & 24 & 14 & 4 & 392 & 61,2 & 35,7 & 10,2 \\
79 & Quintino Bocaiúva & 14 & 9 & 3 & 1.021 & 13,7 & 8,8 & 2,9 \\
80 & Cavalcanti & 3 & 2 & 0 & 597 & 5,0 & 3,4 & 0,0 \\
81 & Engenheiro Leal & 13 & 7 & 2 & 245 & 53,1 & 28,6 & 8,2 \\
82 & Cascadura & 0 & 0 & 0 & 1.184 & 0,0 & 0,0 & 0,0 \\
83 & Madureira & 0 & 0 & 0 & 1.766 & 0,0 & 0,0 & 0,0 \\
84 & Vaz Lobo & 10 & 8 & 1 & 582 & 17,2 & 13,7 & 1,7 \\
85 & Turiaçú & 13 & 8 & 0 & 649 & 20,0 & 12,3 & 0,0 \\
86 & Rocha Miranda & 5 & 3 & 1 & 1.724 & 2,9 & 1,7 & 0,6 \\
87 & Honório Gurgel & 5 & 1 & 0 & 933 & 5,4 & 1,1 & 0,0 \\
88 & Osvaldo Cruz & 7 & 4 & 1 & 1.289 & 5,4 & 3,1 & 0,8 \\
89 & Bento Ribeiro & 7 & 6 & 1 & 1.541 & 4,5 & 3,9 & 0,6 \\
90 & Marechal Hermes & 2 & 2 & 1 & 1.963 & 1,0 & 1,0 & 0,5 \\
91 & Ribeira & 8 & 4 & 1 & 129 & 62,0 & 31,0 & 7,8 \\
92 & Zumbi & 25 & 18 & 0 & 78 & 320,5 & 230,8 & 0,0 \\
93 & Cacuia & 0 & 0 & 1 & 439 & 0,0 & 0,0 & 2,3 \\
94 & Pitangueiras & 0 & 0 & 0 & 503 & 0,0 & 0,0 & 0,0 \\
95 & Praia da Bandeira & 3 & 2 & 2 & 205 & 14,6 & 9,8 & 9,8 \\
96 & Cocotá & 1 & 3 & 1 & 174 & 5,7 & 17,2 & 5,7 \\
97 & Bancários & 0 & 0 & 0 & 600 & 0,0 & 0,0 & 0,0 \\
98 & Freguesia (Ilha) & 0 & 1 & 0 & 828 & 0,0 & 1,2 & 0,0 \\
99 & Jardim Guanabara & 2 & 2 & 0 & 967 & 2,1 & 2,1 & 0,0 \\
100 & Jardim Carioca & 2 & 1 & 2 & 790 & 2,5 & 1,3 & 2,5 \\
101 & Tauá & 2 & 1 & 0 & 1.230 & 1,6 & 0,8 & 0,0 \\
102 & Moneró & 3 & 4 & 0 & 181 & 16,6 & 22,1 & 0,0 \\
103 & Portuguesa & 4 & 0 & 2 & 839 & 4,8 & 0,0 & 2,4 \\
104 & Galeão & 0 & 0 & 0 & 1.448 & 0,0 & 0,0 & 0,0 \\
105 & Cidade Universitária & 3 & 1 & 0 & 65 & 46,2 & 15,4 & 0,0 \\
106 & Guadalupe & 9 & 2 & 3 & 2.073 & 4,3 & 1,0 & 1,4 \\
107 & Anchieta & 0 & 0 & 2 & 2.866 & 0,0 & 0,0 & 0,7 \\
108 & Parque Anchieta & 22 & 20 & 5 & 962 & 22,9 & 20,8 & 5,2 \\
109 & Ricardo de Albuquerque & 38 & 35 & 8 & 1.345 & 28,3 & 26,0 & 5,9 \\
110 & Coelho Neto & 9 & 8 & 6 & 1.190 & 7,6 & 6,7 & 5,0 \\
111 & Acari & 17 & 16 & 4 & 2.112 & 8,0 & 7,6 & 1,9 \\
112 & Barros Filho & 19 & 17 & 1 & 1.400 & 13,6 & 12,1 & 0,7 \\
113 & Costa Barros & 27 & 15 & 3 & 1.952 & 13,8 & 7,7 & 1,5 \\
114 & Pavuna & 7 & 6 & 6 & 5.892 & 1,2 & 1,0 & 1,0 \\
115 & Jacarepaguá & 16 & 9 & 6 & 11.499 & 1,4 & 0,8 & 0,5 \\
116 & Anil & 27 & 19 & 5 & 1.504 & 18,0 & 12,6 & 3,3 \\
117 & Gardênia Azul & 33 & 25 & 6 & 1.227 & 26,9 & 20,4 & 4,9 \\
118 & Cidade de Deus & 10 & 7 & 3 & 1.968 & 5,1 & 3,6 & 1,5 \\
119 & Curicica & 12 & 9 & 4 & 1.485 & 8,1 & 6,1 & 2,7 \\
120 & Freguesia (Jacarepaguá) & 15 & 10 & 6 & 3.266 & 4,6 & 3,1 & 1,8 \\
121 & Pechincha & 5 & 6 & 7 & 1.656 & 3,0 & 3,6 & 4,2 \\
122 & Taquara & 1 & 1 & 2 & 5.078 & 0,2 & 0,2 & 0,4 \\
123 & Tanque & 1 & 3 & 1 & 1.818 & 0,6 & 1,7 & 0,6 \\
124 & Praça Seca & 10 & 10 & 8 & 3.399 & 2,9 & 2,9 & 2,4 \\
125 & Vila Valqueire & 2 & 1 & 0 & 1.382 & 1,4 & 0,7 & 0,0 \\
126 & Joá & 11 & 10 & 2 & 22 & 500,0 & 454,5 & 90,9 \\
127 & Itanhangá & 4 & 4 & 0 & 4.220 & 0,9 & 0,9 & 0,0 \\
128 & Barra da Tijuca & 0 & 0 & 0 & 5.310 & 0,0 & 0,0 & 0,0 \\
129 & Camorim & 10 & 5 & 2 & 83 & 120,5 & 60,2 & 24,1 \\
130 & Vargem Pequena & 5 & 3 & 3 & 1.827 & 2,7 & 1,6 & 1,6 \\
131 & Vargem Grande & 0 & 0 & 0 & 1.191 & 0,0 & 0,0 & 0,0 \\
132 & Recreio dos Bandeirantes & 4 & 3 & 1 & 6.450 & 0,6 & 0,5 & 0,2 \\
133 & Grumari & 0 & 1 & 1 & 15 & 0,0 & 66,7 & 66,7 \\
134 & Deodoro & 6 & 7 & 2 & 742 & 8,1 & 9,4 & 2,7 \\
135 & Vila Militar & 0 & 1 & 0 & 798 & 0,0 & 1,3 & 0,0 \\
136 & Campo dos Afonsos & 4 & 6 & 0 & 178 & 22,5 & 33,7 & 0,0 \\
137 & Jardim Sulacap & 1 & 2 & 1 & 696 & 1,4 & 2,9 & 1,4 \\
138 & Magalhães Bastos & 1 & 1 & 0 & 968 & 1,0 & 1,0 & 0,0 \\
139 & Realengo & 9 & 9 & 0 & 8.532 & 1,1 & 1,1 & 0,0 \\
140 & Padre Miguel & 13 & 12 & 1 & 3.101 & 4,2 & 3,9 & 0,3 \\
141 & Bangu & 127 & 121 & 11 & 11.418 & 11,1 & 10,6 & 1,0 \\
142 & Senador Camará & 60 & 53 & 3 & 5.702 & 10,5 & 9,3 & 0,5 \\
143 & Santíssimo & 165 & 132 & 26 & 2.482 & 66,5 & 53,2 & 10,5 \\
144 & Campo Grande & 49 & 50 & 16 & 18.505 & 2,6 & 2,7 & 0,9 \\
145 & Senador Vasconcelos & 23 & 22 & 3 & 1.638 & 14,0 & 13,4 & 1,8 \\
146 & Inhoaíba & 81 & 92 & 26 & 4.495 & 18,0 & 20,5 & 5,8 \\
147 & Cosmos & 4 & 5 & 3 & 5.662 & 0,7 & 0,9 & 0,5 \\
148 & Paciência & 13 & 11 & 4 & 6.333 & 2,1 & 1,7 & 0,6 \\
149 & Santa Cruz & 23 & 20 & 8 & 16.093 & 1,4 & 1,2 & 0,5 \\
150 & Sepetiba & 27 & 23 & 15 & 3.682 & 7,3 & 6,2 & 4,1 \\
151 & Guaratiba & 116 & 100 & 50 & 10.088 & 11,5 & 9,9 & 5,0 \\
152 & Barra de Guaratiba & 43 & 34 & 18 & 208 & 206,7 & 163,5 & 86,5 \\
153 & Pedra de Guaratiba & 50 & 37 & 18 & 599 & 83,5 & 61,8 & 30,1 \\
154 & Rocinha & 4 & 4 & 2 & 4.208 & 1,0 & 1,0 & 0,5 \\
155 & Jacarezinho & 2 & 3 & 5 & 2.375 & 0,8 & 1,3 & 2,1 \\
156 & Complexo do Alemão & 27 & 14 & 5 & 3.477 & 7,8 & 4,0 & 1,4 \\
157 & Maré & 6 & 5 & 1 & 8.553 & 0,7 & 0,6 & 0,1 \\
158 & Vasco da Gama & 2 & 1 & 4 & 745 & 2,7 & 1,3 & 5,4 \\
159 & Parque Colúmbia & 16 & 12 & 10 & 483 & 33,1 & 24,8 & 20,7 \\
160 & Gericinó & 0 & 0 & 0 & 178 & 0,0 & 0,0 & 0,0 \\
161 & Lapa & 0 & 0 & 0 & 206 & 0,0 & 0,0 & 0,0 \\
162 & Vila Kennedy & 0 & 0 & 0 & 1.199 & 0,0 & 0,0 & 0,0 \\
163 & Jabour & 0 & 0 & 0 & 173 & 0,0 & 0,0 & 0,0 \\
164 & Ilha de Guaratiba & 0 & 0 & 0 & 309 & 0,0 & 0,0 & 0,0 \\
165 & Barra Olímpica & 0 & 0 & 0 & 3.251 & 0,0 & 0,0 & 0,0 \\
166 & Argentino & 0 & 0 & 0 & 33 & 0,0 & 0,0 & 0,0 \\
\end{longtable}}
\vspace{-6pt}\fonte{Elaboração IPP com dados de \citeonline{ipp_datario_censo}; \citeonline{ms_ripsa_populacao}; \citeonline{sms_rio_sinan}.}
\end{landscape}


{\small\setlength{\tabcolsep}{5pt}
\begin{longtable}{@{}llr@{}}
\caption{Taxa de notificações de violência (0 a 6 anos) (violencia familiar taxa top bairros 2025)}\label{tab:violencia-familiar-taxa-top-bairros-2025}\\
\toprule \textbf{Bairro} & \textbf{Vinculo} & \textbf{Taxa por 1.000} \\ \midrule \endfirsthead
\multicolumn{3}{@{}l}{\footnotesize\itshape (continuação)}\\ \toprule \textbf{Bairro} & \textbf{Vinculo} & \textbf{Taxa por 1.000} \\ \midrule \endhead
\midrule \multicolumn{3}{r@{}}{\footnotesize\itshape (continua)}\\ \endfoot
\bottomrule \endlastfoot
Barra de Guaratiba & Mãe & 206,7 \\
Pedra de Guaratiba & Mãe & 83,5 \\
Santíssimo & Mãe & 66,5 \\
Ribeira & Mãe & 62,0 \\
Campinho & Mãe & 61,2 \\
Engenheiro Leal & Mãe & 53,1 \\
Parque Colúmbia & Mãe & 33,1 \\
Higienópolis & Mãe & 29,8 \\
Encantado & Mãe & 29,4 \\
Ricardo de Albuquerque & Mãe & 28,3 \\
Barra de Guaratiba & Pai & 163,5 \\
Pedra de Guaratiba & Pai & 61,8 \\
Santíssimo & Pai & 53,2 \\
Ribeira & Pai & 31,0 \\
Campinho & Pai & 35,7 \\
Engenheiro Leal & Pai & 28,6 \\
Parque Colúmbia & Pai & 24,8 \\
Higienópolis & Pai & 15,9 \\
Encantado & Pai & 27,1 \\
Ricardo de Albuquerque & Pai & 26,0 \\
\end{longtable}}
\vspace{-6pt}\fonte{Elaboração IPP com dados de \citeonline{ipp_datario_censo}; \citeonline{ms_ripsa_populacao}; \citeonline{sms_rio_sinan}.}


\section*{Tabelas disponíveis em formato digital}

As tabelas abaixo têm mais de 200 linhas e não são impressas; estão em \texttt{tabelas\_finais/} no repositório do projeto.

\begin{itemize}\item \texttt{violencia\_familiar\_por\_bairro.csv} --- Violência familiar (0 a 5 anos, por vínculo do provável autor) (2.490 linhas)\item \texttt{notif\_autoprovocada\_por\_bairro\_ano.csv} --- Notificações de violência interpessoal/autoprovocada (menores de 1 ano, 1 a 5 anos) (1.494 linhas)\end{itemize}

\chapter{Tabelas do eixo Alimentação}\label{ap:eixo-5}

{\small\setlength{\tabcolsep}{5pt}
\begin{longtable}{@{}rrr@{}}
\caption{Baixo peso ao nascer (número) (nascidos abaixo peso por ano)}\label{tab:nascidos-abaixo-peso-por-ano}\\
\toprule \textbf{Ano} & \textbf{Nascidos abaixo peso} & \textbf{Percentual abaixo do peso} \\ \midrule \endfirsthead
\multicolumn{3}{@{}l}{\footnotesize\itshape (continuação)}\\ \toprule \textbf{Ano} & \textbf{Nascidos abaixo peso} & \textbf{Percentual abaixo do peso} \\ \midrule \endhead
\midrule \multicolumn{3}{r@{}}{\footnotesize\itshape (continua)}\\ \endfoot
\bottomrule \endlastfoot
2006 & 9.306 & 10,0 \\
2007 & 9.285 & 10,1 \\
2008 & 9.280 & 10,1 \\
2009 & 9.639 & 10,2 \\
2010 & 9.192 & 10,0 \\
2011 & 7.950 & 9,3 \\
2012 & 8.094 & 9,4 \\
2013 & 8.017 & 9,2 \\
2014 & 8.253 & 9,2 \\
2015 & 8.393 & 9,3 \\
2016 & 7.684 & 9,3 \\
2017 & 7.716 & 9,1 \\
2018 & 7.646 & 9,3 \\
2019 & 7.155 & 9,4 \\
2020 & 6.842 & 9,4 \\
2021 & 6.509 & 9,5 \\
2022 & 7.588 & 10,5 \\
2023 & 7.498 & 10,6 \\
2024 & 6.533 & 10,2 \\
2025 & 6.723 & 10,3 \\
\end{longtable}}
\vspace{-6pt}\fonte{Elaboração IPP com dados de \citeonline{sms_rio_sim}; \citeonline{sms_rio_sinasc}.}


\begin{landscape}
{\scriptsize\setlength{\tabcolsep}{3pt}
\begin{longtable}{@{}rrrrrrrrrrr@{}}
\caption{Desnutrição SISVAN (número) (sisvan desnutricao por ano)}\label{tab:sisvan-desnutricao-por-ano}\\
\toprule \textbf{Ano} & \textbf{Peso muito baixo bruto} & \textbf{Peso muito baixo (\%)} & \textbf{Peso baixo bruto} & \textbf{Peso baixo (\%)} & \textbf{Peso adequado bruto} & \textbf{Peso adequado (\%)} & \textbf{Peso elevado bruto} & \textbf{Peso elevado (\%)} & \textbf{Total} & \textbf{Percent. baixo peso total} \\ \midrule \endfirsthead
\multicolumn{11}{@{}l}{\footnotesize\itshape (continuação)}\\ \toprule \textbf{Ano} & \textbf{Peso muito baixo bruto} & \textbf{Peso muito baixo (\%)} & \textbf{Peso baixo bruto} & \textbf{Peso baixo (\%)} & \textbf{Peso adequado bruto} & \textbf{Peso adequado (\%)} & \textbf{Peso elevado bruto} & \textbf{Peso elevado (\%)} & \textbf{Total} & \textbf{Percent. baixo peso total} \\ \midrule \endhead
\midrule \multicolumn{11}{r@{}}{\footnotesize\itshape (continua)}\\ \endfoot
\bottomrule \endlastfoot
2008 & 169 & 1,0 & 438 & 2,5 & 15.910 & 90,0 & 1.154 & 6,5 & 17.671 & 3,4 \\
2009 & 196 & 0,8 & 569 & 2,5 & 20.672 & 89,1 & 1.757 & 7,6 & 23.194 & 3,3 \\
2010 & 262 & 1,1 & 610 & 2,6 & 20.507 & 87,8 & 1.992 & 8,5 & 23.371 & 3,7 \\
2011 & 181 & 1,0 & 492 & 2,7 & 15.851 & 86,0 & 1.905 & 10,3 & 18.429 & 3,6 \\
2012 & 294 & 1,1 & 628 & 2,2 & 24.441 & 87,7 & 2.495 & 9,0 & 27.858 & 3,3 \\
2013 & 964 & 2,0 & 1.462 & 3,1 & 40.112 & 84,2 & 5.098 & 10,7 & 47.636 & 5,1 \\
2014 & 748 & 1,3 & 1.610 & 2,8 & 50.380 & 87,0 & 5.203 & 9,0 & 57.941 & 4,1 \\
2015 & 951 & 1,4 & 2.261 & 3,2 & 61.127 & 87,4 & 5.583 & 8,0 & 69.922 & 4,6 \\
2016 & 986 & 1,8 & 2.187 & 4,0 & 46.908 & 86,7 & 4.009 & 7,4 & 54.090 & 5,9 \\
2017 & 1.614 & 2,6 & 2.222 & 3,6 & 53.827 & 87,2 & 4.078 & 6,6 & 61.741 & 6,2 \\
2018 & 9.707 & 12,1 & 2.689 & 3,4 & 62.619 & 78,2 & 5.070 & 6,3 & 80.085 & 15,5 \\
2019 & 2.469 & 2,6 & 4.131 & 4,3 & 81.422 & 85,6 & 7.122 & 7,5 & 95.144 & 6,9 \\
2020 & 2.684 & 1,9 & 5.332 & 3,8 & 121.472 & 86,7 & 10.623 & 7,6 & 140.111 & 5,7 \\
2021 & 2.795 & 2,0 & 5.007 & 3,6 & 121.050 & 86,9 & 10.506 & 7,5 & 139.358 & 5,6 \\
2022 & 3.484 & 2,2 & 5.787 & 3,6 & 137.472 & 86,8 & 11.684 & 7,4 & 158.427 & 5,8 \\
2023 & 4.074 & 0,0 & 7.575 & 3,7 & 177.453 & 87,2 & 14.517 & 7,1 & 203.619 & 3,7 \\
2024 & 3.639 & 1,7 & 7.749 & 3,6 & 189.359 & 88,0 & 14.479 & 6,7 & 215.226 & 5,3 \\
2025 & 7.115 & 3,3 & 9.100 & 4,2 & 186.262 & 86,6 & 12.655 & 5,9 & 215.132 & 7,5 \\
2026 & 3.332 & 2,8 & 4.526 & 3,7 & 106.781 & 88,1 & 6.530 & 5,4 & 121.169 & 6,5 \\
\end{longtable}}
\vspace{-6pt}\fonte{Elaboração IPP com dados de \citeonline{ms_sisvan}.}
\end{landscape}


\begin{landscape}
{\scriptsize\setlength{\tabcolsep}{3pt}
\begin{longtable}{@{}rrrrrrrrrrrrrrr@{}}
\caption{Sobrepeso SISVAN (número) (sisvan sobrepeso por ano)}\label{tab:sisvan-sobrepeso-por-ano}\\
\toprule \textbf{Ano} & \textbf{Magreza acentuada bruto} & \textbf{Magreza acentuada (\%)} & \textbf{Magreza bruto} & \textbf{Magreza (\%)} & \textbf{Eutrofia bruto} & \textbf{Eutrofia (\%)} & \textbf{Risco sobrepeso bruto} & \textbf{Risco sobrepeso (\%)} & \textbf{Sobrepeso bruto} & \textbf{Sobrepeso (\%)} & \textbf{Obesidade bruto} & \textbf{Obesidade (\%)} & \textbf{Total} & \textbf{Percent. sobrepeso total} \\ \midrule \endfirsthead
\multicolumn{15}{@{}l}{\footnotesize\itshape (continuação)}\\ \toprule \textbf{Ano} & \textbf{Magreza acentuada bruto} & \textbf{Magreza acentuada (\%)} & \textbf{Magreza bruto} & \textbf{Magreza (\%)} & \textbf{Eutrofia bruto} & \textbf{Eutrofia (\%)} & \textbf{Risco sobrepeso bruto} & \textbf{Risco sobrepeso (\%)} & \textbf{Sobrepeso bruto} & \textbf{Sobrepeso (\%)} & \textbf{Obesidade bruto} & \textbf{Obesidade (\%)} & \textbf{Total} & \textbf{Percent. sobrepeso total} \\ \midrule \endhead
\midrule \multicolumn{15}{r@{}}{\footnotesize\itshape (continua)}\\ \endfoot
\bottomrule \endlastfoot
2008 & 526 & 3,0 & 336 & 1,9 & 10.835 & 61,3 & 3.505 & 19,8 & 1.419 & 8,0 & 1.050 & 5,9 & 17.671 & 14,0 \\
2009 & 679 & 2,9 & 514 & 2,2 & 13.825 & 59,6 & 4.551 & 19,6 & 2.004 & 8,6 & 1.623 & 0,1 & 23.196 & 8,7 \\
2010 & 848 & 3,6 & 642 & 2,8 & 13.652 & 58,4 & 4.526 & 19,4 & 1.948 & 8,3 & 1.756 & 7,5 & 23.372 & 15,8 \\
2011 & 713 & 3,9 & 540 & 2,9 & 10.154 & 55,1 & 3.599 & 19,5 & 1.638 & 8,9 & 1.785 & 9,7 & 18.429 & 18,6 \\
2012 & 1.013 & 3,6 & 738 & 2,6 & 15.669 & 56,2 & 5.443 & 19,5 & 2.599 & 9,3 & 2.409 & 8,6 & 27.871 & 18,0 \\
2013 & 2.371 & 5,0 & 1.432 & 3,0 & 24.549 & 51,5 & 9.052 & 0,2 & 4.819 & 10,1 & 5.413 & 11,4 & 47.636 & 21,5 \\
2014 & 2.650 & 4,6 & 1.744 & 3,0 & 30.849 & 53,2 & 11.581 & 20,0 & 5.665 & 9,8 & 5.452 & 9,4 & 57.941 & 19,2 \\
2015 & 2.160 & 3,1 & 1.917 & 2,7 & 40.054 & 57,3 & 13.823 & 19,8 & 6.527 & 9,3 & 5.441 & 7,8 & 69.922 & 17,1 \\
2016 & 1.950 & 3,6 & 1.524 & 2,8 & 30.532 & 56,5 & 10.046 & 18,6 & 4.897 & 9,1 & 5.141 & 9,5 & 54.090 & 18,6 \\
2017 & 1.856 & 3,0 & 1.389 & 2,2 & 34.727 & 56,2 & 12.365 & 20,0 & 5.651 & 9,2 & 5.753 & 9,3 & 61.741 & 18,5 \\
2018 & 7.705 & 9,6 & 2.334 & 2,9 & 41.305 & 51,6 & 13.717 & 17,1 & 6.383 & 8,0 & 8.623 & 10,8 & 80.067 & 18,7 \\
2019 & 4.145 & 4,4 & 2.789 & 2,9 & 56.215 & 59,1 & 17.025 & 17,9 & 7.640 & 8,0 & 7.330 & 7,7 & 95.144 & 15,7 \\
2020 & 3.748 & 2,7 & 3.758 & 2,7 & 89.019 & 63,5 & 25.857 & 18,4 & 10.224 & 7,3 & 7.505 & 5,4 & 140.111 & 12,7 \\
2021 & 4.204 & 3,0 & 4.052 & 2,9 & 88.592 & 63,6 & 25.368 & 18,2 & 9.611 & 6,9 & 7.531 & 5,4 & 139.358 & 12,3 \\
2022 & 3.559 & 2,3 & 3.679 & 2,4 & 100.188 & 65,8 & 25.412 & 16,7 & 10.336 & 6,8 & 9.154 & 6,0 & 152.328 & 12,8 \\
2023 & 3.079 & 1,6 & 4.153 & 2,1 & 130.109 & 66,5 & 33.045 & 16,9 & 14.095 & 7,2 & 11.289 & 5,8 & 195.770 & 13,0 \\
2024 & 2.921 & 1,4 & 4.323 & 2,0 & 138.845 & 65,4 & 37.423 & 17,6 & 16.503 & 7,8 & 12.223 & 5,8 & 212.238 & 13,5 \\
2025 & 3.542 & 1,7 & 4.791 & 2,2 & 127.884 & 59,9 & 41.227 & 19,3 & 20.388 & 9,6 & 15.674 & 7,3 & 213.506 & 16,9 \\
2026 & 1.949 & 1,6 & 3.266 & 2,7 & 68.972 & 56,9 & 25.982 & 21,4 & 12.333 & 10,2 & 8.667 & 7,2 & 121.169 & 17,3 \\
\end{longtable}}
\vspace{-6pt}\fonte{Elaboração IPP com dados de \citeonline{ms_sisvan}.}
\end{landscape}


\chapter{Tabelas do eixo Moradia}\label{ap:eixo-6}

Este eixo não tem tabelas de dados nesta edição.

\end{apendicesenv}


\end{document}
